\documentclass[11pt,reqno]{amsart}
\usepackage[letterpaper,margin=1in]{geometry}
\usepackage[T1]{fontenc}
\usepackage{lmodern,microtype,amsmath,amssymb,mathtools,mathrsfs}
\usepackage{enumitem,tikz,aliascnt,float,array,longtable}
\usetikzlibrary{arrows.meta,calc,decorations.markings,positioning}
\usepackage[colorlinks=true,linkcolor=blue!45!black,citecolor=blue!45!black,urlcolor=blue!45!black]{hyperref}
\usepackage[nameinlink,capitalise,noabbrev]{cleveref}
\numberwithin{equation}{section}
\newtheorem{theorem}{Theorem}[section]
\newaliascnt{proposition}{theorem}
\newtheorem{proposition}[proposition]{Proposition}
\aliascntresetthe{proposition}
\newaliascnt{lemma}{theorem}
\newtheorem{lemma}[lemma]{Lemma}
\aliascntresetthe{lemma}
\newaliascnt{corollary}{theorem}
\newtheorem{corollary}[corollary]{Corollary}
\aliascntresetthe{corollary}
\theoremstyle{definition}
\newaliascnt{definition}{theorem}
\newtheorem{definition}[definition]{Definition}
\aliascntresetthe{definition}
\newaliascnt{remark}{theorem}
\newtheorem{remark}[remark]{Remark}
\aliascntresetthe{remark}
\crefname{proposition}{Proposition}{Propositions}
\crefname{lemma}{Lemma}{Lemmas}
\crefname{corollary}{Corollary}{Corollaries}
\crefname{definition}{Definition}{Definitions}
\crefname{remark}{Remark}{Remarks}
\newcommand{\R}{\mathbb R}
\newcommand{\eps}{\varepsilon}
\newcommand{\GammaI}{\Gamma(C_0)}
\newcommand{\distH}{d_{\mathrm H}}
\newcommand{\dd}{\,\mathrm d}
\newcommand{\e}{\mathrm e}
\newcommand{\cD}{\mathcal D}
\newcommand{\cH}{\mathcal H}
\newcommand{\cP}{\mathcal P}
\newcommand{\cV}{\mathcal V}
\DeclareMathOperator{\PV}{PV}
\DeclareMathOperator{\Cycl}{Cycl}
\title[Sharp local cyclicity of the $I_2^1$ graphic]
  {Sharp local cyclicity of the $I_2^1$ graphic in a five-parameter quadratic family}
\author{Haibo Lu}
\address{Shanghai Institute of Technology}
\email{luhaibo1985@gmail.com}
\date{}
\subjclass[2020]{34C07, 34C23, 34E15}
\keywords{local cyclicity, quadratic vector field, nilpotent saddle--node, return map}
\hypersetup{pdftitle={Sharp local cyclicity of the I\unichar{"2082}\unichar{"00B9} graphic in a five-parameter quadratic family},
  pdfauthor={Haibo Lu},
  pdfsubject={Local cyclicity exactly three for each fixed C\unichar{"2080}>1}}
\begin{document}
\begin{abstract}
The DRR graphic $I_2^1$ is a loop through a nilpotent saddle--node at
infinity. We prove that its local cyclicity is exactly three in a
five-parameter quadratic family, for every fixed organizing value
$C_0>1$. The upper bound holds throughout a full parameter neighborhood
and counts periodic orbits close to the whole graphic in Hausdorff distance.
We parameterize both halves of the singular passage by one central height
and compare their endpoint derivatives after the exterior transition.
Drift-weighted errors preserve this comparison when boundary roots merge
or passage times diverge. With $C=C_0$ fixed, direct action matching and a
transverse variation of the connection splitting yield three limit cycles
converging to the whole graphic. They are hyperbolic and persist on open
parameter sets approaching the organizing value. Thus the local upper
bound is attained within the same quadratic family.
\end{abstract}
\maketitle
\setcounter{tocdepth}{1}
\tableofcontents

\section{Local cyclicity of \texorpdfstring{\(I_2^1\)}{I21}}
\label{sec:introduction}

The Dumortier--Roussarie--Rousseau program reduces the finiteness
problem for quadratic limit cycles to local cyclicity questions for
specified limit-periodic sets~\cite{DRR1994}. We study the loop
\(I_2^1\) through a nilpotent saddle--node at infinity. Our aim is
to count the cycles approaching the entire loop, including parameter
directions in which boundary roots merge or coefficient scales vanish
independently, and to attain the resulting bound in the original family.

Fix \(C_0>1\). We consider the five-parameter quadratic family
\begin{equation}
\label{eq:quadratic-family}
\begin{aligned}
 \dot x&=Ax-y+x^2+(\mu_2+\mu_3)xy+\mu_1y^2,\\
 \dot y&=Cx+x^2+xy+\mu_3y^2
\end{aligned}
\end{equation}
near the parameter
\begin{equation}
\label{eq:organizing-parameter}
 (A,C,\mu_1,\mu_2,\mu_3)=(1,C_0,0,0,0).
\end{equation}
At this parameter the invariant parabola
\begin{equation}
\label{eq:organizing-parabola}
 y=(x^2-C_0)/2
\end{equation}
and its resolved nilpotent saddle--node passage at infinity form
the directed graphic \(\GammaI\), denoted \(I_2^1\) in the DRR
classification \cite[Proposition~4.2.2(6), Figure~8(6)]{DRR1994}
(see \cref{sec:graphic-central-return,fig:organizing-graphic-new}).

We measure proximity to this graphic on the Poincar\'e sphere
\(\mathbb S^2\), where the Poincar\'e compactification of
\eqref{eq:quadratic-family} is an analytic vector field. Let
\(\distH\) denote Hausdorff distance for a fixed metric on
\(\mathbb S^2\). When resolved or scaled coordinates are used below,
the sets are projected back to \(\mathbb S^2\) before their Hausdorff
distance is taken.

For \(\varepsilon>0\) and a neighborhood \(V\subset\mathbb R^5\) of the
organizing parameter \eqref{eq:organizing-parameter}, let
\(N(V,\varepsilon)\) be the supremum over \(p\in V\) of the number of
limit cycles \(\gamma\) of \eqref{eq:quadratic-family} with parameter
\(p\) such that \(\distH(\overline\gamma,\GammaI)<\varepsilon\).
Following the Hausdorff formulation of cyclicity in
\cite[Definition~2.2(2)]{DRR1994} (see also
\cite[\S2.2.1]{Roussarie1998}), the cyclicity of
\(\GammaI\) in the family \eqref{eq:quadratic-family} is
\begin{equation}
\label{eq:cyclicity-definition}
 \Cycl(\GammaI)=\inf_{V,\,\varepsilon}N(V,\varepsilon).
\end{equation}
The Hausdorff condition requires proximity to the whole graphic.

\begin{theorem}[Local upper bound]
\label{thm:direct-finite-cyclicity}
For every fixed \(C_0>1\), there exist a neighborhood
\(V\subset\mathbb R^5\) of \((1,C_0,0,0,0)\) and a number
\(\varepsilon>0\) such that, for every \(p\in V\), the vector field
\eqref{eq:quadratic-family} with parameter \(p\) has at most three
periodic orbits \(\gamma\) with
\[
 \distH(\overline\gamma,\Gamma(C_0))<\varepsilon .
\]
Each of these periodic orbits is a limit cycle. In particular,
\(\Cycl(\GammaI)\le3\).
\end{theorem}

\begin{proposition}[Three periodic orbits]
\label{prop:three-cycle-lower-bound}
For every fixed \(C_0>1\), there is a sequence of parameters
\(p_n\to(1,C_0,0,0,0)\), with \(C=C_0\), for which
\eqref{eq:quadratic-family} has three distinct isolated periodic
orbits \(\gamma_{j,n}\), \(j=1,2,3\), satisfying
\[
 \distH(\overline{\gamma_{j,n}},\Gamma(C_0))\longrightarrow0.
\]
These orbits eventually lie in every fixed sufficiently small
neighborhood of the whole graphic, and in particular satisfy the
Hausdorff condition of \cref{thm:direct-finite-cyclicity}.
\end{proposition}

\begin{corollary}[Cyclicity three]
\label{cor:exact-local-cyclicity-three}
For every fixed \(C_0>1\), \(\Cycl(\GammaI)=3\).
\end{corollary}

The three orbits of \cref{prop:three-cycle-lower-bound} are also
hyperbolic and persist on open sets in the original five-parameter space
arbitrarily close to the organizing parameter
(\cref{cor:hyperbolic-triple}).

The neighborhood \(V\) and the radius \(\varepsilon\) depend on
\(C_0\). Their construction uses, in particular, the positive
speed \eqref{eq:parabola-speed-new} of the organizing orbit and the
limiting contraction multiplier \(\Lambda_0<1\) of
\cref{lem:exterior-multiplier-new}.
The speed bound \((C_0-1)/2\) tends to zero as \(C_0\downarrow1\),
and the margin \(1-\Lambda_0\) tends to zero as \(C_0\to\infty\).

\subsection*{Relation to previous work}

The DRR program develops normal forms near saddle--nodes and applies them to
finite-cyclicity questions; see Dumortier, Ilyashenko, and Rousseau
\cite{DIR2002}. Related parts of the program treat the degenerate graphics
$DF_{1a}$ and $DF_{2a}$ \cite{DR2009}, with the latter study completed in the
remaining chart by Huzak \cite{Huzak2018}. For graphics through
triple nilpotent points, Roussarie and Rousseau study the center graphic
$I_{14}^1$ and distinguish it from the blown-up boundary limit-periodic sets
appearing in the same analysis \cite{RR2015}; Rousseau, Shan, and Zhu treat
the saddle-type graphics $I_{12}^1$ and $I_{13}^1$ \cite{RSZ2016}.

Huzak and Kristiansen develop entry--exit formulas for degenerate turning
points and apply them to the present five-parameter family
\cite[Theorem~2.4 and \S6]{HK2026}. Their theorem constructs a local Dulac
map, with finite-order remainder estimates in a logarithmic parameterization,
in a principal root-free region where the scaled boundary drift has a uniform
sign margin. Their Section~6 derives the entry--exit relation for this
family and leaves the cyclicity analysis to further work. The passage
problem here also includes merging boundary
roots, roots approaching the central section, and independently
vanishing coefficient scales. We compare the two endpoint maps on the
full interval of admissible central heights, compose them with the
exterior transition, and use geometric capture to count the periodic
orbits satisfying the whole-graphic Hausdorff condition.

The root-free calculation enters our buffered compact-exit analysis
through the leading relation between normalized exit radii,
\[
 f_\Lambda(\rho)=\frac{\Lambda\rho}{1+(1+\Lambda)\rho}.
\]
Here \(\rho\) is the normalized positive exit radius and \(\Lambda\)
is the boundary-action multiplier. This relation is the inverse,
after a sign change and positive rescaling, of the base-point relation
in Huzak--Kristiansen \cite[(2.13) and (6.8)]{HK2026}.
The change of variables is given after
\eqref{eq:direct-buffered-fractional-relation}.

Uniform \(C^3\) matching transfers this relation to the endpoint maps
used for counting. Drift-weighted estimates treat the other passage
regions, and a scattering comparison treats contact with the central
section. The parameter cover combines these comparisons into the
upper bound. A direct matching construction and variation of the
connection splitting then attain it with \(C=C_0\).

Finite-order passage analysis for canard cycles with singular slow dynamics,
including divergence comparisons and exponential-flatness estimates, was
developed by De Maesschalck and Dumortier \cite{DMD2008}. For limit-periodic
sets associated with singular slow dynamics, their later work obtains
cycle bounds by studying suitable derivatives of the slow-divergence integral,
even when that integral is unbounded \cite{DMD2011}.
Mar\'in and Villadelprat give parameter-dependent
hyperbolic-saddle Dulac expansions, including the near-resonant case
\cite{MV2024}. These works provide the background for the passage
estimates and derivative comparisons developed here.

\subsection*{Proof strategy}

The proof has two counting steps. For fixed parameters, we control
one return displacement \(\Delta_p\) on its full interval
\(\mathcal H_p\) of central heights. A finite cover of parameter
space then selects an applicable estimate throughout a neighborhood
of the organizing value. Geometric capture connects these zeros to
the periodic orbits in the theorem.

\emph{A common height for the two halves.}
In normal-rate coordinates \((\eta,z)\), the field is
\[
 \dot\eta=q(\eta,p)-zB(\eta,z,p),\qquad \dot z=-\eta z.
\]
Here \(q(\eta,p)=q_0+q_1\eta+q_2\eta^2+\eta^3R_q(\eta,p)\)
is the boundary drift. At fixed nearby \(A,C\), its coefficients
\((q_0,q_1,q_2)\) are local coordinates in
\((\mu_1,\mu_2,\mu_3)\) (Appendix~\ref{app:coordinates}).
The central height \(h\) is the coordinate \(z\) on \(\eta=0\).
From \((0,h)\), follow the positive half-orbit backward and the negative
half-orbit forward. Their endpoint functions \(I_\pm(h)\) are defined
on the same interval \(\mathcal H_p\). The regular exterior map
\(\Pi_p\) maps the negative endpoint window to the positive one.
The return displacement is
\[
 \Delta_p(h)=\Pi_p(I_-(h))-I_+(h).
\]
A zero expresses the matching of the two endpoints through the exterior
transition.

\emph{From endpoint derivatives to a zero count.}
Both endpoint derivatives with respect to \(h\) are positive. The
logarithmic slope ratio
\(\mathcal L=\log[\Pi_p'(I_-)I_{-,h}/I_{+,h}]\) satisfies
\[
 \Delta_{p,h}=I_{+,h}(\e^{\mathcal L}-1).
\]
The zeros of \(\mathcal L\) are therefore the critical points of
\(\Delta_p\). To bound them on the full height interval, we combine
crossing directions on the outer height ranges with negative
curvature at critical points on the middle range. The order of these
ranges allows at most two zeros of \(\mathcal L\); Rolle's theorem
then gives at most three zeros of \(\Delta_p\).

\emph{Uniformity near degenerate passages.}
The comparison must remain uniform when roots degenerate and coefficient
scales vanish independently. Each endpoint derivative retains the tail
multiplier \((z_*/t)^{-2q_2}\), even though the common quadratic term
cancels from their ratio; this allows unbounded \(q_2\log t\).
Near a boundary root, the leading curvature vanishes with the drift.
The boundary-flow coordinate of \cref{lem:direct-boundary-flow}
absorbs the height-dependent radial term, and
\cref{lem:direct-weighted-phase} makes the phase correction smaller
than any prescribed finite drift power after a fixed positive
boundary-time advance. The resulting drift-weighted curvature formula
\eqref{eq:direct-rate-curvature} preserves every fixed sign margin.

The exterior transition supplies the remaining contraction when the
singular slopes nearly agree. Its organizing derivative tends to
\[
 \Lambda_0=\exp\!\left(-\frac{4\pi}{\sqrt{C_0-1}}\right)<1
\]
as the sections approach infinity. Fixed sections retain a contraction
margin, while the nonconstant exterior term carries the same drift
weights as the passage terms. At central contact, a decreasing scattering
derivative ratio with a small \(C^2\) error gives at most one middle
critical point. The outer derivative signs then give at most two
contact zeros on the full height interval.

\emph{Realizing three cycles.}
For the lower bound, keep \(C=C_0\) and choose a coefficient path in
the root-free region with negative quadratic drift. Comparing the two
boundary actions gives a compact exit profile with a strict negative
minimum. The parameter \(A\) leaves the boundary drift unchanged and
splits the organizing connection transversely, so it can adjust the
additive constant of this profile. Three test exits then give two
displacement zeros.

This compact exit comparison alone does not locate the third cycle.
Returning to the common central height, the organizing contraction gives
the opposite displacement sign at a larger height. This supplies a third
zero beyond the heights tested by the compact profile. Graph order and
regular continuation then show that all three orbits converge to the
whole graphic. The coefficient path, action matching, and splitting
calculation are developed in \cref{sec:sharpness}; hyperbolicity and the
separation of the third height follow in
\cref{cor:hyperbolic-triple,cor:third-height-separation}.

The proof is organized around these comparisons. The geometry and the
common-height return problem are established in
\cref{sec:graphic-central-return}, and
\cref{sec:standard-tools,sec:direct-passage-estimates} give the counting
principles and uniform passage estimates. Contact and positive offsets
are treated in \cref{sec:contact-positive-offsets}. For negative offsets,
\cref{sec:recurring-ranges} develops the height-range comparisons used on
both coefficient scales; the primary-scale analysis is in
\cref{sec:primary-regimes}, and the secondary-scale cases begin in
\cref{sec:quadratic-regimes}. These are estimates on overlapping height
intervals for fixed parameters. Their coverage of a full parameter
neighborhood is proved separately in \cref{sec:direct-parameter-cover},
where \cref{tab:direct-cover-guide} identifies the applicable comparison
and bound in each region.


\section{The graphic and its central displacement}
\label{sec:graphic-central-return}

To count periodic orbits near the whole graphic, we first show that each
of them makes one passage through a fixed system of sections. Its central
height then parameterizes both halves of the singular passage. We use
this common parameter to construct the displacement and its interval of
admissible heights.

\subsection{The organizing orbit and the projective field}

At the parameter \eqref{eq:organizing-parameter}, set
\[
 H_0(x,y)=y-\frac{x^2}{2}+\frac{C_0}{2}.
\]
A direct calculation gives
\begin{equation}
\label{eq:parabola-invariance-new}
 \dot H_0=2xH_0.
\end{equation}
Thus \eqref{eq:organizing-parabola} is invariant.  On this parabola,
\begin{equation}
\label{eq:parabola-speed-new}
 \dot x=\frac{(x+1)^2+C_0-1}{2}>0.
\end{equation}
It follows that the parabola contains one regular orbit, directed from
$x=-\infty$ to $x=+\infty$.

Near the positive $y$-direction at infinity, put
\begin{equation}
\label{eq:projective-coordinates-new}
 x=\frac vz,\qquad y=\frac1z,\qquad z>0,
\end{equation}
and multiply the transformed vector field by the positive factor $z$.  We
obtain
\begin{equation}
\label{eq:projective-family-new}
\begin{aligned}
 v'&=P(v)-(1-Av+Cv^2)z,\\
 z'&=-z\eta(v,z),
\end{aligned}
\end{equation}
where
\begin{equation}
\label{eq:projective-factors-new}
\begin{aligned}
 P(v)&=\mu _1+\mu _2v-v^3,\\
 \eta(v,z)&=\mu _3+v+v^2+Cvz.
\end{aligned}
\end{equation}
The boundary drift \(q\) below therefore depends on the three
\(\mu\)-parameters, whereas the positive factor \(B\) retains the
dependence on \((A,C)\). This separation is used later when \(A\) is
varied to change the connection without changing the boundary action.
After shrinking the $(A,C)$-neighborhood of $(1,C_0)$,
$A^2<4C$ and hence
\begin{equation}
\label{eq:K-positive-new}
 1-Av+Cv^2>0\qquad(v\in\R).
\end{equation}
At the organizing parameter, $(v,z)=(0,0)$ is the nilpotent saddle--node
where the two ends of the affine orbit meet. Together with the affine
orbit, it forms the compact directed graphic $\GammaI$.

\begin{figure}[t]
\centering
% Generated by scripts/figure_organizing_graphic.py; do not edit.
% C_0 = 10, R = 3.
\newcommand{\figOneSpiral}[1]{\draw[#1] (0.016,0.000) -- (0.017,0.015) -- (0.016,0.030) -- (0.014,0.045) -- (0.010,0.059) -- (-0.003,0.067) -- (-0.013,0.056) -- (-0.018,0.042) -- (-0.022,0.027) -- (-0.025,0.013) -- (-0.027,-0.002) -- (-0.028,-0.017) -- (-0.028,-0.032) -- (-0.027,-0.047) -- (-0.025,-0.063) -- (-0.022,-0.077) -- (-0.018,-0.091) -- (-0.010,-0.104) -- (0.004,-0.108) -- (0.015,-0.099) -- (0.023,-0.086) -- (0.029,-0.072) -- (0.033,-0.058) -- (0.037,-0.042) -- (0.040,-0.027) -- (0.042,-0.012) -- (0.044,0.003) -- (0.045,0.019) -- (0.046,0.035) -- (0.046,0.051) -- (0.046,0.068) -- (0.045,0.084) -- (0.043,0.099) -- (0.041,0.114) -- (0.038,0.128) -- (0.034,0.143) -- (0.029,0.158) -- (0.022,0.171) -- (0.012,0.182) -- (-0.003,0.186) -- (-0.016,0.179) -- (-0.027,0.168) -- (-0.034,0.155) -- (-0.041,0.142) -- (-0.047,0.126) -- (-0.052,0.112) -- (-0.057,0.095) -- (-0.060,0.081) -- (-0.063,0.065) -- (-0.066,0.049) -- (-0.069,0.032) -- (-0.071,0.015) -- (-0.073,-0.003) -- (-0.074,-0.018) -- (-0.075,-0.032) -- (-0.075,-0.047) -- (-0.076,-0.062) -- (-0.076,-0.077) -- (-0.075,-0.092) -- (-0.075,-0.106) -- (-0.074,-0.121) -- (-0.072,-0.138) -- (-0.070,-0.156) -- (-0.067,-0.172) -- (-0.064,-0.188) -- (-0.061,-0.203) -- (-0.057,-0.217) -- (-0.052,-0.233) -- (-0.046,-0.247) -- (-0.039,-0.261) -- (-0.029,-0.274) -- (-0.018,-0.284) -- (-0.004,-0.289) -- (0.010,-0.287) -- (0.024,-0.280) -- (0.035,-0.270) -- (0.044,-0.259) -- (0.053,-0.245) -- (0.060,-0.232) -- (0.068,-0.216) -- (0.073,-0.202) -- (0.079,-0.188) -- (0.084,-0.172) -- (0.089,-0.155) -- (0.094,-0.136) -- (0.098,-0.122) -- (0.101,-0.107) -- (0.105,-0.092) -- (0.108,-0.076) -- (0.111,-0.059) -- (0.113,-0.043) -- (0.116,-0.025) -- (0.118,-0.007) -- (0.120,0.011) -- (0.122,0.029) -- (0.123,0.048) -- (0.125,0.067) -- (0.126,0.086) -- (0.126,0.105) -- (0.127,0.125) -- (0.127,0.144) -- (0.127,0.164) -- (0.126,0.183) -- (0.125,0.203) -- (0.124,0.222) -- (0.123,0.241) -- (0.121,0.260) -- (0.119,0.279) -- (0.117,0.297) -- (0.114,0.315) -- (0.111,0.332) -- (0.108,0.349) -- (0.104,0.365) -- (0.100,0.381) -- (0.096,0.396) -- (0.092,0.410) -- (0.087,0.424) -- (0.080,0.441) -- (0.073,0.456) -- (0.065,0.470) -- (0.057,0.483) -- (0.046,0.496) -- (0.035,0.506) -- (0.021,0.514) -- (0.006,0.518) -- (-0.009,0.518) -- (-0.025,0.513) -- (-0.038,0.505) -- (-0.051,0.494) -- (-0.062,0.483) -- (-0.072,0.469) -- (-0.083,0.454) -- (-0.091,0.441) -- (-0.098,0.427) -- (-0.106,0.412) -- (-0.113,0.395) -- (-0.121,0.378) -- (-0.128,0.359) -- (-0.135,0.340) -- (-0.142,0.319) -- (-0.146,0.305) -- (-0.150,0.290) -- (-0.154,0.275) -- (-0.158,0.259) -- (-0.162,0.243) -- (-0.166,0.227) -- (-0.169,0.210) -- (-0.173,0.193) -- (-0.176,0.176) -- (-0.179,0.158) -- (-0.182,0.141) -- (-0.185,0.123) -- (-0.188,0.104) -- (-0.190,0.086) -- (-0.192,0.067) -- (-0.194,0.048) -- (-0.196,0.029) -- (-0.198,0.010) -- (-0.199,-0.010) -- (-0.201,-0.029) -- (-0.202,-0.048) -- (-0.203,-0.068) -- (-0.203,-0.087) -- (-0.204,-0.107) -- (-0.204,-0.126) -- (-0.204,-0.146) -- (-0.204,-0.165) -- (-0.203,-0.184) -- (-0.203,-0.203) -- (-0.202,-0.222) -- (-0.201,-0.241) -- (-0.200,-0.259) -- (-0.199,-0.278) -- (-0.197,-0.296) -- (-0.195,-0.314) -- (-0.194,-0.331) -- (-0.191,-0.349) -- (-0.189,-0.366) -- (-0.187,-0.383) -- (-0.184,-0.399) -- (-0.181,-0.415) -- (-0.178,-0.431) -- (-0.175,-0.447) -- (-0.172,-0.462) -- (-0.168,-0.476) -- (-0.165,-0.491) -- (-0.161,-0.505) -- (-0.155,-0.525) -- (-0.149,-0.544) -- (-0.142,-0.563) -- (-0.135,-0.580) -- (-0.128,-0.596) -- (-0.121,-0.612) -- (-0.113,-0.626) -- (-0.105,-0.639) -- (-0.097,-0.652) -- (-0.085,-0.666) -- (-0.074,-0.679) -- (-0.062,-0.690) -- (-0.049,-0.699) -- (-0.037,-0.706) -- (-0.020,-0.711) -- (-0.004,-0.714) -- (0.014,-0.713) -- (0.031,-0.708) -- (0.045,-0.702) -- (0.060,-0.694) -- (0.074,-0.683) -- (0.086,-0.673) -- (0.097,-0.662) -- (0.108,-0.649) -- (0.119,-0.634) -- (0.131,-0.619) -- (0.142,-0.601) -- (0.150,-0.588) -- (0.157,-0.575) -- (0.165,-0.561) -- (0.172,-0.546) -- (0.180,-0.531) -- (0.188,-0.514) -- (0.195,-0.497) -- (0.203,-0.479) -- (0.210,-0.460) -- (0.217,-0.440) -- (0.225,-0.420) -- (0.232,-0.399) -- (0.239,-0.376) -- (0.246,-0.353) -- (0.253,-0.329) -- (0.259,-0.305) -- (0.266,-0.279) -- (0.273,-0.253) -- (0.279,-0.226) -- (0.285,-0.198) -- (0.288,-0.184) -- (0.291,-0.169) -- (0.294,-0.154) -- (0.296,-0.140) -- (0.299,-0.124) -- (0.302,-0.109) -- (0.304,-0.094) -- (0.307,-0.078) -- (0.310,-0.063) -- (0.312,-0.047) -- (0.314,-0.031) -- (0.317,-0.014) -- (0.319,0.002) -- (0.321,0.018) -- (0.323,0.035) -- (0.325,0.052) -- (0.327,0.069) -- (0.329,0.086) -- (0.330,0.103) -- (0.332,0.120) -- (0.333,0.138) -- (0.335,0.155) -- (0.336,0.173) -- (0.338,0.190) -- (0.339,0.208) -- (0.340,0.226) -- (0.341,0.243) -- (0.342,0.261) -- (0.342,0.279) -- (0.343,0.297) -- (0.344,0.315) -- (0.344,0.333) -- (0.344,0.351) -- (0.345,0.369) -- (0.345,0.387) -- (0.345,0.405) -- (0.345,0.422) -- (0.345,0.440) -- (0.344,0.458) -- (0.344,0.476) -- (0.344,0.494) -- (0.343,0.511) -- (0.342,0.529) -- (0.342,0.546) -- (0.341,0.563) -- (0.340,0.581) -- (0.339,0.598) -- (0.337,0.615) -- (0.336,0.632) -- (0.335,0.649) -- (0.333,0.665) -- (0.331,0.682) -- (0.330,0.698) -- (0.328,0.714) -- (0.326,0.730) -- (0.324,0.746) -- (0.322,0.762) -- (0.320,0.777) -- (0.318,0.793) -- (0.315,0.808) -- (0.313,0.823) -- (0.310,0.838) -- (0.308,0.852) -- (0.305,0.867) -- (0.302,0.881) -- (0.296,0.909) -- (0.290,0.935) -- (0.284,0.962) -- (0.277,0.987) -- (0.270,1.011) -- (0.263,1.034) -- (0.256,1.057) -- (0.248,1.079) -- (0.240,1.099) -- (0.232,1.119) -- (0.224,1.138) -- (0.216,1.156) -- (0.207,1.173) -- (0.198,1.190) -- (0.189,1.205) -- (0.180,1.220) -- (0.171,1.233) -- (0.162,1.246) -- (0.152,1.258) -- (0.142,1.269) -- (0.128,1.285) -- (0.113,1.298) -- (0.097,1.310) -- (0.082,1.319) -- (0.066,1.328) -- (0.050,1.334) -- (0.034,1.338) -- (0.018,1.341) -- (0.001,1.342) -- (-0.015,1.341) -- (-0.032,1.339) -- (-0.050,1.335) -- (-0.067,1.328) -- (-0.085,1.320) -- (-0.102,1.310) -- (-0.120,1.298) -- (-0.132,1.289) -- (-0.144,1.279) -- (-0.157,1.268) -- (-0.169,1.256) -- (-0.181,1.243) -- (-0.194,1.229) -- (-0.206,1.214) -- (-0.219,1.197) -- (-0.231,1.180) -- (-0.244,1.162) -- (-0.256,1.142) -- (-0.269,1.122) -- (-0.281,1.100) -- (-0.294,1.077) -- (-0.306,1.053) -- (-0.318,1.027) -- (-0.325,1.014) -- (-0.331,1.000) -- (-0.337,0.986) -- (-0.343,0.972) -- (-0.349,0.958) -- (-0.355,0.943) -- (-0.361,0.928) -- (-0.367,0.912) -- (-0.373,0.896) -- (-0.379,0.880) -- (-0.385,0.864) -- (-0.390,0.847) -- (-0.396,0.830) -- (-0.402,0.812) -- (-0.407,0.794) -- (-0.413,0.776) -- (-0.418,0.758) -- (-0.423,0.739) -- (-0.429,0.720) -- (-0.434,0.701) -- (-0.439,0.681) -- (-0.444,0.661) -- (-0.449,0.641) -- (-0.454,0.620) -- (-0.458,0.599) -- (-0.463,0.578) -- (-0.467,0.557) -- (-0.472,0.535) -- (-0.476,0.513) -- (-0.480,0.491) -- (-0.484,0.469) -- (-0.488,0.447) -- (-0.491,0.424) -- (-0.495,0.401) -- (-0.498,0.378) -- (-0.502,0.355) -- (-0.505,0.331) -- (-0.508,0.308) -- (-0.511,0.284) -- (-0.513,0.260) -- (-0.516,0.236) -- (-0.518,0.212) -- (-0.520,0.188) -- (-0.522,0.164) -- (-0.524,0.140) -- (-0.526,0.116) -- (-0.527,0.091) -- (-0.528,0.067) -- (-0.530,0.043) -- (-0.531,0.019) -- (-0.531,-0.006) -- (-0.532,-0.030) -- (-0.533,-0.054) -- (-0.533,-0.078) -- (-0.533,-0.101) -- (-0.533,-0.125) -- (-0.533,-0.149) -- (-0.533,-0.172) -- (-0.532,-0.196) -- (-0.532,-0.219) -- (-0.531,-0.242) -- (-0.530,-0.265) -- (-0.529,-0.287) -- (-0.528,-0.310) -- (-0.526,-0.332) -- (-0.525,-0.354) -- (-0.523,-0.376) -- (-0.521,-0.398) -- (-0.519,-0.419) -- (-0.517,-0.440) -- (-0.515,-0.461) -- (-0.513,-0.482) -- (-0.510,-0.502) -- (-0.508,-0.522) -- (-0.505,-0.542) -- (-0.503,-0.562) -- (-0.500,-0.581) -- (-0.497,-0.600) -- (-0.494,-0.619) -- (-0.491,-0.638) -- (-0.487,-0.656) -- (-0.484,-0.674) -- (-0.481,-0.692) -- (-0.477,-0.709) -- (-0.474,-0.726) -- (-0.470,-0.743) -- (-0.466,-0.760) -- (-0.463,-0.776) -- (-0.459,-0.792) -- (-0.455,-0.808) -- (-0.451,-0.823) -- (-0.447,-0.838) -- (-0.443,-0.853) -- (-0.439,-0.868) -- (-0.435,-0.882) -- (-0.430,-0.896) -- (-0.426,-0.910) -- (-0.417,-0.937) -- (-0.409,-0.963) -- (-0.400,-0.988) -- (-0.390,-1.012) -- (-0.381,-1.035) -- (-0.372,-1.057) -- (-0.362,-1.079) -- (-0.353,-1.099) -- (-0.343,-1.118) -- (-0.334,-1.137) -- (-0.324,-1.155) -- (-0.314,-1.172) -- (-0.304,-1.189) -- (-0.294,-1.204) -- (-0.284,-1.219) -- (-0.274,-1.233) -- (-0.264,-1.247) -- (-0.254,-1.260) -- (-0.244,-1.272) -- (-0.234,-1.284) -- (-0.224,-1.295) -- (-0.213,-1.305) -- (-0.198,-1.320) -- (-0.182,-1.333) -- (-0.167,-1.345) -- (-0.151,-1.356) -- (-0.135,-1.366) -- (-0.119,-1.374) -- (-0.103,-1.382) -- (-0.087,-1.388) -- (-0.071,-1.394) -- (-0.055,-1.398) -- (-0.038,-1.401) -- (-0.021,-1.403) -- (-0.004,-1.404) -- (0.013,-1.403) -- (0.030,-1.402) -- (0.048,-1.399) -- (0.066,-1.395) -- (0.084,-1.390) -- (0.102,-1.384) -- (0.121,-1.376) -- (0.140,-1.367) -- (0.153,-1.360) -- (0.166,-1.352) -- (0.180,-1.343) -- (0.193,-1.334) -- (0.207,-1.324) -- (0.221,-1.313) -- (0.235,-1.301) -- (0.249,-1.288) -- (0.264,-1.274) -- (0.279,-1.259) -- (0.294,-1.243) -- (0.309,-1.226) -- (0.325,-1.208) -- (0.341,-1.188) -- (0.357,-1.167) -- (0.373,-1.145) -- (0.382,-1.133) -- (0.390,-1.121) -- (0.399,-1.108) -- (0.407,-1.095) -- (0.416,-1.081) -- (0.425,-1.067) -- (0.433,-1.053) -- (0.442,-1.038) -- (0.451,-1.023) -- (0.460,-1.007) -- (0.469,-0.990) -- (0.479,-0.973) -- (0.488,-0.956) -- (0.497,-0.938) -- (0.506,-0.919) -- (0.516,-0.900) -- (0.525,-0.880) -- (0.535,-0.860) -- (0.544,-0.838) -- (0.554,-0.817) -- (0.564,-0.794) -- (0.573,-0.771) -- (0.583,-0.747) -- (0.593,-0.722) -- (0.603,-0.697) -- (0.612,-0.671) -- (0.622,-0.644) -- (0.632,-0.616) -- (0.642,-0.588) -- (0.651,-0.559) -- (0.661,-0.528) -- (0.670,-0.498) -- (0.680,-0.466) -- (0.689,-0.434) -- (0.698,-0.400) -- (0.707,-0.366) -- (0.716,-0.331) -- (0.725,-0.296) -- (0.734,-0.259) -- (0.742,-0.222) -- (0.751,-0.185) -- (0.758,-0.146) -- (0.766,-0.107) -- (0.774,-0.067) -- (0.781,-0.027) -- (0.788,0.014) -- (0.794,0.055) -- (0.800,0.097) -- (0.806,0.139) -- (0.811,0.182) -- (0.816,0.225) -- (0.821,0.268) -- (0.825,0.311) -- (0.828,0.354) -- (0.832,0.397) -- (0.834,0.440) -- (0.837,0.483) -- (0.839,0.526) -- (0.840,0.569) -- (0.841,0.611) -- (0.842,0.653) -- (0.842,0.694) -- (0.841,0.735) -- (0.840,0.776) -- (0.839,0.816) -- (0.838,0.855) -- (0.836,0.893) -- (0.833,0.931) -- (0.830,0.968) -- (0.827,1.005) -- (0.824,1.040) -- (0.820,1.075) -- (0.816,1.109) -- (0.812,1.142) -- (0.807,1.174) -- (0.802,1.206) -- (0.797,1.237) -- (0.792,1.266) -- (0.787,1.295) -- (0.781,1.323) -- (0.775,1.351) -- (0.769,1.377) -- (0.763,1.403) -- (0.757,1.428) -- (0.751,1.452) -- (0.745,1.475) -- (0.738,1.498) -- (0.732,1.520) -- (0.725,1.541) -- (0.719,1.562) -- (0.712,1.582) -- (0.705,1.601) -- (0.699,1.620) -- (0.692,1.638) -- (0.685,1.656) -- (0.678,1.673) -- (0.672,1.689) -- (0.665,1.705) -- (0.658,1.720) -- (0.651,1.735) -- (0.645,1.750) -- (0.638,1.764) -- (0.631,1.777) -- (0.625,1.790) -- (0.612,1.815) -- (0.599,1.839) -- (0.586,1.861) -- (0.573,1.882) -- (0.561,1.901) -- (0.548,1.919) -- (0.536,1.937) -- (0.524,1.953) -- (0.512,1.968) -- (0.500,1.983) -- (0.489,1.996) -- (0.477,2.009) -- (0.466,2.021) -- (0.455,2.033) -- (0.444,2.044) -- (0.434,2.054) -- (0.423,2.064) -- (0.407,2.078) -- (0.392,2.091) -- (0.377,2.102) -- (0.362,2.113) -- (0.348,2.123) -- (0.334,2.133) -- (0.320,2.141) -- (0.306,2.150) -- (0.292,2.157) -- (0.279,2.164) -- (0.265,2.171) -- (0.252,2.177) -- (0.235,2.184) -- (0.218,2.191) -- (0.201,2.196) -- (0.185,2.202) -- (0.168,2.207) -- (0.152,2.211) -- (0.135,2.215) -- (0.119,2.218) -- (0.103,2.221) -- (0.087,2.223) -- (0.070,2.225) -- (0.054,2.227) -- (0.038,2.228) -- (0.021,2.228) -- (0.005,2.229) -- (-0.012,2.229) -- (-0.028,2.228) -- (-0.045,2.227) -- (-0.062,2.226) -- (-0.079,2.224) -- (-0.097,2.222) -- (-0.115,2.219) -- (-0.133,2.216) -- (-0.151,2.213) -- (-0.170,2.208) -- (-0.184,2.205) -- (-0.199,2.201) -- (-0.213,2.197) -- (-0.228,2.192) -- (-0.244,2.187) -- (-0.259,2.182) -- (-0.275,2.176) -- (-0.292,2.170) -- (-0.308,2.163) -- (-0.325,2.155) -- (-0.343,2.147) -- (-0.360,2.138) -- (-0.379,2.129) -- (-0.397,2.119) -- (-0.410,2.111) -- (-0.423,2.104) -- (-0.436,2.096) -- (-0.450,2.087) -- (-0.463,2.078) -- (-0.477,2.069) -- (-0.492,2.059) -- (-0.506,2.048) -- (-0.521,2.037) -- (-0.536,2.025) -- (-0.552,2.013) -- (-0.567,2.000) -- (-0.583,1.986) -- (-0.600,1.971) -- (-0.617,1.956) -- (-0.634,1.939) -- (-0.651,1.922) -- (-0.669,1.903) -- (-0.687,1.883) -- (-0.706,1.863) -- (-0.715,1.852) -- (-0.724,1.840) -- (-0.734,1.829) -- (-0.744,1.817) -- (-0.753,1.804) -- (-0.763,1.792) -- (-0.773,1.778) -- (-0.783,1.765) -- (-0.793,1.751) -- (-0.804,1.736) -- (-0.814,1.721) -- (-0.824,1.706) -- (-0.835,1.690) -- (-0.845,1.674) -- (-0.856,1.657) -- (-0.866,1.639) -- (-0.877,1.621) -- (-0.888,1.603) -- (-0.899,1.584) -- (-0.909,1.564) -- (-0.920,1.544) -- (-0.931,1.523) -- (-0.942,1.501) -- (-0.953,1.479) -- (-0.964,1.456) -- (-0.975,1.432) -- (-0.985,1.407) -- (-0.996,1.382) -- (-1.007,1.356) -- (-1.018,1.330) -- (-1.028,1.302) -- (-1.039,1.274) -- (-1.049,1.245) -- (-1.059,1.215) -- (-1.069,1.184) -- (-1.079,1.153) -- (-1.089,1.120) -- (-1.098,1.087) -- (-1.107,1.053) -- (-1.116,1.019) -- (-1.125,0.983) -- (-1.133,0.947) -- (-1.141,0.910) -- (-1.149,0.872) -- (-1.156,0.834) -- (-1.163,0.795) -- (-1.170,0.755) -- (-1.176,0.715) -- (-1.182,0.674) -- (-1.187,0.632) -- (-1.191,0.590) -- (-1.196,0.548) -- (-1.199,0.505) -- (-1.202,0.462) -- (-1.205,0.418) -- (-1.207,0.375) -- (-1.208,0.331) -- (-1.209,0.287) -- (-1.210,0.243) -- (-1.210,0.199) -- (-1.209,0.155) -- (-1.208,0.111) -- (-1.206,0.068) -- (-1.203,0.025) -- (-1.201,-0.018) -- (-1.197,-0.061) -- (-1.193,-0.103) -- (-1.189,-0.145) -- (-1.184,-0.187) -- (-1.179,-0.227) -- (-1.174,-0.268) -- (-1.168,-0.307) -- (-1.161,-0.346) -- (-1.155,-0.385) -- (-1.148,-0.422) -- (-1.140,-0.459) -- (-1.133,-0.496) -- (-1.125,-0.531) -- (-1.117,-0.566) -- (-1.109,-0.600) -- (-1.100,-0.634) -- (-1.091,-0.666) -- (-1.083,-0.698) -- (-1.074,-0.729) -- (-1.064,-0.760) -- (-1.055,-0.789) -- (-1.046,-0.818) -- (-1.036,-0.847) -- (-1.027,-0.874) -- (-1.017,-0.901) -- (-1.007,-0.927) -- (-0.998,-0.953) -- (-0.988,-0.977) -- (-0.978,-1.002) -- (-0.968,-1.025) -- (-0.959,-1.048) -- (-0.949,-1.070) -- (-0.939,-1.092) -- (-0.929,-1.113) -- (-0.920,-1.134) -- (-0.910,-1.154) -- (-0.900,-1.173) -- (-0.891,-1.192) -- (-0.881,-1.211) -- (-0.872,-1.229) -- (-0.862,-1.246) -- (-0.853,-1.263) -- (-0.843,-1.280) -- (-0.834,-1.296) -- (-0.825,-1.312) -- (-0.816,-1.327) -- (-0.806,-1.342) -- (-0.797,-1.356) -- (-0.788,-1.371) -- (-0.779,-1.384) -- (-0.770,-1.398) -- (-0.762,-1.411) -- (-0.753,-1.424) -- (-0.744,-1.436) -- (-0.736,-1.448) -- (-0.727,-1.460) -- (-0.710,-1.483) -- (-0.693,-1.504) -- (-0.677,-1.525) -- (-0.661,-1.544) -- (-0.645,-1.563) -- (-0.629,-1.581) -- (-0.614,-1.597) -- (-0.599,-1.614) -- (-0.584,-1.629) -- (-0.569,-1.644) -- (-0.555,-1.658) -- (-0.540,-1.671) -- (-0.526,-1.684) -- (-0.512,-1.696) -- (-0.498,-1.708) -- (-0.485,-1.719) -- (-0.472,-1.729) -- (-0.458,-1.740) -- (-0.445,-1.749) -- (-0.432,-1.759) -- (-0.419,-1.768) -- (-0.407,-1.776) -- (-0.394,-1.784) -- (-0.382,-1.792) -- (-0.363,-1.803) -- (-0.345,-1.813) -- (-0.327,-1.823) -- (-0.310,-1.832) -- (-0.292,-1.840) -- (-0.275,-1.848) -- (-0.257,-1.855) -- (-0.240,-1.862) -- (-0.223,-1.868) -- (-0.206,-1.874) -- (-0.189,-1.879) -- (-0.173,-1.884) -- (-0.156,-1.888) -- (-0.139,-1.892) -- (-0.122,-1.895) -- (-0.106,-1.898) -- (-0.089,-1.901) -- (-0.072,-1.903) -- (-0.055,-1.904) -- (-0.038,-1.906) -- (-0.021,-1.906) -- (-0.004,-1.907) -- (0.014,-1.907) -- (0.032,-1.906) -- (0.049,-1.905) -- (0.067,-1.903) -- (0.086,-1.901) -- (0.104,-1.899) -- (0.123,-1.895) -- (0.142,-1.892) -- (0.162,-1.887) -- (0.182,-1.882) -- (0.202,-1.876) -- (0.216,-1.872) -- (0.230,-1.868) -- (0.245,-1.863) -- (0.259,-1.857) -- (0.274,-1.851) -- (0.290,-1.845) -- (0.305,-1.838) -- (0.321,-1.831) -- (0.337,-1.823) -- (0.354,-1.815) -- (0.371,-1.806) -- (0.389,-1.797) -- (0.407,-1.786) -- (0.425,-1.775) -- (0.444,-1.764) -- (0.463,-1.751) -- (0.484,-1.737) -- (0.504,-1.722) -- (0.526,-1.707) -- (0.548,-1.690) -- (0.559,-1.681) -- (0.571,-1.671) -- (0.582,-1.662) -- (0.594,-1.652) -- (0.607,-1.641) -- (0.619,-1.630) -- (0.631,-1.619) -- (0.644,-1.607) -- (0.657,-1.595) -- (0.671,-1.582) -- (0.684,-1.569) -- (0.698,-1.555) -- (0.712,-1.540) -- (0.727,-1.525) -- (0.741,-1.509) -- (0.756,-1.493) -- (0.772,-1.476) -- (0.787,-1.458) -- (0.803,-1.439) -- (0.819,-1.420) -- (0.836,-1.400) -- (0.853,-1.378) -- (0.870,-1.356) -- (0.887,-1.333) -- (0.905,-1.308) -- (0.923,-1.283) -- (0.942,-1.256) -- (0.960,-1.228) -- (0.980,-1.199) -- (0.999,-1.168) -- (1.019,-1.136) -- (1.039,-1.102) -- (1.059,-1.067) -- (1.080,-1.030) -- (1.100,-0.991) -- (1.121,-0.951) -- (1.142,-0.908) -- (1.163,-0.864) -- (1.185,-0.818) -- (1.206,-0.770) -- (1.227,-0.720) -- (1.247,-0.668) -- (1.268,-0.614) -- (1.288,-0.557) -- (1.308,-0.499) -- (1.327,-0.439) -- (1.345,-0.377) -- (1.363,-0.313) -- (1.380,-0.248) -- (1.396,-0.181) -- (1.410,-0.113) -- (1.424,-0.044) -- (1.436,0.027) -- (1.447,0.097) -- (1.456,0.169) -- (1.464,0.240) -- (1.470,0.311) -- (1.474,0.382) -- (1.477,0.452) -- (1.479,0.521) -- (1.479,0.589) -- (1.478,0.656) -- (1.475,0.721) -- (1.471,0.784) -- (1.465,0.846) -- (1.459,0.906) -- (1.451,0.963) -- (1.443,1.019) -- (1.433,1.072) -- (1.423,1.124) -- (1.413,1.173) -- (1.401,1.220) -- (1.390,1.265) -- (1.378,1.308) -- (1.365,1.349) -- (1.352,1.389) -- (1.339,1.426) -- (1.326,1.461) -- (1.313,1.495) -- (1.300,1.528) -- (1.287,1.558) -- (1.274,1.588) -- (1.261,1.615) -- (1.248,1.642) -- (1.235,1.667) -- (1.222,1.691) -- (1.209,1.714) -- (1.197,1.736) -- (1.184,1.757) -- (1.172,1.777) -- (1.160,1.796) -- (1.148,1.814) -- (1.136,1.831) -- (1.125,1.848) -- (1.114,1.863) -- (1.103,1.878) -- (1.092,1.893) -- (1.081,1.907) -- (1.070,1.920) -- (1.060,1.933) -- (1.050,1.945) -- (1.040,1.957) -- (1.030,1.968) -- (1.021,1.979) -- (1.002,1.999) -- (0.984,2.018) -- (0.966,2.035) -- (0.950,2.051) -- (0.934,2.067) -- (0.918,2.081) -- (0.903,2.094) -- (0.888,2.106) -- (0.874,2.118) -- (0.861,2.128) -- (0.848,2.138) -- (0.835,2.148) -- (0.823,2.157) -- (0.811,2.165) -- (0.793,2.177) -- (0.777,2.188) -- (0.761,2.198) -- (0.746,2.208) -- (0.732,2.216) -- (0.718,2.224) -- (0.704,2.232) -- (0.691,2.239) -- (0.675,2.248) -- (0.659,2.256) -- (0.644,2.263) -- (0.630,2.270) -- (0.617,2.276) -- (0.600,2.283) -- (0.585,2.290) -- (0.570,2.296) -- (0.556,2.301) -- (0.540,2.307) -- (0.525,2.313) -- (0.510,2.318) -- (0.496,2.322) -- (0.481,2.327) -- (0.466,2.332) -- (0.452,2.336) -- (0.437,2.340) -- (0.423,2.344) -- (0.407,2.348) -- (0.393,2.352) -- (0.377,2.355) -- (0.362,2.359) -- (0.347,2.362) -- (0.332,2.365) -- (0.317,2.368) -- (0.302,2.370) -- (0.287,2.373) -- (0.273,2.375) -- (0.258,2.377) -- (0.243,2.379) -- (0.229,2.381) -- (0.214,2.383) -- (0.199,2.385) -- (0.184,2.386) -- (0.169,2.388) -- (0.154,2.389) -- (0.140,2.390) -- (0.125,2.391) -- (0.110,2.392) -- (0.095,2.393) -- (0.080,2.394) -- (0.065,2.394) -- (0.050,2.395) -- (0.035,2.395) -- (0.021,2.395) -- (0.006,2.395) -- (-0.009,2.395) -- (-0.024,2.395) -- (-0.039,2.395) -- (-0.053,2.395) -- (-0.068,2.394) -- (-0.083,2.394) -- (-0.098,2.393) -- (-0.113,2.392) -- (-0.127,2.391) -- (-0.142,2.390) -- (-0.157,2.389) -- (-0.171,2.388) -- (-0.186,2.387) -- (-0.201,2.385) -- (-0.216,2.383) -- (-0.230,2.382) -- (-0.245,2.380) -- (-0.260,2.378) -- (-0.275,2.376) -- (-0.290,2.374) -- (-0.305,2.371) -- (-0.320,2.369) -- (-0.334,2.367) -- (-0.349,2.364) -- (-0.364,2.361) -- (-0.378,2.358) -- (-0.393,2.355) -- (-0.409,2.352) -- (-0.423,2.348) -- (-0.438,2.345) -- (-0.453,2.341) -- (-0.468,2.337) -- (-0.482,2.333) -- (-0.498,2.329) -- (-0.514,2.325) -- (-0.528,2.320) -- (-0.543,2.316) -- (-0.558,2.311) -- (-0.572,2.306) -- (-0.586,2.302) -- (-0.601,2.296) -- (-0.616,2.291) -- (-0.631,2.285) -- (-0.645,2.280) -- (-0.659,2.274) -- (-0.673,2.268) -- (-0.688,2.262) -- (-0.704,2.255) -- (-0.720,2.248) -- (-0.733,2.242) -- (-0.746,2.236) -- (-0.760,2.229) -- (-0.775,2.222) -- (-0.790,2.214) -- (-0.805,2.206) -- (-0.821,2.197) -- (-0.837,2.188) -- (-0.850,2.181) -- (-0.863,2.173) -- (-0.876,2.165) -- (-0.890,2.156) -- (-0.904,2.148) -- (-0.919,2.138) -- (-0.934,2.128) -- (-0.949,2.117) -- (-0.965,2.106) -- (-0.981,2.094) -- (-0.998,2.082) -- (-1.009,2.073) -- (-1.021,2.064) -- (-1.033,2.054) -- (-1.045,2.044) -- (-1.058,2.034) -- (-1.071,2.023) -- (-1.084,2.012) -- (-1.097,2.000) -- (-1.110,1.987) -- (-1.124,1.974) -- (-1.139,1.961) -- (-1.153,1.946) -- (-1.168,1.931) -- (-1.183,1.916) -- (-1.198,1.899) -- (-1.214,1.882) -- (-1.230,1.863) -- (-1.247,1.844) -- (-1.263,1.824) -- (-1.280,1.802) -- (-1.298,1.780) -- (-1.307,1.768) -- (-1.316,1.756) -- (-1.325,1.743) -- (-1.334,1.731) -- (-1.343,1.718) -- (-1.352,1.704) -- (-1.361,1.690) -- (-1.370,1.676) -- (-1.380,1.661) -- (-1.389,1.646) -- (-1.399,1.631) -- (-1.408,1.615) -- (-1.418,1.599) -- (-1.428,1.582) -- (-1.437,1.565) -- (-1.447,1.547) -- (-1.457,1.528) -- (-1.466,1.509) -- (-1.476,1.490) -- (-1.486,1.470) -- (-1.496,1.450) -- (-1.505,1.429) -- (-1.515,1.407) -- (-1.525,1.384) -- (-1.534,1.362) -- (-1.544,1.338) -- (-1.553,1.314) -- (-1.563,1.289) -- (-1.572,1.263) -- (-1.581,1.237) -- (-1.590,1.210) -- (-1.599,1.182) -- (-1.608,1.154) -- (-1.616,1.125) -- (-1.625,1.095) -- (-1.633,1.064) -- (-1.640,1.033) -- (-1.648,1.001) -- (-1.655,0.968) -- (-1.662,0.934) -- (-1.668,0.900) -- (-1.675,0.865) -- (-1.680,0.830) -- (-1.686,0.793) -- (-1.690,0.756) -- (-1.695,0.719) -- (-1.699,0.681) -- (-1.702,0.642) -- (-1.705,0.602) -- (-1.708,0.563) -- (-1.709,0.522) -- (-1.711,0.481) -- (-1.711,0.440) -- (-1.711,0.399) -- (-1.711,0.357) -- (-1.710,0.315) -- (-1.708,0.272) -- (-1.706,0.230) -- (-1.703,0.187) -- (-1.699,0.145) -- (-1.695,0.102) -- (-1.690,0.060) -- (-1.685,0.017) -- (-1.679,-0.025) -- (-1.672,-0.067) -- (-1.665,-0.109) -- (-1.657,-0.151) -- (-1.649,-0.192) -- (-1.640,-0.232) -- (-1.631,-0.273) -- (-1.622,-0.313) -- (-1.612,-0.352) -- (-1.601,-0.391) -- (-1.590,-0.429) -- (-1.579,-0.467) -- (-1.568,-0.504) -- (-1.556,-0.540) -- (-1.544,-0.576) -- (-1.532,-0.611) -- (-1.519,-0.645) -- (-1.506,-0.679) -- (-1.494,-0.712) -- (-1.480,-0.744) -- (-1.467,-0.776) -- (-1.454,-0.806) -- (-1.441,-0.837) -- (-1.427,-0.866) -- (-1.413,-0.895) -- (-1.400,-0.923) -- (-1.386,-0.950) -- (-1.373,-0.977) -- (-1.359,-1.003) -- (-1.345,-1.028) -- (-1.332,-1.053) -- (-1.318,-1.077) -- (-1.304,-1.101) -- (-1.291,-1.124) -- (-1.278,-1.146) -- (-1.264,-1.168) -- (-1.251,-1.189) -- (-1.238,-1.209) -- (-1.224,-1.229) -- (-1.211,-1.249) -- (-1.198,-1.268) -- (-1.186,-1.287) -- (-1.173,-1.305) -- (-1.160,-1.322) -- (-1.148,-1.339) -- (-1.135,-1.356) -- (-1.123,-1.372) -- (-1.111,-1.388) -- (-1.099,-1.403) -- (-1.087,-1.418) -- (-1.075,-1.433) -- (-1.063,-1.447) -- (-1.052,-1.461) -- (-1.040,-1.475) -- (-1.029,-1.488) -- (-1.017,-1.501) -- (-1.006,-1.513) -- (-0.995,-1.525) -- (-0.984,-1.537) -- (-0.974,-1.549) -- (-0.963,-1.560) -- (-0.952,-1.571) -- (-0.942,-1.582) -- (-0.932,-1.593) -- (-0.921,-1.603) -- (-0.901,-1.623) -- (-0.881,-1.641) -- (-0.862,-1.659) -- (-0.843,-1.677) -- (-0.825,-1.693) -- (-0.807,-1.709) -- (-0.789,-1.724) -- (-0.771,-1.738) -- (-0.754,-1.752) -- (-0.737,-1.765) -- (-0.721,-1.777) -- (-0.705,-1.789) -- (-0.689,-1.801) -- (-0.673,-1.812) -- (-0.658,-1.822) -- (-0.643,-1.832) -- (-0.628,-1.842) -- (-0.613,-1.852) -- (-0.599,-1.861) -- (-0.585,-1.869) -- (-0.571,-1.878) -- (-0.557,-1.885) -- (-0.544,-1.893) -- (-0.531,-1.901) -- (-0.517,-1.908) -- (-0.504,-1.914) -- (-0.485,-1.924) -- (-0.467,-1.933) -- (-0.448,-1.942) -- (-0.430,-1.950) -- (-0.412,-1.958) -- (-0.394,-1.965) -- (-0.377,-1.972) -- (-0.360,-1.979) -- (-0.343,-1.985) -- (-0.326,-1.991) -- (-0.310,-1.996) -- (-0.293,-2.001) -- (-0.277,-2.006) -- (-0.261,-2.010) -- (-0.245,-2.015) -- (-0.229,-2.019) -- (-0.213,-2.022) -- (-0.198,-2.026) -- (-0.182,-2.029) -- (-0.166,-2.031) -- (-0.151,-2.034) -- (-0.135,-2.036) -- (-0.120,-2.038) -- (-0.104,-2.040) -- (-0.089,-2.042) -- (-0.073,-2.043) -- (-0.057,-2.044) -- (-0.042,-2.045) -- (-0.026,-2.045) -- (-0.010,-2.046) -- (0.006,-2.046) -- (0.022,-2.046) -- (0.039,-2.045) -- (0.055,-2.044) -- (0.072,-2.043) -- (0.089,-2.042) -- (0.106,-2.040) -- (0.123,-2.038) -- (0.141,-2.036) -- (0.159,-2.033) -- (0.177,-2.030) -- (0.196,-2.026) -- (0.215,-2.022) -- (0.234,-2.017) -- (0.254,-2.012) -- (0.275,-2.007) -- (0.289,-2.003) -- (0.303,-1.999) -- (0.318,-1.994) -- (0.332,-1.989) -- (0.348,-1.984) -- (0.363,-1.978) -- (0.379,-1.972) -- (0.395,-1.966) -- (0.412,-1.959) -- (0.429,-1.952) -- (0.446,-1.944) -- (0.464,-1.936) -- (0.483,-1.927) -- (0.502,-1.917) -- (0.521,-1.907) -- (0.541,-1.896) -- (0.562,-1.884) -- (0.584,-1.872) -- (0.606,-1.858) -- (0.629,-1.844) -- (0.653,-1.828) -- (0.666,-1.820) -- (0.678,-1.811) -- (0.691,-1.802) -- (0.704,-1.793) -- (0.717,-1.783) -- (0.731,-1.773) -- (0.745,-1.763) -- (0.759,-1.752) -- (0.773,-1.741) -- (0.788,-1.729) -- (0.803,-1.716) -- (0.819,-1.704) -- (0.834,-1.690) -- (0.850,-1.676) -- (0.867,-1.661) -- (0.884,-1.646) -- (0.901,-1.630) -- (0.919,-1.613) -- (0.937,-1.596) -- (0.955,-1.577) -- (0.974,-1.558) -- (0.993,-1.538) -- (1.013,-1.516) -- (1.033,-1.494) -- (1.054,-1.471) -- (1.075,-1.446) -- (1.096,-1.420) -- (1.118,-1.393) -- (1.141,-1.364) -- (1.164,-1.334) -- (1.187,-1.303) -- (1.211,-1.269) -- (1.235,-1.234) -- (1.260,-1.197) -- (1.285,-1.159) -- (1.310,-1.118) -- (1.336,-1.075) -- (1.362,-1.030) -- (1.388,-0.982) -- (1.414,-0.933) -- (1.440,-0.881) -- (1.466,-0.826) -- (1.491,-0.769) -- (1.517,-0.710) -- (1.541,-0.648) -- (1.565,-0.583) -- (1.589,-0.517) -- (1.611,-0.448) -- (1.632,-0.377) -- (1.652,-0.304) -- (1.670,-0.229) -- (1.686,-0.153) -- (1.701,-0.076) -- (1.714,0.002) -- (1.725,0.081) -- (1.734,0.160) -- (1.740,0.238) -- (1.745,0.316) -- (1.747,0.393) -- (1.748,0.469) -- (1.746,0.543) -- (1.742,0.616) -- (1.737,0.687) -- (1.730,0.755) -- (1.721,0.822) -- (1.712,0.885) -- (1.700,0.947) -- (1.688,1.006) -- (1.675,1.062) -- (1.661,1.116) -- (1.647,1.167) -- (1.632,1.216) -- (1.616,1.262) -- (1.600,1.306) -- (1.584,1.348) -- (1.568,1.388) -- (1.551,1.426) -- (1.535,1.462) -- (1.518,1.496) -- (1.502,1.528) -- (1.486,1.559) -- (1.470,1.588) -- (1.454,1.615) -- (1.438,1.642) -- (1.423,1.666) -- (1.408,1.690) -- (1.392,1.713) -- (1.378,1.734) -- (1.363,1.754) -- (1.349,1.774) -- (1.335,1.792) -- (1.321,1.810) -- (1.308,1.827) -- (1.295,1.843) -- (1.282,1.858) -- (1.269,1.873) -- (1.257,1.887) -- (1.245,1.900) -- (1.233,1.913) -- (1.221,1.925) -- (1.210,1.937) -- (1.199,1.948) -- (1.188,1.959) -- (1.177,1.969) -- (1.167,1.979) -- (1.146,1.998) -- (1.127,2.016) -- (1.108,2.032) -- (1.090,2.047) -- (1.073,2.061) -- (1.056,2.074) -- (1.040,2.086) -- (1.025,2.098) -- (1.010,2.108) -- (0.996,2.118) -- (0.982,2.128) -- (0.969,2.137) -- (0.956,2.145) -- (0.943,2.153) -- (0.926,2.164) -- (0.909,2.175) -- (0.892,2.184) -- (0.877,2.193) -- (0.862,2.201) -- (0.848,2.209) -- (0.834,2.216) -- (0.821,2.223) -- (0.804,2.231) -- (0.789,2.239) -- (0.774,2.246) -- (0.759,2.252) -- (0.746,2.258) -- (0.729,2.265) -- (0.714,2.272) -- (0.700,2.278) -- (0.686,2.283) -- (0.670,2.289) -- (0.655,2.294) -- (0.641,2.299) -- (0.626,2.305) -- (0.611,2.310) -- (0.598,2.314) -- (0.583,2.319) -- (0.569,2.323) -- (0.554,2.327) -- (0.539,2.332) -- (0.525,2.335) -- (0.510,2.339) -- (0.495,2.343) -- (0.480,2.347) -- (0.466,2.350) -- (0.452,2.353) -- (0.437,2.357) -- (0.422,2.360) -- (0.408,2.362) -- (0.393,2.365) -- (0.379,2.368) -- (0.364,2.370) -- (0.350,2.373) -- (0.335,2.375) -- (0.321,2.377) -- (0.307,2.379) -- (0.292,2.381) -- (0.278,2.383) -- (0.263,2.385) -- (0.249,2.387) -- (0.234,2.388) -- (0.220,2.390) -- (0.205,2.391) -- (0.191,2.392) -- (0.176,2.393) -- (0.162,2.394) -- (0.148,2.395) -- (0.133,2.396) -- (0.119,2.397) -- (0.104,2.398) -- (0.090,2.398) -- (0.075,2.399) -- (0.061,2.399) -- (0.047,2.400) -- (0.032,2.400) -- (0.018,2.400) -- (0.003,2.400) -- (-0.011,2.400) -- (-0.025,2.400) -- (-0.040,2.400) -- (-0.054,2.399) -- (-0.069,2.399) -- (-0.083,2.399) -- (-0.098,2.398) -- (-0.112,2.397) -- (-0.126,2.397) -- (-0.141,2.396) -- (-0.155,2.395) -- (-0.170,2.394) -- (-0.184,2.393) -- (-0.199,2.392) -- (-0.213,2.390) -- (-0.227,2.389) -- (-0.242,2.387) -- (-0.256,2.386) -- (-0.271,2.384) -- (-0.285,2.382) -- (-0.299,2.380) -- (-0.314,2.379) -- (-0.328,2.376) -- (-0.343,2.374) -- (-0.357,2.372) -- (-0.371,2.369) -- (-0.386,2.367) -- (-0.400,2.364) -- (-0.415,2.361) -- (-0.429,2.359) -- (-0.444,2.355) -- (-0.458,2.352) -- (-0.473,2.349) -- (-0.488,2.346) -- (-0.502,2.342) -- (-0.516,2.339) -- (-0.530,2.335) -- (-0.544,2.331) -- (-0.559,2.327) -- (-0.573,2.323) -- (-0.587,2.318) -- (-0.602,2.314) -- (-0.616,2.309) -- (-0.632,2.304) -- (-0.646,2.299) -- (-0.660,2.294) -- (-0.674,2.289) -- (-0.689,2.283) -- (-0.703,2.278) -- (-0.718,2.272) -- (-0.731,2.267) -- (-0.745,2.261) -- (-0.760,2.254) -- (-0.775,2.248) -- (-0.789,2.241) -- (-0.803,2.234) -- (-0.818,2.227) -- (-0.831,2.221) -- (-0.845,2.214) -- (-0.859,2.207) -- (-0.873,2.199) -- (-0.888,2.190) -- (-0.901,2.183) -- (-0.914,2.176) -- (-0.927,2.168) -- (-0.941,2.159) -- (-0.956,2.150) -- (-0.970,2.141) -- (-0.986,2.131) -- (-0.998,2.123) -- (-1.010,2.114) -- (-1.023,2.105) -- (-1.036,2.096) -- (-1.049,2.086) -- (-1.063,2.076) -- (-1.078,2.065) -- (-1.092,2.053) -- (-1.108,2.041) -- (-1.123,2.027) -- (-1.139,2.014) -- (-1.150,2.004) -- (-1.162,1.994) -- (-1.173,1.983) -- (-1.185,1.972) -- (-1.197,1.961) -- (-1.209,1.949) -- (-1.222,1.936) -- (-1.235,1.923) -- (-1.248,1.909) -- (-1.262,1.895) -- (-1.275,1.880) -- (-1.289,1.864) -- (-1.304,1.847) -- (-1.318,1.830) -- (-1.333,1.812) -- (-1.348,1.792) -- (-1.364,1.772) -- (-1.380,1.751) -- (-1.396,1.729) -- (-1.412,1.705) -- (-1.420,1.693) -- (-1.429,1.680) -- (-1.437,1.667) -- (-1.446,1.654) -- (-1.454,1.641) -- (-1.463,1.627) -- (-1.471,1.612) -- (-1.480,1.597) -- (-1.489,1.582) -- (-1.498,1.567) -- (-1.507,1.551) -- (-1.515,1.534) -- (-1.524,1.517) -- (-1.533,1.500) -- (-1.542,1.482) -- (-1.551,1.464) -- (-1.560,1.445) -- (-1.569,1.426) -- (-1.578,1.406) -- (-1.587,1.385) -- (-1.596,1.364) -- (-1.605,1.343) -- (-1.613,1.320) -- (-1.622,1.298) -- (-1.631,1.274) -- (-1.639,1.250) -- (-1.648,1.226) -- (-1.656,1.200) -- (-1.664,1.174) -- (-1.672,1.148) -- (-1.680,1.121) -- (-1.688,1.093) -- (-1.695,1.064) -- (-1.703,1.035) -- (-1.710,1.005) -- (-1.716,0.974) -- (-1.723,0.943) -- (-1.729,0.911) -- (-1.735,0.878) -- (-1.740,0.844) -- (-1.745,0.810) -- (-1.750,0.776) -- (-1.754,0.740) -- (-1.758,0.705) -- (-1.762,0.668) -- (-1.765,0.631) -- (-1.767,0.593) -- (-1.769,0.555) -- (-1.771,0.517) -- (-1.772,0.478) -- (-1.772,0.438) -- (-1.772,0.398) -- (-1.771,0.358) -- (-1.770,0.318) -- (-1.768,0.277) -- (-1.765,0.236) -- (-1.762,0.195) -- (-1.759,0.154) -- (-1.754,0.113) -- (-1.749,0.072) -- (-1.744,0.030) -- (-1.738,-0.011) -- (-1.731,-0.051) -- (-1.724,-0.092) -- (-1.717,-0.133) -- (-1.709,-0.173) -- (-1.700,-0.213) -- (-1.691,-0.252) -- (-1.681,-0.291) -- (-1.671,-0.330) -- (-1.661,-0.368) -- (-1.650,-0.406) -- (-1.639,-0.443) -- (-1.628,-0.480) -- (-1.616,-0.516) -- (-1.604,-0.551) -- (-1.591,-0.586) -- (-1.579,-0.620) -- (-1.566,-0.654) -- (-1.553,-0.687) -- (-1.540,-0.719) -- (-1.526,-0.751) -- (-1.513,-0.782) -- (-1.499,-0.812) -- (-1.485,-0.842) -- (-1.472,-0.871) -- (-1.458,-0.899) -- (-1.444,-0.927) -- (-1.430,-0.954) -- (-1.416,-0.980) -- (-1.402,-1.006) -- (-1.388,-1.031) -- (-1.374,-1.056) -- (-1.360,-1.080) -- (-1.346,-1.103) -- (-1.333,-1.126) -- (-1.319,-1.148) -- (-1.305,-1.169) -- (-1.292,-1.190) -- (-1.278,-1.211) -- (-1.265,-1.231) -- (-1.251,-1.250) -- (-1.238,-1.269) -- (-1.225,-1.288) -- (-1.212,-1.306) -- (-1.199,-1.323) -- (-1.186,-1.340) -- (-1.173,-1.357) -- (-1.161,-1.373) -- (-1.148,-1.389) -- (-1.136,-1.405) -- (-1.124,-1.420) -- (-1.112,-1.434) -- (-1.100,-1.448) -- (-1.088,-1.462) -- (-1.076,-1.476) -- (-1.064,-1.489) -- (-1.053,-1.502) -- (-1.041,-1.515) -- (-1.030,-1.527) -- (-1.019,-1.539) -- (-1.008,-1.550) -- (-0.997,-1.562) -- (-0.986,-1.573) -- (-0.975,-1.584) -- (-0.965,-1.594) -- (-0.954,-1.605) -- (-0.944,-1.615) -- (-0.924,-1.634) -- (-0.904,-1.653) -- (-0.884,-1.670) -- (-0.865,-1.687) -- (-0.846,-1.703) -- (-0.828,-1.719) -- (-0.810,-1.734) -- (-0.792,-1.748) -- (-0.775,-1.761) -- (-0.758,-1.774) -- (-0.741,-1.787) -- (-0.725,-1.798) -- (-0.709,-1.810) -- (-0.693,-1.821) -- (-0.677,-1.831) -- (-0.662,-1.841) -- (-0.647,-1.851) -- (-0.632,-1.860) -- (-0.618,-1.869) -- (-0.604,-1.878) -- (-0.590,-1.886) -- (-0.576,-1.894) -- (-0.562,-1.901) -- (-0.549,-1.909) -- (-0.535,-1.916) -- (-0.522,-1.923) -- (-0.509,-1.929) -- (-0.490,-1.939) -- (-0.472,-1.947) -- (-0.453,-1.956) -- (-0.435,-1.964) -- (-0.418,-1.971) -- (-0.400,-1.978) -- (-0.383,-1.985) -- (-0.366,-1.991) -- (-0.349,-1.997) -- (-0.332,-2.003) -- (-0.316,-2.008) -- (-0.300,-2.013) -- (-0.283,-2.018) -- (-0.267,-2.022) -- (-0.252,-2.026) -- (-0.236,-2.030) -- (-0.220,-2.033) -- (-0.205,-2.037) -- (-0.189,-2.040) -- (-0.174,-2.043) -- (-0.158,-2.045) -- (-0.143,-2.047) -- (-0.127,-2.050) -- (-0.112,-2.051) -- (-0.096,-2.053) -- (-0.081,-2.054) -- (-0.066,-2.056) -- (-0.050,-2.056) -- (-0.034,-2.057) -- (-0.019,-2.057) -- (-0.003,-2.058) -- (0.013,-2.058) -- (0.029,-2.057) -- (0.045,-2.057) -- (0.062,-2.056) -- (0.078,-2.055) -- (0.095,-2.053) -- (0.112,-2.051) -- (0.129,-2.049) -- (0.147,-2.047) -- (0.165,-2.044) -- (0.183,-2.041) -- (0.202,-2.037) -- (0.221,-2.033) -- (0.240,-2.029) -- (0.260,-2.024) -- (0.281,-2.018) -- (0.294,-2.014) -- (0.309,-2.010) -- (0.323,-2.006) -- (0.338,-2.001) -- (0.353,-1.996) -- (0.368,-1.990) -- (0.384,-1.984) -- (0.400,-1.978) -- (0.417,-1.971) -- (0.434,-1.964) -- (0.451,-1.957) -- (0.469,-1.949) -- (0.488,-1.940) -- (0.507,-1.931) -- (0.526,-1.921) -- (0.546,-1.910) -- (0.567,-1.899) -- (0.589,-1.886) -- (0.611,-1.873) -- (0.634,-1.859) -- (0.658,-1.844) -- (0.670,-1.836) -- (0.683,-1.828) -- (0.695,-1.819) -- (0.708,-1.810) -- (0.722,-1.801) -- (0.735,-1.791) -- (0.749,-1.781) -- (0.763,-1.770) -- (0.778,-1.759) -- (0.792,-1.748) -- (0.807,-1.736) -- (0.823,-1.723) -- (0.838,-1.710) -- (0.854,-1.696) -- (0.871,-1.682) -- (0.888,-1.667) -- (0.905,-1.652) -- (0.922,-1.635) -- (0.940,-1.618) -- (0.959,-1.600) -- (0.978,-1.582) -- (0.997,-1.562) -- (1.017,-1.542) -- (1.037,-1.520) -- (1.057,-1.497) -- (1.078,-1.473) -- (1.100,-1.448) -- (1.122,-1.422) -- (1.145,-1.394) -- (1.168,-1.365) -- (1.191,-1.334) -- (1.215,-1.302) -- (1.239,-1.268) -- (1.264,-1.232) -- (1.289,-1.195) -- (1.315,-1.155) -- (1.341,-1.113) -- (1.367,-1.070) -- (1.393,-1.023) -- (1.419,-0.975) -- (1.446,-0.924) -- (1.472,-0.871) -- (1.498,-0.815) -- (1.524,-0.757) -- (1.550,-0.696) -- (1.575,-0.633) -- (1.599,-0.567) -- (1.622,-0.500) -- (1.644,-0.429) -- (1.665,-0.357) -- (1.684,-0.283) -- (1.702,-0.207) -- (1.718,-0.130) -- (1.732,-0.052) -- (1.744,0.027) -- (1.754,0.106) -- (1.762,0.185) -- (1.768,0.263) -- (1.772,0.342) -- (1.773,0.419) -- (1.772,0.494) -- (1.770,0.568) -- (1.765,0.641) -- (1.759,0.711) -- (1.751,0.779) -- (1.741,0.844) -- (1.731,0.907) -- (1.719,0.968) -- (1.706,1.025) -- (1.692,1.081) -- (1.678,1.134) -- (1.662,1.184) -- (1.647,1.232) -- (1.631,1.278) -- (1.614,1.321) -- (1.598,1.362) -- (1.581,1.401) -- (1.565,1.438) -- (1.548,1.473) -- (1.531,1.506) -- (1.515,1.538) -- (1.498,1.568) -- (1.482,1.596) -- (1.466,1.623) -- (1.450,1.649) -- (1.435,1.673) -- (1.419,1.696) -- (1.404,1.719) -- (1.389,1.739) -- (1.375,1.759) -- (1.360,1.778) -- (1.346,1.796) -- (1.333,1.814) -- (1.319,1.830) -- (1.306,1.846) -- (1.293,1.861) -- (1.280,1.875) -- (1.268,1.889) -- (1.256,1.902) -- (1.244,1.915) -- (1.232,1.927) -- (1.221,1.938) -- (1.210,1.949) -- (1.199,1.960) -- (1.188,1.970) -- (1.168,1.990) -- (1.148,2.007) -- (1.129,2.024) -- (1.110,2.039) -- (1.093,2.054) -- (1.076,2.067) -- (1.059,2.079) -- (1.044,2.091) -- (1.029,2.102) -- (1.014,2.112) -- (1.000,2.122) -- (0.987,2.131) -- (0.973,2.140) -- (0.961,2.148) -- (0.949,2.156) -- (0.931,2.166) -- (0.914,2.176) -- (0.898,2.186) -- (0.883,2.194) -- (0.868,2.202) -- (0.854,2.210) -- (0.841,2.217) -- (0.828,2.223) -- (0.812,2.231) -- (0.796,2.239) -- (0.781,2.245) -- (0.767,2.252) -- (0.754,2.258) -- (0.738,2.265) -- (0.722,2.271) -- (0.708,2.277) -- (0.695,2.282) -- (0.679,2.288) -- (0.664,2.293) -- (0.651,2.298) -- (0.635,2.303) -- (0.621,2.308) -- (0.606,2.313) -- (0.591,2.318) -- (0.576,2.322) -- (0.561,2.326) -- (0.547,2.331) -- (0.533,2.334) -- (0.518,2.338) -- (0.504,2.342) -- (0.489,2.346) -- (0.474,2.349) -- (0.460,2.352) -- (0.445,2.355) -- (0.431,2.358) -- (0.416,2.361) -- (0.402,2.364) -- (0.388,2.367) -- (0.373,2.369) -- (0.359,2.372) -- (0.345,2.374) -- (0.330,2.376) -- (0.316,2.378) -- (0.302,2.380) -- (0.287,2.382) -- (0.272,2.384) -- (0.258,2.386) -- (0.244,2.387) -- (0.229,2.389) -- (0.215,2.390) -- (0.200,2.391) -- (0.186,2.393) -- (0.172,2.394) -- (0.157,2.395) -- (0.143,2.396) -- (0.128,2.397) -- (0.114,2.397) -- (0.099,2.398) -- (0.085,2.398) -- (0.071,2.399) -- (0.056,2.399) -- (0.042,2.400) -- (0.027,2.400) -- (0.013,2.400) -- (-0.001,2.400) -- (-0.016,2.400) -- (-0.019,2.400);}
\newcommand{\figOneGraphic}[1]{\draw[#1] (-0.000,2.400) -- (-0.010,2.400) -- (-0.019,2.400) -- (-0.029,2.400) -- (-0.039,2.400) -- (-0.048,2.400) -- (-0.058,2.399) -- (-0.068,2.399) -- (-0.078,2.399) -- (-0.087,2.398) -- (-0.097,2.398) -- (-0.107,2.398) -- (-0.117,2.397) -- (-0.126,2.397) -- (-0.136,2.396) -- (-0.146,2.396) -- (-0.156,2.395) -- (-0.166,2.394) -- (-0.175,2.394) -- (-0.185,2.393) -- (-0.195,2.392) -- (-0.205,2.391) -- (-0.215,2.390) -- (-0.224,2.389) -- (-0.234,2.388) -- (-0.244,2.387) -- (-0.254,2.386) -- (-0.264,2.385) -- (-0.274,2.384) -- (-0.284,2.383) -- (-0.294,2.381) -- (-0.304,2.380) -- (-0.313,2.379) -- (-0.323,2.377) -- (-0.333,2.376) -- (-0.343,2.374) -- (-0.353,2.373) -- (-0.363,2.371) -- (-0.373,2.369) -- (-0.382,2.368) -- (-0.392,2.366) -- (-0.402,2.364) -- (-0.412,2.362) -- (-0.421,2.360) -- (-0.431,2.358) -- (-0.441,2.356) -- (-0.451,2.354) -- (-0.460,2.352) -- (-0.470,2.350) -- (-0.480,2.348) -- (-0.490,2.345) -- (-0.499,2.343) -- (-0.509,2.341) -- (-0.519,2.338) -- (-0.529,2.336) -- (-0.538,2.333) -- (-0.548,2.330) -- (-0.557,2.328) -- (-0.568,2.325) -- (-0.578,2.322) -- (-0.587,2.319) -- (-0.597,2.316) -- (-0.606,2.313) -- (-0.616,2.310) -- (-0.626,2.306) -- (-0.636,2.303) -- (-0.645,2.300) -- (-0.654,2.297) -- (-0.664,2.293) -- (-0.674,2.290) -- (-0.684,2.286) -- (-0.694,2.282) -- (-0.703,2.279) -- (-0.712,2.275) -- (-0.721,2.271) -- (-0.730,2.268) -- (-0.740,2.264) -- (-0.750,2.259) -- (-0.760,2.255) -- (-0.770,2.251) -- (-0.780,2.246) -- (-0.789,2.242) -- (-0.798,2.238) -- (-0.807,2.233) -- (-0.816,2.229) -- (-0.826,2.224) -- (-0.835,2.220) -- (-0.845,2.215) -- (-0.855,2.210) -- (-0.865,2.204) -- (-0.875,2.199) -- (-0.883,2.194) -- (-0.892,2.189) -- (-0.901,2.184) -- (-0.909,2.179) -- (-0.918,2.174) -- (-0.927,2.169) -- (-0.937,2.163) -- (-0.946,2.157) -- (-0.955,2.152) -- (-0.965,2.146) -- (-0.974,2.139) -- (-0.984,2.133) -- (-0.994,2.126) -- (-1.002,2.121) -- (-1.010,2.115) -- (-1.018,2.110) -- (-1.026,2.104) -- (-1.035,2.098) -- (-1.043,2.092) -- (-1.052,2.086) -- (-1.060,2.079) -- (-1.069,2.073) -- (-1.077,2.066) -- (-1.086,2.059) -- (-1.095,2.052) -- (-1.104,2.045) -- (-1.113,2.037) -- (-1.122,2.030) -- (-1.131,2.022) -- (-1.141,2.014) -- (-1.150,2.006) -- (-1.159,1.997) -- (-1.167,1.991) -- (-1.174,1.984) -- (-1.181,1.977) -- (-1.188,1.971) -- (-1.196,1.964) -- (-1.203,1.957) -- (-1.210,1.949) -- (-1.218,1.942) -- (-1.225,1.934) -- (-1.233,1.927) -- (-1.240,1.919) -- (-1.248,1.911) -- (-1.256,1.903) -- (-1.263,1.895) -- (-1.271,1.886) -- (-1.279,1.878) -- (-1.287,1.869) -- (-1.294,1.860) -- (-1.302,1.851) -- (-1.310,1.841) -- (-1.318,1.832) -- (-1.326,1.822) -- (-1.334,1.812) -- (-1.342,1.802) -- (-1.350,1.792) -- (-1.358,1.782) -- (-1.367,1.771) -- (-1.375,1.760) -- (-1.383,1.749) -- (-1.391,1.738) -- (-1.399,1.726) -- (-1.407,1.714) -- (-1.413,1.706) -- (-1.418,1.698) -- (-1.424,1.690) -- (-1.429,1.682) -- (-1.435,1.674) -- (-1.440,1.665) -- (-1.446,1.657) -- (-1.451,1.648) -- (-1.457,1.639) -- (-1.462,1.630) -- (-1.468,1.621) -- (-1.473,1.612) -- (-1.479,1.603) -- (-1.484,1.593) -- (-1.490,1.584) -- (-1.495,1.574) -- (-1.501,1.565) -- (-1.506,1.555) -- (-1.512,1.545) -- (-1.517,1.535) -- (-1.522,1.524) -- (-1.528,1.514) -- (-1.533,1.504) -- (-1.538,1.493) -- (-1.544,1.482) -- (-1.549,1.471) -- (-1.554,1.461) -- (-1.560,1.449) -- (-1.565,1.438) -- (-1.570,1.427) -- (-1.575,1.415) -- (-1.581,1.404) -- (-1.586,1.392) -- (-1.591,1.380) -- (-1.596,1.368) -- (-1.601,1.356) -- (-1.606,1.343) -- (-1.611,1.331) -- (-1.616,1.318) -- (-1.621,1.306) -- (-1.626,1.293) -- (-1.630,1.280) -- (-1.635,1.267) -- (-1.640,1.253) -- (-1.645,1.240) -- (-1.649,1.226) -- (-1.654,1.213) -- (-1.658,1.199) -- (-1.663,1.185) -- (-1.667,1.171) -- (-1.671,1.157) -- (-1.676,1.142) -- (-1.680,1.128) -- (-1.684,1.113) -- (-1.688,1.098) -- (-1.692,1.083) -- (-1.696,1.068) -- (-1.700,1.053) -- (-1.704,1.037) -- (-1.707,1.022) -- (-1.711,1.006) -- (-1.714,0.991) -- (-1.718,0.975) -- (-1.721,0.959) -- (-1.725,0.943) -- (-1.728,0.926) -- (-1.731,0.910) -- (-1.734,0.893) -- (-1.737,0.877) -- (-1.740,0.860) -- (-1.742,0.843) -- (-1.745,0.826) -- (-1.747,0.809) -- (-1.750,0.792) -- (-1.752,0.774) -- (-1.754,0.757) -- (-1.756,0.739) -- (-1.758,0.722) -- (-1.760,0.704) -- (-1.762,0.686) -- (-1.763,0.668) -- (-1.765,0.650) -- (-1.766,0.632) -- (-1.768,0.613) -- (-1.769,0.595) -- (-1.770,0.577) -- (-1.771,0.558) -- (-1.772,0.539) -- (-1.772,0.521) -- (-1.773,0.502) -- (-1.773,0.483) -- (-1.773,0.464) -- (-1.774,0.446) -- (-1.774,0.427) -- (-1.774,0.408) -- (-1.773,0.388) -- (-1.773,0.369) -- (-1.772,0.350) -- (-1.772,0.341) -- (-1.772,0.331) -- (-1.772,0.321) -- (-1.771,0.312) -- (-1.771,0.302) -- (-1.770,0.292) -- (-1.770,0.283) -- (-1.769,0.273) -- (-1.769,0.263) -- (-1.768,0.254) -- (-1.767,0.244) -- (-1.767,0.234) -- (-1.766,0.225) -- (-1.765,0.215) -- (-1.765,0.205) -- (-1.764,0.196) -- (-1.763,0.186) -- (-1.762,0.176) -- (-1.761,0.167) -- (-1.760,0.157) -- (-1.759,0.147) -- (-1.758,0.138) -- (-1.757,0.128) -- (-1.756,0.118) -- (-1.755,0.108) -- (-1.754,0.099) -- (-1.753,0.089) -- (-1.752,0.079) -- (-1.751,0.070) -- (-1.750,0.060) -- (-1.748,0.050) -- (-1.747,0.041) -- (-1.746,0.031) -- (-1.744,0.021) -- (-1.743,0.012) -- (-1.741,0.002) -- (-1.740,-0.007) -- (-1.739,-0.017) -- (-1.737,-0.027) -- (-1.735,-0.036) -- (-1.734,-0.046) -- (-1.732,-0.055) -- (-1.731,-0.065) -- (-1.729,-0.075) -- (-1.727,-0.084) -- (-1.726,-0.094) -- (-1.724,-0.103) -- (-1.722,-0.113) -- (-1.720,-0.122) -- (-1.718,-0.132) -- (-1.717,-0.141) -- (-1.715,-0.151) -- (-1.713,-0.160) -- (-1.711,-0.169) -- (-1.709,-0.179) -- (-1.707,-0.188) -- (-1.703,-0.207) -- (-1.699,-0.226) -- (-1.694,-0.244) -- (-1.690,-0.263) -- (-1.685,-0.281) -- (-1.681,-0.300) -- (-1.676,-0.318) -- (-1.671,-0.336) -- (-1.666,-0.354) -- (-1.661,-0.372) -- (-1.656,-0.390) -- (-1.651,-0.408) -- (-1.646,-0.426) -- (-1.640,-0.443) -- (-1.635,-0.461) -- (-1.629,-0.478) -- (-1.624,-0.495) -- (-1.618,-0.513) -- (-1.612,-0.530) -- (-1.607,-0.547) -- (-1.601,-0.564) -- (-1.595,-0.580) -- (-1.589,-0.597) -- (-1.583,-0.613) -- (-1.577,-0.630) -- (-1.570,-0.646) -- (-1.564,-0.662) -- (-1.558,-0.678) -- (-1.551,-0.694) -- (-1.545,-0.710) -- (-1.538,-0.725) -- (-1.532,-0.741) -- (-1.525,-0.756) -- (-1.519,-0.771) -- (-1.512,-0.786) -- (-1.505,-0.801) -- (-1.499,-0.816) -- (-1.492,-0.831) -- (-1.485,-0.846) -- (-1.478,-0.860) -- (-1.471,-0.874) -- (-1.464,-0.888) -- (-1.457,-0.902) -- (-1.451,-0.916) -- (-1.444,-0.930) -- (-1.437,-0.944) -- (-1.430,-0.957) -- (-1.422,-0.971) -- (-1.415,-0.984) -- (-1.408,-0.997) -- (-1.401,-1.010) -- (-1.394,-1.023) -- (-1.387,-1.035) -- (-1.380,-1.048) -- (-1.373,-1.060) -- (-1.366,-1.073) -- (-1.358,-1.085) -- (-1.351,-1.097) -- (-1.344,-1.109) -- (-1.337,-1.121) -- (-1.330,-1.132) -- (-1.323,-1.144) -- (-1.315,-1.155) -- (-1.308,-1.167) -- (-1.301,-1.178) -- (-1.294,-1.189) -- (-1.287,-1.200) -- (-1.280,-1.210) -- (-1.272,-1.221) -- (-1.265,-1.232) -- (-1.258,-1.242) -- (-1.251,-1.252) -- (-1.244,-1.263) -- (-1.237,-1.273) -- (-1.230,-1.283) -- (-1.223,-1.292) -- (-1.216,-1.302) -- (-1.209,-1.312) -- (-1.202,-1.321) -- (-1.195,-1.331) -- (-1.188,-1.340) -- (-1.181,-1.349) -- (-1.174,-1.358) -- (-1.167,-1.367) -- (-1.160,-1.376) -- (-1.153,-1.385) -- (-1.146,-1.394) -- (-1.139,-1.402) -- (-1.132,-1.411) -- (-1.125,-1.419) -- (-1.118,-1.427) -- (-1.112,-1.435) -- (-1.105,-1.443) -- (-1.098,-1.451) -- (-1.091,-1.459) -- (-1.085,-1.467) -- (-1.078,-1.475) -- (-1.071,-1.482) -- (-1.065,-1.490) -- (-1.058,-1.497) -- (-1.051,-1.504) -- (-1.045,-1.512) -- (-1.038,-1.519) -- (-1.032,-1.526) -- (-1.022,-1.536) -- (-1.012,-1.547) -- (-1.003,-1.557) -- (-0.993,-1.567) -- (-0.984,-1.576) -- (-0.974,-1.586) -- (-0.965,-1.595) -- (-0.956,-1.604) -- (-0.946,-1.613) -- (-0.937,-1.622) -- (-0.928,-1.631) -- (-0.919,-1.639) -- (-0.910,-1.648) -- (-0.901,-1.656) -- (-0.892,-1.664) -- (-0.883,-1.672) -- (-0.874,-1.680) -- (-0.865,-1.688) -- (-0.857,-1.695) -- (-0.848,-1.702) -- (-0.840,-1.710) -- (-0.831,-1.717) -- (-0.823,-1.724) -- (-0.814,-1.731) -- (-0.806,-1.737) -- (-0.797,-1.744) -- (-0.789,-1.751) -- (-0.781,-1.757) -- (-0.773,-1.763) -- (-0.765,-1.769) -- (-0.757,-1.775) -- (-0.749,-1.781) -- (-0.741,-1.787) -- (-0.733,-1.793) -- (-0.725,-1.799) -- (-0.715,-1.806) -- (-0.705,-1.813) -- (-0.694,-1.820) -- (-0.684,-1.827) -- (-0.674,-1.834) -- (-0.664,-1.840) -- (-0.655,-1.847) -- (-0.645,-1.853) -- (-0.635,-1.859) -- (-0.626,-1.865) -- (-0.616,-1.871) -- (-0.607,-1.876) -- (-0.597,-1.882) -- (-0.588,-1.887) -- (-0.579,-1.892) -- (-0.570,-1.898) -- (-0.561,-1.903) -- (-0.552,-1.908) -- (-0.543,-1.912) -- (-0.534,-1.917) -- (-0.525,-1.922) -- (-0.517,-1.926) -- (-0.508,-1.930) -- (-0.497,-1.936) -- (-0.487,-1.941) -- (-0.476,-1.946) -- (-0.466,-1.951) -- (-0.455,-1.955) -- (-0.445,-1.960) -- (-0.435,-1.964) -- (-0.425,-1.969) -- (-0.415,-1.973) -- (-0.405,-1.977) -- (-0.395,-1.981) -- (-0.385,-1.984) -- (-0.376,-1.988) -- (-0.366,-1.991) -- (-0.356,-1.995) -- (-0.347,-1.998) -- (-0.337,-2.001) -- (-0.328,-2.004) -- (-0.319,-2.007) -- (-0.310,-2.010) -- (-0.299,-2.014) -- (-0.288,-2.017) -- (-0.277,-2.020) -- (-0.266,-2.023) -- (-0.255,-2.026) -- (-0.245,-2.028) -- (-0.234,-2.031) -- (-0.223,-2.033) -- (-0.213,-2.035) -- (-0.203,-2.038) -- (-0.192,-2.040) -- (-0.182,-2.041) -- (-0.172,-2.043) -- (-0.161,-2.045) -- (-0.151,-2.047) -- (-0.141,-2.048) -- (-0.131,-2.049) -- (-0.121,-2.051) -- (-0.111,-2.052) -- (-0.101,-2.053) -- (-0.091,-2.054) -- (-0.081,-2.055) -- (-0.071,-2.055) -- (-0.061,-2.056) -- (-0.051,-2.057) -- (-0.041,-2.057) -- (-0.031,-2.057) -- (-0.021,-2.058) -- (-0.012,-2.058) -- (-0.002,-2.058) -- (0.008,-2.058) -- (0.018,-2.058) -- (0.028,-2.058) -- (0.038,-2.057) -- (0.048,-2.057) -- (0.058,-2.056) -- (0.068,-2.056) -- (0.078,-2.055) -- (0.088,-2.054) -- (0.097,-2.053) -- (0.108,-2.052) -- (0.118,-2.051) -- (0.128,-2.050) -- (0.138,-2.049) -- (0.148,-2.047) -- (0.158,-2.046) -- (0.168,-2.044) -- (0.178,-2.042) -- (0.189,-2.040) -- (0.199,-2.038) -- (0.210,-2.036) -- (0.220,-2.034) -- (0.230,-2.031) -- (0.241,-2.029) -- (0.252,-2.026) -- (0.262,-2.024) -- (0.273,-2.021) -- (0.284,-2.018) -- (0.295,-2.015) -- (0.306,-2.011) -- (0.315,-2.009) -- (0.324,-2.006) -- (0.334,-2.003) -- (0.343,-1.999) -- (0.353,-1.996) -- (0.362,-1.993) -- (0.372,-1.989) -- (0.381,-1.986) -- (0.391,-1.982) -- (0.401,-1.978) -- (0.411,-1.974) -- (0.421,-1.970) -- (0.431,-1.966) -- (0.441,-1.962) -- (0.451,-1.957) -- (0.462,-1.953) -- (0.472,-1.948) -- (0.482,-1.943) -- (0.493,-1.938) -- (0.504,-1.933) -- (0.512,-1.928) -- (0.521,-1.924) -- (0.530,-1.919) -- (0.538,-1.915) -- (0.547,-1.910) -- (0.556,-1.905) -- (0.565,-1.900) -- (0.574,-1.895) -- (0.584,-1.890) -- (0.593,-1.885) -- (0.602,-1.879) -- (0.611,-1.873) -- (0.621,-1.868) -- (0.630,-1.862) -- (0.640,-1.856) -- (0.650,-1.850) -- (0.660,-1.843) -- (0.669,-1.837) -- (0.679,-1.830) -- (0.689,-1.824) -- (0.699,-1.817) -- (0.710,-1.810) -- (0.720,-1.802) -- (0.730,-1.795) -- (0.738,-1.789) -- (0.746,-1.783) -- (0.754,-1.777) -- (0.762,-1.771) -- (0.770,-1.765) -- (0.778,-1.759) -- (0.786,-1.753) -- (0.795,-1.746) -- (0.803,-1.740) -- (0.811,-1.733) -- (0.820,-1.726) -- (0.828,-1.719) -- (0.837,-1.712) -- (0.845,-1.705) -- (0.854,-1.698) -- (0.863,-1.690) -- (0.871,-1.682) -- (0.880,-1.675) -- (0.889,-1.667) -- (0.898,-1.659) -- (0.907,-1.651) -- (0.916,-1.642) -- (0.925,-1.634) -- (0.934,-1.625) -- (0.943,-1.616) -- (0.952,-1.607) -- (0.962,-1.598) -- (0.971,-1.589) -- (0.981,-1.579) -- (0.990,-1.570) -- (1.000,-1.560) -- (1.009,-1.550) -- (1.019,-1.540) -- (1.029,-1.529) -- (1.038,-1.519) -- (1.045,-1.512) -- (1.051,-1.504) -- (1.058,-1.497) -- (1.065,-1.490) -- (1.071,-1.482) -- (1.078,-1.475) -- (1.085,-1.467) -- (1.091,-1.459) -- (1.098,-1.451) -- (1.105,-1.443) -- (1.112,-1.435) -- (1.118,-1.427) -- (1.125,-1.419) -- (1.132,-1.411) -- (1.139,-1.402) -- (1.146,-1.394) -- (1.153,-1.385) -- (1.160,-1.376) -- (1.167,-1.367) -- (1.174,-1.358) -- (1.181,-1.349) -- (1.188,-1.340) -- (1.195,-1.331) -- (1.202,-1.321) -- (1.209,-1.312) -- (1.216,-1.302) -- (1.223,-1.292) -- (1.230,-1.283) -- (1.237,-1.273) -- (1.244,-1.263) -- (1.251,-1.252) -- (1.258,-1.242) -- (1.265,-1.232) -- (1.272,-1.221) -- (1.280,-1.210) -- (1.287,-1.200) -- (1.294,-1.189) -- (1.301,-1.178) -- (1.308,-1.167) -- (1.315,-1.155) -- (1.323,-1.144) -- (1.330,-1.132) -- (1.337,-1.121) -- (1.344,-1.109) -- (1.351,-1.097) -- (1.358,-1.085) -- (1.366,-1.073) -- (1.373,-1.060) -- (1.380,-1.048) -- (1.387,-1.035) -- (1.394,-1.023) -- (1.401,-1.010) -- (1.408,-0.997) -- (1.415,-0.984) -- (1.422,-0.971) -- (1.430,-0.957) -- (1.437,-0.944) -- (1.444,-0.930) -- (1.451,-0.916) -- (1.457,-0.902) -- (1.464,-0.888) -- (1.471,-0.874) -- (1.478,-0.860) -- (1.485,-0.846) -- (1.492,-0.831) -- (1.499,-0.816) -- (1.505,-0.801) -- (1.512,-0.786) -- (1.519,-0.771) -- (1.525,-0.756) -- (1.532,-0.741) -- (1.538,-0.725) -- (1.545,-0.710) -- (1.551,-0.694) -- (1.558,-0.678) -- (1.564,-0.662) -- (1.570,-0.646) -- (1.577,-0.630) -- (1.583,-0.613) -- (1.589,-0.597) -- (1.595,-0.580) -- (1.601,-0.564) -- (1.607,-0.547) -- (1.612,-0.530) -- (1.618,-0.513) -- (1.624,-0.495) -- (1.629,-0.478) -- (1.635,-0.461) -- (1.640,-0.443) -- (1.646,-0.426) -- (1.651,-0.408) -- (1.656,-0.390) -- (1.661,-0.372) -- (1.666,-0.354) -- (1.671,-0.336) -- (1.676,-0.318) -- (1.681,-0.300) -- (1.685,-0.281) -- (1.690,-0.263) -- (1.694,-0.244) -- (1.699,-0.226) -- (1.703,-0.207) -- (1.707,-0.188) -- (1.709,-0.179) -- (1.711,-0.169) -- (1.713,-0.160) -- (1.715,-0.151) -- (1.717,-0.141) -- (1.718,-0.132) -- (1.720,-0.122) -- (1.722,-0.113) -- (1.724,-0.103) -- (1.726,-0.094) -- (1.727,-0.084) -- (1.729,-0.075) -- (1.731,-0.065) -- (1.732,-0.055) -- (1.734,-0.046) -- (1.735,-0.036) -- (1.737,-0.027) -- (1.739,-0.017) -- (1.740,-0.007) -- (1.741,0.002) -- (1.743,0.012) -- (1.744,0.021) -- (1.746,0.031) -- (1.747,0.041) -- (1.748,0.050) -- (1.750,0.060) -- (1.751,0.070) -- (1.752,0.079) -- (1.753,0.089) -- (1.754,0.099) -- (1.755,0.108) -- (1.756,0.118) -- (1.757,0.128) -- (1.758,0.138) -- (1.759,0.147) -- (1.760,0.157) -- (1.761,0.167) -- (1.762,0.176) -- (1.763,0.186) -- (1.764,0.196) -- (1.765,0.205) -- (1.765,0.215) -- (1.766,0.225) -- (1.767,0.234) -- (1.767,0.244) -- (1.768,0.254) -- (1.769,0.263) -- (1.769,0.273) -- (1.770,0.283) -- (1.770,0.292) -- (1.771,0.302) -- (1.771,0.312) -- (1.772,0.321) -- (1.772,0.331) -- (1.772,0.341) -- (1.772,0.350) -- (1.773,0.369) -- (1.773,0.388) -- (1.774,0.408) -- (1.774,0.427) -- (1.774,0.446) -- (1.773,0.464) -- (1.773,0.483) -- (1.773,0.502) -- (1.772,0.521) -- (1.772,0.539) -- (1.771,0.558) -- (1.770,0.577) -- (1.769,0.595) -- (1.768,0.613) -- (1.766,0.632) -- (1.765,0.650) -- (1.763,0.668) -- (1.762,0.686) -- (1.760,0.704) -- (1.758,0.722) -- (1.756,0.739) -- (1.754,0.757) -- (1.752,0.774) -- (1.750,0.792) -- (1.747,0.809) -- (1.745,0.826) -- (1.742,0.843) -- (1.740,0.860) -- (1.737,0.877) -- (1.734,0.893) -- (1.731,0.910) -- (1.728,0.926) -- (1.725,0.943) -- (1.721,0.959) -- (1.718,0.975) -- (1.714,0.991) -- (1.711,1.006) -- (1.707,1.022) -- (1.704,1.037) -- (1.700,1.053) -- (1.696,1.068) -- (1.692,1.083) -- (1.688,1.098) -- (1.684,1.113) -- (1.680,1.128) -- (1.676,1.142) -- (1.671,1.157) -- (1.667,1.171) -- (1.663,1.185) -- (1.658,1.199) -- (1.654,1.213) -- (1.649,1.226) -- (1.645,1.240) -- (1.640,1.253) -- (1.635,1.267) -- (1.630,1.280) -- (1.626,1.293) -- (1.621,1.306) -- (1.616,1.318) -- (1.611,1.331) -- (1.606,1.343) -- (1.601,1.356) -- (1.596,1.368) -- (1.591,1.380) -- (1.586,1.392) -- (1.581,1.404) -- (1.575,1.415) -- (1.570,1.427) -- (1.565,1.438) -- (1.560,1.449) -- (1.554,1.461) -- (1.549,1.471) -- (1.544,1.482) -- (1.538,1.493) -- (1.533,1.504) -- (1.528,1.514) -- (1.522,1.524) -- (1.517,1.535) -- (1.512,1.545) -- (1.506,1.555) -- (1.501,1.565) -- (1.495,1.574) -- (1.490,1.584) -- (1.484,1.593) -- (1.479,1.603) -- (1.473,1.612) -- (1.468,1.621) -- (1.462,1.630) -- (1.457,1.639) -- (1.451,1.648) -- (1.446,1.657) -- (1.440,1.665) -- (1.435,1.674) -- (1.429,1.682) -- (1.424,1.690) -- (1.418,1.698) -- (1.413,1.706) -- (1.407,1.714) -- (1.399,1.726) -- (1.391,1.738) -- (1.383,1.749) -- (1.375,1.760) -- (1.367,1.771) -- (1.358,1.782) -- (1.350,1.792) -- (1.342,1.802) -- (1.334,1.812) -- (1.326,1.822) -- (1.318,1.832) -- (1.310,1.841) -- (1.302,1.851) -- (1.294,1.860) -- (1.287,1.869) -- (1.279,1.878) -- (1.271,1.886) -- (1.263,1.895) -- (1.256,1.903) -- (1.248,1.911) -- (1.240,1.919) -- (1.233,1.927) -- (1.225,1.934) -- (1.218,1.942) -- (1.210,1.949) -- (1.203,1.957) -- (1.196,1.964) -- (1.188,1.971) -- (1.181,1.977) -- (1.174,1.984) -- (1.167,1.991) -- (1.159,1.997) -- (1.150,2.006) -- (1.141,2.014) -- (1.131,2.022) -- (1.122,2.030) -- (1.113,2.037) -- (1.104,2.045) -- (1.095,2.052) -- (1.086,2.059) -- (1.077,2.066) -- (1.069,2.073) -- (1.060,2.079) -- (1.052,2.086) -- (1.043,2.092) -- (1.035,2.098) -- (1.026,2.104) -- (1.018,2.110) -- (1.010,2.115) -- (1.002,2.121) -- (0.994,2.126) -- (0.984,2.133) -- (0.974,2.139) -- (0.965,2.146) -- (0.955,2.152) -- (0.946,2.157) -- (0.937,2.163) -- (0.927,2.169) -- (0.918,2.174) -- (0.909,2.179) -- (0.901,2.184) -- (0.892,2.189) -- (0.883,2.194) -- (0.875,2.199) -- (0.865,2.204) -- (0.855,2.210) -- (0.845,2.215) -- (0.835,2.220) -- (0.826,2.224) -- (0.816,2.229) -- (0.807,2.233) -- (0.798,2.238) -- (0.789,2.242) -- (0.780,2.246) -- (0.771,2.250) -- (0.761,2.255) -- (0.751,2.259) -- (0.741,2.263) -- (0.732,2.267) -- (0.722,2.271) -- (0.713,2.275) -- (0.704,2.278) -- (0.695,2.282) -- (0.685,2.286) -- (0.675,2.289) -- (0.665,2.293) -- (0.655,2.296) -- (0.646,2.300) -- (0.637,2.303) -- (0.628,2.306) -- (0.617,2.309) -- (0.607,2.312) -- (0.598,2.316) -- (0.588,2.319) -- (0.579,2.321) -- (0.570,2.324) -- (0.559,2.327) -- (0.550,2.330) -- (0.540,2.332) -- (0.530,2.335) -- (0.521,2.337) -- (0.511,2.340) -- (0.501,2.343) -- (0.492,2.345) -- (0.482,2.347) -- (0.472,2.350) -- (0.462,2.352) -- (0.452,2.354) -- (0.443,2.356) -- (0.433,2.358) -- (0.423,2.360) -- (0.413,2.362) -- (0.404,2.364) -- (0.394,2.366) -- (0.384,2.367) -- (0.375,2.369) -- (0.365,2.371) -- (0.355,2.372) -- (0.345,2.374) -- (0.335,2.375) -- (0.326,2.377) -- (0.316,2.378) -- (0.306,2.380) -- (0.296,2.381) -- (0.286,2.382) -- (0.277,2.384) -- (0.267,2.385) -- (0.257,2.386) -- (0.247,2.387) -- (0.238,2.388) -- (0.228,2.389) -- (0.218,2.390) -- (0.209,2.391) -- (0.199,2.392) -- (0.189,2.392) -- (0.179,2.393) -- (0.169,2.394) -- (0.160,2.395) -- (0.150,2.395) -- (0.140,2.396) -- (0.130,2.396) -- (0.121,2.397) -- (0.111,2.397) -- (0.101,2.398) -- (0.091,2.398) -- (0.082,2.399) -- (0.072,2.399) -- (0.062,2.399) -- (0.053,2.399) -- (0.043,2.400) -- (0.033,2.400) -- (0.024,2.400) -- (0.014,2.400) -- (0.004,2.400) -- (0.000,2.400);}
\newcommand{\figOneSectionZero}{\draw[thick,gray] (0.000,-2.160) -- (0.000,-1.884);}
\newcommand{\figOneCentralSection}{\draw[thick,gray] (0.000,2.058) -- (0.000,2.400);}
\newcommand{\figOneVertex}{(0.000,-2.058)}
\newcommand{\figOneChartOrbits}[1]{\draw[#1] (8.852,1.888) -- (8.830,1.864) -- (8.808,1.841) -- (8.787,1.818) -- (8.766,1.797) -- (8.746,1.775) -- (8.727,1.754) -- (8.707,1.734) -- (8.689,1.715) -- (8.670,1.695) -- (8.653,1.677) -- (8.635,1.658) -- (8.618,1.640) -- (8.593,1.614) -- (8.569,1.589) -- (8.546,1.565) -- (8.523,1.542) -- (8.501,1.519) -- (8.480,1.497) -- (8.459,1.476) -- (8.439,1.455) -- (8.419,1.435) -- (8.400,1.415) -- (8.382,1.397) -- (8.363,1.378) -- (8.346,1.360) -- (8.329,1.343) -- (8.306,1.321) -- (8.285,1.299) -- (8.264,1.278) -- (8.244,1.258) -- (8.224,1.239) -- (8.205,1.220) -- (8.186,1.202) -- (8.168,1.184) -- (8.151,1.167) -- (8.134,1.150) -- (8.113,1.130) -- (8.093,1.111) -- (8.074,1.092) -- (8.055,1.074) -- (8.036,1.056) -- (8.019,1.040) -- (8.002,1.023) -- (7.982,1.004) -- (7.962,0.986) -- (7.943,0.969) -- (7.925,0.952) -- (7.908,0.936) -- (7.888,0.917) -- (7.869,0.900) -- (7.850,0.883) -- (7.832,0.867) -- (7.815,0.851) -- (7.796,0.834) -- (7.777,0.817) -- (7.759,0.801) -- (7.742,0.786) -- (7.723,0.769) -- (7.705,0.754) -- (7.687,0.738) -- (7.669,0.722) -- (7.651,0.707) -- (7.633,0.692) -- (7.614,0.676) -- (7.597,0.661) -- (7.579,0.647) -- (7.561,0.632) -- (7.544,0.618) -- (7.526,0.603) -- (7.508,0.589) -- (7.490,0.574) -- (7.473,0.560) -- (7.455,0.546) -- (7.437,0.533) -- (7.420,0.519) -- (7.401,0.505) -- (7.384,0.492) -- (7.366,0.479) -- (7.349,0.466) -- (7.331,0.452) -- (7.313,0.439) -- (7.295,0.427) -- (7.278,0.415) -- (7.260,0.402) -- (7.243,0.390) -- (7.226,0.378) -- (7.208,0.366) -- (7.190,0.355) -- (7.173,0.343) -- (7.155,0.331) -- (7.138,0.320) -- (7.120,0.309) -- (7.102,0.298) -- (7.085,0.288) -- (7.067,0.277) -- (7.050,0.267) -- (7.033,0.257) -- (7.015,0.247) -- (6.997,0.237) -- (6.980,0.227) -- (6.962,0.218) -- (6.944,0.209) -- (6.927,0.200) -- (6.909,0.191) -- (6.892,0.182) -- (6.874,0.174) -- (6.856,0.165) -- (6.839,0.157) -- (6.821,0.150) -- (6.804,0.142) -- (6.786,0.135) -- (6.769,0.128) -- (6.751,0.121) -- (6.734,0.114) -- (6.716,0.108) -- (6.698,0.102) -- (6.681,0.096) -- (6.663,0.090) -- (6.646,0.085) -- (6.628,0.080) -- (6.610,0.075) -- (6.593,0.070) -- (6.575,0.066) -- (6.558,0.062) -- (6.540,0.058) -- (6.522,0.054) -- (6.505,0.051) -- (6.487,0.048) -- (6.470,0.045) -- (6.452,0.042) -- (6.434,0.040) -- (6.417,0.038) -- (6.399,0.037) -- (6.382,0.035) -- (6.364,0.034) -- (6.346,0.033) -- (6.329,0.033) -- (6.311,0.032) -- (6.293,0.032) -- (6.276,0.033) -- (6.258,0.033) -- (6.241,0.034) -- (6.223,0.035) -- (6.205,0.036) -- (6.188,0.037) -- (6.170,0.039) -- (6.153,0.041) -- (6.135,0.043) -- (6.117,0.046) -- (6.100,0.048) -- (6.082,0.051) -- (6.064,0.054) -- (6.047,0.058) -- (6.029,0.061) -- (6.012,0.065) -- (5.994,0.069) -- (5.976,0.073) -- (5.959,0.078) -- (5.941,0.083) -- (5.924,0.088) -- (5.906,0.093) -- (5.888,0.098) -- (5.871,0.104) -- (5.853,0.110) -- (5.836,0.116) -- (5.818,0.122) -- (5.800,0.129) -- (5.783,0.135) -- (5.765,0.142) -- (5.748,0.149) -- (5.730,0.157) -- (5.713,0.164) -- (5.695,0.172) -- (5.678,0.180) -- (5.660,0.188) -- (5.643,0.197) -- (5.625,0.205) -- (5.607,0.214) -- (5.590,0.223) -- (5.572,0.232) -- (5.554,0.241) -- (5.537,0.251) -- (5.519,0.260) -- (5.502,0.270) -- (5.484,0.280) -- (5.466,0.290) -- (5.449,0.301) -- (5.431,0.311) -- (5.414,0.322) -- (5.396,0.333) -- (5.378,0.344) -- (5.361,0.355) -- (5.343,0.366) -- (5.326,0.377) -- (5.308,0.389) -- (5.290,0.401) -- (5.273,0.413) -- (5.255,0.425) -- (5.237,0.437) -- (5.220,0.449) -- (5.202,0.462) -- (5.184,0.474) -- (5.166,0.487) -- (5.148,0.500) -- (5.130,0.513) -- (5.113,0.526) -- (5.096,0.539) -- (5.078,0.552) -- (5.060,0.566) -- (5.043,0.579) -- (5.025,0.593) -- (5.007,0.607) -- (4.989,0.621) -- (4.971,0.635) -- (4.954,0.649) -- (4.936,0.663) -- (4.919,0.677) -- (4.901,0.691) -- (4.883,0.706) -- (4.865,0.721) -- (4.847,0.736) -- (4.829,0.750) -- (4.811,0.765) -- (4.794,0.780) -- (4.776,0.795) -- (4.758,0.810) -- (4.741,0.825) -- (4.723,0.840) -- (4.705,0.856) -- (4.686,0.872) -- (4.669,0.887) -- (4.651,0.903) -- (4.632,0.919) -- (4.615,0.935) -- (4.597,0.950) -- (4.579,0.967) -- (4.560,0.983) -- (4.543,0.999) -- (4.525,1.015) -- (4.507,1.031) -- (4.489,1.048) -- (4.470,1.065) -- (4.453,1.081) -- (4.435,1.097) -- (4.417,1.114) -- (4.398,1.131) -- (4.379,1.149) -- (4.360,1.167) -- (4.343,1.183) -- (4.325,1.200) -- (4.307,1.217) -- (4.288,1.234) -- (4.269,1.253) -- (4.250,1.271) -- (4.229,1.290) -- (4.212,1.307) -- (4.195,1.324) -- (4.177,1.341) -- (4.158,1.359) -- (4.139,1.377) -- (4.120,1.396) -- (4.100,1.415) -- (4.080,1.435) -- (4.059,1.455) -- (4.042,1.472) -- (4.024,1.489) -- (4.006,1.507) -- (3.988,1.525) -- (3.969,1.544) -- (3.950,1.563) -- (3.930,1.582) -- (3.910,1.602) -- (3.889,1.623) -- (3.868,1.644) -- (3.846,1.666) -- (3.824,1.688) -- (3.807,1.705) -- (3.790,1.723) -- (3.772,1.741) -- (3.754,1.759) -- (3.735,1.778) -- (3.716,1.797) -- (3.697,1.817) -- (3.677,1.837) -- (3.663,1.851); \draw[#1] (8.852,2.064) -- (8.828,2.037) -- (8.805,2.011) -- (8.783,1.985) -- (8.761,1.960) -- (8.740,1.936) -- (8.719,1.913) -- (8.699,1.890) -- (8.679,1.868) -- (8.660,1.846) -- (8.641,1.825) -- (8.623,1.805) -- (8.605,1.785) -- (8.587,1.765) -- (8.570,1.746) -- (8.545,1.719) -- (8.521,1.692) -- (8.498,1.666) -- (8.475,1.641) -- (8.453,1.617) -- (8.432,1.594) -- (8.411,1.571) -- (8.391,1.549) -- (8.371,1.528) -- (8.352,1.507) -- (8.333,1.487) -- (8.315,1.468) -- (8.298,1.449) -- (8.280,1.430) -- (8.258,1.407) -- (8.237,1.384) -- (8.216,1.362) -- (8.196,1.341) -- (8.176,1.320) -- (8.157,1.300) -- (8.138,1.281) -- (8.120,1.262) -- (8.103,1.244) -- (8.086,1.226) -- (8.065,1.205) -- (8.045,1.185) -- (8.026,1.165) -- (8.007,1.146) -- (7.989,1.127) -- (7.971,1.110) -- (7.954,1.093) -- (7.934,1.073) -- (7.915,1.053) -- (7.896,1.035) -- (7.878,1.017) -- (7.860,1.000) -- (7.843,0.984) -- (7.824,0.965) -- (7.806,0.947) -- (7.788,0.930) -- (7.770,0.913) -- (7.751,0.895) -- (7.732,0.878) -- (7.714,0.861) -- (7.697,0.845) -- (7.678,0.827) -- (7.660,0.810) -- (7.642,0.794) -- (7.624,0.777) -- (7.605,0.761) -- (7.588,0.745) -- (7.569,0.729) -- (7.551,0.713) -- (7.534,0.698) -- (7.516,0.682) -- (7.498,0.667) -- (7.480,0.651) -- (7.462,0.637) -- (7.444,0.621) -- (7.427,0.607) -- (7.409,0.592) -- (7.391,0.578) -- (7.373,0.563) -- (7.355,0.549) -- (7.338,0.535) -- (7.320,0.521) -- (7.302,0.507) -- (7.285,0.494) -- (7.267,0.481) -- (7.250,0.468) -- (7.232,0.455) -- (7.214,0.442) -- (7.197,0.429) -- (7.179,0.416) -- (7.161,0.404) -- (7.144,0.392) -- (7.126,0.380) -- (7.108,0.368) -- (7.091,0.357) -- (7.073,0.345) -- (7.056,0.334) -- (7.038,0.323) -- (7.020,0.312) -- (7.003,0.302) -- (6.986,0.291) -- (6.968,0.281) -- (6.950,0.271) -- (6.933,0.261) -- (6.915,0.252) -- (6.898,0.243) -- (6.880,0.233) -- (6.862,0.225) -- (6.845,0.216) -- (6.827,0.207) -- (6.810,0.199) -- (6.792,0.191) -- (6.774,0.184) -- (6.757,0.176) -- (6.739,0.169) -- (6.722,0.162) -- (6.704,0.156) -- (6.687,0.149) -- (6.669,0.143) -- (6.651,0.137) -- (6.634,0.132) -- (6.616,0.126) -- (6.598,0.121) -- (6.581,0.117) -- (6.563,0.112) -- (6.546,0.108) -- (6.528,0.104) -- (6.510,0.101) -- (6.493,0.097) -- (6.475,0.094) -- (6.458,0.092) -- (6.440,0.089) -- (6.422,0.087) -- (6.405,0.085) -- (6.387,0.084) -- (6.369,0.083) -- (6.352,0.082) -- (6.334,0.081) -- (6.316,0.081) -- (6.299,0.080) -- (6.281,0.081) -- (6.263,0.081) -- (6.246,0.082) -- (6.228,0.083) -- (6.211,0.084) -- (6.193,0.085) -- (6.175,0.087) -- (6.158,0.089) -- (6.140,0.091) -- (6.122,0.093) -- (6.105,0.096) -- (6.087,0.099) -- (6.069,0.102) -- (6.052,0.105) -- (6.034,0.109) -- (6.016,0.113) -- (5.999,0.117) -- (5.981,0.121) -- (5.963,0.125) -- (5.946,0.130) -- (5.928,0.135) -- (5.911,0.140) -- (5.893,0.145) -- (5.875,0.151) -- (5.858,0.157) -- (5.840,0.163) -- (5.823,0.169) -- (5.805,0.175) -- (5.787,0.182) -- (5.770,0.189) -- (5.752,0.196) -- (5.734,0.203) -- (5.717,0.210) -- (5.699,0.218) -- (5.682,0.226) -- (5.664,0.234) -- (5.646,0.242) -- (5.629,0.250) -- (5.611,0.259) -- (5.594,0.267) -- (5.576,0.276) -- (5.559,0.285) -- (5.541,0.295) -- (5.523,0.304) -- (5.506,0.314) -- (5.488,0.324) -- (5.470,0.334) -- (5.452,0.344) -- (5.435,0.354) -- (5.417,0.365) -- (5.399,0.376) -- (5.382,0.386) -- (5.364,0.397) -- (5.346,0.409) -- (5.329,0.420) -- (5.311,0.431) -- (5.293,0.443) -- (5.276,0.455) -- (5.258,0.466) -- (5.241,0.478) -- (5.223,0.490) -- (5.205,0.503) -- (5.187,0.515) -- (5.170,0.528) -- (5.152,0.540) -- (5.134,0.553) -- (5.117,0.566) -- (5.099,0.579) -- (5.081,0.593) -- (5.064,0.606) -- (5.046,0.620) -- (5.028,0.633) -- (5.010,0.647) -- (4.992,0.661) -- (4.974,0.675) -- (4.956,0.689) -- (4.938,0.704) -- (4.921,0.718) -- (4.903,0.732) -- (4.885,0.746) -- (4.868,0.761) -- (4.849,0.776) -- (4.832,0.790) -- (4.814,0.805) -- (4.795,0.821) -- (4.778,0.835) -- (4.760,0.851) -- (4.741,0.866) -- (4.724,0.881) -- (4.706,0.896) -- (4.688,0.912) -- (4.669,0.929) -- (4.652,0.944) -- (4.634,0.960) -- (4.616,0.976) -- (4.597,0.993) -- (4.580,1.008) -- (4.562,1.024) -- (4.544,1.040) -- (4.525,1.057) -- (4.506,1.075) -- (4.489,1.090) -- (4.471,1.106) -- (4.453,1.123) -- (4.434,1.140) -- (4.415,1.158) -- (4.395,1.176) -- (4.378,1.192) -- (4.360,1.209) -- (4.342,1.226) -- (4.323,1.244) -- (4.304,1.262) -- (4.284,1.280) -- (4.264,1.300) -- (4.247,1.316) -- (4.229,1.333) -- (4.211,1.350) -- (4.193,1.368) -- (4.174,1.386) -- (4.155,1.405) -- (4.135,1.424) -- (4.114,1.444) -- (4.093,1.465) -- (4.076,1.482) -- (4.058,1.499) -- (4.040,1.516) -- (4.022,1.535) -- (4.003,1.553) -- (3.984,1.572) -- (3.964,1.592) -- (3.944,1.612) -- (3.924,1.632) -- (3.902,1.653) -- (3.881,1.675) -- (3.858,1.698) -- (3.841,1.715) -- (3.824,1.732) -- (3.806,1.750) -- (3.788,1.769) -- (3.769,1.788) -- (3.750,1.807) -- (3.731,1.827) -- (3.711,1.847) -- (3.691,1.867) -- (3.670,1.889) -- (3.663,1.896); \draw[#1] (8.852,2.256) -- (8.827,2.225) -- (8.802,2.195) -- (8.778,2.166) -- (8.755,2.138) -- (8.733,2.111) -- (8.711,2.085) -- (8.689,2.059) -- (8.669,2.034) -- (8.648,2.010) -- (8.629,1.986) -- (8.609,1.963) -- (8.590,1.941) -- (8.572,1.919) -- (8.554,1.898) -- (8.537,1.877) -- (8.520,1.857) -- (8.503,1.837) -- (8.478,1.808) -- (8.455,1.781) -- (8.432,1.754) -- (8.410,1.728) -- (8.388,1.703) -- (8.367,1.679) -- (8.347,1.655) -- (8.327,1.633) -- (8.308,1.611) -- (8.289,1.589) -- (8.271,1.568) -- (8.253,1.548) -- (8.236,1.529) -- (8.219,1.509) -- (8.197,1.485) -- (8.176,1.461) -- (8.155,1.438) -- (8.135,1.416) -- (8.116,1.395) -- (8.097,1.374) -- (8.079,1.354) -- (8.061,1.335) -- (8.044,1.316) -- (8.023,1.293) -- (8.003,1.271) -- (7.984,1.250) -- (7.965,1.230) -- (7.946,1.210) -- (7.929,1.191) -- (7.911,1.173) -- (7.891,1.152) -- (7.872,1.131) -- (7.853,1.112) -- (7.835,1.093) -- (7.817,1.074) -- (7.800,1.057) -- (7.780,1.037) -- (7.762,1.018) -- (7.744,1.000) -- (7.726,0.982) -- (7.709,0.965) -- (7.690,0.947) -- (7.672,0.929) -- (7.655,0.911) -- (7.636,0.893) -- (7.617,0.875) -- (7.599,0.858) -- (7.582,0.842) -- (7.564,0.824) -- (7.546,0.808) -- (7.529,0.791) -- (7.510,0.775) -- (7.493,0.758) -- (7.474,0.742) -- (7.457,0.726) -- (7.439,0.710) -- (7.421,0.694) -- (7.404,0.679) -- (7.386,0.664) -- (7.369,0.649) -- (7.351,0.633) -- (7.334,0.619) -- (7.316,0.604) -- (7.298,0.589) -- (7.281,0.575) -- (7.263,0.561) -- (7.245,0.547) -- (7.228,0.533) -- (7.210,0.520) -- (7.193,0.506) -- (7.175,0.493) -- (7.157,0.480) -- (7.140,0.467) -- (7.122,0.454) -- (7.104,0.441) -- (7.087,0.429) -- (7.069,0.417) -- (7.052,0.405) -- (7.034,0.393) -- (7.016,0.382) -- (6.999,0.370) -- (6.981,0.359) -- (6.964,0.349) -- (6.946,0.338) -- (6.929,0.328) -- (6.911,0.318) -- (6.894,0.308) -- (6.876,0.298) -- (6.859,0.289) -- (6.841,0.280) -- (6.823,0.271) -- (6.806,0.262) -- (6.788,0.254) -- (6.770,0.246) -- (6.753,0.238) -- (6.735,0.230) -- (6.718,0.223) -- (6.700,0.216) -- (6.682,0.210) -- (6.665,0.203) -- (6.647,0.197) -- (6.630,0.191) -- (6.612,0.186) -- (6.594,0.181) -- (6.577,0.176) -- (6.559,0.171) -- (6.541,0.167) -- (6.524,0.163) -- (6.506,0.159) -- (6.488,0.156) -- (6.470,0.153) -- (6.453,0.150) -- (6.435,0.148) -- (6.417,0.146) -- (6.400,0.144) -- (6.382,0.142) -- (6.364,0.141) -- (6.347,0.140) -- (6.329,0.140) -- (6.311,0.139) -- (6.294,0.139) -- (6.276,0.139) -- (6.258,0.140) -- (6.241,0.141) -- (6.223,0.142) -- (6.205,0.143) -- (6.188,0.145) -- (6.170,0.147) -- (6.152,0.149) -- (6.135,0.151) -- (6.117,0.154) -- (6.099,0.156) -- (6.082,0.159) -- (6.064,0.163) -- (6.046,0.166) -- (6.029,0.170) -- (6.011,0.174) -- (5.993,0.178) -- (5.975,0.182) -- (5.958,0.187) -- (5.940,0.192) -- (5.922,0.197) -- (5.905,0.202) -- (5.887,0.208) -- (5.870,0.213) -- (5.852,0.219) -- (5.834,0.225) -- (5.817,0.232) -- (5.799,0.238) -- (5.781,0.245) -- (5.764,0.252) -- (5.746,0.259) -- (5.728,0.266) -- (5.711,0.274) -- (5.693,0.281) -- (5.675,0.289) -- (5.658,0.297) -- (5.640,0.306) -- (5.622,0.314) -- (5.604,0.323) -- (5.587,0.332) -- (5.569,0.341) -- (5.551,0.350) -- (5.533,0.359) -- (5.516,0.369) -- (5.498,0.379) -- (5.481,0.388) -- (5.463,0.398) -- (5.446,0.409) -- (5.428,0.419) -- (5.411,0.429) -- (5.393,0.440) -- (5.376,0.451) -- (5.358,0.462) -- (5.340,0.473) -- (5.322,0.485) -- (5.304,0.496) -- (5.287,0.508) -- (5.269,0.519) -- (5.251,0.531) -- (5.234,0.543) -- (5.216,0.556) -- (5.198,0.568) -- (5.181,0.580) -- (5.163,0.593) -- (5.146,0.605) -- (5.128,0.618) -- (5.110,0.631) -- (5.092,0.644) -- (5.075,0.658) -- (5.057,0.671) -- (5.040,0.684) -- (5.022,0.698) -- (5.003,0.712) -- (4.986,0.726) -- (4.968,0.740) -- (4.950,0.754) -- (4.932,0.768) -- (4.914,0.783) -- (4.896,0.797) -- (4.878,0.812) -- (4.860,0.827) -- (4.842,0.841) -- (4.825,0.856) -- (4.806,0.872) -- (4.789,0.886) -- (4.771,0.901) -- (4.753,0.917) -- (4.734,0.933) -- (4.717,0.948) -- (4.699,0.963) -- (4.681,0.979) -- (4.662,0.995) -- (4.643,1.012) -- (4.625,1.028) -- (4.607,1.044) -- (4.589,1.060) -- (4.570,1.077) -- (4.553,1.093) -- (4.535,1.109) -- (4.517,1.125) -- (4.499,1.142) -- (4.480,1.159) -- (4.460,1.177) -- (4.443,1.193) -- (4.426,1.209) -- (4.408,1.226) -- (4.390,1.243) -- (4.371,1.260) -- (4.351,1.278) -- (4.332,1.297) -- (4.311,1.316) -- (4.294,1.333) -- (4.276,1.350) -- (4.258,1.367) -- (4.239,1.385) -- (4.220,1.404) -- (4.201,1.423) -- (4.181,1.442) -- (4.160,1.462) -- (4.139,1.483) -- (4.122,1.500) -- (4.104,1.517) -- (4.086,1.535) -- (4.067,1.553) -- (4.048,1.572) -- (4.029,1.591) -- (4.009,1.610) -- (3.989,1.631) -- (3.968,1.651) -- (3.947,1.673) -- (3.925,1.695) -- (3.903,1.717) -- (3.885,1.734) -- (3.868,1.752) -- (3.850,1.770) -- (3.832,1.789) -- (3.813,1.808) -- (3.794,1.827) -- (3.774,1.847) -- (3.755,1.867) -- (3.734,1.888) -- (3.713,1.910) -- (3.692,1.932) -- (3.670,1.954) -- (3.663,1.962); \draw[#1] (8.852,0.560) -- (8.829,0.550) -- (8.807,0.541) -- (8.786,0.531) -- (8.765,0.522) -- (8.745,0.514) -- (8.725,0.505) -- (8.705,0.497) -- (8.687,0.489) -- (8.668,0.481) -- (8.650,0.474) -- (8.629,0.465) -- (8.608,0.456) -- (8.587,0.447) -- (8.567,0.439) -- (8.548,0.431) -- (8.529,0.423) -- (8.511,0.416) -- (8.493,0.409) -- (8.472,0.400) -- (8.452,0.392) -- (8.432,0.384) -- (8.413,0.377) -- (8.395,0.369) -- (8.377,0.362) -- (8.357,0.354) -- (8.337,0.347) -- (8.318,0.339) -- (8.299,0.332) -- (8.281,0.325) -- (8.261,0.318) -- (8.242,0.310) -- (8.224,0.303) -- (8.205,0.296) -- (8.188,0.290) -- (8.169,0.283) -- (8.150,0.276) -- (8.132,0.270) -- (8.113,0.263) -- (8.095,0.256) -- (8.077,0.250) -- (8.057,0.243) -- (8.039,0.236) -- (8.021,0.230) -- (8.002,0.224) -- (7.984,0.217) -- (7.965,0.211) -- (7.947,0.205) -- (7.929,0.199) -- (7.911,0.193) -- (7.893,0.187) -- (7.875,0.181) -- (7.858,0.176) -- (7.840,0.170) -- (7.821,0.164) -- (7.804,0.159) -- (7.786,0.153) -- (7.767,0.148) -- (7.750,0.143) -- (7.732,0.137) -- (7.714,0.132) -- (7.696,0.127) -- (7.678,0.122) -- (7.660,0.117) -- (7.642,0.112) -- (7.625,0.108) -- (7.607,0.103) -- (7.589,0.098) -- (7.571,0.094) -- (7.553,0.089) -- (7.535,0.085) -- (7.518,0.081) -- (7.500,0.077) -- (7.482,0.073) -- (7.464,0.069) -- (7.446,0.065) -- (7.429,0.061) -- (7.411,0.058) -- (7.393,0.054) -- (7.375,0.051) -- (7.358,0.047) -- (7.340,0.044) -- (7.322,0.041) -- (7.304,0.038) -- (7.287,0.035) -- (7.269,0.032) -- (7.251,0.030) -- (7.234,0.027) -- (7.216,0.025) -- (7.198,0.023) -- (7.181,0.020) -- (7.163,0.018) -- (7.145,0.017) -- (7.128,0.015) -- (7.110,0.013) -- (7.092,0.012) -- (7.075,0.010) -- (7.057,0.009) -- (7.039,0.008) -- (7.022,0.007) -- (7.004,0.006) -- (6.986,0.005) -- (6.969,0.004) -- (6.951,0.003) -- (6.934,0.003) -- (6.916,0.002) -- (6.898,0.002) -- (6.881,0.001) -- (6.863,0.001) -- (6.845,0.001) -- (6.828,0.001) -- (6.810,0.000) -- (6.793,0.000) -- (6.775,0.000) -- (6.757,0.000) -- (6.740,0.000) -- (6.722,0.000) -- (6.705,0.000) -- (6.687,0.000) -- (6.669,0.000) -- (6.652,0.000) -- (6.634,0.000) -- (6.617,0.000) -- (6.599,0.000) -- (6.581,0.000) -- (6.564,0.000) -- (6.546,0.000) -- (6.529,0.000) -- (6.511,0.000) -- (6.493,0.000) -- (6.476,0.000) -- (6.458,0.000) -- (6.441,0.000) -- (6.423,0.000) -- (6.405,0.000) -- (6.388,0.000) -- (6.378,0.000); \draw[#1] (8.852,1.120) -- (8.828,1.103) -- (8.805,1.086) -- (8.782,1.069) -- (8.760,1.053) -- (8.739,1.038) -- (8.718,1.023) -- (8.698,1.008) -- (8.678,0.994) -- (8.659,0.980) -- (8.640,0.967) -- (8.622,0.954) -- (8.604,0.941) -- (8.587,0.929) -- (8.565,0.913) -- (8.543,0.898) -- (8.521,0.883) -- (8.501,0.868) -- (8.481,0.854) -- (8.461,0.841) -- (8.443,0.828) -- (8.424,0.815) -- (8.406,0.803) -- (8.389,0.791) -- (8.368,0.776) -- (8.347,0.762) -- (8.327,0.749) -- (8.308,0.735) -- (8.289,0.723) -- (8.271,0.711) -- (8.254,0.699) -- (8.233,0.685) -- (8.213,0.672) -- (8.194,0.659) -- (8.175,0.647) -- (8.157,0.635) -- (8.140,0.623) -- (8.120,0.610) -- (8.100,0.598) -- (8.082,0.586) -- (8.064,0.574) -- (8.047,0.563) -- (8.027,0.551) -- (8.009,0.539) -- (7.991,0.528) -- (7.973,0.517) -- (7.954,0.505) -- (7.936,0.494) -- (7.918,0.483) -- (7.899,0.471) -- (7.881,0.460) -- (7.863,0.449) -- (7.845,0.438) -- (7.827,0.427) -- (7.809,0.417) -- (7.791,0.406) -- (7.774,0.396) -- (7.755,0.385) -- (7.738,0.375) -- (7.719,0.365) -- (7.702,0.355) -- (7.684,0.345) -- (7.665,0.334) -- (7.647,0.324) -- (7.629,0.314) -- (7.611,0.304) -- (7.593,0.295) -- (7.575,0.286) -- (7.558,0.276) -- (7.540,0.267) -- (7.522,0.257) -- (7.504,0.248) -- (7.486,0.239) -- (7.468,0.230) -- (7.450,0.222) -- (7.432,0.213) -- (7.415,0.204) -- (7.397,0.196) -- (7.379,0.188) -- (7.361,0.179) -- (7.344,0.171) -- (7.326,0.163) -- (7.308,0.155) -- (7.290,0.148) -- (7.272,0.140) -- (7.255,0.133) -- (7.237,0.126) -- (7.220,0.119) -- (7.202,0.112) -- (7.184,0.105) -- (7.166,0.098) -- (7.149,0.092) -- (7.131,0.086) -- (7.114,0.079) -- (7.096,0.074) -- (7.079,0.068) -- (7.061,0.062) -- (7.043,0.057) -- (7.026,0.052) -- (7.008,0.047) -- (6.990,0.042) -- (6.973,0.037) -- (6.955,0.033) -- (6.938,0.029) -- (6.920,0.025) -- (6.902,0.022) -- (6.885,0.018) -- (6.867,0.015) -- (6.850,0.013) -- (6.832,0.010) -- (6.814,0.008) -- (6.797,0.006) -- (6.779,0.005) -- (6.761,0.003) -- (6.744,0.002) -- (6.726,0.002) -- (6.709,0.001) -- (6.691,0.001) -- (6.673,0.000) -- (6.656,0.000) -- (6.638,0.000) -- (6.621,0.000) -- (6.603,0.000) -- (6.585,0.000) -- (6.568,0.000) -- (6.550,0.000) -- (6.533,0.000) -- (6.515,0.000) -- (6.497,0.000) -- (6.480,0.000) -- (6.462,0.000) -- (6.445,0.000) -- (6.427,0.000) -- (6.409,0.000) -- (6.392,0.000) -- (6.377,0.000); \draw[#1] (3.748,0.096) -- (3.768,0.097) -- (3.788,0.098) -- (3.807,0.098) -- (3.825,0.099) -- (3.843,0.100) -- (3.861,0.101) -- (3.881,0.102) -- (3.900,0.102) -- (3.918,0.103) -- (3.936,0.104) -- (3.956,0.105) -- (3.975,0.106) -- (3.994,0.108) -- (4.012,0.109) -- (4.031,0.110) -- (4.050,0.111) -- (4.068,0.112) -- (4.086,0.114) -- (4.104,0.115) -- (4.123,0.117) -- (4.140,0.118) -- (4.159,0.120) -- (4.177,0.121) -- (4.195,0.123) -- (4.213,0.125) -- (4.232,0.127) -- (4.250,0.129) -- (4.269,0.131) -- (4.287,0.134) -- (4.305,0.136) -- (4.323,0.139) -- (4.341,0.142) -- (4.359,0.145) -- (4.377,0.148) -- (4.395,0.151) -- (4.413,0.155) -- (4.431,0.159) -- (4.449,0.163) -- (4.467,0.168) -- (4.485,0.173) -- (4.503,0.179) -- (4.520,0.186) -- (4.538,0.193) -- (4.556,0.202) -- (4.573,0.212) -- (4.591,0.225) -- (4.608,0.242) -- (4.624,0.268) -- (4.631,0.326) -- (4.619,0.373) -- (4.604,0.407) -- (4.588,0.436) -- (4.572,0.462) -- (4.555,0.486) -- (4.539,0.509) -- (4.522,0.532) -- (4.505,0.553) -- (4.488,0.574) -- (4.471,0.594) -- (4.454,0.614) -- (4.437,0.635) -- (4.420,0.654) -- (4.403,0.673) -- (4.386,0.692) -- (4.368,0.712) -- (4.350,0.731) -- (4.333,0.750) -- (4.316,0.768) -- (4.298,0.787) -- (4.281,0.805) -- (4.264,0.824) -- (4.246,0.842) -- (4.229,0.860) -- (4.210,0.880) -- (4.192,0.898) -- (4.174,0.917) -- (4.157,0.935) -- (4.139,0.953) -- (4.121,0.972) -- (4.104,0.989) -- (4.086,1.007) -- (4.068,1.026) -- (4.049,1.045) -- (4.032,1.063) -- (4.015,1.080) -- (3.997,1.099) -- (3.978,1.118) -- (3.959,1.138) -- (3.942,1.155) -- (3.924,1.173) -- (3.906,1.191) -- (3.887,1.210) -- (3.868,1.230) -- (3.848,1.250) -- (3.831,1.267) -- (3.814,1.285) -- (3.796,1.303) -- (3.777,1.322) -- (3.758,1.341) -- (3.739,1.361) -- (3.718,1.381) -- (3.698,1.402) -- (3.681,1.419) -- (3.663,1.437); \draw[#1] (3.748,0.320) -- (3.766,0.325) -- (3.785,0.331) -- (3.803,0.337) -- (3.821,0.343) -- (3.839,0.350) -- (3.858,0.358) -- (3.876,0.367) -- (3.894,0.376) -- (3.912,0.387) -- (3.930,0.400) -- (3.948,0.415) -- (3.965,0.433) -- (3.981,0.457) -- (3.996,0.494) -- (4.000,0.556) -- (3.990,0.608) -- (3.976,0.649) -- (3.961,0.682) -- (3.945,0.713) -- (3.929,0.742) -- (3.912,0.769) -- (3.896,0.795) -- (3.878,0.820) -- (3.861,0.845) -- (3.843,0.870) -- (3.826,0.893) -- (3.808,0.916) -- (3.790,0.939) -- (3.773,0.961) -- (3.756,0.982) -- (3.738,1.004) -- (3.719,1.026) -- (3.702,1.047) -- (3.684,1.068) -- (3.665,1.090) -- (3.662,1.093); \draw[#1] (3.748,0.720) -- (3.734,0.758) -- (3.718,0.792) -- (3.702,0.823) -- (3.685,0.853) -- (3.669,0.880) -- (3.661,0.893);}
\newcommand{\figOneChartLeft}[1]{\draw[#1] (3.660,1.832) -- (3.673,1.818) -- (3.686,1.805) -- (3.700,1.791) -- (3.713,1.778) -- (3.726,1.765) -- (3.739,1.751) -- (3.752,1.738) -- (3.766,1.725) -- (3.779,1.711) -- (3.792,1.698) -- (3.805,1.685) -- (3.818,1.672) -- (3.832,1.658) -- (3.845,1.645) -- (3.858,1.632) -- (3.871,1.619) -- (3.884,1.606) -- (3.898,1.592) -- (3.911,1.579) -- (3.924,1.566) -- (3.937,1.553) -- (3.950,1.540) -- (3.964,1.527) -- (3.977,1.514) -- (3.990,1.501) -- (4.003,1.488) -- (4.016,1.475) -- (4.030,1.462) -- (4.043,1.449) -- (4.056,1.436) -- (4.069,1.423) -- (4.082,1.410) -- (4.096,1.398) -- (4.109,1.385) -- (4.122,1.372) -- (4.135,1.359) -- (4.148,1.346) -- (4.162,1.334) -- (4.175,1.321) -- (4.188,1.308) -- (4.201,1.296) -- (4.214,1.283) -- (4.228,1.270) -- (4.241,1.258) -- (4.254,1.245) -- (4.267,1.233) -- (4.280,1.220) -- (4.294,1.208) -- (4.307,1.195) -- (4.320,1.183) -- (4.333,1.170) -- (4.346,1.158) -- (4.360,1.146) -- (4.373,1.133) -- (4.386,1.121) -- (4.399,1.109) -- (4.412,1.097) -- (4.426,1.084) -- (4.439,1.072) -- (4.452,1.060) -- (4.465,1.048) -- (4.478,1.036) -- (4.492,1.024) -- (4.505,1.012) -- (4.518,1.000) -- (4.531,0.988) -- (4.544,0.976) -- (4.558,0.964) -- (4.571,0.952) -- (4.584,0.940) -- (4.597,0.929) -- (4.610,0.917) -- (4.624,0.905) -- (4.637,0.894) -- (4.650,0.882) -- (4.663,0.870) -- (4.676,0.859) -- (4.690,0.847) -- (4.703,0.836) -- (4.716,0.824) -- (4.729,0.813) -- (4.742,0.802) -- (4.756,0.790) -- (4.769,0.779) -- (4.782,0.768) -- (4.795,0.757) -- (4.808,0.746) -- (4.822,0.735) -- (4.835,0.724) -- (4.848,0.713) -- (4.861,0.702) -- (4.874,0.691) -- (4.888,0.680) -- (4.901,0.669) -- (4.914,0.658) -- (4.927,0.648) -- (4.940,0.637) -- (4.954,0.626) -- (4.967,0.616) -- (4.980,0.605) -- (4.993,0.595) -- (5.006,0.585) -- (5.020,0.574) -- (5.033,0.564) -- (5.046,0.554) -- (5.059,0.544) -- (5.072,0.534) -- (5.086,0.524) -- (5.099,0.514) -- (5.112,0.504) -- (5.125,0.494) -- (5.138,0.484) -- (5.152,0.475) -- (5.165,0.465) -- (5.178,0.455) -- (5.191,0.446) -- (5.204,0.436) -- (5.218,0.427) -- (5.231,0.418) -- (5.244,0.409) -- (5.257,0.399) -- (5.270,0.390) -- (5.284,0.381) -- (5.297,0.373) -- (5.310,0.364) -- (5.323,0.355) -- (5.336,0.346) -- (5.350,0.338) -- (5.363,0.329) -- (5.376,0.321) -- (5.389,0.312) -- (5.402,0.304) -- (5.416,0.296) -- (5.429,0.288) -- (5.442,0.280) -- (5.455,0.272) -- (5.468,0.264) -- (5.482,0.256) -- (5.495,0.249) -- (5.508,0.241) -- (5.521,0.234) -- (5.534,0.226) -- (5.548,0.219) -- (5.561,0.212) -- (5.574,0.205) -- (5.587,0.198) -- (5.600,0.191) -- (5.614,0.184) -- (5.627,0.177) -- (5.640,0.171) -- (5.653,0.164) -- (5.666,0.158) -- (5.680,0.152) -- (5.693,0.146) -- (5.706,0.140) -- (5.719,0.134) -- (5.732,0.128) -- (5.746,0.122) -- (5.759,0.117) -- (5.772,0.111) -- (5.785,0.106) -- (5.798,0.101) -- (5.812,0.096) -- (5.825,0.091) -- (5.838,0.086) -- (5.851,0.081) -- (5.864,0.077) -- (5.878,0.072) -- (5.891,0.068) -- (5.904,0.064) -- (5.917,0.059) -- (5.930,0.055) -- (5.944,0.052) -- (5.957,0.048) -- (5.970,0.044) -- (5.983,0.041) -- (5.996,0.038) -- (6.010,0.034) -- (6.023,0.031) -- (6.036,0.029) -- (6.049,0.026) -- (6.062,0.023) -- (6.076,0.021) -- (6.089,0.018) -- (6.102,0.016) -- (6.115,0.014) -- (6.128,0.012) -- (6.142,0.010) -- (6.155,0.009) -- (6.168,0.007) -- (6.181,0.006) -- (6.194,0.005) -- (6.208,0.004) -- (6.221,0.003) -- (6.234,0.002) -- (6.247,0.001) -- (6.260,0.001) -- (6.274,0.000) -- (6.287,0.000) -- (6.300,0.000);}
\newcommand{\figOneChartRight}[1]{\draw[#1] (6.300,0.000) -- (6.313,0.000) -- (6.326,0.000) -- (6.340,0.001) -- (6.353,0.001) -- (6.366,0.002) -- (6.379,0.003) -- (6.392,0.004) -- (6.406,0.005) -- (6.419,0.006) -- (6.432,0.007) -- (6.445,0.009) -- (6.458,0.010) -- (6.472,0.012) -- (6.485,0.014) -- (6.498,0.016) -- (6.511,0.018) -- (6.524,0.021) -- (6.538,0.023) -- (6.551,0.026) -- (6.564,0.029) -- (6.577,0.031) -- (6.590,0.034) -- (6.604,0.038) -- (6.617,0.041) -- (6.630,0.044) -- (6.643,0.048) -- (6.656,0.052) -- (6.670,0.055) -- (6.683,0.059) -- (6.696,0.064) -- (6.709,0.068) -- (6.722,0.072) -- (6.736,0.077) -- (6.749,0.081) -- (6.762,0.086) -- (6.775,0.091) -- (6.788,0.096) -- (6.802,0.101) -- (6.815,0.106) -- (6.828,0.111) -- (6.841,0.117) -- (6.854,0.122) -- (6.868,0.128) -- (6.881,0.134) -- (6.894,0.140) -- (6.907,0.146) -- (6.920,0.152) -- (6.934,0.158) -- (6.947,0.164) -- (6.960,0.171) -- (6.973,0.177) -- (6.986,0.184) -- (7.000,0.191) -- (7.013,0.198) -- (7.026,0.205) -- (7.039,0.212) -- (7.052,0.219) -- (7.066,0.226) -- (7.079,0.234) -- (7.092,0.241) -- (7.105,0.249) -- (7.118,0.256) -- (7.132,0.264) -- (7.145,0.272) -- (7.158,0.280) -- (7.171,0.288) -- (7.184,0.296) -- (7.198,0.304) -- (7.211,0.312) -- (7.224,0.321) -- (7.237,0.329) -- (7.250,0.338) -- (7.264,0.346) -- (7.277,0.355) -- (7.290,0.364) -- (7.303,0.373) -- (7.316,0.381) -- (7.330,0.390) -- (7.343,0.399) -- (7.356,0.409) -- (7.369,0.418) -- (7.382,0.427) -- (7.396,0.436) -- (7.409,0.446) -- (7.422,0.455) -- (7.435,0.465) -- (7.448,0.475) -- (7.462,0.484) -- (7.475,0.494) -- (7.488,0.504) -- (7.501,0.514) -- (7.514,0.524) -- (7.528,0.534) -- (7.541,0.544) -- (7.554,0.554) -- (7.567,0.564) -- (7.580,0.574) -- (7.594,0.585) -- (7.607,0.595) -- (7.620,0.605) -- (7.633,0.616) -- (7.646,0.626) -- (7.660,0.637) -- (7.673,0.648) -- (7.686,0.658) -- (7.699,0.669) -- (7.712,0.680) -- (7.726,0.691) -- (7.739,0.702) -- (7.752,0.713) -- (7.765,0.724) -- (7.778,0.735) -- (7.792,0.746) -- (7.805,0.757) -- (7.818,0.768) -- (7.831,0.779) -- (7.844,0.790) -- (7.858,0.802) -- (7.871,0.813) -- (7.884,0.824) -- (7.897,0.836) -- (7.910,0.847) -- (7.924,0.859) -- (7.937,0.870) -- (7.950,0.882) -- (7.963,0.894) -- (7.976,0.905) -- (7.990,0.917) -- (8.003,0.929) -- (8.016,0.940) -- (8.029,0.952) -- (8.042,0.964) -- (8.056,0.976) -- (8.069,0.988) -- (8.082,1.000) -- (8.095,1.012) -- (8.108,1.024) -- (8.122,1.036) -- (8.135,1.048) -- (8.148,1.060) -- (8.161,1.072) -- (8.174,1.084) -- (8.188,1.097) -- (8.201,1.109) -- (8.214,1.121) -- (8.227,1.133) -- (8.240,1.146) -- (8.254,1.158) -- (8.267,1.170) -- (8.280,1.183) -- (8.293,1.195) -- (8.306,1.208) -- (8.320,1.220) -- (8.333,1.233) -- (8.346,1.245) -- (8.359,1.258) -- (8.372,1.270) -- (8.386,1.283) -- (8.399,1.296) -- (8.412,1.308) -- (8.425,1.321) -- (8.438,1.334) -- (8.452,1.346) -- (8.465,1.359) -- (8.478,1.372) -- (8.491,1.385) -- (8.504,1.398) -- (8.518,1.410) -- (8.531,1.423) -- (8.544,1.436) -- (8.557,1.449) -- (8.570,1.462) -- (8.584,1.475) -- (8.597,1.488) -- (8.610,1.501) -- (8.623,1.514) -- (8.636,1.527) -- (8.650,1.540) -- (8.663,1.553) -- (8.676,1.566) -- (8.689,1.579) -- (8.702,1.592) -- (8.716,1.606) -- (8.729,1.619) -- (8.742,1.632) -- (8.755,1.645) -- (8.768,1.658) -- (8.782,1.672) -- (8.795,1.685) -- (8.808,1.698) -- (8.821,1.711) -- (8.834,1.725) -- (8.848,1.738) -- (8.861,1.751) -- (8.874,1.765) -- (8.887,1.778) -- (8.900,1.791) -- (8.914,1.805) -- (8.927,1.818) -- (8.940,1.832);}

\begin{tikzpicture}[>=Stealth,font=\small,
    mid/.style={postaction={decorate},decoration={markings,
      mark=at position #1 with {\arrow{Stealth}}}},
    midrev/.style={postaction={decorate},decoration={markings,
      mark=at position #1 with {\arrowreversed{Stealth}}}}]
  % (a) the Poincare disk
  \draw[gray!60] (0,0) circle (2.4);
  \figOneSpiral{gray,thin,mid=.12}
  \figOneGraphic{very thick,
    postaction={decorate},decoration={markings,
      mark=at position .3 with {\arrow{Stealth}},
      mark=at position .7 with {\arrow{Stealth}}}}
  \figOneSectionZero
  \figOneCentralSection
  \fill (0,0) circle (1.5pt) (0,2.4) circle (1.8pt);
  \node[above] at (0,2.4) {saddle--node};
  \node[right,xshift=4pt,yshift=-14pt] at \figOneVertex {$\ \Sigma_0$};
  \node[right] at (0,2.02) {$\Sigma_c$};
  \node[right] at (2.55,-.8) {$\Gamma(C_0)$};
  \node at (0,-2.85) {(a)};
  % (b) the projective chart near the saddle--node
  \draw[->,gray] (3.4,0)--(9.25,0) node[right] {$v$};
  \draw[->,gray] (6.3,0)--(6.3,2.75) node[above] {$z$};
  \figOneChartOrbits{gray,thin,mid=.55}
  \draw[very thick,mid=.55] (3.66,0)--(6.3,0);
  \draw[very thick,mid=.55] (8.94,0)--(6.3,0);
  \figOneChartLeft{very thick,midrev=.5}
  \figOneChartRight{very thick,midrev=.5}
  \draw[dashed] (6.3,0)--(6.3,2.4) node[above right] {$\eta=0$};
  \fill (6.3,0) circle (1.8pt);
  \node[below] at (6.3,-.05) {$z=0$, \ $v'=-v^3$};
  \node[above] at (3.95,1.9) {leaving};
  \node[above] at (8.65,1.9) {entering};
  \node at (6.3,-.85) {(b)};
\end{tikzpicture}
\caption{The organizing field of \eqref{eq:quadratic-family} for
\(C_0=10\), computed. (a) The Poincar\'e disk, through
\((x,y)\mapsto(x,y)/\sqrt{9+x^2+y^2}\): the graphic \(\Gamma(C_0)\), the
unstable focus at the origin with an orbit spiralling to
\(\Gamma(C_0)\), the section \(\Sigma_0\), and the central section
\(\Sigma_c\) below the nilpotent saddle--node at infinity. (b) The chart
\eqref{eq:projective-coordinates-new}--\eqref{eq:projective-family-new}
near the saddle--node, with the vertical scale enlarged: the line at
infinity \(z=0\), the branches of \(\Gamma(C_0)\) entering the
saddle--node from \(v>0\) and leaving it into \(v<0\), the central
section \(\eta=0\), and orbits in the three sectors of \(z\ge0\).}
\label{fig:organizing-graphic-new}
\end{figure}

\subsection{A fixed neighborhood of the whole graphic}

The Hausdorff condition must force passage through the sections as well
as proximity to the graphic. Removing a neighborhood of the finite
repelling equilibrium gives an isolating block and sections valid for
every nearby parameter.

\begin{proposition}[Isolation and whole-graphic capture]
\label{prop:isolation-capture-new}
There exist a compact set $B_{\mathrm{iso}}\subset\mathbb S^2$, a compact annular
neighborhood $B_{\mathrm{cap}}\Subset\operatorname{int}B_{\mathrm{iso}}$ of $\GammaI$,
a parameter neighborhood $V_{\mathrm{cap}}$, a transverse affine section $\Sigma_0$,
and $\rho_{\mathrm{cap}}>0$ with the following properties.

\begin{enumerate}[label=\textup{(\roman*)}]
\item The maximal invariant subset of $B_{\mathrm{iso}}$ for each field with parameter in
  $V_{\mathrm{cap}}$ is contained in $\operatorname{int}B_{\mathrm{iso}}$.
\item If a periodic orbit $\gamma$ satisfies
  $\distH(\overline\gamma,\GammaI)<\rho_{\mathrm{cap}}$, then
  $\overline\gamma\subset B_{\mathrm{cap}}$ and $\gamma$ crosses
  $\Sigma_0$ and the entrance and exit faces of a fixed projective box
  exactly once per minimal period, in the order of $\GammaI$.
\end{enumerate}
\end{proposition}

A periodic orbit as in (ii) is called \emph{captured}. The reverse
Hausdorff inclusion forces section crossings; in a thin tubular annulus,
their common orientation and the intersection number of an embedded
circle make each crossing unique. Appendix~\ref{app:graphic-capture-details}
gives the neighborhood construction and proof.

\begin{remark}
\label{rem:cyclicity-formulation}
The bound in \cref{thm:direct-finite-cyclicity} also applies after
restricting the orbits to a compact set: for each \(p\in V\) and each
compact \(B\subset\mathbb S^2\), there are at most three limit cycles
\(\gamma\) with \(\overline\gamma\subset B\) and
\(\distH(\overline\gamma,\GammaI)<\varepsilon\).
One such set is \(B_{\mathrm{iso}}\) of
\cref{prop:isolation-capture-new}, whose maximal invariant subset lies
in its interior.
\end{remark}

\subsection{The normal-rate coordinate}

The negative logarithmic rate of \(z\) identifies the central section
and the two directions of passage. Following \cite[Lemma~2.1]{HK2026},
we take this rate, \(\eta=-\dot z/z\), as a coordinate. In the projective time of
\eqref{eq:projective-family-new}, it is
\begin{equation}
\label{eq:normal-rate-coordinate-new}
 \eta=\mu _3+v+v^2+Cvz.
\end{equation}
Since $\partial_v\eta=1+2v+Cz$, this is an analytic coordinate together with
$z$ on a fixed projective box and a sufficiently small parameter
neighborhood.

\begin{lemma}[Normal form]
\label{lem:physical-normal-form-new}
In the coordinates $(\eta,z)$, the projective field has the exact form
\begin{equation}
\label{eq:physical-normal-form-new}
 \dot\eta=q(\eta,p)-zB(\eta,z,p),
 \qquad \dot z=-\eta z,
\end{equation}
where $p=(A,C,\mu _1,\mu _2,\mu _3)$, $B$ and $q$ are analytic, and
\begin{equation}
\label{eq:physical-normal-form-signs-new}
 B>0,\qquad \partial_z\{zB(0,z,p)\}>0
\end{equation}
on a fixed smaller box.  On $z=0$,
\begin{equation}
\label{eq:boundary-drift-conjugacy-new}
 q(\mu _3+v+v^2,p)=(1+2v)P(v).
\end{equation}
At the organizing parameter,
\begin{equation}
\label{eq:organizing-normal-coefficients-new}
 B(0,0)=1,\qquad B_\eta(0,0)=1,\qquad
 q(\eta)=-\eta^3+O(\eta^4).
\end{equation}
Consequently, after one more reduction of the box and parameter
neighborhood, there are positive constants $b_0,b_1,b_2$ such that
\begin{equation}
\label{eq:normal-sign-margins-new}
 B\ge b_0,\qquad B_\eta\ge b_1,\qquad
 3R_q+\eta\partial_\eta R_q\le-b_2,
\end{equation}
where
\begin{equation}
\label{eq:q-expansion-new}
 q(\eta,p)=q_0+q_1\eta+q_2\eta^2+\eta^3R_q(\eta,p).
\end{equation}
\end{lemma}

\begin{proof}
The equation for \(z\) in \eqref{eq:physical-normal-form-new} follows directly
from \eqref{eq:projective-family-new} and
\eqref{eq:normal-rate-coordinate-new}. Differentiating $\eta$ along the
projective vector field \eqref{eq:projective-family-new} gives an analytic first
component.  Its restriction to $z=0$ is
$(1+2v)P(v)$, which proves \eqref{eq:boundary-drift-conjugacy-new}.  The
remainder vanishes on $z=0$ and therefore factors as $-zB$.  Direct
evaluation at the organizing parameter gives
\eqref{eq:organizing-normal-coefficients-new}.  The inequalities follow by
continuity on a compact smaller box.
\end{proof}

For fixed $(A,C)$, the coefficients in \eqref{eq:q-expansion-new} are local
coordinates on $(\mu _1,\mu _2,\mu _3)$, as shown in
Appendix~\ref{app:coordinates}.

\subsection{Regular transitions on fixed flow boxes}

\begin{lemma}[Regular transitions]
\label{lem:regular-transitions}
Fix a finite order \(k\). Let \(\mathcal X_p\) be a \(C^{k+1}\) family
with parameters in a neighborhood of a compact subset of a
finite-dimensional parameter space. Consider transitions between
transverse sections, defined on neighborhoods of the closed section
windows under consideration. Suppose that their orbit segments remain
in a common compact regular region and have flight times bounded by
a fixed \(t_{\max}\). Suppose also that the section parameterizations,
their inverse charts, and their defining functions have uniform
\(C^{k+1}\) bounds, and that the input and output transverse fluxes
have absolute value at least a fixed \(c>0\).
Then the transition maps and their local inverses have uniform \(C^k\)
bounds on the section windows.
\end{lemma}

\begin{proof}
Smooth flow dependence \cite[Theorem~2.11]{Teschl2012}, on finitely
many flow neighborhoods, gives uniform derivatives up to time \(t_{\max}\).
The implicit function theorem for the hitting equation
\(g_p(\varphi_p^t(x))=0\), whose \(t\)-derivative is the nonzero
transverse flux, gives the transition bounds; reversed time gives
the inverse bounds.
\end{proof}


\subsection{The exterior transition}

Fix a small level $z=z_*$.  At the organizing parameter the two branches of
the graphic meet this level in two points with opposite signs of $\eta$.
Let $W_+$ and $W_-$ be compact transverse windows around the positive and
negative endpoints.  The finite part of the graphic defines, for nearby
parameters, a regular first-hit map
\begin{equation}
\label{eq:exterior-map-new}
 \Pi_p:W_-\longrightarrow W_+.
\end{equation}
On both windows we use the auxiliary energy coordinate
\begin{equation}
\label{eq:physical-energy-new}
 I=z_*-\frac{\eta^2}{2}.
\end{equation}
Choose a bounded normal-rate graph box which is downward closed in
the energy variable \(U=\eta^2/2\), for example
\(0\le U\le U_{\max}\), \(0<z\le z_*\), with the organizing
branches away from its upper radius face. It contains the two strips
of \(B_{\mathrm{cap}}\) along the organizing branches with
\(0<z\le z_*\). Choose the organizing endpoints in the interiors
of \(W_\pm\). After reducing the capture data of
\cref{prop:isolation-capture-new} and the parameter neighborhood of
\eqref{eq:exterior-map-new}, every captured periodic orbit stays inside
the graph box, away from its upper radius face, during its central passage
and meets the interiors of both windows. This follows by narrowing
the tubular neighborhood at the two regular section crossings.

\begin{lemma}[Exterior multiplier]
\label{lem:exterior-multiplier-new}
Put
\[
 r_*=\sqrt{2z_*+C_0z_*^2},\qquad a_*=1+C_0z_*,
 \qquad R_*=\frac{r_*}{z_*},
\]
and take $z_*$ small enough that $a_*>2r_*$.  Define the finite-section
transverse multiplier
\[
 m_{\mathrm{aff}}(R_*)=
 \exp\!\left(\int_{-R_*}^{R_*}
       \frac{4x}{x^2+2x+C_0}\,\dd x\right).
\]
Let $I_-^0$ and $I_+^0$ be the two endpoint states of the organizing orbit
in the coordinate \eqref{eq:physical-energy-new}.  The derivative of the
organizing exterior transition at that orbit is
\begin{equation}
\label{eq:exterior-multiplier-new}
 (\Pi_0)'(I_-^0)=
 m_{\mathrm{aff}}(R_*)\frac{(a_*+r_*)(a_*+2r_*)}
                    {(a_*-r_*)(a_*-2r_*)},
 \qquad
 \Lambda_0=\exp\!\left(-\frac{4\pi}{\sqrt{C_0-1}}\right).
\end{equation}
In particular, $(\Pi_0)'(I_-^0)=\Lambda_0\{1+O(\sqrt{z_*})\}$ as
$z_*\downarrow0$.
Let $\Lambda_0<\Lambda_+<1$ and fix $z_*$ sufficiently small. There are
enlarged compact transverse windows \(\widetilde W_\pm\) at \(z=z_*\)
around the organizing endpoints, in whose interiors the windows \(W_\pm\)
are chosen, and \(\Lambda_->0\) such that, for nearby parameters,
\(\Pi_p\) extends to a neighborhood of \(\widetilde W_-\) with values in
the positive section, has uniformly bounded derivatives of every fixed
finite order there, and satisfies
\begin{equation}
\label{eq:exterior-uniform-contraction-new}
 0<\Lambda_-\le\Pi_p'(I)\le\Lambda_+<1\qquad(I\in\widetilde W_-).
\end{equation}
In particular, endpoints uniformly \(o(1)\)-close to \(W_\pm\) lie in
\(\widetilde W_\pm\) for sufficiently small parameters.
\end{lemma}

\begin{proof}
Along the invariant parabola, use $x$ as time and let $\mathfrak k$ be a transverse
variation measured by $H_0$.  Equations
\eqref{eq:parabola-invariance-new}--\eqref{eq:parabola-speed-new} give
\[
 \frac{d}{dx}\log\mathfrak k
 =\frac{4x}{x^2+2x+C_0}.
\]
The finite-section multiplier in this affine transverse coordinate is
$m_{\mathrm{aff}}(R_*)$.  Since
\[
 \PV\int_{-\infty}^{\infty}
 \frac{4x}{x^2+2x+C_0}\,\dd x
 =-\frac{4\pi}{\sqrt{C_0-1}},
\]
and direct integration on the symmetric finite interval gives
\[
 m_{\mathrm{aff}}(R_*)
 =\Lambda_0\{1+O(R_*^{-1})\},
\]
we have $m_{\mathrm{aff}}(R_*)=\Lambda_0\{1+O(\sqrt{z_*})\}$.
We now convert the affine multiplier to the endpoint coordinate \(I\).
At fixed $z=z_*$ the organizing parabola meets the section at $v=\pm r_*$.
Direct differentiation of the
energy coordinate with respect to the affine transverse coordinate $H_0$
gives the two section factors
\[
 z_*^2(a_*+r_*)(a_*+2r_*),\qquad
 z_*^2(a_*-r_*)(a_*-2r_*),
\]
respectively.  Their ratio, multiplied by $m_{\mathrm{aff}}(R_*)$, yields
\eqref{eq:exterior-multiplier-new}; the ratio tends to one as
$z_*\downarrow0$.
Since $\Lambda_0<\Lambda_+$, fix $z_*$ with this derivative below
\(\Lambda_+\). Smooth dependence then gives the asserted extension and
\eqref{eq:exterior-uniform-contraction-new} on sufficiently small
windows \(\widetilde W_\pm\) around the organizing endpoints, for nearby
parameters. The nearby
exterior arcs have uniformly bounded flight times through regular flow
boxes and uniformly transverse endpoints, so
\cref{lem:regular-transitions} gives the finite-order derivative bounds
on \(\widetilde W_-\). The last assertion holds because \(W_\pm\) lie in
the interiors of \(\widetilde W_\pm\).
\end{proof}

\subsection{One central height}

The central section separates the incoming and outgoing halves of the
projective passage. Put
\begin{equation}
\label{eq:central-section-new}
 \Sigma_c(p)=\{\eta=0:0<z<z_*\}.
\end{equation}
The height $h=z|_{\Sigma_c(p)}$ gives a common analytic coordinate on
this parameter-dependent section in the original $(v,z)$ chart.
We call the point
$(\eta,z)=(0,h)$ of $\Sigma_c(p)$ a \emph{central point} and $h$ its
\emph{central height}.

\begin{lemma}[Unique central crossing]
\label{lem:unique-central-crossing-new}
Every periodic orbit captured by \cref{prop:isolation-capture-new} meets
$\Sigma_c(p)$ exactly once during its projective passage, and this
intersection is transverse.
\end{lemma}

\begin{proof}
At $\eta=0$, set
\[
 f_p(z)=q(0,p)-zB(0,z,p).
\]
By \eqref{eq:physical-normal-form-signs-new}, $f_p$ is strictly decreasing.
If $f_p(z)=0$, both components of
\eqref{eq:physical-normal-form-new} vanish, so the point is an equilibrium
and cannot lie on a periodic orbit.  Hence every central crossing of a
captured periodic orbit is transverse.

Suppose a crossing at $z_1$ goes from $\eta<0$ to $\eta>0$.  Until a
hypothetical second crossing, $z$ decreases.  The second crossing would
therefore occur at $z_2<z_1$ and would require $f_p(z_2)<0$, whereas strict
decrease gives $f_p(z_2)>f_p(z_1)>0$.  The opposite orientation is treated
by reversing all inequalities.  There is at most one crossing.  Existence
follows because the entrance and exit faces have opposite signs of $\eta$.
\end{proof}

For the directed passage considered here the crossing is from $\eta>0$ to
$\eta<0$; hence a passing height satisfies
\begin{equation}
\label{eq:central-passing-condition-new}
 f_p(h)<0.
\end{equation}

\subsection{The interval of central returns}

For a passing height $h$, parameterize both half-orbits by increasing
$z$: follow the positive half-orbit backwards and the negative half-orbit
forwards. Put
\begin{equation}
\label{eq:squared-graph-new}
 U_\sigma(z;h)=\frac{\eta_\sigma(z;h)^2}{2},\qquad
 \eta_\sigma=\sigma\sqrt{2U_\sigma},\qquad
 \sigma\in\{+,-\}.
\end{equation}
With $z$ increasing, both functions satisfy
\begin{equation}
\label{eq:physical-graph-equation-new}
 \frac{dU_\sigma}{dz}
 =\mathcal G_{\sigma,p}(U_\sigma,z)
 :=B(\sigma\sqrt{2U_\sigma},z,p)
   -\frac{q(\sigma\sqrt{2U_\sigma},p)}{z},
 \qquad U_\sigma(h;h)=0.
\end{equation}
The smooth field selects a unique branch through the initial
square-root singularity of this coordinate.

\begin{proposition}[Central displacement on one interval]
\label{prop:central-displacement-interval-new}
After reducing the capture data, there is, for every parameter
$p$, an interval $\cH_p\subset(0,z_*)$, possibly empty, such that both graphs
\eqref{eq:physical-graph-equation-new} reach $W_+$ and $W_-$.  The endpoint
functions
\begin{equation}
\label{eq:endpoint-functions-new}
 I_\sigma(h)=z_*-U_\sigma(z_*;h)
\end{equation}
are continuous and strictly increasing on $\cH_p$.  The central first-return
domain is an interval $\cD_p\subset\cH_p$, and its fixed points are exactly
the zeros on $\cH_p$ of
\begin{equation}
\label{eq:central-displacement-new}
 \Delta_p(h)=\Pi_p(I_-(h))-I_+(h).
\end{equation}
For each fixed \(p\), the maps \(I_\pm\) and \(\Delta_p\) are real
analytic on \(\operatorname{int}\cH_p\).
Every captured periodic orbit gives one zero of
\eqref{eq:central-displacement-new}, and every zero gives one periodic
orbit in the selected return neighborhood.  This correspondence preserves
isolation.
\end{proposition}

We call a central height \(h\in\cH_p\) \emph{admissible}, and use the
same adjective for its half-orbits and their endpoints in \(W_\pm\).

\begin{proof}
We first compare the graphs to determine which heights reach the
endpoint windows. We then determine the domain and fixed points of the
first return. Let
\[
 \mathcal Z_p=\{0<h<z_*:f_p(h)<0\}.
\]
This is an interval because $f_p$ is strictly decreasing.  For $h\in
\mathcal Z_p$ and $z>h$, condition
\eqref{eq:central-passing-condition-new} and the strict
decrease of $f_p$ give
\begin{equation}
\label{eq:graph-positive-at-axis-new}
 \mathcal G_{\sigma,p}(0,z)
 =\frac{zB(0,z,p)-q(0,p)}{z}>0.
\end{equation}
Thus a graph which has become positive cannot return to $U=0$ at a larger
height.

Let $h_1<h_2$ be passing heights and suppose first that the graph starting at
$h_1$ exists up to $h_2$.  Then
$U_\sigma(h_2;h_1)>0=U_\sigma(h_2;h_2)$.  Scalar uniqueness away from
$U=0$ gives
\begin{equation}
\label{eq:central-graph-order-new}
 U_\sigma(z;h_1)>U_\sigma(z;h_2)
\end{equation}
for every common later height. If the first graph stays in the fixed
downward-closed graph box until $z_*$, the second stays between it and
zero and also reaches $z_*$ by standard continuation. Consequently, for
each sign, the set $\mathcal Z_{\sigma,p}\subset\mathcal Z_p$ of heights
whose graph reaches $z_*$ inside that box is an interval which
contains, with each of its heights, every larger height of
$\mathcal Z_p$.

To show that \(zB-q>0\) along these graphs, note that initially
\(\mathcal G_{\sigma,p}>0\). At the first zero of
\(\mathcal G_{\sigma,p}(U(z),z)\), if one exists,
one has \(U>0\) and \(q=zB\), so
\begin{equation}
\label{eq:graph-slope-first-zero}
 \frac{d}{dz}\mathcal G_{\sigma,p}(U(z),z)
   =B_z+q/z^2=\frac{B+zB_z}{z}>0.
\end{equation}
This contradicts a first crossing from positive values. The strict
inequality \(B+zB_z>0\) holds on the fixed normal-rate box after
it is reduced.

To determine how the endpoints vary with the central height, differentiate
the moving initial condition at an interior height. This gives
\begin{equation}
\label{eq:moving-height-variation-new}
 \partial_hU_\sigma(z_*;h)
 =\frac{f_p(h)}h\exp\!\left(
  \int_h^{z_*}\partial_U\mathcal G_{\sigma,p}
       (U_\sigma(z;h),z)\,\dd z\right)<0,
\end{equation}
where $f_p(h)/h=q(0,p)/h-B(0,h,p)<0$.
The possible initial $O((z-h)^{-1/2})$ singularity is integrable; equivalently
one may differentiate the smooth flow before passing to $U$.
Hence $U_\sigma(z_*;h)$ is continuous and
strictly decreasing, and $I_\sigma(h)$ is continuous and strictly increasing,
on $\mathcal Z_{\sigma,p}$.

We can now impose the endpoint windows. Each $W_\sigma$ is a compact
interval in the energy coordinate, so
\[
 \mathcal H_{\sigma,p}
 =\{h\in\mathcal Z_{\sigma,p}:I_\sigma(h)\in W_\sigma\}
\]
is an interval, because it is the inverse image of an interval under a
strictly monotone map.  Define
\[
 \mathcal H_p=\mathcal H_{+,p}\cap\mathcal H_{-,p}.
\]
Every captured periodic orbit has its unique central height in
$\mathcal H_p$, by the choice of these strips and of the windows above.
In fact its height belongs to \(\operatorname{int}\mathcal H_p\).
Its central crossing is regular, its endpoint crossings are transverse
and lie in \(\operatorname{int}W_\pm\), and its two compact half-orbits
avoid the radius face. Smooth flow dependence therefore continues
both half-orbits for all central heights in an open neighborhood
of that height, with endpoints still in the two windows.

For fixed \(p\) and \(h_0\in\operatorname{int}\cH_p\), the
analytic field has a regular central point, since \(f_p(h_0)<0\).
The selected half-orbits reach \(z=z_*\) in finite time:
\(\ddot z(0)=-h_0f_p(h_0)>0\) near the central point, and
\(U_\sigma>0\) on the remaining compact graph. At the endpoint,
\(\dot z=-\eta z_*\ne0\). Analytic flow dependence and the analytic
implicit function theorem for these hitting equations therefore give
analytic \(I_\sigma\). The same argument applies to the exterior
transition, whose section coordinate has derivative \(-\eta\ne0\).
Thus \(\Delta_p\) is analytic on \(\operatorname{int}\cH_p\).

It remains to identify the return domain within the admissible interval.
We restrict the central return to heights whose return height lies in
$\cH_p$. The images $I_\sigma(\cH_p)$ are intervals, and this return is defined
precisely when $\Pi_p(I_-(h))$ belongs to $I_+(\cH_p)$. Since $\Pi_p\circ I_-$
is strictly increasing, its inverse image of this interval is an interval
$\cD_p$.  On $\cD_p$ the return is
\[
 \cP_p(h)=I_+^{-1}\bigl(\Pi_p(I_-(h))\bigr).
\]
The equality $\cP_p(h)=h$ is equivalent to
\eqref{eq:central-displacement-new}. Every zero on $\cH_p$ belongs to
$\cD_p$, since its common endpoint lies
in $I_+(\cH_p)$.  Joining the negative half-orbit, the exterior arc, and the
positive half-orbit proves the correspondence with periodic orbits.  The
unique central crossing proves injectivity, and a transverse flow box
preserves isolation in both directions.
\end{proof}

\begin{figure}[t]
\centering
\begin{tikzpicture}[>=Stealth,x=1.05cm,y=.9cm]
  \draw[->,gray!70] (-4.4,0)--(4.4,0) node[right] {$\eta$};
  \draw[->,gray!70] (0,-.2)--(0,4.5) node[above] {$z$};
  \draw[dashed,gray!70] (-3.6,3.8)--(3.6,3.8) node[right] {$z=z_*$};
  \draw[very thick,blue!55!black,
    postaction={decorate},decoration={markings,
      mark=at position .58 with {\arrowreversed{Stealth}}}]
    (0,1) .. controls (1.0,1.25) and (1.8,2.45) .. (2.8,3.8);
  \draw[very thick,blue!55!black,
    postaction={decorate},decoration={markings,
      mark=at position .58 with {\arrow{Stealth}}}]
    (0,1) .. controls (-1.0,1.25) and (-1.8,2.45) .. (-2.8,3.8);
  \draw[very thick,red!65!black,->]
    (-2.8,4.02) to[bend left=28] node[above] {$\Pi_p$} (2.8,4.02);
  \fill (0,1) circle (2pt) node[below right] {$h$};
  \fill (-2.8,3.8) circle (2pt) node[below left] {$I_-(h)$};
  \fill (2.8,3.8) circle (2pt) node[below right] {$I_+(h)$};
  \node[above left] at (-2.8,3.8) {$W_-$};
  \node[above right] at (2.8,3.8) {$W_+$};
  \node[right] at (.08,2.55) {$\Sigma_c(p)$};
\end{tikzpicture}
\caption{The central return in normal-rate coordinates. The
blue curves are the two halves of the orbit through $(0,h)$; the
red arrow is the regular exterior transition.  Positions and curvature are
schematic.  The closing condition is
$\Pi_p(I_-(h))=I_+(h)$.}
\label{fig:central-return-new}
\end{figure}

Returns based at other transverse sections are conjugate to
\(\cP_p\); they have the same number of fixed points, and the
conjugacy preserves isolation.

\subsection{Paired variations of the central passage}
\label{sec:paired-variation}

The two graphs start at the same central height. Pairing their variations
at equal \(z\)-levels cancels the quadratic boundary term from the
endpoint derivative ratio. We use the remaining terms to prove
contraction when \(q_1\le0\), and then at sufficiently large central
heights when \(q_1>0\).

Reduce the capture data of \cref{prop:isolation-capture-new} so that
\begin{equation}
\label{eq:small-central-height-new}
 0<h<h_0<\frac{z_*}{2}.
\end{equation}
This is possible because the central crossing of the organizing graphic is
$(\eta,z)=(0,0)$.

Let $L_p$ denote the local map from the positive endpoint window to the
negative endpoint window obtained by following the central passage.
Expressing both endpoints in terms of the same central height gives,
in the coordinate $I=z_*-\eta^2/2$,
\begin{equation}
\label{eq:local-map-derivative-new}
 L_p'(I_+(h))=\frac{I_-'(h)}{I_+'(h)}>0.
\end{equation}

\begin{lemma}[Paired-variation identity]
\label{lem:paired-variation-identity-new}
Let $U_\pm$ be the graph solutions in
\eqref{eq:physical-graph-equation-new}.  Define
\begin{equation}
\label{eq:J-sigma-new}
 J_\sigma(z)=B_\eta(\eta_\sigma,z,p)-\frac{q_1}{z}
 -\frac{\eta_\sigma^2}{z}
  \bigl(3R_q(\eta_\sigma,p)
       +\eta_\sigma\partial_\eta R_q(\eta_\sigma,p)\bigr).
\end{equation}
Then
\begin{equation}
\label{eq:paired-log-derivative-new}
 \log L_p'(I_+(h))
 =-\int_h^{z_*}
 \left(\frac{J_-(z)}{\sqrt{2U_-(z;h)}}
       +\frac{J_+(z)}{\sqrt{2U_+(z;h)}}\right)\dd z.
\end{equation}
The integral is improper only at $z=h$ and is convergent.
\end{lemma}

\begin{proof}
Differentiate the right-hand side of
\eqref{eq:physical-graph-equation-new} with respect to $U$.  Since
$\eta_\sigma=\sigma\sqrt{2U}$, direct differentiation of
\eqref{eq:q-expansion-new} gives
\begin{equation}
\label{eq:F-U-derivative-new}
 \partial_U\mathcal G_{\sigma,p}
 =-\frac{2q_2}{z}
 +\frac{\sigma}{\sqrt{2U}}\left\{
 B_\eta-\frac{q_1}{z}
 -\frac{\eta^2}{z}(3R_q+\eta R_{q,\eta})\right\}.
\end{equation}
We compare the endpoint variations using
\eqref{eq:moving-height-variation-new}. Their initial factor $f_p(h)/h$
is common to both signs and cancels in the quotient. The term $-2q_2/z$
in \eqref{eq:F-U-derivative-new} cancels at each common \(z\)-level.
Subtracting the two remaining integrands gives
\eqref{eq:paired-log-derivative-new}. As in
\eqref{eq:moving-height-variation-new}, the singularity at \(z=h\)
is integrable.
\end{proof}

\begin{proposition}[Signed contraction]
\label{prop:signed-contraction-new}
If $q_1\le0$, then
\begin{equation}
\label{eq:signed-local-contraction-new}
 0<L_p'<1
\end{equation}
on the whole central return interval. The return is strictly
contracting in the endpoint coordinate $I$, and $\Delta_p$ is strictly
decreasing in $h$. In particular, $\Delta_p$ has at most one zero.
\end{proposition}

\begin{proof}
The sign conditions \eqref{eq:normal-sign-margins-new} and $q_1\le0$ give
$J_\sigma\ge b_1$ along both halves. The right-hand side of
\eqref{eq:paired-log-derivative-new} is therefore strictly negative,
which proves \eqref{eq:signed-local-contraction-new}. Composing with
the exterior transition in the endpoint coordinate and using
\eqref{eq:exterior-uniform-contraction-new} gives
$0<(\Pi_p\circ L_p)'<1$. In terms of the central height, this is
\[
 \Delta_p'(h)=I_+'(h)
   \left\{\Pi_p'(I_-(h))\frac{I_-'(h)}{I_+'(h)}-1\right\}<0.
\]
Thus the displacement has at most one zero.
\end{proof}

For $q_1>0$, the term $-q_1/z$ in \eqref{eq:J-sigma-new} may reverse
the sign of $J_\sigma$ near the central crossing. The next lemma controls
this contribution at central heights that are large relative to the
boundary coefficients.

\begin{lemma}[Large central heights]
\label{lem:relative-height-paired-new}
After fixing $z_*$ sufficiently small, there are constants
$\theta_0,c_*>0$ such that the following holds.  Suppose $q_1>0$,
$0<h<h_0$, and
\begin{equation}
\label{eq:relative-height-paired-condition-new}
 \frac{|q_0|}{h}+\frac{q_1}{\sqrt h}+|q_2|\le\theta_0.
\end{equation}
Then both halves through $(0,h)$ reach the endpoint windows and
\begin{equation}
\label{eq:relative-height-local-bound-new}
 0<L_p'\le\exp(-c_*\sqrt{z_*})<1.
\end{equation}
The return is strictly contracting in $I$, and $\Delta_p'(h)<0$ if
$h\in\mathcal H_p$.
\end{lemma}

\begin{proof}
We first continue both graphs in a cone, then place their endpoints in
the windows, and finally estimate the derivative ratio. A first-contact
comparison in
\eqref{eq:physical-graph-equation-new}, first in a provisional fixed box
and then in the resulting cone, gives constants $c,C>0$ such that
\begin{equation}
\label{eq:large-height-cone-new}
 c(z-h)\le U_\sigma(z;h)\le C(z-h),\qquad
 \mathcal G_{\sigma,p}(U_\sigma,z)\ge c,
 \qquad h\le z\le z_*.
\end{equation}
Indeed, in this cone the constant, linear, quadratic, and cubic parts of
$q/z$ are bounded, respectively, by
\[
 \frac{|q_0|}{h},\qquad
 C\frac{q_1}{\sqrt h},\qquad
 C|q_2|,\qquad C\sqrt{z_*}.
\]
Choose $z_*$ and then $\theta_0$ so that these quantities fit the strict
cone margins. This proves continuation to $z=z_*$ and
\eqref{eq:large-height-cone-new}.

To place the endpoints in the prescribed windows, choose a fixed
intermediate height \(\zeta\ll z_*\), then reduce \(h_0\) so that
\(h_0<\zeta/2\). The cone gives \(U_\sigma(\zeta)=O(\zeta)\),
and the organizing boundary graph has the same bound. On
\([\zeta,z_*]\), continuous dependence and the first-variation
estimate show that the endpoint difference between these graphs is
\(O(\zeta)+o_p(1)\). The organizing variational coefficient on its
positive cone is \(O(z^{-1/2})\), whose integral is uniformly bounded;
for fixed \(\zeta\), the perturbed coefficients converge uniformly on
this interval. Choose \(\zeta\) small relative to the fixed window
widths, and then reduce the parameter neighborhood.

With continuation and the endpoint windows established, we return to
the derivative ratio. The sign margins \eqref{eq:normal-sign-margins-new} give
$J_\sigma\ge b_1-q_1/z$.
In \eqref{eq:paired-log-derivative-new}, use the upper cone bound for
the negative contribution and the lower cone bound for the positive
contribution. Since
\[
 \int_h^{z_*}\frac{\dd z}{\sqrt{z-h}}=2\sqrt{z_*-h},
 \qquad
 \int_h^{z_*}\frac{\dd z}{z\sqrt{z-h}}
 \le\frac{\pi}{\sqrt h},
\]
we obtain
\begin{equation}
\label{eq:relative-height-log-bound-new}
 \log L_p'
 \le-c\sqrt{z_*-h}+C\frac{q_1}{\sqrt h}
 \le-c'\sqrt{z_*}+C\theta_0,
\end{equation}
where $h<h_0<z_*/2$. Decreasing $\theta_0$ absorbs the positive term and proves
\eqref{eq:relative-height-local-bound-new}.  Composition with
\eqref{eq:exterior-uniform-contraction-new} proves the last assertion.
\end{proof}

\section{Logarithmic slopes and zero-counting tools}
\label{sec:standard-tools}

To pass from local derivative estimates to a bound on the whole height
interval, we must retain the order and direction of zero crossings.
We first give the scalar counting tools for this step. The coefficient
scales then identify which passage estimates supply their hypotheses.
\subsection{Counting zeros of the return displacement}
\label{sec:direct-zero-counting}

We count distinct zeros of \(\Delta_p\) on an interval of central
heights \(h\). Multiplication by a nonvanishing factor or a regular
monotone change of height coordinate preserves their number. To relate
the critical points of the displacement to the endpoint slopes, put
\(m_\sigma=hI_{\sigma,h}>0\) and define
\begin{equation}
\label{eq:direct-completed-log-ratio}
 \mathcal L=\log\left[\Pi_p'(I_-)\frac{m_-}{m_+}\right],
 \qquad
 \Delta_p'(h)=I_+'(h)(\e^{\mathcal L}-1).
\end{equation}
We express the crossing directions in a logarithmic coordinate that
increases as the central height decreases:
\begin{equation}
\label{eq:direct-source-clock-convention}
 T=-c_T\log(h/h_{\rm ref}),\qquad
 \partial_T=-\frac{h}{c_T}\partial_h,\qquad
 m_\sigma=-c_T I_{\sigma,T},
\end{equation}
where \(c_T>0\) and \(h_{\rm ref}>0\) depend only on the
parameters and are held fixed when differentiating in \(h\). The
height-range covers are ordered by increasing \(T\), hence decreasing \(h\).
Changing these scales replaces \(T\) by \(c_1T+c_2\), with
\(c_1>0\) constant for fixed parameters. This preserves zero-crossing
directions and the sign of curvature at a critical point. We write
\(\mathcal L(T)\) for the same slope ratio in this coordinate.

By \eqref{eq:direct-completed-log-ratio}, a bound of \(n\) for the
zeros of \(\mathcal L\) on one interval of central heights gives at
most \(n+1\) zeros of \(\Delta_p\), by Rolle's theorem. This bound
is uniform in the additive constant of \(\Delta_p\). It remains to
combine estimates from adjacent height ranges without losing their
crossing directions. The next lemma also allows a zero at a strict
maximum.

\begin{lemma}[Ordered zero crossings]
\label{lem:direct-oriented-zero-count}
Let \(f\in C^2(J)\), where \(J\) is an interval. If every
critical point of \(f\) satisfies \(f''<0\), then \(f\) has
at most one critical point and at most two distinct zeros.
More generally, suppose \(J\) has an ordered cover by an initial,
middle, and final interval, with overlaps allowed and any of these
intervals possibly empty. Assume that
\begin{enumerate}[label=(\roman*)]
\item every zero in the initial interval has \(f'>0\);
\item every critical point in the middle interval has \(f''<0\);
\item every zero in the final interval has \(f'<0\).
\end{enumerate}
Then \(f\) has at most two distinct zeros on \(J\).
If every zero on \(J\) has the same strict derivative sign,
there is at most one. An interval on which \(f\) never vanishes
may be included in either end interval.
\end{lemma}

\begin{proof}
At a zero of \(f'\) with \(f''<0\), the derivative crosses
strictly from positive to negative. Two such crossings would require
an intervening zero of \(f'\) that is not a downward crossing.
Thus \(f\) is monotone or increases and then decreases. It has at
most two zeros; when there are two, an upward crossing is followed by
a downward crossing. A zero at the critical point is a strict maximum
and is the only zero. The same argument gives at most one zero when
\(f'\) has the same strict sign at every zero of \(f\).

For the ordered cover, choose cuts in the overlaps and assign each
shared endpoint to only one adjacent interval. In the initial interval
zeros cross upward; in the middle interval they occur in the order
just proved; in the final interval they cross downward. Two upward
crossings require an intervening downward crossing, and two downward
crossings require an intervening upward crossing. Either configuration
contradicts the permitted crossing directions and their order. A zero
at a middle maximum likewise excludes zeros in both end intervals.
Hence the only possible pair consists of an upward crossing followed
by a downward crossing.
\end{proof}

\begin{samepage}
\begin{corollary}[Ordered cover count]
\label{cor:ordered-cover-zero-count}
Let \(J\) be an interval of central heights on which \(\mathcal L\) is
\(C^2\), and let \(T\) be a logarithmic height coordinate as in
\eqref{eq:direct-source-clock-convention}. If \(J\) has an ordered cover,
in increasing \(T\), satisfying the hypotheses of
\cref{lem:direct-oriented-zero-count} for \(f=\mathcal L\), then
\(\Delta_p\) has at most three distinct zeros on \(J\). If every zero of
\(\mathcal L\) on \(J\) has the same strict crossing direction, then
\(\Delta_p\) has at most two.
\end{corollary}
\end{samepage}

\begin{proof}
By \cref{lem:direct-oriented-zero-count}, \(\mathcal L\) has at most two
zeros on \(J\), and at most one in the second case. Since
\(\Delta_T=-(m_+/c_T)(\e^{\mathcal L}-1)\) with \(m_+>0\), the critical
points of \(\Delta_p\) are the zeros of \(\mathcal L\). Rolle's theorem
gives the assertion.
\end{proof}

\begin{figure}[t]
\centering
\begin{tikzpicture}[x=1.05cm,y=.8cm,>=Stealth,font=\small]
  % logarithmic slope ratio
  \draw[->] (0,0)--(10.5,0) node[right] {$T$};
  \draw[->] (0,-1.25)--(0,1.75) node[above] {$\mathcal L$};
  \draw[thick] plot[smooth,tension=.7] coordinates
    {(0.2,-1) (1.5,-.6) (2.6,0) (3.7,.82) (5,1.18) (6.3,.82) (7.4,0)
     (8.5,-.6) (10,-.95)};
  \fill (2.6,0) circle (1.7pt) (7.4,0) circle (1.7pt) (5,1.18) circle (1.7pt);
  \node[above left] at (2.55,0) {upward};
  \node[above right] at (7.45,0) {downward};
  \node[above] at (5,1.2) {maximum};
  % ordered cover
  \draw[line width=2.4pt,gray!55] (0.2,-1.8)--(4.2,-1.8);
  \draw[line width=2.4pt,gray!55] (3.3,-2.3)--(6.9,-2.3);
  \draw[line width=2.4pt,gray!55] (6.1,-2.8)--(10,-2.8);
  \node[right] at (4.25,-1.8) {initial};
  \node[right] at (7.65,-2.3) {middle};
  \node[left] at (6.05,-2.8) {final};
  % displacement
  \draw[->] (0,-4.6)--(10.5,-4.6) node[right] {$T$};
  \draw[->] (0,-5.85)--(0,-3.4) node[above] {$\Delta_p$};
  \draw[thick,domain=0.3:9.8,samples=80] plot
    (\x,{-4.6+0.0814*((\x-5)*(\x-5)*(\x-5)/3-5.76*(\x-5))});
  \draw[dashed,gray] (2.6,0)--(2.6,-3.85);
  \draw[dashed,gray] (7.4,0)--(7.4,-5.35);
\end{tikzpicture}
\caption{Schematic of \cref{cor:ordered-cover-zero-count}, in increasing
\(T\), that is, decreasing central height: \(\mathcal L\) with an upward
zero, a strict maximum and a downward zero; the initial, middle and final
intervals of an ordered cover; and \(\Delta_p\), whose critical points
(dashed) are the zeros of \(\mathcal L\).}
\label{fig:ordered-cover}
\end{figure}


\Cref{fig:ordered-cover} shows the configuration allowed by
\cref{cor:ordered-cover-zero-count}. In the applications, thresholds
on the half-orbits define the pieces of the cover. The following
lemma uses the graph order \eqref{eq:central-graph-order-new} to
place these pieces in their order in \(T\).

\begin{lemma}[Order of central heights]
\label{lem:central-height-order}
Let \(h_1<h_2\) be central heights whose half-orbits of the same sign
\(\sigma\) are defined up to a common height \(z\). Then
\(|\eta_\sigma(z;h_1)|>|\eta_\sigma(z;h_2)|\), and at every common radius
the half-orbit through \(h_2\) has the larger height. Consequently, for
fixed \(z_e,R>0\), the central heights with \(h\ge z_e\) or
\(|\eta_+(z_e;h)|\le R\) form an interval that contains all larger
heights, that is, an initial interval in \(T\). For a fixed height
\(z_o\) and an increasing function \(g\), the central heights with
\(g(|\eta_+(z_o;h)|)\ge c_{\rm th}\) form a final interval in \(T\). The
intersection of such intervals with \(\cH_p\) is an interval.
\end{lemma}

\begin{proof}
The first assertion is \eqref{eq:central-graph-order-new}, since
\(|\eta_\sigma|=\sqrt{2U_\sigma}\). Each graph is strictly increasing in
\(z\): the first-zero argument \eqref{eq:graph-slope-first-zero}
gives \(zB-q>0\), using \(B+zB_z>0\) on the reduced normal-rate box.
Hence, at a common radius, the graph through the larger central
height has the larger height. The remaining assertions follow because
\(T=-c_T\log(h/h_{\rm ref})\) decreases with \(h\) and
\(\cH_p\) is an interval.
\end{proof}

\Cref{tab:ordered-covers} lists the ordered covers used in
\cref{sec:contact-positive-offsets,sec:recurring-ranges,sec:primary-regimes,sec:quadratic-regimes,sec:direct-negative-annulus-buffer,sec:direct-negative-root-regimes}.
The remaining regional statements reduce to
\cref{prop:weak-root-passage}.
\begin{table}[ht]
\centering
\small
\caption{Ordered covers of \(\cH_p\), in increasing \(T\). ``Up'' and ``down'' mean that every zero of \(\mathcal L\) in the
range crosses upward, respectively downward; ``max'' means that every
critical point of \(\mathcal L\) in the range is a strict maximum. The
bounds follow from \cref{cor:ordered-cover-zero-count} unless another
lemma is named.}
\label{tab:ordered-covers}
\begin{tabular}{@{}>{\raggedright\arraybackslash}p{\dimexpr.22\textwidth-8pt}
                >{\raggedright\arraybackslash}p{\dimexpr.66\textwidth-8pt}
                >{\raggedright\arraybackslash}p{\dimexpr.12\textwidth-8pt}@{}}
\hline
Proposition & Ranges and the property of \(\mathcal L\) & Zero bound for
  \(\Delta_p\)\\
\hline
\ref{prop:direct-bounded-contact}
  & large heights, \(\mathcal L<0\); compact heights, up, by
    \cref{lem:direct-contact-scattering-ratio,lem:direct-derivative-ratio-count};
    small heights, \(\mathcal L>0\)
  & 2\\
\ref{prop:direct-positive-offset}
  & heights above the saddle collar, \(\mathcal L<0\); the collar, up, by
    \cref{lem:one-logarithm-zero-count}
  & 2\\
\ref{prop:weak-root-passage}
  & early heights and the root interval, up
  & 2\\
\ref{prop:direct-primary-negative}
  & early heights, \(\mathcal L<0\); bounded radii, max, by
    \cref{lem:curvature-criterion}; large exits, \(\mathcal L<0\)
  & 3\\
\ref{prop:direct-positive-quadratic}, compact drift
  & early heights, up; the root interval, max, by
    \cref{lem:curvature-criterion}
  & 3\\
\ref{prop:direct-positive-quadratic}, small \(b\)
  & early heights and the root collar, up, by
    \cref{lem:one-logarithm-zero-count} on the collar
  & 2\\
\ref{prop:direct-buffered-negative}
  & early heights, up; annular exits, max, or \(\mathcal L<0\) for
    small \(b\); later exits, down
  & 3\\
\ref{prop:direct-fold-bounded}, \ref{prop:direct-negative-fold}
  & early heights, up; bounded outputs and escaped heights, max
  & 3\\
\ref{prop:direct-negative-separated}
  & early heights, up; the root interval, max
  & 3\\
\hline
\end{tabular}
\end{table}

Near a boundary root, taking the logarithm of a power-law derivative
reduces the slope comparison to a constant, one logarithm, and a
remainder. The next lemma uses the strict negative constant to control
the zeros even as the coefficient of the logarithm tends to zero.

\begin{samepage}
\begin{lemma}[One logarithm]
\label{lem:one-logarithm-zero-count}
For \(c_0,C_e,\vartheta>0\) there is \(x_c\in(0,1)\) with the following
property. Let \(G\in C^1((0,x_c))\) have the form
\[
 G(x)=c+\alpha\log x+e(x),\qquad c\le-c_0,\quad \alpha\in\R,\quad
 |e(x)|+|xe'(x)|\le C_e x^\vartheta .
\]
If \(\alpha\ge0\), then \(G<0\). In general \(xG'(x)<0\) at every zero of
\(G\), so \(G\) has at most one zero. If \(f\in C^1((0,x_c))\)
and \(f'\) has the sign of \(G\), then \(f\) has at most
two zeros.
\end{lemma}
\end{samepage}

\begin{proof}
Choose \(x_c\) so that \(C_e x^\vartheta<c_0/2\) and
\(C_e x^\vartheta|\log x|<c_0/4\) for \(0<x<x_c\). If \(\alpha\ge0\), then
\(G\le c+e<-c_0/2\). If \(\alpha<0\), every zero satisfies
\(|\alpha||\log x|=|c+e(x)|\ge c_0/2\), and hence
\(xG'(x)=\alpha+xe'(x)<-|\alpha|+|\alpha|/2<0\). All zeros are crossings
in the same direction, so there is at most one. Rolle's theorem gives the
bound for \(f\).
\end{proof}

\begin{remark}
\label{lem:direct-frozen-section-sign}
The auxiliary sections used in a local representation are held fixed
during differentiation. If each point of an interval \(J\) has a
neighborhood on which \(f\in C^1(J)\) is represented with auxiliary
sections fixed for that neighborhood, then a sign of \(f'\) proved in
these representations holds on all of \(J\).
Endpoint preparations formed from the same \(I_\pm\) and \(\Pi_p\) on
\(\widetilde W_\pm\) of \cref{lem:exterior-multiplier-new} represent the
same \(\mathcal L\) in \eqref{eq:direct-completed-log-ratio}, including
shared cuts.
\end{remark}

\subsection{Derivative-ratio comparisons on compact intervals}

For analytic functions, the classical Wronskian criterion
\cite[Lemma~2.3(b)]{GMV2011} identifies the extended complete
Chebyshev property by nonvanishing successive Wronskians.
For the three functions \(1,\phi,\psi\), the relevant identities are
\[
 W(1,\phi)=\phi',\qquad
 W(1,\phi,\psi)=(\phi')^2(\psi'/\phi')'.
\]
For the compact-interval comparisons below, we need this argument to
remain valid after adding a small \(C^2\) remainder. The next lemma
records the strict margins on the denominator and the derivative
of the ratio that make this possible.

\begin{lemma}[Stable derivative-ratio count]
\label{lem:direct-derivative-ratio-count}
On a fixed compact interval \(J\), let \(\phi,\psi\in C^2(J)\),
with \(\phi'\ge c_*>0\) and a uniform \(C^2\) bound on \(\phi\).
Let \(0<\lambda_0\le\lambda\le\lambda_1\). If
\[
 |(\psi'/\phi')'|\ge c_*
 \quad\hbox{on }J,
\]
then, for every constant \(D\) and every sufficiently small
\(C^2\) remainder \(R\), the function
\(g=D+\phi-\lambda\psi+R\) has at most two zeros on \(J\), and
\((g'/\phi')'\) has the sign of \(-(\psi'/\phi')'\) on \(J\).
If instead \(1-\lambda\psi'/\phi'\ge c_*\), then \(g'/\phi'>0\), and
\(g\) has at most one zero.
The remainder threshold depends only on the displayed uniform bounds.
\end{lemma}

\begin{proof}
Dividing the derivative by \(\phi'>0\) gives
\(g'/\phi'=1-\lambda\psi'/\phi'+R'/\phi'\).
The derivative of the remainder term is bounded by
\(C\|R\|_{C^2}\). In the first case, the continuous function
\((\psi'/\phi')'\) has one sign and
\(|\lambda(\psi'/\phi')'|\ge\lambda_0c_*\). Consequently,
\((g'/\phi')'=-\lambda(\psi'/\phi')'+O(\|R\|_{C^2})\) has the sign of
\(-(\psi'/\phi')'\) when the remainder is sufficiently small.
The quotient is strictly monotone, so Rolle's theorem gives at most
two zeros of \(g\). In the second case, the quotient remains positive,
which gives at most one zero.
\end{proof}


\subsection{Coefficient scales}
There are two covers in the argument. The ordered height cover above
combines estimates for one fixed parameter on subintervals of
\(\mathcal H_p\). The parameter cover selects a regional result
that already bounds the zeros on that whole interval, so its bound
is the maximum of the regional bounds. Signed contraction treats
\(q_1\le0\); the scales below distinguish the regions for \(q_1>0\).
We use the coefficient scales
\[
 H=\max\{|q_0|^{1/3},\sqrt{q_1},|q_2|\},\qquad
 \rho_1=\sqrt{q_1}.
\]
For \(q_0<0\) and \(q_2\ne0\), the secondary parameters are
\[
 s=|q_2|,\qquad
 r=\max\{\sqrt{|q_0|/s^3},q_1/s^2\}.
\]
The central-height unit reflects the balance between the boundary drift
\(q\) and \(zB\) in \eqref{eq:physical-normal-form-new}; the output-height
unit reflects the comparability of \(U=\eta^2/2\) and \(z\) in the common
cone. On the primary scale, \(q(HX)=O(H^3)\) for bounded \(X\), whereas
\(|\eta|=O(H)\) gives \(U=O(H^2)\), yielding the height units \(H^3\)
and \(H^2\), respectively.
\Cref{tab:direct-scale-guide} lists the coordinate normalizations;
the rescaled central height ranges over the corresponding image of
\(\mathcal H_p\). Near a weak root, whose radius tends to zero relative
to the current coefficient scale, we rescale again.

For \(q_1>0\), the contact strip \(|q_0|\le Lq_1^2\) is treated by
the scattering comparison. Positive offsets outside this strip are
treated by the interior saddle.

For negative offsets, we use the primary scale when \(|q_0|/H^3\)
has a fixed positive lower bound. At smaller ratios, the secondary
scale treats small \(r\), separately for the two signs of \(q_2\).
The remaining parameters are treated either on the linear scale
\(\rho_1\), by the weak-root comparison, or in the contact strip.
After the regional estimates, \cref{sec:direct-parameter-cover} fixes
\(L\) and the overlaps needed to make this classification exhaustive.

\begin{table}[!htbp]
\centering
\small
\caption{Coordinate scales for \(|\eta|\), \(h\) and \(z\), respectively.
  Output heights are fixed positive multiples of the listed scale.}
\label{tab:direct-scale-guide}
\begin{tabular}{@{}lccc@{}}
\hline
Region & Radius scale & Height scale & Output scale\\
\hline
Contact & \(q_1\) & \(q_1^2\) & \(q_1^2\)\\
Positive offset & \(\sqrt{q_0}\) & \(q_0\) & \(q_0\)\\
Primary negative drift & \(H\) & \(H^3\) & \(H^2\)\\
Secondary negative offset & \(sr\) & \(s^3r^2\) & \(s^2r^2\)\\
\hline
\end{tabular}
\end{table}

\Cref{tab:direct-cover-guide} in \cref{sec:direct-parameter-cover}
lists the parameter regions and the bound each provides on the whole
height interval \(\mathcal H_p\). Appendix~\ref{app:notation-terminology} collects the recurring notation
and passage terminology.

\section{Uniform passage estimates}
\label{sec:direct-passage-estimates}

Two degeneracies govern the endpoint comparison. At exponentially small
incoming heights, the quadratic boundary coefficient can produce an
unbounded power-law multiplier in each endpoint derivative. Near a
boundary root, the drift and the leading curvature both tend to zero,
so an unweighted error estimate cannot preserve the curvature sign.
We retain the common tail multiplier explicitly, remove the
height-dependent term by a boundary-flow coordinate, and fix the phase
of both projected orbits using the same central height. After a fixed
positive boundary-time advance, the phase error carries arbitrarily
high powers of the drift. These estimates give
the curvature criterion of \cref{lem:curvature-criterion}, including
the exterior contribution, and the one-sided slope comparison in the
weak-root range of \cref{cor:direct-weak-root-comparison}.

In all passage estimates below, derivatives are taken in the stated state
variables, with parameters and all parameter-only scales and actions
fixed. Bounds are uniform on the indicated parameter regions, and each
finite order \(k\) is fixed before choosing the small parameter threshold.

\subsection{Weighted tail estimates}

We first estimate the map from an incoming height \(t\) to the fixed
section \(z=z_*\), keeping its power-law multiplier explicit.
Write \(\mathcal G_\sigma(U,z)=B(\sigma\sqrt{2U},z)
-q(\sigma\sqrt{2U})/z\), as in
\eqref{eq:physical-graph-equation-new}. Fix a compact interval of
positive incoming energies \(\mathcal U=U(t)/t\). For a sufficiently small
fixed \(z_*\), the conditions
\begin{equation}
\label{eq:direct-tail-smallness}
 \frac{|q_0|}{t}+\frac{|q_1|}{\sqrt t}+|q_2|
 \quad\hbox{sufficiently small},\qquad 0<t<z_*,
\end{equation}
put all those solutions in a slightly enlarged invariant cone
\(c_-z\le U(z)\le c_+z\). Indeed, on this cone,
\[
 \frac{|q(\sigma\sqrt{2U})|}{z}
 \le C\left(\frac{|q_0|}{z}+\frac{|q_1|}{\sqrt z}
                  +|q_2|+\sqrt z\right).
\]
Choose \(c_-<b_0/4\) and \(c_+>4\sup B\), containing the incoming
energies in their interior. Reduction of \(z_*\) and
\eqref{eq:direct-tail-smallness} makes
\(b_0/2<\mathcal G_\sigma<2\sup B\). The vector field points into the
cone along both boundary lines.
This proves continuation to \(z_*\), also for nearby incoming
energies needed for differentiation.

To make the coefficient of \(\sqrt{z_*}\) independent of later scaled
height thresholds, fix a common cone using coefficient bounds on a
preliminary box in \((\eta,z)\). Enlarge it once to contain the
moving-entrance energy \(1/(2m_{\rm land})\) used in the adaptive and
annular landings, with \(m_{\rm land}\) fixed before the coefficient
and height thresholds. A scaled entrance with a larger energy-to-height
ratio is first continued to this cone. Constants for this finite
continuation may depend on the thresholds; the coefficient
\(C_{\rm cone}\) remains independent of them.

For a fixed positive scaled entrance outside the common cone, put
\(z=\ell^2Z,\ U=\ell^2\mathcal U\),
with fixed \(Z_0>0\) and \(\mathcal U(Z_0)=\mathcal U_0\) in a compact positive
interval, and set
\[
 \epsilon_{\rm in}=\ell+|q_0|/\ell^2+|q_1|/\ell+|q_2|.
\]
Write \(B_0=B(0,0,p)\). On any fixed positive-height interval the graph
equation is
\(\mathcal U_Z=B_0+O_{C^k}(\epsilon_{\rm in})\) on a fixed enlarged
positive energy interval. Integration and the finite variational
equations give
\begin{equation}
\label{eq:direct-finite-common-cone-entry}
 \mathcal U(Z_1;\mathcal U_0)=\mathcal U_0+B_0(Z_1-Z_0)
                +O_{C^k}(\epsilon_{\rm in}).
\end{equation}
Choose \(Z_1\) large after the entrance interval, and then
\(\epsilon_{\rm in}\) small; the ratio \(\mathcal U(Z_1)/Z_1\) lies
in the common cone. The constants in this finite continuation may
depend on \(Z_0,Z_1\) and the entrance interval. Its derivative with
respect to \(\mathcal U_0\) is \(1+O(\epsilon_{\rm in})\), and replacing the tail multiplier
based at \(\ell^2Z_0\) by that based at \(\ell^2Z_1\) introduces
only \((Z_1/Z_0)^{2q_2}=1+O(|q_2|)\), with \(Z_0\) and \(Z_1\) fixed.

\begin{lemma}[Weighted tail]
\label{lem:direct-weighted-tail}
For every fixed \(k\), set
\[
 \varpi_t(z)=(z/t)^{-2q_2},\qquad \varpi_t=\varpi_t(z_*),\qquad
 e(t)=\sqrt t+|q_1|/\sqrt t .
\]
If \(|q_2|<1/8\) and \eqref{eq:direct-tail-smallness} holds, then
the centered endpoint map satisfies
\begin{equation}
\label{eq:direct-centered-tail}
 \frac{U(z_*;t\mathcal U)-U(z_*;t\mathcal U_\circ)}{t\varpi_t}
 =a_{t,p}(\mathcal U-\mathcal U_\circ)+R_{t,p}(\mathcal U),\qquad
 c<a_{t,p}<C,\quad \|R_{t,p}\|_{C^k}\le C_ke(t).
\end{equation}
Its inverse in these centered coordinates has uniformly bounded
derivatives through order \(k\). If the incoming energies lie in the
common cone, then
\begin{equation}
\label{eq:direct-tail-first-slope}
 \left|\log a_{t,p}\right|
       \le C_{\rm cone}\left(\sqrt{z_*}+\frac{|q_1|}{\sqrt t}\right).
\end{equation}
Here \(C_{\rm cone}\) depends only on the common cone and the
uniform bounds for \(B\), \(R_q\) and their first derivatives
on the preliminary normal-rate box. The same constant applies
for all admissible \(t,z_*\) in that box, independently of the
later coefficient and height thresholds. For another fixed cone
the same bound holds with a constant depending on that cone.
For a moving entrance \(U(t)=ct\), with fixed
\(c>4\sup B\), let \(U_{\rm end}(t)\) be its energy at \(z_*\).
Then
\begin{equation}
\label{eq:direct-moving-tail}
 \begin{gathered}
 U_{\rm end}'(t)=\varpi_t a(t),\qquad
 c_1<a(t)<C_1,\\
 |(t\partial_t)^j a(t)|+|(t\partial_t)^j\log a(t)|
 \le C_j\left(e(t)+\frac{|q_0|}{t}\right),\quad 1\le j\le k .
 \end{gathered}
\end{equation}
Here the cone is enlarged to contain \(c\) before the smallness
constants are chosen.
\end{lemma}

\begin{proof}
We separate the quadratic term in the first energy derivative from the
coefficients that control the higher variations. On the cone,
\begin{align}
\label{eq:direct-tail-coefficients}
 \mathcal G_U&=-2q_2/z+R_1,\qquad
 |R_1|\le C(z^{-1/2}+|q_1|z^{-3/2}),\\
 |\mathcal G_{U^j}|&\le C_j(z^{1/2-j}+|q_1|z^{-j-1/2}),
 \qquad j\ge2.\nonumber
\end{align}
The constant drift disappears on differentiation, and the constant in
the first line depends only on the stated cone and coefficient bounds.
Hence the first variation is
\begin{equation}
\label{eq:direct-tail-variation}
 J(z;t)=\varpi_t(z)
     \exp\int_t^z R_1(U(\zeta),\zeta)\,d\zeta .
\end{equation}
Its quotient by \(\varpi_t(z)\) is positive and bounded; integrating the
bound for \(R_1\) in \eqref{eq:direct-tail-coefficients} bounds its
logarithm by the right side of \eqref{eq:direct-tail-first-slope}.

To estimate the higher variations, factor out the same multiplier. Put
\(V_j=t^{-1}\partial_{\mathcal U}^jU=\varpi_t(z)W_j\).
Here \(W_1(t)=1\) and \(W_j(t)=0\) for \(j\ge2\).
After cancellation of the quadratic linear term, their equations are
\[
 (W_j)_z=R_1W_j+
 \sum_{m=2}^j t^{m-1}\varpi_t(z)^{m-1}\mathcal G_{U^m}
 \sum_{\substack{i_1+\cdots+i_m=j\\i_\ell\ge1}}
       c_{i_1,\ldots,i_m}W_{i_1}\cdots W_{i_m}.
\]
Every variation in the inner sum has order less than \(j\), so it remains
to bound the coefficients of these products. The two required integrals obey
\begin{align}
\label{eq:direct-weighted-integrals}
 t^{m-1}\int_t^{z_*}z^{1/2-m}(z/t)^{2|q_2|(m-1)}\,dz
       &\le C_m\sqrt t,\\
 |q_1|t^{m-1}\int_t^{z_*}z^{-m-1/2}
                       (z/t)^{2|q_2|(m-1)}\,dz
       &\le C_m|q_1|/\sqrt t.\nonumber
\end{align}
After \(z=t\zeta\), the first integral has decay margin
\(m-3/2-2|q_2|(m-1)\ge1/4\); the second has a larger margin.
Variation of constants and induction give bounded \(W_1\) and
\(W_j=O(e(t))\) for \(j\ge2\). Choose \(a_{t,p}\) as the first
normalized derivative at \(\mathcal U_\circ\) and integrate the second
derivative bound. Constructing one extra derivative if necessary
proves \eqref{eq:direct-centered-tail}. Positivity of the first
derivative gives the inverse bounds.

For the moving entrance, differentiation includes the change of initial
height as well as that of initial energy. It gives
\[
 U_{\rm end}'(t)=J(z_*;t)\,f_c(t),\qquad
 f_c(t)=c-\mathcal G_\sigma(ct,t)>c/2.
\]
Thus \(a(t)=f_c(t)J(z_*;t)/\varpi_t\). To justify its derivatives,
write \(\mathcal E_t=t\partial_t\), holding \(z\) fixed inside the flow.
The differentiated equations and
\eqref{eq:direct-weighted-integrals} give
\begin{equation}
\label{eq:direct-moving-flow-jets}
 |\mathcal E_t^jU(z;t)|\le C_jt\varpi_t(z),\qquad j\ge1.
\end{equation}
For \(j=1\) this is \(\mathcal E_tU=tJf_c\). To compute higher moving
boundary data, set \(\mathcal B_j(t)=(\mathcal E_t^jU)(t;t)\). Then
\[
 \mathcal B_0=ct,\qquad
 \mathcal B_j=t\frac{d\mathcal B_{j-1}}{dt}
       -t(\partial_z\mathcal E_t^{j-1}U)(t;t).
\]
In the local variables \(z=t\zeta,\ U=tW\), near
\((\zeta,W)=(1,c)\), the equation has bounded finite jets in
\(\sqrt t,q_1/\sqrt t,q_0/t,q_2\). The positive-order Euler
derivatives of these quantities are bounded under
\eqref{eq:direct-tail-smallness}, so \(\mathcal B_j=O_j(t)\).
The same triangular induction then proves
\eqref{eq:direct-moving-flow-jets}.

These moving-flow bounds allow us to differentiate the logarithm in
\eqref{eq:direct-tail-variation}. The lower-boundary terms are
\(O(e(t))\). At order \(j\), the chain rule gives a finite sum
of terms with \(m\) flow variations, \(1\le m\le j\).
By \eqref{eq:direct-weighted-integrals} with \(m+1\) in place of
\(m\), each is bounded by \(C_je(t)\).
Repeated boundary differentiation has the same estimate.
Finally \(\mathcal E_t^jf_c=O(e(t)+|q_0|/t)\) for \(j\ge1\).
In particular its quadratic contribution \(2q_2c\) is constant
in \(t\). These estimates bound the Euler derivatives of \(a\); since
\(a\) has a positive lower bound, they also bound those of \(\log a\).
This proves \eqref{eq:direct-moving-tail}.
\end{proof}

For \(|q_2|<1/8\) and \(0<t\le z_*\), the tail amplitude satisfies
\begin{equation}
\label{eq:direct-tail-amplitude}
 t\varpi_t=z_*^{1/4}t^{3/4}(t/z_*)^{1/4+2q_2}\le z_*^{1/4}t^{3/4}.
\end{equation}

When the scaled entrance height also tends to zero, we first continue the
graph to a fixed scaled height where the weighted-tail estimate applies.
The following lemma supplies this continuation and the resulting endpoint
expansion. Its additional symbols and those of its proof are local.

\begin{lemma}[Continuation from a small scaled height]
\label{lem:direct-small-height-continuation}
Suppose \(0<\ell\le\delta\) and
\(|q_0|/(\delta\ell^2)+|q_1|/(\delta\ell)+|q_2|/\delta\le C\).
Put \(z=\ell^2Z\), \(U=\ell^2\mathcal U\) and
\(K_\sigma(x)=-q(\sigma\ell x)/(\delta\ell^2)\). Let a graph enter at
\(Z=\delta\zeta_0\) with \(\mathcal U=\mathcal U_0\in[0,\mathcal U_{\max}]\), where
\(0<c\le\zeta_0\le C\) and \(K_\sigma(\sqrt{2\mathcal U_0})\ge k_0>0\).
For every fixed \(Z_1>0\) and sufficiently small \(\delta\), it continues
with positive flux to
\[
 \mathcal U(Z_1)=\mathcal U_0+B_0Z_1+O(\delta(1+|\log\delta|)).
\]
A large fixed \(Z_1\), chosen after \(\mathcal U_{\max}\) and before \(\delta\),
puts \(\mathcal U(Z_1)/Z_1\) in the common cone.
If \(\mathcal U_0=x^2/2\), with \(x\) in a fixed compact positive interval and
\(\|\zeta_0(x)\|_{C^3}\le C\), set
\(t=\ell^2Z_1\), \(E=\ell^2\varpi_t\) and
\[
 b_\sigma=z_*-U(z_*;t,B_0t),\qquad
 A_\sigma=\varpi_t^{-1}
       \left.\partial_uU(z_*;t,u)\right|_{u=B_0t}.
\]
Then the endpoint satisfies
\begin{equation}
\label{eq:direct-small-height-endpoint}
 I_\sigma(x)=b_\sigma-\tfrac12EA_\sigma x^2+Ee_\sigma(x),\qquad
 c<A_\sigma<C,\quad
 \|e_\sigma\|_{C^3}\le C\delta(1+|\log\delta|).
\end{equation}
For two branches with the same \(\ell,Z_1\),
\(|\log(A_-/A_+)|\le2C_{\rm cone}\sqrt{z_*}+C\delta\).
The constants depend only on the displayed bounds, the compact radius
interval and \(Z_1\).
\end{lemma}

\begin{proof}
We first prove continuation, then retain the radius derivatives needed
for the endpoint expansion. The coefficient bounds and \(\ell\le\delta\)
bound all radius derivatives
of \(K_\sigma\) on compact intervals. Choose
\(\mathcal U_1>\mathcal U_{\max}+2Z_1\sup B\), then \(\zeta_1>\sup\zeta_0\) so large that
\(|K_\sigma|\le b_0\zeta_1/2\) for \(0\le \mathcal U\le \mathcal U_1\).
In \(Z=\delta\zeta\), the graph equation is
\[
 \mathcal U_\zeta=\delta\{B(\sigma\ell\sqrt{2\mathcal U},\delta\ell^2\zeta)
                       +K_\sigma(\sqrt{2\mathcal U})/\zeta\}.
\]
Up to \(\zeta_1\), the energy changes by \(O(\delta)\) and the radius by
\(O(\sqrt\delta)\), preserving \(K_\sigma\ge k_0/2\) and positive flux.
For \(Z\ge\delta\zeta_1\), \(|\delta K_\sigma/Z|\le b_0/2\) while
\(\mathcal U\le \mathcal U_1\), so \(b_0/2\le \mathcal U_Z\le2\sup B\), and this energy interval
persists to \(Z_1\). Integrating \(B-B_0=O(\ell)\) and
\(\delta K_\sigma/Z\) gives the displayed continuation estimate.
Since \(B_0\) is interior to the common cone, large \(Z_1\) gives entry.

For the radius-dependent entrance, the energies stay bounded away from
zero, so the preceding continuation can be differentiated in the entrance
radius. The flow in \(\zeta\) has uniformly bounded finite jets;
composition with \((\zeta_0(x),x^2/2)\) gives an \(O_{C^3}(\delta)\)
change up to \(\zeta_1\). Thereafter the coefficient error and its energy
derivatives are bounded by \(C_j(\delta+\delta/Z)\). Integration of the
finite variational equations gives the same continuation estimate in
\(C^3\). At \(t\), the smallness expression in
\eqref{eq:direct-tail-smallness} and \(e(t)\) are \(O(\delta)\).
Apply \eqref{eq:direct-centered-tail} with center
\(\mathcal U_\circ=B_0\) to the incoming energy \(\mathcal U(Z_1;x)/Z_1\).
The coefficient of this centered map is \(A_\sigma\), the normalized center derivative,
and \(t\varpi_t=Z_1E\); this proves
\eqref{eq:direct-small-height-endpoint}.
Finally \eqref{eq:direct-tail-first-slope}, with
\(|q_1|/\sqrt t=O(\delta)\), gives the slope-ratio bound.
\end{proof}

\subsection{A coordinate following the boundary flow}

The passage to a fixed scaled height combines boundary motion with a
height-dependent term. Dividing by the boundary drift would be singular
at its roots. We instead construct a coordinate that absorbs the
height-dependent term and follows the boundary flow, uniformly through
those roots. On a fixed positive radius interval, write
the scaled height equation as
\begin{equation}
\label{eq:direct-height-system}
 \frac{dX}{d\log Y}=\delta F(X)+YG(X,Y),\qquad \delta>0,
\end{equation}
where \(F\) has bounded derivatives through the required order and
\(G(X,Y)=B_0/X+O_{C^k}(\delta)\), with \(B_0\) in a compact
positive interval.

\begin{lemma}[Boundary-flow coordinate]
\label{lem:direct-boundary-flow}
On a smaller positive radius interval and \(0\le Y\le Y_0\),
for every fixed finite \(k\) and all sufficiently small \(\delta\),
there is a uniformly \(C^k\) coordinate
\[
 \mathcal W(X,Y)=X+Y\mathcal W_1(X,Y),\qquad
 \mathcal W_X>0,\quad \mathcal W_Y<0,
\]
which satisfies \(\mathcal W_{\log Y}=\delta F(\mathcal W)\) along
\eqref{eq:direct-height-system}. Its inverse \(X=\chi(w,Y)\)
has
\begin{equation}
\label{eq:direct-projection-energy}
 \frac{\chi(w,Y)^2}{2}
   =\frac{w^2}{2}+B_0Y+\delta YE_\chi(w,Y),\qquad
 \|E_\chi\|_{C^k}\le C_k,\quad \chi_Y>0.
\end{equation}
The construction is uniform through simple or multiple zeros of \(F\).
The height bound \(Y_0\) depends only on the bounds for \(B_0\) and on the
lower ends of the two radius intervals; the threshold for \(\delta\) and
the constants \(C_k\) may depend on the whole interval.
\end{lemma}

\begin{proof}
We construct the coordinate by following the full passage backwards and
then undoing the boundary motion. Extend the coefficients smoothly to a
larger interval and then to bounded functions, preserving positive \(G\). Let \(\phi^u\)
be the flow of \(\delta F\). Follow
\eqref{eq:direct-height-system} backward for logarithmic time \(u\),
writing its radius as \(X_b(u)\) and its height as \(Y\e^{-u}\).
Define
\[
 \mathcal W(X,Y)=\lim_{u\to\infty}\phi^u(X_b(u)).
\]
The derivative of the expression inside this limit is
\begin{equation}
\label{eq:direct-projection-integrand}
 -Y\e^{-u}(\phi^u)'(X_b(u))G(X_b(u),Y\e^{-u}).
\end{equation}
To differentiate the limit, we compare growth of the flow derivatives
with the height decay in this integrand. For each fixed order,
boundary-flow derivatives grow at most as
a polynomial in \(u\) times \(\e^{C_k\delta u}\).
The backward equation is
\[
 \partial_uX_b=-\delta F(X_b)-Y\e^{-u}G(X_b,Y\e^{-u}).
\]
Its state variations have linear coefficient
\(O(\delta)+O(Y\e^{-u})\). In each higher variational equation,
the remaining terms contain lower variations; coefficients from
\(F\) carry \(\delta\), and those from \(G\) carry \(\e^{-u}\).
Variation of constants therefore gives the same
polynomial-times-exponential bounds at every fixed order.
Consequently
\[
 \mathcal W=X-Y\int_0^\infty \e^{-u}
       (\phi^u)'(X_b(u))G(X_b(u),Y\e^{-u})\,du.
\]
Every prescribed derivative of this integral is bounded by an
integrable function
\(C_k(1+u)^{N_k}\e^{-(1-C_k\delta)u}\), after reducing
\(\delta\). This proves convergence and \(\mathcal W=X+Y\mathcal W_1\)
with uniform derivative bounds for \(\mathcal W_1\).
The flow identity gives \(\mathcal W_{\log Y}=\delta F(\mathcal W)\).

The signs and the inverse energy expansion follow by comparison with
\(\delta=0\). At that parameter, the energy identity is
\(\mathcal W_0^2=X^2-2B_0Y\). Choose \(Y_0\) so that this backward
graph remains inside the positive interval; since \(\mathcal W_0\le X\), this
constrains only the lower ends of the intervals. Write
\(G=B_0/X+\delta g\), and interpolate by
\(\lambda\delta F,\ B_0/X+\lambda\delta g\), \(0\le\lambda\le1\),
using compatible extensions. A \(\lambda\)-variation has
forcing of size \(\delta\) and introduces only additional
polynomials in \(u\). Differentiating the preceding integral
therefore gives
\[
 \mathcal W-\mathcal W_0=\delta YR(X,Y),\qquad \|R\|_{C^k}\le C_k.
\]
The construction gives two additional derivatives whenever needed.
Since \(\mathcal W_{0,X}>c>0\) and
\(\mathcal W_{0,Y}<-c<0\), the same signs and uniform inverse bounds
hold for \(\mathcal W\). If \(\chi_0=(w^2+2B_0Y)^{1/2}\), the identity
\[
 \mathcal W_0(\chi,Y)-\mathcal W_0(\chi_0,Y)=-\delta YR(\chi,Y)
\]
and the positive lower bound on \(\mathcal W_{0,X}\) give
\(\chi-\chi_0=\delta YO_{C^k}(1)\).
Squaring proves \eqref{eq:direct-projection-energy}, while
\(\chi_Y=-\mathcal W_Y/\mathcal W_X>0\).
\end{proof}

\begin{remark}[Continuation of the projected state]
\label{rem:direct-projected-continuation}
If the projected orbit starts in a component where \(F>0\), it
cannot cross the first zero of \(F\) in finite boundary-flow time.
Its inverse is an outward graph of \eqref{eq:direct-height-system}, since
\begin{equation}
\label{eq:direct-inverse-flux}
 X_{\log Y}=\delta F(w)\chi_w+Y\chi_Y>0.
\end{equation}
This remains valid as the projected state approaches a double root.
\end{remark}

The projected state moves with bounded speed in the boundary-flow time
\(\delta\log Y\). This bounds the height at which a fixed radial advance
can start.

\begin{lemma}[Height before a radial advance]
\label{lem:direct-radial-advance}
In \cref{lem:direct-boundary-flow}, let \(|F|\le C_F\),
\(\delta\sup|E_\chi|\le B_0\) and \(X_1<X_2\). Suppose a solution of
\eqref{eq:direct-height-system} lies in the domain of \(\mathcal W\) at
the heights \(Y_1\le Y\le Y_2\), with radius at most \(X_1\) at \(Y_1\)
and at least \(X_2\) at \(Y_2\). If \(8B_0Y_2\le X_2^2-X_1^2\), then
\[
 Y_1\le Y_2\exp\Bigl(-\frac{X_2-X_1}{2C_F\delta}\Bigr).
\]
\end{lemma}

\begin{proof}
At a point of the solution with radius \(X\), height \(Y\) and projected
state \(w\), \eqref{eq:direct-projection-energy} gives
\(0\le X^2-w^2=2Y\{B_0+\delta E_\chi(w,Y)\}\le4B_0Y\).
Hence \(w\le X_1\) at \(Y_1\). At \(Y_2\) the hypothesis gives
\(w^2\ge(X_1^2+X_2^2)/2\), so \(w\ge(X_1+X_2)/2\). Since
\(w_{\log Y}=\delta F(w)\) along the solution, \(w\) changes by at most
\(C_F\delta\log(Y_2/Y_1)\) between the two heights, and this change is at
least \((X_2-X_1)/2\).
\end{proof}

\subsection{Contracting height coordinates and the exact phase}

The projected boundary-flow times must now be expressed in the common
central height \(h\). We first linearize the height equation over a
possibly multidimensional base flow. We then normalize the passage from
the central section and use that normalization to fix the phase of each
projected orbit.
The following height normalization uses contraction in backward
\(\tau\)-time.

\begin{lemma}[Contracting height coordinate]
\label{lem:direct-contracting-height-coordinate}
Consider a skew system
\[
 x_\tau=g(x),\qquad y_\tau=\mathfrak a(x)y+y^2\mathcal N(x,y)
\]
on a base domain \(\mathcal O\) and \(|y|\le2y_0\), with \(\mathfrak a,\mathcal N\) and
their \(C^{k+1}\) derivatives bounded for a prescribed finite
\(k\). Suppose \(\partial_y(\mathfrak a y+y^2\mathcal N)\ge2c_y>0\) there. Fix a set
of initial base points \(\mathcal O_0\subset\mathcal O\), and suppose
that for \(x\) in a neighborhood of \(\mathcal O_0\) the backward base
flow \(x_b(t;x)\) stays in \(\mathcal O\) with uniform bounds
\[
 |D_x^j x_b(t;x)|\le C_j(1+t)^{M_j}\e^{jc_x t},
 \qquad 1\le j\le k+1,
\]
where \((k+1)c_x<c_y\). After reducing \(y_0\), the
function
\[
 \Theta(x,y)=-\int_0^\infty
      y_b(t;x,y)\mathcal N(x_b(t;x),y_b(t;x,y))\,dt
\]
is uniformly \(C^{k+1}\) on \(\mathcal O_0\times[-y_0,y_0]\), is divisible
by \(y\) with a bounded
\(C^k\) quotient, and
\(\Psi(x,y)=y \e^{\Theta(x,y)}\) has positive height derivative and
uniformly bounded inverse height derivatives through order \(k\).
Along the system, \(\Psi_\tau=\mathfrak a(x)\Psi\).
\end{lemma}

\begin{proof}
The backward height satisfies \(|y_b(t)|\le |y|\e^{-2c_y t}\).
In its variational equations the highest height jet has a negative
coefficient bounded away from zero. Each pure base derivative of
\(\mathfrak a y+y^2\mathcal N\) contains the decaying height; the other forcings are
products of lower height jets and base-flow jets. Variation of
constants and induction therefore bound every derivative of the
integrand through order \(k+1\) by
\(C_j(1+t)^{m_j}\e^{-(2c_y-jc_x)t}\), after a
reduction of the decay margin. The stated inequality makes these
bounds integrable. Differentiation under the integral gives
\(\Theta=O(y)\) and its uniform derivatives; one additional derivative
gives the bounded quotient \(\Theta/y\). Thus \(\Psi_y>c>0\) on a
smaller height interval, with the asserted inverse bounds.
Finally the flow property gives \(\Theta_\tau=-y\mathcal N\), and hence
\(\Psi_\tau=\mathfrak a\Psi\).
\end{proof}

For the passage from the central section, we next normalize the height
at a fixed regular radius. On a regular central prefix, the scaled
height equation is
\begin{equation}
\label{eq:direct-source-prefix}
 y_x=\frac{xy}{\delta(K(x)+B(x,y)y)},\qquad
 y(0)=h/Q,\qquad K(x)\ge k_0>0
\end{equation}
on a fixed interval. Here \(Q>0\) depends only on parameters.
The required derivatives of \(K,B\) are uniformly bounded.
Use \(t=\delta^{-1}\mathcal A(x)\), where
\(\mathcal A(x)=\int_0^x u/K(u)\,du\). The height equation becomes
\[
 y_t=y+y^2\mathcal N_t(y),\qquad
 \mathcal N_t(y)=-B/(K+By).
\]
The inverse of \(x\mapsto t\) is continuous at zero; the normalization below
differentiates only the height.

\begin{lemma}[Exact central-height normalization]
\label{lem:direct-source-height}
At each fixed terminal radius \(x_i>0\), on a sufficiently
small fixed height interval there is a coordinate
\[
 \Psi_i(y)=y\exp \Theta_i(y),\quad \Theta_i(0)=0,\quad
 \Psi_i'(y)>c,\quad
 \|\Psi_i\|_{C^k}+\|\Psi_i^{-1}\|_{C^k}\le C_k
\]
such that
\begin{equation}
\label{eq:direct-prefix-action}
 \Psi_i(y(x_i))=\e^{\mathcal A(x_i)/\delta}h/Q .
\end{equation}
These bounds hold on arbitrarily long intervals of boundary time.
The same conclusion holds if
\(K(x)\asymp K_*+x\), \(K'(x)\ge c\), \(K_*/\delta\ge1\),
with \(0<K_*\le K_{\max}\) and fixed \(K_{\max}\),
and the other coefficient derivative bounds are uniform.
\end{lemma}

\begin{proof}
Let \(T_i=\mathcal A(x_i)/\delta\), and let \(x_b(u)\) and \(y_b(u;y)\)
be the radius and height at backward time \(u\).
Set
\[
 Z(u)=\e^uy_b(u;y),\qquad
 a_\delta(u)=\frac{\e^{-u}}{K(x_b(u))}.
\]
The backward height equation becomes
\[
 Z'=\frac{a_\delta(u)Z^2B(x_b(u),\e^{-u}Z)}
           {1+a_\delta(u)ZB(x_b(u),\e^{-u}Z)}
       =:\mathcal N_u(Z).
\]
The long-time estimate reduces to a uniform bound on this coefficient
and its integral. In both cases of the lemma,
\[
 \sup_{0\le u\le T_i}a_\delta(u)
       +\int_0^{T_i}a_\delta(u)\,du\le C.
\]
For \(K\ge k_0\), this follows from \(a_\delta\le \e^{-u}/k_0\).
For \(K\asymp K_*+x\), let \(u_\delta\) be the backward
time to \(x_i/2\). The fixed terminal segment gives
\(u_\delta\ge c/\delta\) and \(K(x_b)\ge c_i>0\) before
that time. Afterwards \(K\ge cK_*\ge c\delta\), so
\[
 \sup_{u\ge u_\delta}a_\delta(u)
       +\int_{u_\delta}^{T_i}a_\delta(u)\,du
       \le C\delta^{-1}\e^{-c/\delta}.
\]
These bounds are independent of the positive lower endpoint
of the allowed \(K_*\)-range.

On a sufficiently small fixed \(Z\)-interval the denominator
in \(\mathcal N_u\) is at least \(1/2\), and
\[
 |\mathcal N_u(Z)|\le Ca_\delta|Z|^2,\qquad
 |\partial_Z\mathcal N_u(Z)|\le Ca_\delta|Z|,\qquad
 |\partial_Z^j\mathcal N_u(Z)|\le C_ja_\delta\quad(j\ge2).
\]
Gr\"onwall's inequality and the finite variational equations, with
the uniform integral bound for \(a_\delta\), give \(|Z(u;y)|\le2|y|\)
and uniform height derivatives of every prescribed finite order on a
smaller fixed interval. The first variation is positive and bounded
above and below. Set \(\Psi_i(y)=Z(T_i;y)\); then
\(y_b(T_i;y)=h/Q\) gives \eqref{eq:direct-prefix-action}, and
\[
 \Theta_i(y)=\log\frac{\Psi_i(y)}y
 =\int_0^{T_i}
   \frac{a_\delta ZB(x_b,\e^{-u}Z)}
        {1+a_\delta ZB(x_b,\e^{-u}Z)}\,du .
\]
The same bounds, with one extra derivative, give \(\Theta_i=O(y)\)
with uniformly bounded derivatives. Thus \(\Psi_i'(0)=1\), and the
inverse function theorem gives the asserted inverse bounds.
\end{proof}

The preceding normalization fixes the terminal radius. We next control
the dependence of the central height on that radius.
Suppose \(x\) ranges in a fixed compact positive
interval with \(K\ge k_0\), and fix a small terminal height \(y_f>0\).
Then the corresponding central height satisfies
\begin{equation}
\label{eq:direct-terminal-radius-jets}
 \delta\log(h(x)/Q)=-\mathcal A(x)+\delta M_p(x),\qquad
 \|M_p\|_{C^k}\le C_k .
\end{equation}
Choose a smaller fixed \(x_0>0\) with a positive action gap to
all terminal radii. On \(x\ge x_0\), use time \(\tau=x/\delta\).
The system is
\[
 x_\tau=\delta,\qquad
 y_\tau=\mathfrak a(x)y+y^2\mathcal N(x,y),\qquad \mathfrak a(x)=x/K(x)>c.
\]
Extend the coefficients by fixed smooth cutoffs preserving the positive
rate, and reduce the height box to preserve its derivative margin.
The backward base \(x_b(u)=x-\delta u\) has bounded state derivatives,
so \cref{lem:direct-contracting-height-coordinate} applies with
\(c_x=0\). It gives a uniformly regular coordinate
\(\Phi(x,y)=y\e^{\Theta(x,y)}\), with the required state derivative
and inverse bounds, satisfying \(\Phi_\tau=\mathfrak a(x)\Phi\).
Compose it at \(x_0\) with
\cref{lem:direct-source-height}, whose coordinate there is
\(\Psi_0\). The unit
\[
 W_p(\xi)=\frac{\Phi(x_0,\Psi_0^{-1}(\xi))}{\xi},
 \qquad W_p(0)=1,
\]
has uniform derivative bounds, using one extra order at zero.
The composition gives
\[
 \Phi(x,y_f)=\e^{\mathcal A(x)/\delta}\frac hQ\,W_p(\xi_0),\qquad
 \xi_0=\e^{\mathcal A(x_0)/\delta}h/Q,\qquad W_p(0)=1.
\]
Rearrange this as
\[
 \xi_0W_p(\xi_0)=
 \e^{-[\mathcal A(x)-\mathcal A(x_0)]/\delta}\Phi(x,y_f).
\]
The positive action gap bounds the right side and its \(j\)-th
terminal-radius derivative by \(C_j\delta^{-j}\e^{-c/\delta}\).
The uniformly regular inverse of \(\xi\mapsto\xi W_p(\xi)\) preserves
these finite-order bounds. Taking logarithms gives
\eqref{eq:direct-terminal-radius-jets} with
\(M_p=\log\Phi(x,y_f)-\log W_p(\xi_0)\).

The central-height normalization and the boundary-flow coordinate now
give the phase of the projected orbit. At a fixed regular entry
\(x_L\), combine \cref{lem:direct-source-height,lem:direct-boundary-flow}.
In the notation \(Y_i=\delta y_i\), let
\(\xi=\e^{\mathcal A(x_L)/\delta}h/Q\), so that
\(Y_i=\delta\xi \Upsilon(\xi)\), where \(\Upsilon(0)=1\) and \(\Upsilon\) has
uniform derivative bounds. Because \(F(x_L)>c>0\), the
coordinate shift at entry corresponds to the boundary-flow time
\[
 \alpha_i(\xi)
 =\int_{x_L}^{\mathcal W(x_L,\delta\xi \Upsilon(\xi))}\frac{dx}{F(x)}
 =\delta O(\xi).
\]
The integral lies in a regular entry collar. At a fixed
terminal height \(Y_o\), the resulting formula is
\begin{equation}
\label{eq:direct-common-phase}
 w=\Phi_F^{\tau(T)}(x_L),\quad
 \tau=T+C_p+\delta\Xi(\xi),\quad
 T=-\delta\log(h/Q),\quad
 C_p=\delta\log(Y_o/\delta)-\mathcal A(x_L).
\end{equation}
Here \(\Xi=\alpha_i/\delta-\log \Upsilon\) vanishes at zero and has
uniform derivative bounds. Both phase origins are now fixed by the same
central height \(h\). Near a boundary root, \(F(w)\) may tend to zero;
the next lemma bounds the derivatives of the phase correction relative
to this drift.

\begin{lemma}[Phase error at a boundary root]
\label{lem:direct-weighted-phase}
Suppose the phase in \eqref{eq:direct-common-phase} satisfies
\(\tau\ge\tau_0>0\), its orbit has \(F>0\), and
\(F(x_L)\ge f_0>0\), \(|F'|\le L_F\).
For every fixed \(N,j\ge0\), writing \(\theta=\tau_T\), one has
\begin{equation}
\label{eq:direct-phase-weight}
 |\theta-1|+\sum_{\ell=1}^j|\partial_T^\ell\theta|
       \le\varepsilon_{N,j}(\delta)F(w)^N,\qquad
 \varepsilon_{N,j}(\delta)\longrightarrow0 .
\end{equation}
\end{lemma}

\begin{proof}
In terms of \(\xi\), \eqref{eq:direct-common-phase} gives
\[
 \xi=(Y_o/\delta)\exp[-\tau/\delta+\Xi(\xi)],\qquad
 \xi_T=-\xi/\delta,\qquad \theta=1-\xi\Xi'(\xi).
\]
The left side of \eqref{eq:direct-phase-weight} is at most
\(C_j\delta^{-j}\xi\). To compare this bound with the drift, use
\(d\log F/d\tau=F'\) along the boundary flow; hence
\(F(w)\ge f_0\e^{-L_F\tau}\).
After division by \(F(w)^N\), the error is bounded by
\[
 C_{N,j}\delta^{-j-1}
       \exp[-(1/\delta-NL_F)\tau]
 \le C_{N,j}\delta^{-j-1}\e^{-\tau_0/(2\delta)}.
\]
Reduction of \(\delta\) proves the assertion, uniformly up to
simple or multiple boundary roots.
\end{proof}

For a no-root bottleneck, we also need height derivatives that remain
uniform as its action grows. On a fixed interval
\([x_L,x_R]\) suppose
\(F>0\), with positive margins only in fixed collars of its
endpoints. Define
\[
 \alpha_j(Y)=\int_{x_j}^{\mathcal W(x_j,Y)}\frac{dx}{F(x)},\qquad
 \Psi_j(y)=y\exp[-\alpha_j(\delta y)/\delta],\quad j=L,R.
\]
These are uniformly regular height coordinates, since
\(\alpha_j(Y)=Y\mathcal A_j(Y)\). The exact identity is
\begin{equation}
\label{eq:direct-bottleneck-action}
 \Psi_R(y_R)=
 \exp\left(\frac1\delta\int_{x_L}^{x_R}\frac{dx}{F(x)}\right)
 \Psi_L(y_L).
\end{equation}
For a small outgoing height, the backward projected flow and
\eqref{eq:direct-inverse-flux} reconstruct the outward passage.
Let \(\mathcal A_L\) be the action of the preceding
regular prefix and write its inverse height coordinate as
\(y_L=P_p(\zeta)\), where
\(\zeta=\e^{\mathcal A_L/\delta}h/Q\). By
\cref{lem:direct-source-height}, \(P_p(0)=0\), \(P_p'(0)=1\),
and its derivatives through the prescribed order are uniformly
bounded on a fixed small interval. Put
\[
 \begin{gathered}
 \mathcal A_f=\int_{x_L}^{x_R}\frac{dx}{F(x)},\qquad
 \beta=\e^{-\mathcal A_f/\delta}\in(0,1],\\
 \xi=\e^{(\mathcal A_L+\mathcal A_f)/\delta}h/Q,
 \qquad \mathcal C_p=\Psi_L\circ P_p.
 \end{gathered}
\]
Thus \(\zeta=\beta\xi\) and
\(\mathcal C_p(x)=x\Upsilon_p(x)\), with \(\Upsilon_p(0)=1\).
Composition with the prefix gives
\begin{equation}
\label{eq:direct-bottleneck-inverse-composition}
 \begin{aligned}
 \Psi_R(y_R)&=\beta^{-1}\mathcal C_p(\beta\xi)
             =\xi \Upsilon_p(\beta\xi),\\
 y_R&=\Psi_R^{-1}\bigl(\xi \Upsilon_p(\beta\xi)\bigr).
 \end{aligned}
\end{equation}
The factor \(\beta^{-1}\) can be large, but differentiation of the
argument \(\beta\xi\) contributes powers of \(\beta\). With \(\beta\)
fixed, the derivatives in the expanded height \(\xi\) are
\[
 \partial_\xi^j[\beta^{-1}\mathcal C_p(\beta\xi)]
       =\beta^{j-1}\mathcal C_p^{(j)}(\beta\xi),\qquad j\ge1.
\]
Its first derivative has a uniform positive lower bound, and
all prescribed higher derivatives are bounded because \(\beta\le1\).
Choose a fixed \(\xi_* >0\) so that \(\beta\xi\) lies in the
prefix domain and \(\xi \Upsilon_p(\beta\xi)\) lies in the common
domain of \(\Psi_R^{-1}\) for \(0\le\xi\le\xi_*\).
The resulting map \(\xi\mapsto y_R\) has a uniformly positive
first derivative; its image contains a fixed small terminal-height interval,
and its inverse has uniformly bounded derivatives there.
Constructing one extra derivative before dividing by \(\xi\) gives
\(y_R=\xi V_p(\xi)\), where \(V_p(0)=1\), and both \(V_p\)
and its reciprocal have uniform derivatives through the prescribed
order. These bounds for derivatives in the expanded height \(\xi\)
are uniform in \(\mathcal A_f\).

\subsection{Endpoint derivatives and the logarithmic slope ratio}

We now transfer the drift estimates to the endpoints. Composing the
inverse boundary projection with the weighted tail gives the derivative
in the projected radius; the phase formula then converts it to a
derivative in \(T\). Keeping the drift factor through this conversion
allows comparison with the exterior contribution to the logarithmic
slope ratio.

\pagebreak[3]
\begin{lemma}[Endpoint after a boundary-flow passage]
\label{lem:boundary-flow-endpoint}
Let \(\ell>0\), and write \(|\eta|=\ell X\) and \(z=\ell^2Y\) on one
branch. Suppose that on a fixed positive radius interval the height
equation has the form \eqref{eq:direct-height-system} with small
parameter \(\delta\), and let \(w\) be the projected state of
\cref{lem:direct-boundary-flow}. Fix a small height \(Y_o>0\) and a
compact interval \(J_w\subset(0,\infty)\) on which
\eqref{eq:direct-projection-energy} holds at \(Y=Y_o\), and put
\[
 t=\ell^2Y_o,\qquad
 \varepsilon_o=\ell+\frac{|q_0|}t+\frac{|q_1|}{\sqrt t}+|q_2|,\qquad
 E=\ell^2\Bigl(\frac{z_*}{t}\Bigr)^{-2q_2}.
\]
For \(\delta\) and \(\varepsilon_o\) sufficiently small, the endpoint
\(I\) of the graph whose projected state at height \(t\) is
\(w\in J_w\) satisfies
\begin{equation}
\label{eq:boundary-flow-endpoint}
 -I_w=EA_p w\,U_p(w),\qquad c<A_p<C,\qquad
 U_p=1+O_{C^3}(\delta+\varepsilon_o),
\end{equation}
where \(A_p\) depends only on the parameters and \(U_p\) is a
dimensionless positive unit factor. The constants depend on
\(J_w\), \(Y_o\) and the fixed boundary-flow data. If the two branches
use the same \(\ell\) and \(Y_o\), then
\begin{equation}
\label{eq:boundary-flow-amplitude-ratio}
 \left|\log\frac{A_-}{A_+}\right|
 \le2C_{\rm cone}\sqrt{z_*}+O(\delta+\varepsilon_o),
\end{equation}
where \(A_\pm\) are the values of \(A_p\) on the negative and positive
branches, respectively.  The branch labels are fixed throughout the
passage; later contraction estimates use only the upper one-sided
consequence of this absolute comparison.
\end{lemma}

\begin{proof}
By \eqref{eq:direct-projection-energy}, the energy at height \(t\) is
\(U(t)=\ell^2\{w^2/2+B_0Y_o+\delta Y_oE_\chi(w,Y_o)\}\), and its
\(w\)-derivative is \(\ell^2w\{1+O_{C^3}(\delta)\}\) on \(J_w\). The
normalized energies \(U(t)/t\) lie in a fixed compact positive
interval. If this interval is not inside the common cone, continue the
graphs first by \eqref{eq:direct-finite-common-cone-entry} from \(Y_o\)
to a fixed larger \(Y_1\), whose energies lie in that cone. There the
error is \(O_{C^k}(\varepsilon_o)\) and the first derivative is
\(1+O(\varepsilon_o)\), since \(\varepsilon_{\rm in}\le C\varepsilon_o\)
for fixed \(Y_o,Y_1\). \Cref{lem:direct-weighted-tail} at
\(t_1=\ell^2Y_1\) gives a centered endpoint map with slope
\(t_1(z_*/t_1)^{-2q_2}a\), where \(c<a<C\) and, by
\eqref{eq:direct-tail-first-slope},
\(|\log a|\le C_{\rm cone}(\sqrt{z_*}+|q_1|/\sqrt{t_1})\), and with a
relative \(C^3\) error \(O(e(t_1))=O(\varepsilon_o)\). Since
\((t_1/t)^{2q_2}=(Y_1/Y_o)^{2q_2}=1+O(\varepsilon_o)\), the chain rule gives
\eqref{eq:boundary-flow-endpoint}. The factors \(\ell^2\),
\((z_*/t)^{-2q_2}\) and \((Y_1/Y_o)^{2q_2}\) are common to the two
branches, which gives the bound for \(A_-/A_+\).
\end{proof}

To compute the logarithmic rate in the common height coordinate,
consider an endpoint map on a projected-radius interval contained in a
fixed compact subset of \((0,\infty)\), of the form
\begin{equation}
\label{eq:direct-endpoint-unit}
 -I_w=EA_p w\,U_p(w),\qquad
 c<A_p<C,\quad U_p=1+O_{C^3}(\delta),\quad 0<E\longrightarrow0 ,
\end{equation}
as given by \cref{lem:boundary-flow-endpoint} when
\(\varepsilon_o=O(\delta)\); its amplitude satisfies
\(E\le C\ell^{3/2}\) by \eqref{eq:direct-tail-amplitude} with
\(t=\ell^2Y_o\). Here \(E=E(p)\) and \(A_p\) are independent of the
central height.
Use the phase \eqref{eq:direct-common-phase}, so that
\(w_T=\theta F(w)\), with \(\theta=\tau_T\) and \(F=K/w>0\).
Put \(m=hI_h=-\delta I_T\),
\(\mathfrak H=\partial_T\log m\).
In the identities below, primes on \(\theta,\mathfrak H\)
denote \(T\)-derivatives, and primes on \(K,F,\log U_p,\mathfrak H_0\) denote
\(w\)-derivatives:
\begin{align}
\label{eq:direct-rate-identities}
 m&=\delta EA_pK(w)\theta U_p(w),\\
 \mathfrak H&=\theta\mathfrak H_0(w)+\theta'/\theta,\qquad
 \mathfrak H_0=K'/w+F(\log U_p)',\nonumber\\
 \mathfrak H'&=\theta^2F\mathfrak H_0'+\theta'\mathfrak H_0
                 +\theta''/\theta-(\theta'/\theta)^2,\nonumber\\
 \mathfrak H_0'&=(wK''-K')/w^2+F'(\log U_p)'+F(\log U_p)''.\nonumber
\end{align}
By \cref{lem:direct-weighted-phase}, these identities imply
\begin{equation}
\label{eq:direct-rate-curvature}
 \mathfrak H'=F(w)\left\{\frac{wK''-K'}{w^2}+o(1)\right\},\qquad
 |I_T|+|I_{TT}|\le CEF(w).
\end{equation}
The errors are uniform up to a limiting boundary root, with the drift
factor retained in both estimates.

For the two endpoints use the same logarithmic height coordinate
\(T=-\delta\log(h/Q)\) in the logarithmic slope ratio
\eqref{eq:direct-completed-log-ratio}.
The exterior contribution is
\begin{equation}
\label{eq:direct-exterior-curvature}
 \partial_T^2[(\log\Pi_p')(I_-)]
       =(\log\Pi_p')''(I_-)I_{-,T}^2+(\log\Pi_p')'(I_-)I_{-,TT}.
\end{equation}
By \eqref{eq:direct-rate-curvature}, this term is
\(O(EF_-+E^2F_-^2)\).

The exterior contribution has the same drift weight as the negative
endpoint estimate and carries the vanishing amplitude \(E\). This leaves
the coefficients \(Q_\sigma=(wK_\sigma''-K_\sigma')/w^2\) of the two
drifts to determine the principal curvature of \(\mathcal L\).

\begin{lemma}[Curvature criterion]
\label{lem:curvature-criterion}
On an interval of central heights, suppose that both endpoints have the
form \eqref{eq:direct-endpoint-unit} in the common logarithmic height
coordinate \(T=-\delta\log(h/Q)\), that their phases satisfy the
hypotheses of \cref{lem:direct-weighted-phase}, and that the negative
projected state stays in a fixed compact subset of \((0,\infty)\). If
\(Q_-\le-c\) and \(Q_+\ge c\) there, then \(\mathcal L''<0\) for all
sufficiently small parameters, uniformly up to simple or multiple
boundary roots.
\end{lemma}

\begin{proof}
By \eqref{eq:direct-rate-curvature},
\(\mathfrak H_\sigma'=F_\sigma(w_\sigma)\{Q_\sigma+o(1)\}\). The exterior
term \eqref{eq:direct-exterior-curvature} is \(O(EF_-+E^2F_-^2)=o(1)F_-\),
since \(F_-\) is bounded and \(E\to0\). Hence
\[
 \mathcal L''=\mathfrak H_-'-\mathfrak H_+'+\partial_T^2[(\log\Pi_p')(I_-)]
 \le-\tfrac c2(F_-+F_+)<0.\qedhere
\]
\end{proof}

For a drift
\(K(w)=K(0)+K'(0)w+\tfrac12K''(0)w^2+w^3R_K(w)\), direct
differentiation gives
\begin{equation}
\label{eq:direct-curvature-coefficient-cubic}
 \frac{wK''-K'}{w^2}
 =-\frac{K'(0)}{w^2}+3R_K(w)+5wR_K'(w)+w^2R_K''(w).
\end{equation}

In the weak-root range, the first rate identity gives a one-sided slope
comparison. The next result expresses it in the common height coordinate,
retaining the change of scale from the positive branch's coordinate.

\begin{corollary}[Weak-root slope comparison]
\label{cor:direct-weak-root-comparison}
Let \(I_\pm(h)\) be the \(C^2\) endpoint maps on one interval of
central heights, with \(I_{\pm,h}>0\), and let \(a,\kappa>0\).
Suppose the positive endpoint map satisfies
\(-I_{+,w}=EA_pwU_p(w)\) with \(E>0\), \(c<A_p<C\) and
\(U_p=1+O_{C^2}(\kappa)\), where \(E,A_p\) are independent of the
central height. The radius
interval satisfies \(0<w_*\le w_a\le w<\alpha\le w^*<\infty\), with
fixed \(w_*,w^*\), and its exact phase uses \(\widehat T=-\kappa\log h\).
Assume the weighted-phase lemma
applies there, \(F=K/w>0\), the \(C^2\) bounds for \(K\) are
uniform, and \(K'(w)/w\le-c_0<0\).
In the common logarithmic height coordinate \(T=a\widehat T\), put
\(\mathfrak H_\sigma=\partial_T\log(hI_{\sigma,h})\).
Then, for sufficiently small \(\kappa\),
\begin{equation}
\label{eq:direct-weak-root-rate}
 \mathfrak H_+\le-\frac{c_0}{2a}.
\end{equation}
If also \(\mathfrak H_-\ge-C_n\), \(|I_{-,T}|\le\varepsilon_p\to0\),
and the \(C^2\) exterior map on these endpoint windows has positive slope with
\(|(\log\Pi_p')'|\le M\), then, for sufficiently small \(a\),
\(\mathcal L_T>0\) on this interval, and the return displacement has at
most two zeros there.
All bounds, the phase-lemma constants, and the convergence
\(\varepsilon_p\to0\) are uniform on a fixed preliminary parameter region
with small \(a\).
\end{corollary}

\begin{proof}
With \(\theta=\partial_{\widehat T}\tau\),
\eqref{eq:direct-rate-identities} gives
\[
 a\mathfrak H_+
 =\theta\{K'/w+F(\log U_p)'\}
                   +\theta_{\widehat T}/\theta.
\]
The bounded radius interval and the stated phase estimates make
the error \(o(1)\), proving \eqref{eq:direct-weak-root-rate}.
In the coordinate \(T\),
\[
 \mathcal L_T
 =(\log\Pi_p')'(I_-)I_{-,T}
                      +\mathfrak H_--\mathfrak H_+
 \ge \frac{c_0}{2a}-C_n-M\varepsilon_p>0.
\]
Fix the constants on the preliminary region, reduce \(a\),
and then choose the phase, tail, and parameter thresholds.
Since \(\mathcal L_T>0\), \eqref{eq:direct-completed-log-ratio} gives the
two-zero count.
\end{proof}

\section{Contact with the central section and positive offsets}
\label{sec:contact-positive-offsets}

We begin with two regions in which \(\Delta_p\) has at most two zeros: the contact scale
\(q_0=\widehat q_0q_1^2\) with bounded \(\widehat q_0\), and positive
offsets with small \(q_1/\sqrt{q_0}\). In the first, a scattering
comparison controls the derivative ratio on a compact range of central
heights, while the paired identity controls the two ends of the passing
interval. In the second, the Dulac map of an interior saddle determines
the crossing direction of the logarithmic slope ratio near the
separatrices; exterior contraction controls the remaining heights.

In this section and in
\cref{sec:recurring-ranges,sec:primary-regimes,sec:quadratic-regimes,sec:direct-negative-annulus-buffer,sec:direct-negative-root-regimes,sec:sharpness},
the height \(z_*\) of the fixed section, the exterior windows, the common
cone, and the landing level \(m_{\rm land}\) are those chosen in the first
row of \cref{tab:direct-constant-order}. All scaled thresholds and parameter
bounds below are chosen after these fixed data.

The constant \(C_{\rm cone}\) in \eqref{eq:direct-tail-first-slope} is
independent of the contact width and the later coefficient margins and
regional cuts.
The continuation constants and higher exterior derivatives may depend
on these cuts and on \(z_*\); they enter only the final parameter
thresholds, with the sections fixed. The section requirements are justified
in the proof of \cref{prop:direct-uniform-displacement} and the derivation
of \eqref{eq:sharpness-lambda-range}.

\subsection{The contact strip}
\label{sec:direct-bounded-contact}

We treat the contact strip, with fixed \(L<\infty\),
\begin{equation}
\label{eq:direct-contact-region}
 q_1>0,\qquad q_0=\widehat q_0 q_1^2,\qquad |\widehat q_0|\le L .
\end{equation}
The limits of \(q_1\) and \(q_2\) are independent, including
\(q_2=0\); constants may depend on \(L\) and \(C_0\).
The difficulty is to compare the two passages uniformly through
contact, without relating the sizes of \(q_1\) and \(q_2\).
The paired identity gives opposite derivative signs at the two ends
of the height interval. On the compact interval between them,
\cref{lem:direct-contact-scattering-ratio} gives a strictly decreasing
ratio of the limiting scattering derivatives. We transfer its derivative
margin to the endpoint maps by a \(C^2\) estimate. Every middle zero
of \(\mathcal L\) then crosses upwards in \(T\), and the outer signs
complete the count on \(\mathcal H_p\).

\begin{proposition}[Bounded contact offsets]
\label{prop:direct-bounded-contact}
For every fixed \(L\), the return displacement
\eqref{eq:central-displacement-new} has at most two
distinct zeros on the height interval \(\cH_p\) in
\eqref{eq:direct-contact-region}, for all sufficiently small parameters.
\end{proposition}

\subsubsection{Barrier and tail estimates}

With the contact scales of \cref{tab:direct-scale-guide}, put
\[
 d_{\rm c}=\max\{|q_2|,\sqrt{q_1}\},\qquad
 z=q_1^2 Z,\qquad h=q_1^2 \widehat h,\qquad
 |\eta_\sigma|=q_1u_\sigma,\qquad
 \mathcal U_\sigma=u_\sigma^2/2 .
\]
The rescaled equations are
\begin{equation}
\label{eq:direct-contact-equations}
 (\mathcal U_\sigma)_Z
 =B(\sigma q_1u_\sigma,q_1^2Z)
  +\frac{-\widehat q_0-\sigma u_\sigma-q_2u_\sigma^2
               -\sigma q_1u_\sigma^3R_q(\sigma q_1u_\sigma)}{Z}.
\end{equation}
They are the graphs of the analytic outward time equations
\begin{equation}
\label{eq:direct-contact-time}
 (u_\sigma)_\tau=-\widehat q_0-\sigma u_\sigma-q_2u_\sigma^2
               -\sigma q_1u_\sigma^3R_q+ZB,\qquad
 Z_\tau=u_\sigma Z .
\end{equation}
The central point is \(u_\sigma=0,\ Z=\widehat h\). Its radial flux is
\[
 f(\widehat h)=\widehat h B(0,q_1^2\widehat h)-\widehat q_0,\qquad
 f_{\widehat h}=B(0,q_1^2\widehat h)
   +q_1^2\widehat h B_z(0,q_1^2\widehat h)\ge b_0.
\]
Thus all \(\widehat h>0\) pass when \(\widehat q_0\le0\); when \(\widehat q_0>0\)
the passing interval is \(\widehat h>\widehat h_0\), where \(f(\widehat h_0)=0\).

The negative graph has the uniform upper barrier
\begin{equation}
\label{eq:direct-contact-cap}
 U_-(z)\le C(d_{\rm c}^2+z),\qquad |\eta_-(z)|\le C(d_{\rm c}+\sqrt z).
\end{equation}
Indeed, for \(x\ge Kd_{\rm c}\), the coefficient bounds give
\[
 q(-x)\ge cx^3-d_{\rm c}x^2-q_1x-Lq_1^2\ge0.
\]
After division by \(x^3\), the three adverse terms are at most
\(K^{-1},K^{-2},Ld_{\rm c}/K^3\). Choose \(K\) large and then \(d_{\rm c}\)
small. On the line \(U=C(d_{\rm c}^2+z)\), with \(\sqrt{2C}>K\),
the graph derivative \(B-q(-\sqrt{2U})/z\) is at most
\(\sup B<C\). The line starts above the central point, so first contact
proves \eqref{eq:direct-contact-cap}.

To continue from bounded rescaled central heights to the fixed section,
we first obtain the following cone bounds. For each fixed bounded set
of rescaled central heights, they hold for a sufficiently large fixed
\(Z_0\):
\begin{equation}
\label{eq:direct-contact-cones}
 cz\le U_\pm(z),\qquad U_+(z)\le Cz,\qquad
 z_0:=q_1^2Z_0\le z\le z_*.
\end{equation}
For continuation on a fixed interval \(\widehat h\le Z\le Z_1\),
first bound the positive radius by \(M_{L,Z_1}\). At a fixed
\(M>2(L+(\sup B)Z_1+1)\), the first right side of
\eqref{eq:direct-contact-time} is negative, after reducing
\(|q_2|\) and \(q_1\). Hence the graph cannot first reach \(M\).
Choose \(Z_1\) larger than twice these rescaled central heights and
so large that \(L/Z_1<b_0/4\). On the lower comparison line
\(\mathcal U=m(Z-Z_1)\), the adverse terms in
\eqref{eq:direct-contact-equations} are bounded by
\[
 L/Z_1+C\sqrt{m/Z_1}+2d_{\rm c}m+C m^{3/2}\sqrt{z_*}.
\]
Choose \(m>0\) small after \(Z_1\), and use the common-cone choice of
\(z_*\) already fixed before the regional thresholds.  Reducing
\(d_{\rm c}\) and \(q_1\) after these choices makes the graph derivative
greater than \(m\).  This proves a lower bound for both
signs. For the positive upper bound use \(W=C_2Z+C_3\), choosing
\(C_3\) to dominate the energy at \(Z_1\) and to have
\(\sqrt{2W}\ge2L\). The term \(-\widehat q_0-\sqrt{2W}\) is then
nonpositive, and the remaining derivative is at most
\[
 \sup B+2d_{\rm c}W/Z+Cq_1W^{3/2}/Z<C_2
\]
for \(C_2\) fixed large and the same successive reductions.
Enlarge \(Z_0\) relative to \(Z_1,C_3\). These comparisons give
\eqref{eq:direct-contact-cones}, with fixed cone constants, up
to \(z_*\). Together with \eqref{eq:direct-contact-cap}
they keep both graphs inside the box in the coordinates \((\eta,z)\).

The negative graph may remain above the upper cone for part of the
passage. To estimate the resulting loss in the paired identity, set
\(v_0=|\eta_-(z_0)|/q_1\) and
\(t_h=\log(Z_0/\widehat h)\). For \(z\ge z_0\), discard the favorable
negative cubic term and use
\(q_1\sqrt{2U}\le U/8+4q_1^2\). Taking \(Z_0\ge L+4\)
and then \(d_{\rm c}\) small yields
\begin{equation}
\label{eq:direct-contact-upper-comparison}
 U_-'\le C+\tfrac12 U_-/z,\qquad
 U_-(z)\le Cz+U_-(z_0)(z/z_0)^{1/2}.
\end{equation}
The second term is at most \(z\) once
\(z\ge z_0\max\{1,(U_-(z_0)/z_0)^2\}\).
Below that height use \eqref{eq:direct-contact-cap}; above it use
the resulting upper cone. Splitting the integral, also when this
height exceeds \(z_*\), gives
\begin{equation}
\label{eq:direct-contact-tail-loss}
 \int_{z_0}^{z_*}\frac{|\eta_-(z)|}{z}\,dz
       \le C\sqrt{z_*}+Cd_{\rm c}\log(1+v_0).
\end{equation}

The tail loss must be smaller than the gain in the paired identity
at small central heights. We therefore express its last logarithm in
terms of the central height. In the
negative equation use \(t=\log(Z/\widehat h)\) and
\(\mathcal V=\sqrt{1+u_-^2}\). Dropping the favorable cubic
term, and using \(|\widehat q_0|\le L\), gives
\[
 \mathcal V_t\le C_L+d_{\rm c}\mathcal V+(\sup B)\widehat h \e^t,\qquad
 \mathcal V(t)\le \e^{d_{\rm c}t}(1+C_Lt)+C\widehat h \e^t .
\]
In particular,
\begin{equation}
\label{eq:direct-contact-source-growth}
 \log(1+v_0)\le C_{L,Z_0}+d_{\rm c}t_h+\log(1+t_h).
\end{equation}
This bounds the tail loss in terms of \(t_h\), uniformly
in \(d_{\rm c}|\log q_1|\).

\subsubsection{The two ends of the height interval}

For the fixed contact width \(L\), choose \(Z_0\) after the bounded-height
comparison, then fix the cone and barrier constants. Reduce \(q_1\)
so that
\begin{equation}
\label{eq:direct-contact-prefix-absorption-choice}
 q_1Z_0\le \frac{b_1}{4C_B},
 \qquad z_0=q_1^2Z_0\le \frac{q_1}{4C_B},
\end{equation}
where \(C_B\) bounds \(\lvert B_\eta\rvert\) on the comparison box and
\(b_1\) is the positive coefficient margin in
\eqref{eq:normal-sign-margins-new}. With \(z_*\) already fixed by the
common-cone construction, reduce \(\lvert q_2\rvert\) (and hence
\(d_{\rm c}\)) relative to this \(z_*\), after \(q_1\). The absorptions
below are then uniform in the independent limits of \(q_1\) and \(q_2\).

For \(\widehat q_0\le0\), choose initially \(\widehat h\le1/2\).
On \(h\le z\le z_0\), the terms \(B_\eta/|\eta_\sigma|\) in the
paired identity \eqref{eq:paired-log-derivative-new} are absorbed
by half of the positive terms \(q_1/(|\eta_\sigma|z)\): indeed,
\eqref{eq:direct-contact-prefix-absorption-choice} gives
\(z\le z_0\le q_1/(4C_B)\), so
\( |B_\eta|/|\eta_\sigma|\le q_1/(4|\eta_\sigma|z)\).
The fixed positive-radius bound
gives a gain at least \(c_Lt_h\). By
\eqref{eq:direct-contact-cap}, the cubic loss in this prefix is at most
\(C_Ld_{\rm c} t_h\). On \(z\ge z_0\), the lower cones bound the
\(B_\eta/|\eta_\sigma|\) integrals; the positive upper cone bounds
the cubic integral on the positive branch, and
\eqref{eq:direct-contact-tail-loss} bounds it on the negative branch.
Therefore
\[
 \log\frac{I_{-,h}}{I_{+,h}}
 \ge (c_L-C_Ld_{\rm c})t_h-C\sqrt{z_*}-Cd_{\rm c}\log(1+v_0)
 \ge \tfrac12c_Lt_h-C_L .
\]
The last step uses \eqref{eq:direct-contact-source-growth} and
\(\log(1+t_h)\le t_h\), reducing \(d_{\rm c}\) after \(L,Z_0\) and
\eqref{eq:direct-contact-prefix-absorption-choice} so that the
coefficient of \(t_h\) remains at least \(c_L/2\).
Choose a fixed \(\widehat h_{\rm lo}>0\)
so small that this bound exceeds \(\log(1/\Lambda_-)\) for
\(\widehat h\le \widehat h_{\rm lo}\). Then \(\Delta_h>0\) there.

For \(\widehat q_0\ge0\), the lower end of the passing interval is
controlled by the radial flux. Consider central heights with small positive flux
\(0<f\le1\). Then
\[
 \widehat h\le (L+1)/b_0,\qquad \widehat h\ge f/\sup B .
\]
For \(0\le t\le1\), subtracting the value of \(ZB\)
at the central point in \eqref{eq:direct-contact-equations} gives
\[
 (u_-^2/2)_t\le C_Lu_-+d_{\rm c}u_-^2+f+C_Lt .
\]
The explicit upper comparison
\[
 \bar v=K(t+\sqrt{ft}),\qquad
 (\bar v^2/2)'=K^2(t+\tfrac32\sqrt{ft}+f/2)
\]
proves \(u_-\le\bar v\). Choose \(K\) large first to dominate the
linear, \(f\), and \(t\) terms, then \(d_{\rm c}\) small to absorb
\(d_{\rm c}\bar v^2\). A fixed provisional radius box larger than
\(4K\) completes the first-contact continuation argument.
Thus the favorable negative-branch term in the paired identity
satisfies
\[
 \int_0^1\frac{dt}{u_-(t)}
 \ge \frac2K\log(1+f^{-1/2})
 \ge c_L\log(1/f)-C_L .
\]
Choose \(Z_0>\e(L+1)/b_0\), so this interval lies in the prefix.
The other prefix losses and the whole tail are bounded
as above. Since \(t_h\le\log(C_L/f)\), the result is
\begin{equation}
\label{eq:direct-contact-flux-gain}
 \log\frac{I_{-,h}}{I_{+,h}}
 \ge \tfrac12c_L\log(1/f)-C_L .
\end{equation}
Choose \(f_{\rm lo}>0\) so small that the right side exceeds
\(\log(1/\Lambda_-)\) for \(f\le f_{\rm lo}\). Thus \(\Delta_h>0\)
on this interval, whose connectedness follows from \(f_{\widehat h}>0\).
This estimate is uniform as \(f/\widehat h\to0\).

At the other end let \(\theta_0\) be the constant of
\cref{lem:relative-height-paired-new}, and choose
\(\widehat h_{\rm hi}>(L+1)/b_0+\widehat h_{\rm lo}\) so large that
\(L/\widehat h_{\rm hi}+\widehat h_{\rm hi}^{-1/2}\le\theta_0/2\).
For \(h\ge q_1^2\widehat h_{\rm hi}\),
\[
 |q_0|/h\le L/\widehat h_{\rm hi},\qquad
 q_1/\sqrt h\le \widehat h_{\rm hi}^{-1/2},
\]
and \(|q_2|\le\theta_0/2\) for small parameters. Hence
\cref{lem:relative-height-paired-new} applies and gives
\(\Delta_h<0\) at these central heights.

Write \(\widehat h_{\rm a}=\widehat h_{\rm lo}\) when \(\widehat q_0\le0\), and otherwise
let \(\widehat h_{\rm a}\) be the unique solution of \(f(\widehat h_{\rm a})=f_{\rm lo}\).
Put \(\widehat h_{\rm b}=\widehat h_{\rm hi}\). Taking \(f_{\rm lo}\le1\), the choice
of \(\widehat h_{\rm b}\) gives \(f(\widehat h_{\rm b})>1\) for \(\widehat q_0>0\), so \(\widehat h_{\rm a}<\widehat h_{\rm b}\) in both
cases. At central heights in \(\cH_p\), the displacement has
\(\Delta_h>0\) for \(\widehat h\le \widehat h_{\rm a}\) and \(\Delta_h<0\) for
\(\widehat h\ge \widehat h_{\rm b}\), including the height thresholds whenever they
are admissible.
The middle interval is \([\widehat h_{\rm a},\widehat h_{\rm b}]\). On bounded heights,
\[
 f(\widehat h)-(B_0\widehat h-\widehat q_0)
       =O(q_1^2\widehat h^2),\qquad B_0=B(0,0,p).
\]
Consequently the central heights of the middle interval belong, after
parameter reduction, to
\begin{equation}
\label{eq:direct-contact-compact-passing}
 \mathcal K_{\rm c}=\left\{(\widehat h,\widehat q_0,B_0):
 \begin{array}{c}
 b_0\le B_0\le B_{\max},\quad |\widehat q_0|\le L,\\
 h_*\le\widehat h\le\widehat h_{\rm hi},\quad
 B_0\widehat h-\widehat q_0\ge f_*
 \end{array}\right\},
\end{equation}
where
\[
 h_*:=\min\{\widehat h_{\rm lo},f_{\rm lo}/(2B_{\max})\}>0,
 \qquad f_*:=\min\{b_0\widehat h_{\rm lo},f_{\rm lo}/2\}>0.
\]
Here \(B_{\max}\) is a fixed upper bound for \(B_0\).
Every fixed \((\widehat q_0,B_0)\) fibre of \(\mathcal K_{\rm c}\) is an interval,
possibly empty. This compact set lies strictly inside the model passing domain
\begin{equation}
\label{eq:direct-contact-model-domain}
 |\widehat q_0|\le L,\qquad b_0\le B_0\le B_{\max},\qquad
 \widehat h>\max\{0,\widehat q_0/B_0\}.
\end{equation}
For the estimates below, choose one fixed
\(\widehat h_{\rm ref}>\widehat h_{\rm hi}+1+2L/b_0\),
and let $\mathcal K_{\rm c}^+$ be the compact set obtained from $\mathcal K_{\rm c}$
by replacing its upper height bound by $\widehat h_{\rm ref}+1$
and its lower height and flux bounds by $h_*/2$ and $f_*/2$.
Every fibre contains the reference height $\widehat h_{\rm ref}$ and
is still an interval. By the estimate of \(f\) above, the radial flux
at the central point is positive on a neighborhood of this enlarged
set for small parameters.

\subsubsection{Scattering and a strictly decreasing derivative ratio}

The outer signs leave only the compact passing set
\(\mathcal K_{\rm c}\). On this set we compare the derivatives of the
limiting scattering maps, and then transfer their strict ratio
margin to the endpoint maps.
For fixed contact heights the limiting equations are
\begin{equation}
\label{eq:direct-contact-model}
 (\mathcal U_\sigma^0)_Z=B_0-\frac{\widehat q_0+\sigma\sqrt{2\mathcal U_\sigma^0}}{Z},
 \qquad \mathcal U_\sigma^0(\widehat h;\widehat h)=0 .
\end{equation}
Define
\begin{equation}
\label{eq:direct-contact-scattering}
 S_\sigma(\widehat h;\widehat q_0,B_0)=
 \lim_{Z\to\infty}
 [\mathcal U_\sigma^0(Z;\widehat h_{\rm ref})-\mathcal U_\sigma^0(Z;\widehat h)] .
\end{equation}
The limit exists with local $C^2$ convergence of order $O(Z^{-1/2})$
in $\widehat h$, uniformly on compact passing sets. Fix a height beyond
all the rescaled central heights of such a set. The time equations
\eqref{eq:direct-contact-time} with
$q_1=q_2=0$ have positive initial radial flux, and their solutions
meet this height section transversely. The resulting
finite-height maps are analytic.
The limiting barriers put both graphs in $cZ\le \mathcal U_\sigma^0\le CZ$
above a sufficiently large fixed cut. There
\[
 \partial_{\mathcal U}\mathcal G_\sigma^0=-\frac{\sigma}{Z\sqrt{2\mathcal U_\sigma^0}},\qquad
 |\partial_{\mathcal U}\mathcal G_\sigma^0|\le CZ^{-3/2},\qquad
 |\partial_{\mathcal U}^2\mathcal G_\sigma^0|\le CZ^{-5/2}.
\]
The mean-value equation for two graphs and the triangular first- and
second-variation equations have these integrable coefficients. Hence
the graph difference and its first two \(\widehat h\)-derivatives remain
bounded, and their \(Z\)-derivatives are \(O(Z^{-3/2})\). Integration to
infinity gives the asserted \(C^2\) convergence and rate. The same
majorants give continuous dependence on \((\widehat q_0,B_0)\), and
\(S_\sigma(\widehat h_{\rm ref})=0\).

The variation formula for a moving initial point gives
\begin{equation}
\label{eq:direct-contact-model-derivative}
 S_\sigma'(\widehat h)=
 \left(B_0-\frac{\widehat q_0}{\widehat h}\right)
 \exp\left[-\sigma\int_{\widehat h}^\infty
          \frac{dZ}{Z\sqrt{2\mathcal U_\sigma^0(Z;\widehat h)}}\right]>0 .
\end{equation}
Indeed, the initial factor of $\partial_{\widehat h}\mathcal U_\sigma^0$
is $-(B_0-\widehat q_0/\widehat h)$ for both signs. The integral in
\eqref{eq:direct-contact-model-derivative} converges both at the
central point and at infinity.

\begin{samepage}
\begin{lemma}[Scattering derivative ratio]
\label{lem:direct-contact-scattering-ratio}
For fixed $(\widehat q_0,B_0)$, the function
\[
 \mathscr R(\widehat h):=S_-'(\widehat h)/S_+'(\widehat h)
\]
is strictly decreasing throughout the passing interval
\eqref{eq:direct-contact-model-domain}. On each fixed compact
passing set there are constants $c_1,c_2>0$ such that
\begin{equation}
\label{eq:direct-contact-ratio-margins}
 S_+'\ge c_1,\qquad -\mathscr R'\ge c_2.
\end{equation}
\end{lemma}
\end{samepage}

\begin{proof}
Put $t=\log(Z/\widehat h)$ and
$W_\sigma(t;\widehat h)=\mathcal U_\sigma^0(\widehat h \e^t;\widehat h)$.
The initial time is now fixed, and
\[
 (W_\sigma)_t=B_0\widehat h \e^t-\widehat q_0-\sigma\sqrt{2W_\sigma},
 \qquad W_\sigma(0;\widehat h)=0.
\]
Take $V_\sigma=\partial_{\widehat h}W_\sigma$ at fixed $t$.
With $f^0=B_0\widehat h-\widehat q_0>0$,
the regular time equations at the central point give, uniformly on compact
passing sets,
\[
 W_\sigma=f^0t+O(t^{3/2}),\qquad
 V_\sigma=B_0t+O(t^{3/2}).
\]
For example, in the analytic time variable $\tau$,
$u_\sigma=f^0\tau+O(\tau^2)$ and
$t=f^0\tau^2/2+O(\tau^3)$; inversion in $\sqrt t$ gives these
expansions with their derivatives in $\widehat h$.
For $t>0$ the variation equation is
\[
 (V_\sigma)_t=B_0\e^t-\frac{\sigma V_\sigma}{\sqrt{2W_\sigma}},
 \qquad V_\sigma(0)=0.
\]
Its coefficient is $O(t^{-1/2})$ at zero. Variation of constants,
first from positive time and then by passage to zero, yields
\begin{equation}
\label{eq:direct-contact-fixed-time-variation}
 V_\sigma(t)=B_0\int_0^t \e^{\bar t}
 \exp\left[-\sigma\int_{\bar t}^t\frac{du}{\sqrt{2W_\sigma(u)}}\right]d\bar t>0.
\end{equation}
For large $t$, the cone estimates give $W_\sigma\asymp \e^t$.
The tail coefficient of this variation equation is integrable,
so \eqref{eq:direct-contact-fixed-time-variation} gives
$V_\sigma=O(\e^t)$. Thus both
$(2W_\sigma)^{-1/2}$ and
$V_\sigma(2W_\sigma)^{-3/2}$ are bounded by
$Ct^{-1/2}$ near zero and by $C\e^{-t/2}$ at infinity,
uniformly on every compact passing set. Dominated differentiation gives
\[
 \mathcal T_\sigma(\widehat h):=\int_0^\infty
            \frac{dt}{\sqrt{2W_\sigma(t;\widehat h)}},\qquad
 \mathcal T_\sigma'=-\int_0^\infty
       \frac{V_\sigma(t;\widehat h)}{(2W_\sigma(t;\widehat h))^{3/2}}\,dt<0,
\]
and $\mathcal T_\sigma'$ is continuous in the passing parameters.
The common positive factor $B_0-\widehat q_0/\widehat h$ in
\eqref{eq:direct-contact-model-derivative} cancels, giving
\[
 \mathscr R=\exp(\mathcal T_-+\mathcal T_+),\qquad
 \mathscr R'=\mathscr R(\mathcal T_-'+\mathcal T_+')<0.
\]
Continuity and compactness give \eqref{eq:direct-contact-ratio-margins}.
\end{proof}

\subsubsection{Endpoint preparation and the two-zero bound}

To transfer the scattering comparison to the endpoint maps, we must
retain the quadratic tail factor even when its logarithm is unbounded.
The common endpoint amplitude is
\begin{equation}
\label{eq:direct-contact-amplitude}
 E=q_1^2(z_*/q_1^2)^{-2q_2}.
\end{equation}
By \eqref{eq:direct-tail-amplitude} with \(t=q_1^2\),
\(E\le z_*^{1/4}q_1^{3/2}\to0\) for \(|q_2|<1/8\), while
\(q_2\log q_1\) may be unbounded.
Let \(b_\sigma=I_\sigma(q_1^2\widehat h_{\rm ref})\) and set
\[
 A_\sigma(p)=
 \frac{\partial_{\widehat h} I_\sigma(q_1^2\widehat h)|_{\widehat h=\widehat h_{\rm ref}}}
      {E S_\sigma'(\widehat h_{\rm ref})}.
\]
On the enlarged compact passing set above,
\begin{equation}
\label{eq:direct-contact-actual-preparation}
 I_\sigma(q_1^2\widehat h)=b_\sigma+
 E A_\sigma(p)[S_\sigma(\widehat h)+R_\sigma(\widehat h,p)],\qquad
 \|R_\sigma\|_{C^2}\longrightarrow0,
\end{equation}
where \(A_\sigma\) and \(A_\sigma^{-1}\) are bounded and positive.
The signs \(\sigma=-\) and \(\sigma=+\) label the negative and positive
branches, as in \eqref{eq:direct-contact-equations}. The coefficient
ratio satisfies
\begin{equation}
\label{eq:direct-contact-slope-quotient}
 A_-/A_+\le \exp(2C_{\rm cone}\sqrt{z_*})+o(1).
\end{equation}
Choose a large fixed matching height \(Z_c>2\widehat h_{\rm ref}\), and put
\[
 D_{\sigma,c}(\widehat h)
 =\mathcal U_\sigma(Z_c;\widehat h_{\rm ref})-\mathcal U_\sigma(Z_c;\widehat h).
\]
At each fixed \(Z_c\), the finite-height analytic flow
\eqref{eq:direct-contact-time} converges, with all state derivatives
through any prescribed finite order, to \eqref{eq:direct-contact-model}. The
scattering estimate therefore gives
\[
 \|D_{\sigma,c}-S_\sigma\|_{C^2}
 \le CZ_c^{-1/2}+C(Z_c)(q_1+|q_2|).
\]
After fixing \(Z_c\) and reducing the parameters, these differences
have bounded $C^2$ norms. Integration of the comparison equation gives
\(\mathcal U_\sigma^0(Z;\widehat h)=B_0Z+O(\sqrt Z+\log Z)\), uniformly on
each fixed compact set of rescaled central heights. Since
\(B_0\in[b_0,\sup B]\) lies in the interior of the common cone, a
sufficiently large \(Z_c\) puts the normalized energies
\(\mathcal U_\sigma^0(Z_c;\widehat h)/Z_c\), and after parameter reduction also
\(\mathcal U_\sigma(Z_c;\widehat h)/Z_c\), in that cone. At \(t=q_1^2Z_c\),
moreover,
\[
 |q_0|/t\le L/Z_c,\qquad q_1/\sqrt t=Z_c^{-1/2}.
\]
Apply \cref{lem:direct-weighted-tail} at \(t=q_1^2Z_c\), retaining
\((z_*/t)^{-2q_2}\). Its centered map is
\(a_{\sigma,c}x+N_\sigma(x)\), where
\(N_\sigma(0)=N_\sigma'(0)=0\),
\(\|N_\sigma\|_{C^4}\le Ce(t)\), and
\(e(t)=q_1\sqrt{Z_c}+Z_c^{-1/2}\). Here
\(x=-D_{\sigma,c}/Z_c=O_{C^2}(Z_c^{-1})\), so Taylor's formula gives
\(\|N_\sigma(x)\|_{C^2}\le Ce(t)/Z_c^2\). Since
\(t(z_*/t)^{-2q_2}/E=Z_c^{1+2q_2}\), the resulting endpoint is
\[
 \frac{I_\sigma(q_1^2\widehat h)-b_\sigma}{E}
 =Z_c^{2q_2}a_{\sigma,c}
       [D_{\sigma,c}(\widehat h)+n_{\sigma,c}(\widehat h)],\qquad
 \|n_{\sigma,c}\|_{C^2}\le Ce(t)/Z_c,
\]
with \(n_{\sigma,c}=n_{\sigma,c}'=0\) at \(\widehat h_{\rm ref}\).
Normalize by the positive reference derivative
\[
 A_\sigma=Z_c^{2q_2}a_{\sigma,c}
 \frac{D_{\sigma,c}'(\widehat h_{\rm ref})}
      {S_\sigma'(\widehat h_{\rm ref})}.
\]
The scattering margin bounds the denominator below, and the preparation
\eqref{eq:direct-contact-actual-preparation} follows with
\[
 \|R_\sigma\|_{C^2}
 \le CZ_c^{-1/2}+C(Z_c)(q_1+|q_2|)
           +C(q_1/\sqrt{Z_c}+Z_c^{-3/2}).
\]
Choose \(Z_c\) large and then \(q_1,|q_2|\) small. With this fixed
\(Z_c\), \(q_2\log Z_c\to0\), so \(A_\sigma^{\pm1}\) are bounded.

Since the normalized energies at \(t\) lie in the common cone,
\eqref{eq:direct-tail-first-slope} gives
\[
 |\log a_{\sigma,c}|\le C_{\rm cone}\bigl(\sqrt{z_*}+Z_c^{-1/2}\bigr).
\]
The factor \(Z_c^{2q_2}\) cancels from \(A_-/A_+\), and
the two reference derivative ratios tend to one in the
successive choices above, proving \eqref{eq:direct-contact-slope-quotient}.

We now combine this preparation with the outer derivative signs on
the same admissible height interval. Let
\(\widehat{\mathcal H}_p=q_1^{-2}\cH_p\) be its rescaling in the
coordinate \(\widehat h\). If \(\widehat{\mathcal H}_p\cap[\widehat h_{\rm a},\widehat h_{\rm b}]\) is empty, connectedness puts
\(\widehat{\mathcal H}_p\) in one outer interval, where the strict derivative sign already
gives at most one zero. Otherwise, the endpoint graphs are defined on
the whole fibre of $\mathcal K_{\rm c}^+$, and every endpoint on this fibre, including
the reference endpoint, is within $O(E)$ of the corresponding
admissible endpoint. By \cref{lem:exterior-multiplier-new}, all these
endpoints lie in \(\widetilde W_\pm\) for small parameters.
By \eqref{eq:direct-contact-slope-quotient} and the choice of \(z_*\),
a reduction of parameters gives
\[
 0<\lambda_0\le\lambda(p):=
 \Pi_p'(b_-)A_-/A_+\le\bar\lambda<1.
\]
Put $G(\widehat h;\widehat q_0,B_0,\lambda)=\lambda S_-(\widehat h)-S_+(\widehat h)$.
Taylor's formula for the exterior map and
\eqref{eq:direct-contact-actual-preparation} give
\begin{equation}
\label{eq:direct-contact-completed-preparation}
 \frac{\Delta(q_1^2\widehat h)}{EA_+}
 =D(p)+G(\widehat h;\widehat q_0,B_0,\lambda(p))+R_\Delta(\widehat h,p),
 \qquad \|R_\Delta\|_{C^2}\to0 .
\end{equation}
Here \(D=(\Pi_p(b_-)-b_+)/(EA_+)\).
The exterior Taylor remainder is
\(O_{C^2}(E)\), since the normalized endpoint functions have
bounded \(C^2\) norms and the exterior map has bounded
derivatives through order four.

Write \(\widehat\Delta=\Delta(q_1^2\widehat h)/(EA_+)\).
Apply \cref{lem:direct-derivative-ratio-count} on each fibre of
\(\mathcal K_{\rm c}\) to \(-\widehat\Delta\), with
\(\phi=S_+\), \(\psi=S_-\) and
\(\lambda\in[\lambda_0,\bar\lambda]\). The uniform margins of
\cref{lem:direct-contact-scattering-ratio} give
\begin{equation}
\label{eq:direct-contact-completed-ratio}
 \left(\frac{\widehat\Delta'}{S_+'}\right)'<0 .
\end{equation}
In increasing \(T\), the large-height, compact and small-height
intervals are ordered by the cuts
\(\widehat h_{\rm b}\) and \(\widehat h_{\rm a}\). The outer signs
exclude slope zeros there, including at the cuts. In the middle,
\eqref{eq:direct-completed-log-ratio} writes
\(\widehat\Delta'/S_+'\) as \(\e^{\mathcal L}-1\) times the positive
factor \(q_1^2I_{+,h}/(EA_+S_+')\). Thus the displayed inequality gives
\(\mathcal L_{\widehat h}<0\), or \(\mathcal L_T>0\), at every slope
zero. All preparations represent the same displacement on the enlarged
windows. The second case of \cref{cor:ordered-cover-zero-count} gives
at most two zeros on the whole interval.

This proves \cref{prop:direct-bounded-contact}, with
\[
 N_{\rm ct}(L)=2.
\]

\subsection{Positive offsets and the interior saddle}
\label{sec:direct-positive-offset}

For positive offsets large relative to \(q_1^2\), rescaling reveals
a uniformly hyperbolic interior saddle. Its exponent determines the
single logarithmic term of \(\mathcal L\) near the separatrices, even
as that exponent crosses one. Continuing to the fixed sections adds
a common, possibly unbounded tail multiplier, which cancels from the
endpoint derivative ratio. Away from the saddle, exterior contraction
makes \(\mathcal L\) negative. Together these estimates control all
its zeros on \(\mathcal H_p\).

\begin{proposition}[Positive offset]
\label{prop:direct-positive-offset}
There is a fixed \(\kappa_{\rm p}>0\) such that the return displacement
has at most two distinct zeros when
\[
 q_0>0,\qquad 0\le q_1/\sqrt{q_0}\le\kappa_{\rm p},
\]
for all sufficiently small parameters.
\end{proposition}

\subsubsection{The saddle scale and its local passage}

Put
\[
 \ell=\sqrt{q_0},\qquad \kappa=q_1/\ell,\qquad
 \eta=\ell X,\quad z=\ell^2S,\quad \tau=\ell t .
\]
The scaled field is
\begin{equation}
\label{eq:direct-positive-scaled-field}
 \begin{aligned}
 X_\tau&=1+\kappa X+q_2X^2+\ell X^3R_q(\ell X)
              -SB(\ell X,\ell^2S),\\
 S_\tau&=-XS.
 \end{aligned}
\end{equation}
On every fixed state box it is
\(O_{C^k}(\kappa+|q_2|+\ell)\) close to
\begin{equation}
\label{eq:direct-positive-comparison}
 X_\tau=1-B_0S,\qquad S_\tau=-XS,\qquad B_0=B(0,0,p).
\end{equation}
The comparison field depends on the independent parameters \(A,C\)
through \(B_0\), which lies in a compact positive interval. Its first
integral is
\begin{equation}
\label{eq:direct-positive-first-integral}
 \mathscr H=B_0S-\log S-X^2/2 .
\end{equation}

The equilibrium of the scaled field is \(X=0,\ S=S_s\), where
\(S_sB(0,\ell^2S_s)=1\). It is unique at positive height
in the box in the coordinates \((\eta,z)\), since \(B+zB_z>0\).
Its Jacobian is
\[
 \begin{pmatrix}
 \kappa-\ell S_sB_\eta(0,\ell^2S_s)&
       -(B+\ell^2S_sB_z)(0,\ell^2S_s)\\
 -S_s&0
 \end{pmatrix}.
\]
The positive quantities \(S_s\), \(\lambda_+\) and \(-\lambda_-\), where
\(\lambda_+>0>\lambda_-\) are its eigenvalues, are bounded
above and below uniformly. Moreover,
\begin{equation}
\label{eq:direct-positive-saddle-ratio}
 \nu_p=-\lambda_-/\lambda_+=1+O(\kappa+\ell).
\end{equation}
For passing central heights, the normalized central height is
\(\widehat h_+=h/q_0>S_s\).

Choose small fixed incoming and outgoing sections in the passing
quadrant of this saddle. On each, choose a positive transverse
coordinate equal to \(\mathscr H-\mathscr H_s\) for the comparison
field, and translate its zero to the perturbed separatrix. The
local passage between these coordinates has the form
\begin{equation}
\label{eq:direct-positive-dulac}
 \begin{gathered}
 D_p(x)=u_px^{\nu_p}(1+r_p(x)),\qquad
 u_p=1+O(\kappa+|q_2|+\ell),\\
 |(x\partial_x)^jr_p(x)|\le C_jx^\vartheta,
 \qquad 0\le j\le2,
 \end{gathered}
\end{equation}
for a fixed \(\vartheta>0\). Parameter-uniform expansions of
hyperbolic Dulac maps are developed in \cite{MV2024}. We derive the
finite-order estimate and its leading coefficient below, uniformly as
\(\nu_p\) crosses one.

We first derive the power law on saddle sections where the nonlinear
term is integrable along the passage. We then transfer it to the
chosen sections, keeping the leading coefficient as well as the
Euler remainder.
The parameter-dependent stable manifold theorem
\cite[Theorem~9.6]{Teschl2012}, applied and normalized in
Appendix~\ref{app:positive-saddle-preparation}, gives on a common
saddle box
\begin{equation}
\label{eq:direct-positive-saddle-preparation}
 x'=x,\qquad y'=-\nu_p y+xy^2c_2(x,y).
\end{equation}
The required state derivatives of \(c_2\) are uniformly bounded
for \(\nu_p\) near one.

From input \((x,y_i)\) to output \(x=x_o\), set
\(\tilde y(\xi;x)=y(\xi)(\xi/x)^{\nu_p}\), and write
\[
 k(t,x)=(x/t)^{\nu_p},\qquad
 \mathscr G(t,x,\tilde y)=k(t,x)\tilde y^2c_2(t,k(t,x)\tilde y).
\]
The solution satisfies the Volterra equation
\[
 \tilde y(\xi;x)=y_i+\int_x^\xi
           \mathscr G(t,x,\tilde y(t;x))\,dt .
\]
For \(3/4\le\nu_p\le5/4\) and fixed \(0<\vartheta<1/2\),
\[
 \int_x^{x_o}k(t,x)\,dt
 \le x^{3/4}\int_x^{x_o}t^{-3/4}\,dt
 \le Cx^{3/4}\le Cx^\vartheta .
\]
Fixed small sections therefore keep
\(y_i/2\le\tilde y\le2y_i\). For \(\mathcal E_x=x\partial_x\),
with \(t,\tilde y\) fixed in the integrand, one has
\(\mathcal E_xk=\nu_pk\), hence
\(\lvert\mathcal E_x^i\partial_{\tilde y}^j\mathscr G\rvert
\le C_{i,j}k\) at every required finite order. Euler differentiation
of the Volterra equation gives scalar linear variational equations
with coefficient \(\mathscr G_{\tilde y}\); their forcing terms are
bounded by \(k\) times lower variations. Moving-lower-limit terms
are \(O(x)\): for the first derivative the term is
\(-xy_i^2c_2(x,y_i)\), and another Euler derivative has the same bound.
Variation of constants and the integral estimate give
\[
 |\tilde y(x_o;x)-y_i|
       +\sum_{j=1}^2|\mathcal E_x^j\tilde y(x_o;x)|\le Cx^\vartheta.
\]
Division by the fixed positive \(y_i\) proves the relative Euler
estimates and the leading power on the prepared sections.

The prepared sections and the chosen incoming and outgoing sections
lie in common regular collars at a fixed positive distance from
the saddle. By \cref{lem:regular-transitions}, the transitions
between them, followed by the changes of section coordinates,
contribute positive units, with remainders \(O(x)\) at input and
\(O(x^{\nu_p})\) at output. They preserve the stated Euler estimates. For
\eqref{eq:direct-positive-comparison} the passage in these coordinates
is exactly the identity, by its first integral. Uniform dependence
of the graph constructions, coordinate changes, and fixed section
maps gives \(u_p=1+O(\kappa+|q_2|+\ell)\).
This proves \eqref{eq:direct-positive-dulac}.

The orbits through central points with \(\widehat h_+\downarrow S_s\)
meet the incoming and outgoing sections in their directed order. The
two eigen-coordinates of these central points are positive multiples
of \(\widehat h_+-S_s\), bounded above and below uniformly. Following
backward to the incoming section decreases the unstable coordinate,
so its incoming value tends uniformly to zero. The endpoint order in
\cref{prop:central-displacement-interval-new}, together with the
positive regular section derivative, makes this rebasing injective.
Thus a sufficiently small preliminary interval
\(0<\widehat h_+-S_s<\rho_0\) maps into a fixed collar \(0<x<x_c\).

\subsubsection{Continuation to the fixed sections}

We next continue the maps to the fixed endpoint windows. This step
must include central heights tending to the saddle, whose local passage
time is unbounded. Choose a large fixed \(M_h\), followed by
\(M_t>2M_h+2\), beyond all saddle heights. For
\(S_s<\widehat h_+\le M_h\), both outward graphs reach \(S=M_t\),
with scaled energies \(\mathcal U_\sigma=X_\sigma^2/2
=U_\sigma/q_0\) bounded above and below there. As
\(\widehat h_+\downarrow S_s\), the local construction above makes
the half-orbits converge to the upper separatrices, whose passage from
the fixed saddle box to \(M_t\) is regular.
For the comparison field their energies are exactly
\[
 \mathcal U_\sigma(M_t;\widehat h_+)
   =B_0(M_t-\widehat h_+)-\log(M_t/\widehat h_+),
\]
including the continuous separatrix value at \(\widehat h_+=1/B_0\).
For large \(M_t\) their ratios to \(M_t\) are close to \(B_0\),
uniformly for \(\widehat h_+\le M_h\). Since
\(B_0\in[b_0,\sup B]\) lies in the interior of the common cone, these
ratios lie in that cone.
Outside the saddle box these segments, including the limiting
separatrices, form a compact regular family with uniformly transverse
endpoints at \(S=M_t\). For a small fixed \(\rho_0>0\), use the local
Dulac construction on \(0<\widehat h_+-S_s<\rho_0\), followed by
\cref{lem:regular-transitions} from the saddle sections. On the compact
complement \(S_s+\rho_0\le\widehat h_+\le M_h\), the same lemma applies
from the central section, where the flux has a positive margin.
These continuations preserve the energy bounds and common-cone entry
at \(M_t\). The regular maps have positive bounded first derivatives
and the required higher derivative bounds; the local Dulac map
accounts for the unbounded time.

At the common cut \(t=M_tq_0\),
\[
 q_0/t=1/M_t,\qquad
 q_1/\sqrt t=\kappa/\sqrt{M_t},\qquad
 \sqrt t=\ell\sqrt{M_t}.
\]
Taking \(M_t\) large and the parameters small permits
\cref{lem:direct-weighted-tail} to continue both graphs to \(z_*\).
Keep the tail multiplier in the common endpoint amplitude:
\begin{equation}
\label{eq:direct-positive-amplitude}
 \varpi_{M_tq_0}=(z_*/(M_tq_0))^{-2q_2},\qquad E=q_0\varpi_{M_tq_0} .
\end{equation}
By \eqref{eq:direct-tail-amplitude} with \(t=M_tq_0\),
\(E\le z_*^{1/4}q_0^{3/4}\to0\) as \(q_0\to0\), uniformly for
\(|q_2|<1/8\).

Let \(e_\sigma\) denote the chosen saddle-section coordinates,
with \(e_+=x\) and \(e_-=D_p(x)\). For the comparison field the first
derivatives of the regular maps from these sections to \(S=M_t\) are
both one in the normalized \(I\) coordinate:
at fixed \(S\), a change of \(\mathscr H\) is exactly a change
of \(S-X^2/2\). Their quotient is consequently
\(1+O_{M_t}(\kappa+|q_2|+\ell)\).
The weighted tail gives the endpoint maps
\begin{equation}
\label{eq:direct-positive-endpoints}
 I_\sigma(e)=b_\sigma+E\mathcal F_\sigma(e),\qquad \mathcal F_\sigma(0)=0,\qquad
 c\le\mathcal F_\sigma'\le C,\qquad \|\mathcal F_\sigma\|_{C^3}\le C,
\end{equation}
and
\begin{equation}
\label{eq:direct-positive-slope-quotient}
 \frac{\mathcal F_-'(0)}{\mathcal F_+'(0)}
 \le \exp\left(2C_{\rm cone}\bigl(\sqrt{z_*}+\kappa/\sqrt{M_t}\bigr)\right)
         [1+O_{M_t}(\kappa+|q_2|+\ell)].
\end{equation}
The tail normalization error is
\(O_{C^3}(\ell\sqrt{M_t}+\kappa/\sqrt{M_t})\), and converting from the
cut scale \(M_tq_0\) to \(q_0\) only changes fixed constants. Both tails
start at the common cut, with normalized energies in the common cone
after parameter reduction; their common multiplier \(\varpi_{M_tq_0}\) is
contained in \(E\), and \eqref{eq:direct-tail-first-slope} bounds the
logarithm of each normalized tail slope by
\(C_{\rm cone}(\sqrt{z_*}+\kappa/\sqrt{M_t})\).

If the saddle collar contains a central height in \(\cH_p\), all its
endpoints differ from the admissible endpoints by \(O(E)\), so they lie
in \(\widetilde W_\pm\) for small parameters, by
\cref{lem:exterior-multiplier-new}.

\subsubsection{The return displacement and the remaining central heights}

In the incoming coordinate \(x\), the closing equation is
\[
 \widetilde\Delta_p(x)=\Pi_p(I_-(D_p(x)))-I_+(x)=0.
\]
Its derivative is
\begin{equation}
\label{eq:direct-positive-completed-derivative}
 \widetilde\Delta_p'(x)=E\mathcal F_+'(x)\,[\exp\mathcal L_p(x)-1],
\end{equation}
where
\[
 \mathcal L_p(x)=
 \log \Pi_p'(I_-(D_p(x)))+\log\mathcal F_-'(D_p(x))
             -\log\mathcal F_+'(x)+\log D_p'(x).
\]
The common amplitude \(E\) has disappeared from this ratio. The local
Dulac derivative supplies the logarithmic term, while the endpoint and
exterior factors contribute a bounded constant and a small remainder.
By \eqref{eq:direct-positive-dulac} and
\eqref{eq:direct-positive-slope-quotient},
\begin{equation}
\label{eq:direct-positive-logarithm}
 \mathcal L_p(x)=c_p+(\nu_p-1)\log x+e_p(x),\qquad
 |e_p|+|xe_p'|\le Cx^\vartheta,
\end{equation}
and, after the fixed choices and parameter reductions,
\[
 c_p=\log\left[
 \Pi_p'(b_-)\frac{\mathcal F_-'(0)}{\mathcal F_+'(0)}u_p\nu_p\right]\le-c_0<0.
\]
For the local factor differentiate
\eqref{eq:direct-positive-dulac}: its relative derivative is
\(1+r_p+(x/\nu_p)r_p'\), with the required Euler bound.
The endpoint factors contribute \(O(x)+O(D_p(x))\);
the nonlinear exterior term contributes \(O(ED_p(x))\).
Their first Euler derivatives have the same bounds, by
\eqref{eq:direct-positive-endpoints}. Decrease \(\vartheta\)
if necessary to obtain \eqref{eq:direct-positive-logarithm}.

Using the full constants in \eqref{eq:direct-positive-logarithm}, choose
\(x_c<\min\{1,e^{-1/\vartheta}\}\) such that
\(Cx_c^\vartheta<c_0/2\) and
\(Cx_c^\vartheta|\log x_c|<c_0/4\).
Decrease the preliminary \(\rho_0\) so that its incoming coordinates
lie in \((0,x_c)\); the preceding estimates restrict to this collar.
Use this final \(\rho_0>0\) in the compact complement below, fixing its
continuation constants before the final parameter reduction.
Then \cref{lem:one-logarithm-zero-count} and the negative bound for \(c_p\)
give \(x\mathcal L_p'(x)<0\) at every zero of \(\mathcal L_p\).
To return to the central height, observe that the ratio agrees with
\eqref{eq:direct-completed-log-ratio}: the factor \(dx/dh\) cancels
from the quotient of endpoint derivatives. Moreover,
\(dx/dh=I_{+,h}/(E\mathcal F_+'(x))>0\). Thus the central heights in \(\cH_p\)
with \(0<\widehat h_+-S_s<\rho_0\) form one interval on which
every zero of \(\mathcal L\) has \(\mathcal L_h<0\), that is,
\(\mathcal L_T>0\).

For the remaining central heights we prove contraction. First, for
\(h\ge M_hq_0\),
\[
 |q_0|/h\le1/M_h,\qquad q_1/\sqrt h\le\kappa/\sqrt{M_h}.
\]
For a sufficiently large fixed \(M_h\) and small \(|q_2|\),
\cref{lem:relative-height-paired-new} therefore gives
\(\Delta_h<0\) at these central heights when \(q_1>0\); when
\(q_1=0\), \cref{prop:signed-contraction-new} gives the same sign.
Second, on \(S_s+\rho_0\le\widehat h_+\le M_h\) the radial flux at
the central point has a fixed positive margin. On its compact inner
portion, the scaled linear integral in \eqref{eq:paired-log-derivative-new}
is bounded by \(C_{\rho_0,M_h}\kappa\), and the other nonquadratic integrals
are \(O_{\rho_0,M_h}(\ell)\). To make the tail constant
independent of these height thresholds, continue first to a fixed
large \(S=S_1\). For the comparison field its energy is
\[
 U/q_0=B_0(S-\widehat h_+)-\log(S/\widehat h_+).
\]
Uniformly on \(S_s+\rho_0\le\widehat h_+\le M_h\), its ratio to
\(S\) tends to \(B_0\). Choose \(S_1\) to enter the interior
of the common cone; the positive radial flux at the central point and
finite-height continuation preserve this property for small
\(\kappa,|q_2|,\ell\). The continuation constants depend on the height
thresholds, while the tail constant is fixed by the common cone. By
\eqref{eq:direct-tail-first-slope}, the two tails from \(S=S_1\)
contribute at most \(2C_{\rm cone}(\sqrt{z_*}+\kappa/\sqrt{S_1})\), and
hence
\[
 \left|\log(I_{-,h}/I_{+,h})\right|
 \le 2C_{\rm cone}\sqrt{z_*}
       +C_{\rho_0,M_h,S_1}(\kappa+\ell).
\]
The fixed common-cone choice of \(z_*\), followed by small
\(\kappa_{\rm p},|q_2|,\ell\),
makes this bound smaller than the exterior contraction margin.
These two ranges join into the single interval of large central
heights \(\widehat h_+\ge S_s+\rho_0\), on which \(\mathcal L\) is
strictly negative, including the shared height thresholds.
In increasing \(T\), this interval and the saddle collar, intersected
with \(\cH_p\), cover \(\cH_p\). Since the endpoints of the collar lie in
\(\widetilde W_\pm\), the collar represents this same displacement.
Hence every zero of \(\mathcal L\) on \(\cH_p\) lies in the collar and
crosses upward, and the second case of
\cref{cor:ordered-cover-zero-count} gives at most two zeros of
\(\Delta_p\), proving \cref{prop:direct-positive-offset}.

\section{Early heights, large exits, and weak roots}
\label{sec:recurring-ranges}

For negative constant coefficient, the height interval may extend
from nearly regular passages to long passages near a boundary root.
We first establish contraction at early heights, then compare the
branches while both drifts remain positive. The other two estimates
control positive half-orbits that leave every fixed multiple of the
coefficient scale and boundary roots that approach the central line
relative to that scale. The applications in
\cref{sec:primary-regimes,sec:quadratic-regimes,sec:direct-negative-annulus-buffer,sec:direct-negative-root-regimes}
verify the hypotheses on the primary and secondary scales and join
the resulting intervals for the same logarithmic slope ratio.

\subsection{Early central heights}
\label{sec:early-central-heights}

The first part of the height interval is controlled by exterior
contraction. Besides heights above a coefficient-dependent threshold,
it includes smaller central heights whose positive half-orbits still
have small normalized radius at that threshold. The estimates below
keep the tail error independent of the threshold, so that the common
fixed sections suffice for later applications. We use the scales
\begin{equation}
\label{eq:direct-quadratic-prefix-scales}
 q_0=-a\delta\ell^2,\qquad q_1=b\delta\ell,\qquad
 q_2=-\bar c\delta,\qquad \epsilon_3=\ell/\delta\le1,
\end{equation}
where \(a\ge a_{\min}>0,\ b\ge0,\ \bar c\) lie in a fixed compact
coefficient set and \(\delta>0\) is sufficiently small. Put
\(Q=\delta\ell^2\). In \(x=|\eta|/\ell\) the outward drifts are
\[
 K_\sigma(x)=-q(\sigma\ell x)/(\delta\ell^2)
 =a-\sigma bx+\bar c x^2-\sigma\epsilon_3 x^3R_q(\sigma\ell x).
\]

\begin{lemma}[Early central heights]
\label{lem:early-central-heights}
Fix \(d_f>0\) and put \(z_e=d_fQ\). There is \(r_0>0\), depending only
on \(d_f\) and the compact coefficient set, such that for all
sufficiently small \(\delta\) the central heights in \(\cH_p\) with
\begin{equation}
\label{eq:direct-primary-early}
 h\ge z_e,\qquad\hbox{or}\qquad
 0<h<z_e,\quad |\eta_+(z_e;h)|/\ell\le r_0,
\end{equation}
form an initial interval in every logarithmic height coordinate
\(T=-c_T\log h\), and
\(\Pi_p'(I_-)I_{-,h}/I_{+,h}<1\) on it. In particular
\(\mathcal L<0\) there. The coefficient of \(\sqrt{z_*}\) in the
estimates is \(2C_{\rm cone}\); the constants that depend on
\(d_f\) and \(r_0\) enter only the final threshold for \(\delta\).
\end{lemma}

\begin{proof}
By \cref{lem:central-height-order}, the set
\eqref{eq:direct-primary-early} is an initial interval in \(T\).
There are three ranges: large heights, the intermediate height band,
and smaller heights with small radius at the prescribed cut.
The large-height estimate handles the first. In the other two, we bound the paired
logarithm \(\log(I_{-,h}/I_{+,h})\) of
\cref{lem:paired-variation-identity-new}. The common-cone tail
contributes \(2C_{\rm cone}\sqrt{z_*}\), already below a fixed
fraction of the exterior contraction margin. The remaining errors are
controlled by choosing the radius cut small where needed, and then
reducing \(\delta\).

\emph{Large heights.} Let \(\theta_0\) be the constant of
\cref{lem:relative-height-paired-new}, and fix \(K\ge2a_{\max}/\theta_0\).
For \(h\ge KQ\),
\[
 |q_0|/h\le a/K,\qquad q_1/\sqrt h\le b\sqrt{\delta/K},\qquad
 |q_2|\le C\delta,
\]
so \cref{lem:relative-height-paired-new} gives
\(0<L_p'<1\) for small \(\delta\).

\emph{The band \(d_fQ\le h\le KQ\).} Set \(z=Q\zeta\), \(U=Q\mathcal U\).
Since \(q_0<0\), the constant term of the graph equation is
favorable, and the centered lower barrier \(U_\sigma\ge m(z-h)\) holds
for a small fixed \(m\). On a fixed interval ending at
\(\zeta=K_1>2K\), the graph equation converges uniformly, with its
finite variations, to
\[
 \mathcal U_\zeta=B_0+a/\zeta,\qquad
 \mathcal U(\zeta)=B_0(\zeta-h/Q)+a\log(\zeta Q/h).
\]
Choose \(K_1\) large after \(d_f,K\). Since
\(\log(K_1/d_f)/K_1\to0\), the energy ratio at \(K_1Q\) lies in the
interior of the common cone, uniformly for \(d_f\le h/Q\le K\).
Before this height, the positive flux and the centered lower barrier
bound the three paired integrals by
\(C_{d_f,K_1}(\sqrt Q+q_1/\sqrt Q)=C_{d_f,K_1}(\sqrt Q+b\sqrt\delta)\).
Their square-root singularities at \(\zeta=h/Q\) are integrable.
Beyond this height, \cref{lem:direct-weighted-tail} applies on the
common cone. Combining the two estimates gives
\begin{equation}
\label{eq:direct-primary-early-uniform-section}
 \left|\log(I_{-,h}/I_{+,h})\right|
 \le 2C_{\rm cone}\sqrt{z_*}+C_{d_f,K_1}\sqrt\delta .
\end{equation}

\emph{Small radius at \(z_e\).} Write
\(x_\sigma=|\eta_\sigma(z_e;h)|/\ell\). Take a fixed large \(L_0\), and
then \(r_0\) so small that both \(K_\sigma\) lie between \(a/2\) and
\(2a\) on \(0\le x\le L_0r_0\). Up to a hypothetical first hit
of \(L_0r_0\) by the negative branch, the graph equations give
\(\mathcal G_{-,p}(U_-,z)\le C\mathcal G_{+,p}(U_+,z)\): compare
\(\sup B+2a\delta\ell^2/z\) with \(b_0+a\delta\ell^2/(2z)\), the
supremum being taken on the fixed box. Integration from the common zero
energy gives \(U_-\le CU_+\), and \(L_0^2>C\) excludes that hit.
Thus \(x_+\le r_0\) implies \(x_-\le L_0r_0\).

On the prefix up to \(z_e\), use
\(dz/d|\eta|=|\eta|z/(zB-q)\). The lower bounds
\(zB-q\ge a_{\min}\delta\ell^2/2\) and \(zB-q\ge b_0z\) give
\[
 \int_h^{z_e}\frac{dz}{|\eta_\sigma|}\le C\ell r_0,\qquad
 q_1\int_h^{z_e}\frac{dz}{|\eta_\sigma|z}\le Cbr_0,\qquad
 \int_h^{z_e}\frac{|\eta_\sigma|}{z}\,dz\le C\epsilon_3 r_0^3 .
\]
From \(z_e\) onward there are barriers
\begin{equation}
\label{eq:direct-primary-shifted-cones}
 m(z-z_e)\le U_\sigma(z)
               \le U_\sigma(z_e)+M(z-z_e).
\end{equation}
The initial adverse ratios
\(q_1\sqrt{U_\sigma(z_e)}/z_e\), \(|q_2|U_\sigma(z_e)/z_e\) and
\(U_\sigma(z_e)^{3/2}/z_e\) are \(O(r_0/d_f)\), \(O(r_0^2/d_f)\) and
\(O(\epsilon_3 r_0^3/d_f)\). Choose \(r_0\) small after \(d_f\).
Use the strip between the barriers only up to \(t=\ell^2Z_1\),
with fixed \(Z_1\). Its added linear, quadratic and cubic losses are
\(O(\sqrt\delta)\), \(O(\delta M)\) and \(O(M^{3/2}\ell\sqrt{Z_1})\),
so small \(\delta\) gives positive flux.

These shifted barriers preserve positive flux, but the endpoint
estimate also requires entry into the common cone. For this, apply
the continuation part of
\cref{lem:direct-small-height-continuation} with \(\zeta_0=d_f\) and
\(0\le \mathcal U_0=U_\sigma(z_e)/\ell^2\le Cr_0^2\).
The preceding radius bounds give
\(K_\sigma(\sqrt{2\mathcal U_0})\ge a_{\min}/2\).
Choose the height \(Z_1\) for entry into the cone large after
\(r_0,d_f\). The shifted barriers
bound the paired integrals up to this height by
\(C_{d_f,r_0,Z_1}\{\sqrt\delta+\ell(1+|\log\delta|)\}\):
the linear integral is \(O(\sqrt\delta)\), and splitting the square
root in the upper barrier bounds the cubic integral. This tends to zero
uniformly for \(0<\epsilon_3\le1\), since \(\ell=\epsilon_3\delta\).
The remaining common-cone tail contributes
\(2C_{\rm cone}\sqrt{z_*}+O(\delta)\). Together with the prefix
bounds, this proves contraction for small \(r_0\) and then small
\(\delta\).
\end{proof}

\subsection{A common-height estimate before a boundary root}
\label{sec:direct-quadratic-prefix}

Beyond the contraction range, suppose both drifts still have positive
margins. Equal leading actions then relate the two radii at the same
central height. The derivative of this relation has a useful
cancellation: the quadratic coefficients drop out of the derivative
ratio. This determines the crossing direction of \(\mathcal L\)
before the positive branch reaches a boundary root, for either sign
of the quadratic coefficient.

Suppose the scales \eqref{eq:direct-quadratic-prefix-scales} hold with
\(\delta\) and \(\epsilon_3=\ell/\delta\) sufficiently small. The drifts
\(K_\sigma\) of \cref{sec:early-central-heights} have the quadratic
comparisons
\[
 K_\sigma^0(x)=a-\sigma bx+\bar c x^2.
\]
Choose a positive terminal radius \(R\), possibly varying
over the compact coefficient set, such that \(K_+^0\)
has a fixed positive margin on \(0\le x\le R\).
Assume also that the equal-action negative radii for these
terminal actions lie in a fixed compact interval on which
\(K_-^0\) has a fixed positive margin.

On these comparison intervals set
\[
 \mathcal A_\sigma^0(x)=\int_0^x\frac{u\,du}{K_\sigma^0(u)},\qquad
 \mathcal A_-^0(y)=\mathcal A_+^0(x),\qquad y=f(x).
\]
The drift margins make \(\mathcal A_-^0\) strictly increasing, so its
inverse at the prescribed actions defines \(f\) on the chosen compact
intervals.

\begin{samepage}
\begin{lemma}[Equal-action comparison]
\label{lem:direct-equal-action-comparison}
For \(x,y=f(x)>0\) in these intervals, put
\(R_1(x)=K_-^0(y)/K_+^0(x)\). Then
\begin{equation}
\label{eq:direct-quadratic-critical-identity}
 \frac{yy'}x=R_1,\qquad
 \frac{R_1'}{R_1}
       =\frac{b(1+x/y)}{K_+^0(x)},\qquad
 \frac{(D+x^2-\lambda y^2)'}{2x}=1-\lambda R_1
\end{equation}
for any constants \(D,\lambda\). In particular, \(R_1\) is
nondecreasing when \(b\ge0\).
\end{lemma}

\begin{proof}
Differentiate \(\mathcal A_-^0(y)=\mathcal A_+^0(x)\) to obtain the first identity.
Since \((K_-^0)'(y)=b+2\bar c y\) and
\((K_+^0)'(x)=-b+2\bar c x\), logarithmic differentiation
of \(R_1\) and substitution of the first identity cancel
the two quadratic terms and give the second. The third
follows by differentiating \(D+x^2-\lambda y^2\).
\end{proof}
\end{samepage}

\begin{lemma}[Quadratic prefix]
\label{lem:direct-quadratic-prefix}
Fix a small \(d_f>0\) and put \(z_e=\delta\ell^2d_f\).
On the central heights in \(\cH_p\) satisfying
\[
 h\ge z_e,\qquad\hbox{or}\qquad
 0<h<z_e,\quad |\eta_+(z_e;h)|/\ell\le R,
\]
the return displacement has at most two distinct zeros.
Moreover, for any logarithmic height coordinate \(T=-c_T\log h\),
where \(c_T>0\)
is independent of the central height, every zero of
\[
 \mathcal L(T)=\log\!\left[
       \Pi_p'(I_-(h))\frac{I_{-,h}(h)}{I_{+,h}(h)}\right]
\]
on this interval satisfies \(\mathcal L_T>0\).
The constants are uniform on the compact coefficient and
action domains just specified.
\end{lemma}

\begin{proof}
By \cref{lem:early-central-heights}, there is \(r_0\in(0,R)\) such that
the central heights with \(h\ge z_e\), or with \(h<z_e\) and
\(x_+=|\eta_+(z_e;h)|/\ell\le r_0\), form an initial interval in \(T\) on
which the return contracts; in particular \(\mathcal L<0\) there.

\emph{Matching the branches.}
It remains to consider \(r_0\le x=x_+\le R\). On the prescribed
compact drift intervals, the central-height identity
\eqref{eq:direct-terminal-radius-jets}, with \(Q=\delta\ell^2\), gives
\[
 \delta\log(h_\sigma(x)/Q)
 =-\mathcal A_\sigma(x)+\delta M_\sigma(x),\qquad
 \mathcal A_\sigma(x)=\int_0^x\frac{u\,du}{K_\sigma(u)},\qquad
 \|M_\sigma\|_{C^4}\le C .
\]
The fixed terminal height here is \(z_e/Q=d_f\).

The leading action relation must first give two passages with exactly
the same central height. The assumed margins and compact negative
inverse provide this matching: on the relevant terminal interval,
\((\mathcal A_-^0)'(y)\) is bounded below. At the endpoints of a slightly
enlarged negative radius
interval, the trial passages have central heights bracketing the
positive passage's central height. The order of central heights
(\cref{lem:central-height-order}) gives a matching negative passage.
Implicit differentiation then gives
\begin{equation}
\label{eq:direct-quadratic-action-relation}
 x_-(x)=f(x)+O_{C^3}(\epsilon_3+\delta).
\end{equation}

\emph{Preparing the endpoints.}
Apply \cref{lem:direct-small-height-continuation} with the scales
\eqref{eq:direct-quadratic-prefix-scales},
\(\mathcal U_0=x_\sigma^2/2\) and the constant entrance height factor
\(\zeta_0=d_f\). Both radius intervals are compact and positive, and
the assumed drift margins persist for small \(\epsilon_3\).
Choose the common-cone height \(Z_1=Z_0\ge1\). The endpoints are
\[
 I_\sigma(x_\sigma)=b_\sigma-\tfrac12EA_\sigma x_\sigma^2
                        +Ee_\sigma(x_\sigma),\qquad
 A_\sigma\asymp1,\quad
 \|e_\sigma\|_{C^3}\le C\delta(1+|\log\delta|),
\]
with \(A_-/A_+\le\exp(2C_{\rm cone}\sqrt{z_*}+o(1))\).
Here \(E=\ell^2(z_*/(\ell^2Z_0))^{-2q_2}\); by
\eqref{eq:direct-tail-amplitude} with \(t=\ell^2Z_0\),
\(E\le z_*^{1/4}\ell^{3/2}\), since \(Z_0\ge1\) and
\(|q_2|<1/8\) for small \(\delta\).

Whenever the interval contains a central height in \(\cH_p\), the
endpoints are within \(O(E)\) of admissible endpoints and therefore
lie in \(\widetilde W_\pm\), by \cref{lem:exterior-multiplier-new}.
Composing with the exterior transition gives a Taylor remainder of
\(O(E)\) after division by \(E\). Inserting
\eqref{eq:direct-quadratic-action-relation} gives
\begin{equation}
\label{eq:direct-quadratic-completed-model}
 \frac{2\Delta(x)}{EA_+}
 =D+x^2-\lambda f(x)^2+R_\Delta(x),\qquad
 \|R_\Delta\|_{C^3}
       \le C\{\epsilon_3+\delta(1+|\log\delta|)+E\},
\end{equation}
where \(D\) depends only on the parameters and
\(0<\lambda_0\le\lambda\le\bar\lambda<1\).

\emph{Controlling the crossing direction.}
Apply \cref{lem:direct-derivative-ratio-count} on \(r_0\le x\le R\) to
\(\mathcal F(x)=2\Delta(x)/(EA_+)\), with \(\phi=x^2\), \(\psi=f^2\),
and the \(\lambda\in[\lambda_0,\bar\lambda]\) and remainder
\(R_\Delta\) of \eqref{eq:direct-quadratic-completed-model}; put
\(\mathscr Q=\mathcal F'/(2x)\). Here \(\phi'=2x\ge2r_0\),
\(\phi''=2\), and \(\psi'/\phi'=R_1\) by
\eqref{eq:direct-quadratic-critical-identity}.

At \(b=0\), equal actions give \(y=x\) and \(R_1=1\).
Uniform smooth dependence on the compact drift domains gives
\(|R_1-1|\le Cb\). Fix \(b_*>0\) so small that
\(1-\lambda R_1\ge(1-\bar\lambda)/2\) for \(0\le b\le2b_*\).
The second case of the lemma gives \(\mathscr Q>0\), and hence
\(\Delta_x>0\) in this range. For \(b\ge b_*\),
\cref{lem:direct-equal-action-comparison} gives \(R_1'\ge c_*>0\),
so the first case of the derivative-ratio lemma gives
\begin{equation}
\label{eq:direct-quadratic-weighted-derivative}
 \mathscr Q'<0 .
\end{equation}
These bounds and cuts are fixed before \(\delta,\epsilon_3\) are reduced
to meet the remainder threshold.

To express these signs in the common logarithmic coordinate, use the
positive central-height formula:
\[
 \frac{h_x}{h}=-\frac{\mathcal A_+'(x)}\delta+M_+'(x)<0.
\]
Thus \(x_T>0\) for every such coordinate \(T\). The contraction
interval of large central heights and small radii has
\(\mathcal L<0\), including the cut \(x=r_0\). For \(b\le2b_*\),
\(\Delta_x>0\) and \(h_x<0\) give \(\mathcal L<0\) on the
remaining interval as well. For \(b\ge b_*\), a zero of
\(\mathcal L\) in the compact-radius part has \(\mathcal F'=0\).
By \eqref{eq:direct-quadratic-weighted-derivative},
\[
 \Delta_{xx}=EA_+x\mathscr Q'<0,\qquad
 \Delta_{TT}=\Delta_{xx}x_T^2<0.
\]
With \(m_+=hI_{+,h}>0\), the exact identity
\[
 \Delta_T=-\frac{m_+}{c_T}(\e^{\mathcal L}-1)
\]
therefore gives
\begin{equation}
\label{eq:direct-quadratic-prefix-zero-orientation}
 \mathcal L=0\quad\Longrightarrow\quad
 \mathcal L_T=-\frac{c_T}{m_+}\Delta_{TT}>0.
\end{equation}

Finally, the central heights considered here form one interval in
\(\cH_p\), by \cref{lem:central-height-order}. In increasing \(T\),
the contraction interval and the compact-radius part cover it.
The two estimates apply to the same displacement and include their
shared cut. The compact-radius estimate also includes the terminal cut
\(x=R\), where the endpoints lie in \(\widetilde W_\pm\) and the
exterior map is smooth.
Thus every zero of \(\mathcal L\) on this interval satisfies
\(\mathcal L_T>0\). The second case of
\cref{cor:ordered-cover-zero-count} gives the two-zero bound.
\end{proof}

The compact comparison ends before the positive drift vanishes.
The next lemma gives a strict overlap with the boundary-flow estimate
near that root. It then combines both estimates for the same
logarithmic slope ratio.

\begin{lemma}[Joining the early and root intervals]
\label{lem:prefix-root-cover}
Under \eqref{eq:direct-quadratic-prefix-scales}, suppose \(R\) ranges
over a compact positive interval
and \(K_+^0\) has a fixed positive margin on \(0\le x\le R\).
The negative equal-action radii then lie in a fixed compact interval
where \(K_-^0\) has a positive margin, as required in
\cref{lem:direct-quadratic-prefix}.

Suppose further that \(K_+\) has a first positive root \(\alpha\) in a
fixed compact positive interval. Choose \(0<w_a<R<\alpha\), with fixed
positive lower bounds for \(w_a\) and \(R-w_a\), and a fixed small
\(Y_f>0\) such that the inverse boundary-flow coordinate is defined on
\([w_a,\alpha]\) up to \(Y_f\). Fix a small \(d_f>0\), put
\(z_e=\delta\ell^2d_f\), and use a logarithmic height coordinate \(T\)
as in \eqref{eq:direct-source-clock-convention}. For sufficiently small
\(\delta,\epsilon_3\), the early heights in \(\cH_p\)
\[
 h\ge z_e,\qquad\hbox{or}\qquad
 h<z_e,\quad |\eta_+(z_e;h)|/\ell\le R
\]
form an initial interval in \(T\); every zero of \(\mathcal L\) there
satisfies \(\mathcal L_T>0\). Every remaining height has its positive
projected state at \(z=\ell^2Y_f\) in \([w_a,\alpha)\), with a strict overlap at the
early cut. If every critical point of \(\mathcal L\) in this root range
satisfies \(\mathcal L_{TT}<0\), then \(\Delta_p\) has at most three distinct
zeros on \(\cH_p\); if \(\mathcal L_T>0\) at every zero there, it has at
most two.
\end{lemma}

\begin{proof}
The fixed radius and drift margins bound the positive terminal actions.
To obtain a compact negative inverse, note that
for \(\bar c\ge0\), the action of \(a+by+\bar c y^2\)
increases to infinity; for \(\bar c<0\), it diverges at the first
positive root. Its inverse at each bounded positive action therefore
lies strictly inside the positive-drift interval. Smooth dependence
at a larger fixed action and compactness of the coefficient set give
the claimed compact interval. Thus
\cref{lem:direct-quadratic-prefix} applies to the early interval.

For every remaining height, the positive graph hits \(R\) below
\(Y=\delta d_f\). By \eqref{eq:direct-projection-energy}, its projected
state there lies between \(R-C\delta d_f>w_a\) and \(R<\alpha\).
\Cref{rem:direct-projected-continuation} keeps it increasing and below
\(\alpha\), and continues the inverse projection to \(Y_f\), also at a double
root. The same strict bounds hold at the early cut and give an overlap.
\Cref{lem:central-height-order} orders the two intervals, and
\cref{cor:ordered-cover-zero-count} gives the stated zero bounds.
\end{proof}

\subsection{Positive passages beyond a coefficient scale}
\label{sec:direct-adaptive-positive}

The bounded-radius estimates do not control positive graphs that leave
every fixed multiple of the coefficient scale. To treat these graphs
in the primary and secondary regions, we follow them within a small
unscaled radius box. An adapted height coordinate gives exact action
transport and estimates for the moving exit radius. Continuation from
that exit to a cone then permits the weighted-tail estimate to give
the positive endpoint slope. The corollary below compares it with the
negative endpoint slope to prove contraction after a sufficiently
large exit.

Write \(\varrho=\eta>0\) for the positive radius; in the negative
quadratic-drift application below, \(\varrho=s\rho\). Suppose that
\begin{equation}
\label{eq:direct-adaptive-drift}
 -q(\varrho)=\varrho^2(s+\varrho)\mathcal K(u,v,\varrho,p),\qquad
 u=\varrho_0/\varrho,\quad v=\varrho/(s+\varrho),\qquad s,\varrho_0\ge0,
\end{equation}
where \(\mathcal K\) has fixed positive upper and lower bounds
and uniform derivative bounds on an enlarged compact box of
the displayed base variables. Both \(s\) and the unscaled radius
are small. The coefficient \(\varrho_0\) is held fixed along a graph.
A restriction of the box in \(v\) is handled by a smooth
extension preserving positivity.
The factor \(\varrho^2(s+\varrho)=s\varrho^2+\varrho^3\) joins the
quadratic range \(\varrho\ll s\) to the cubic range
\(\varrho\gg s\).

Put
\[
 \varrho^+=s+\varrho,\qquad y=\frac{z}{\varrho^2\varrho^+},\qquad \mathfrak l=\log\varrho,
 \qquad \frac{d\tau}{d\mathfrak l}=\frac1{\varrho^+},\qquad
 \mathcal B=B(\varrho,\varrho^2\varrho^+y,p).
\]
The graph equation gives the extended system
\begin{equation}
\label{eq:direct-adaptive-system}
 \begin{aligned}
 u_\tau&=-(s+\varrho)u,& \varrho_\tau&=(s+\varrho)\varrho,& v_\tau&=sv,\\
 y_\tau&=\mathfrak a\,y+y^2\mathcal N,&
 \mathfrak a&=\mathcal K^{-1}-2s-3\varrho,&
 \mathcal N&=-\frac{\mathcal B}{\mathcal K(\mathcal K+\mathcal B y)} .
 \end{aligned}
\end{equation}
The defining relations among the base variables are invariant.
For a small fixed \(y_0>0\), the quantities
\(\mathfrak a,\ y_\tau/y,\ \partial_y y_\tau\) have fixed
positive lower bounds on \(|y|\le2y_0\).

\begin{lemma}[Adaptive height and cone landing]
\label{lem:direct-adaptive-positive}
For every fixed finite order \(k\), after reducing \(s\) and the
radius box, there is a uniformly regular coordinate
\[
 \Psi(u,v,\varrho,y,p)=y\exp\Theta(u,v,\varrho,y,p),\qquad
 \Theta=y\widetilde\Theta,\quad
 \|\widetilde\Theta\|_{C^k}\le C_k,
\]
with positive height derivative and bounded inverse derivatives,
such that \(\Psi_\tau=\mathfrak a\Psi\).
The adapted unscaled height \(\widehat z=\varrho^2(s+\varrho)\Psi\)
therefore satisfies exactly
\begin{equation}
\label{eq:direct-adaptive-action}
 \widehat z(\varrho)=
 \exp\left(\int_{\varrho_i}^{\varrho}\frac{t\,dt}{-q(t)}\right)
 \widehat z(\varrho_i)
\end{equation}
as long as the passage is thin, \(y\le y_0\).

For an exit at \(y=y_0\), let \(\varrho_e\) be its radius and
\(\lambda_{\rm in}=\log\widehat z(\varrho_i)\). Then
\begin{equation}
\label{eq:direct-adaptive-exit}
 -C(s+\varrho_e)\le\frac{d\log\varrho_e}{d\lambda_{\rm in}}
                 \le-c(s+\varrho_e),\qquad
 \left|\frac{d^j\log\varrho_e}{d\lambda_{\rm in}^j}\right|
                 \le C_j(s+\varrho_e)^j,\quad 1\le j\le k .
\end{equation}
There is a subsequent first landing on \(z=m_{\rm land}\varrho^2\), for a
small fixed \(m_{\rm land}>0\), with radius \(\varrho_l\) satisfying
\begin{equation}
\label{eq:direct-adaptive-landing}
 \begin{gathered}
 \varrho_e\le\varrho_l\le2\varrho_e,\qquad
 \partial_{\log\varrho_e}\log\varrho_l=1+O(\epsilon_e)>c>0,\\
 \bigl|\partial_{\log\varrho_e}^j\log\varrho_l\bigr|\le C_j\epsilon_e
       \quad(2\le j\le k),\qquad
 \epsilon_e=m_{\rm land}+(s+\varrho_e)|\log(s+\varrho_e)|,
 \end{gathered}
\end{equation}
with constants independent of \(m_{\rm land}\).
If a graph has \(y\ge y_0\) and \(z<m_{\rm land}\varrho^2\) at a radius
\(\bar\varrho\) of the box, its first landing after \(\bar\varrho\) has
radius at most \(2\bar\varrho\).
The construction applies to admissible graphs before their
endpoint.
\end{lemma}

\begin{proof}
\emph{The adapted height coordinate.}
At backward time \(t\ge0\), the base flow of
\eqref{eq:direct-adaptive-system} is explicitly
\[
 g_s(t)=\frac{1-\e^{-st}}s,\quad
 \varrho_b(t)=\frac{\varrho \e^{-st}}{1+\varrho g_s(t)},\quad
 u_b(t)=u\e^{st}(1+\varrho g_s(t)),\quad v_b(t)=v\e^{-st},
\]
where \(g_0(t)=t\). Since \(0\le g_s(t)\le t\) and the
denominator is at least one, each fixed-order state derivative
is bounded by a polynomial in \(t\) times \(\e^{jst}\).
The bounds are uniform for initial base points in the original compact
box. To cover backward trajectories that leave it, extend the
coefficients by smooth bounded retractions of the base variables into
the enlarged positive box, preserving the positive height-rate and
height-derivative margins. The extension agrees with the original
coefficients and makes \(\mathfrak a,\mathcal N\), together with their
required derivatives, bounded on the entire nonnegative base domain
and a small negative-height strip. Along the backward trajectories,
\(\varrho\) decreases.

Choose \(c_y>0\) so that \(\partial_y y_\tau\ge2c_y\)
on the extended box, and reduce \(s\) so that
\((k+1)s<c_y\). The base-flow bounds above verify
\cref{lem:direct-contracting-height-coordinate} with \(c_x=s\).
It supplies \(\Theta\) with the asserted bounds on
\(\Theta/y\) and the inverse, and \(\Psi_\tau=\mathfrak a\Psi\).
The logarithmic derivative of \(\varrho^2(s+\varrho)\) in \(\tau\) is
\(2s+3\varrho\), so
\((\log\widehat z)_{\mathfrak l}=1/(\varrho^+\mathcal K)\).
This is \eqref{eq:direct-adaptive-action}, by
\eqref{eq:direct-adaptive-drift}.

\emph{The exit radius.}
At the thin exit, \(\mathfrak l_e=\log\varrho_e\) solves
\[
 \mathcal M(\mathfrak l_e)=\lambda_{\rm in},\qquad
 \mathcal M(\mathfrak l)=
 \log[\varrho^2(s+\varrho)\Psi(u,v,\varrho,y_0,p)]
          -\int_{\varrho_i}^{\varrho}\frac{t\,dt}{-q(t)} .
\]
The transport equation gives
\[
 \mathcal M'(\mathfrak l)
 =-\frac{f_0}{\varrho^+}\,\partial_y\log\Psi(u,v,\varrho,y_0,p)
 =-\frac{g(u,v,\varrho,p)}{\varrho^+},\qquad c\le g\le C,
\]
where \(f_0=\mathfrak a y_0+y_0^2\mathcal N(u,v,\varrho,y_0,p)\) is the
height velocity at the exit. Fixed logarithmic radial derivatives of \(g\) are
bounded. The corresponding base vector field is
\(-u\partial_u+v(1-v)\partial_v+\varrho\partial_\varrho\), whose
fixed iterates applied to \(\varrho^+\) are \(O(\varrho^+)\). Set
\(\phi_0=-\varrho^+/g\), so that \(\partial_{\mathfrak l}^j\phi_0=O_j(\varrho^+)\).
Inverse differentiation gives
\[
 \frac{d^j\mathfrak l_e}{d\lambda_{\rm in}^j}
 =\bigl[(\phi_0\partial_{\mathfrak l})^{j-1}\phi_0\bigr]_{\mathfrak l=\mathfrak l_e}.
\]
Each expanded term contains exactly \(j\) factors, each a
radial derivative of \(\phi_0\) of some nonnegative order.
It is therefore \(O_j((s+\varrho_e)^j)\). The sign and two-sided
bound at first order follow from \(c\le g\le C\), proving
\eqref{eq:direct-adaptive-exit}.
Since the thin height increases and scalar order holds, the
exit is unique and its existence set is one
interval of incoming heights.

\emph{Passage to the cone.}
After exit use \(\mathcal Y=z/\varrho^2=\varrho^+y\) as the independent height.
The equations become
\[
 \mathcal Y_{\mathfrak l}=\frac{\mathcal Y}{B\mathcal Y+\varrho^+\mathcal K}-2\mathcal Y,\qquad
 \frac{d\mathfrak l}{d\mathcal Y}
 =\frac{\varrho^+\mathcal K+B\mathcal Y}
        {\mathcal Y(1-2\varrho^+\mathcal K-2B\mathcal Y)} .
\]
Choose \(m_{\rm land}\) small and then the radius box and \(s\) small,
so that the last denominator factor is at least \(1/2\)
on \(\mathcal Y\le m_{\rm land}\). Starting at \(\mathcal Y_e=(s+\varrho_e)y_0\), bootstrap
\(\varrho\le2\varrho_e\). With \(\varrho_e^+=s+\varrho_e\), integration gives
\[
 0\le\log(\varrho_l/\varrho_e)
 \le C\{m_{\rm land}+\varrho_e^+\log(m_{\rm land}/(y_0\varrho_e^+))\}<\log2 .
\]
This closes the bootstrap and proves arrival at \(\mathcal Y=m_{\rm land}\).
For the derivatives, write \(x=\log\varrho_e\) and
\(\mathfrak l(\mathcal Y;x)=\log\varrho(\mathcal Y;x)\). Denote by \(R(\mathfrak l,\mathcal Y)\) the
right side of the equation for \(\mathfrak l\) in the variable \(\mathcal Y\), with
parameters fixed, so that
\[
 \mathfrak l_{\mathcal Y}=R(\mathfrak l,\mathcal Y),\qquad \mathfrak l(\mathcal Y_e(x);x)=x.
\]
On the preceding bootstrap range,
\[
 |\partial_{\mathfrak l}^jR|\le C_j(\varrho_e^+/\mathcal Y+1),\qquad
 \int_{\mathcal Y_e}^{m_{\rm land}}(\varrho_e^+/\mathcal Y+1)\,d\mathcal Y
       \le C\epsilon_e,\quad
 \epsilon_e=m_{\rm land}+\varrho_e^+|\log \varrho_e^+|.
\]
Since the last denominator factor is at least \(1/2\), the constants
\(C_j\) depend only on the bounds for \(\mathcal K\) and \(B\), and not on
\(m_{\rm land}\).
The moving initial height is treated as in the moving-entrance part
of \cref{lem:direct-weighted-tail}: near the entrance set
\(\mathcal Y=\mathcal Y_e(x)\zeta\) and
\(\mathfrak l=x+\tilde{\mathfrak l}\). The rescaled equation has
right side \(\mathcal Y_eR\), whose required mixed derivatives are
\(O_j(\varrho_e^+)\) on a fixed \(\zeta\)-interval. Conversion back to
fixed \(\mathcal Y\) uses
\(\partial_x|_{\mathcal Y}=\partial_x|_\zeta-v_e\zeta\partial_\zeta\),
where \(v_e=\varrho_e/\varrho_e^+\) has bounded derivatives. Thus the
initial data for \(P_j=\partial_x^j\mathfrak l|_{\mathcal Y}\) are
\[
 P_1(\mathcal Y_e)=1-R(x,\mathcal Y_e)y_0\varrho_e
       =1+O(\varrho_e^+),\qquad
 P_j(\mathcal Y_e)=O_j(\varrho_e^+)\quad(j\ge2).
\]
Their triangular variational equations have linear coefficient
\(R_{\mathfrak l}\) and higher forcing coefficients
\(\partial_{\mathfrak l}^mR\). The displayed integral bound and
variation of constants therefore give
\(P_1(m_{\rm land})=1+O(\epsilon_e)\) and
\(P_j(m_{\rm land})=O_j(\epsilon_e)\) for \(j\ge2\), proving
\eqref{eq:direct-adaptive-landing}.

For an input above the thin threshold at a radius \(\bar\varrho\), the
same equation in \(\mathcal Y\), started at a larger value of \(\mathcal Y\), gives radial
growth to \(\mathcal Y=m_{\rm land}\) no larger than for an input at that
threshold; hence its landing radius is at most \(2\bar\varrho\).
An input with \(\mathcal Y\ge m_{\rm land}\) already has the required cone height.

\emph{Reaching the cone before the endpoint.}
The organizing endpoints satisfy
\(\eta_\pm^0=\pm a_*r_*+r_*^2\), in the notation of
\cref{lem:exterior-multiplier-new}, by
\eqref{eq:normal-rate-coordinate-new}. Hence
\(z_*/(\eta_\pm^0)^2\to1/2\) as \(z_*\downarrow0\).
For each sufficiently small fixed \(z_*\), the endpoint windows and
parameter neighborhood can therefore be chosen so that
\(\mathcal Y_*=z_*/\varrho_*^2\ge c_{\rm end}>0\), with one
\(c_{\rm end}\) independent of \(z_*\). Thus
\(m_{\rm land}<c_{\rm end}/2\) can be fixed before \(z_*\). The endpoint's normalized thin height is
\(y_*=\mathcal Y_* /(s+\varrho_*)>y_0\) on the small box in the coordinates
\((\eta,z)\).
Thus every admissible graph that starts thin exits and lands before
the fixed endpoint section. Monotonicity of the positive
radius keeps the construction in the chosen radius box.
\end{proof}

The preceding lemma controls the moving exit and its first cone
landing. To obtain a return-map estimate, we now differentiate the
endpoint reached from this landing. At
\(t=m_{\rm land}\varrho_l^2\), the entrance energy is \(U=t/(2m_{\rm land})\).
By the choice \(1/(2m_{\rm land})>4\sup B\), the moving-cone
part of \cref{lem:direct-weighted-tail} applies whenever
\eqref{eq:direct-tail-smallness} holds there. The endpoint energy
then has
\[
 \frac{dU(z_*;t/(2m_{\rm land}))}{dt}
       =(z_*/t)^{-2q_2}a(t),\qquad c<a(t)<C.
\]
Under this tail condition, suppose also that the adapted
incoming height satisfies
\(0<c\le d\lambda_{\rm in}/d\log h\le C\), where
\(\lambda_{\rm in}=\log\widehat z(\varrho_i)\). Then
\begin{equation}
\label{eq:direct-adaptive-source-slope}
 hI_{+,h}\asymp \varrho_e^2(s+\varrho_e)
                    (z_*/(m_{\rm land}\varrho_l^2))^{-2q_2}.
\end{equation}
Indeed \eqref{eq:direct-adaptive-exit} and
\eqref{eq:direct-adaptive-landing} give
\(-dt/d\log h\asymp t(s+\varrho_e)\), and \(I_+=z_*-U(z_*)\).

\begin{corollary}[Contraction after a large exit]
\label{cor:large-exit-contraction}
In the setting of \cref{lem:direct-adaptive-positive}, let
\(\mathcal J\subset\cH_p\) be an interval of central heights whose positive
half-orbits are thin at \(\varrho_i\), satisfy
\eqref{eq:direct-tail-smallness} at their landings, and have
\(c\le d\lambda_{\rm in}/d\log h\le C\). Let \(c_1\) be the constant of the
lower bound in \eqref{eq:direct-adaptive-source-slope}. Suppose that
\(q_2\le C_2s\) and \(0<m_-\le C_1s^3\) on \(\mathcal J\), with fixed
\(C_1,C_2\), and fix \(M_1\ge1\) with
\(c_1M_1^2(1+M_1)>2\Lambda_+C_1\). Then, for sufficiently small \(s\),
the central heights in \(\mathcal J\) with \(\varrho_e\ge M_1s\) form an
interval on which \(\mathcal L<0\). In particular \(\Delta_p\) is strictly
decreasing there and has at most one zero.
\end{corollary}

\begin{proof}
By \eqref{eq:direct-adaptive-exit} and the lower bound for
\(d\lambda_{\rm in}/d\log h\), the exit radius decreases with \(h\), so
these central heights form an interval. The landing precedes the
endpoint, so \(t=m_{\rm land}\varrho_l^2<z_*\), and the multiplier in
\eqref{eq:direct-adaptive-source-slope} is at least one if \(q_2\le0\).
If \(0<q_2\le C_2s\), then \(\varrho_l\ge\varrho_e\ge s\) gives
\(t\ge m_{\rm land}s^2\), and
\[
 (z_*/t)^{-2q_2}\ge\exp[-2C_2s\log(z_*/(m_{\rm land}s^2))]\ge\tfrac12
\]
for small \(s\). In both cases
\(m_+\ge\frac12c_1\varrho_e^2(s+\varrho_e)\), and therefore
\[
 \frac{m_-}{m_+}\le\frac{2C_1}{c_1}\,
        \frac1{(\varrho_e/s)^2(1+\varrho_e/s)} .
\]
The endpoints of central heights in \(\cH_p\) lie in the windows, where
\(\Pi_p'\le\Lambda_+\) by \cref{lem:exterior-multiplier-new}. For
\(\varrho_e\ge M_1s\), the choice of \(M_1\) gives
\(\e^{\mathcal L}=\Pi_p'(I_-)m_-/m_+<1\), and
\(\Delta_p'=I_+'(\e^{\mathcal L}-1)<0\) by
\eqref{eq:direct-completed-log-ratio}.
\end{proof}

\subsection{A boundary root approaching the central line}
\label{sec:weak-root-passage}

When the constant coefficient is small compared with the linear one,
the first boundary root on the positive branch approaches the central
line relative to the current coefficient scale. Rescaling by its radius
keeps this root simple and near \(x=1\). In the resulting logarithmic
height coordinate, the positive endpoint has a large negative
logarithmic rate. A uniform lower bound for the negative endpoint's
rate then makes \(\mathcal L\) increase throughout the root interval.

Let \(a,\kappa>0\) be small and \(\hat r\in(0,1]\), and suppose
\begin{equation}
\label{eq:weak-root-scales}
 q_0=-\kappa\ell^2,\qquad q_1=\kappa\ell,\qquad q_2=-\epsilon_2\kappa,
 \qquad \ell=\hat r\kappa a^2,\qquad |\epsilon_2|\le C_2a,
\end{equation}
with a fixed \(C_2\). In \(|\eta|=\ell x\), \(z=\ell^2Y\) the outward
drifts and height equations are
\[
 K_\sigma(x)=-\frac{q(\sigma\ell x)}{\kappa\ell^2}
 =1-\sigma x+\epsilon_2x^2-\sigma\hat ra^2x^3R_q(\sigma\ell x),\qquad
 x_{\log Y}=\kappa K_\sigma(x)/x+YB(\sigma\ell x,\ell^2Y)/x .
\]
For a fixed small preliminary upper bound on \(a\), the first positive
root of \(K_+\) satisfies \(\alpha=1+O(a)\) and \(K_+'(\alpha)<-c<0\).
Put \(x_L=\alpha/4\), \(w_a=\alpha/2\) and \(R_p=3\alpha/4\), and use
the logarithmic height coordinates \(T=-\kappa a\log h\) and
\(\widehat T=T/a=-\kappa\log h\).
We compare the branches in \(T\), after estimating the positive branch
in \(\widehat T\).
For a fixed small height \(Y_f>0\), the root interval is
\begin{equation}
\label{eq:weak-root-interval}
 w=\mathcal W_+(x(\ell^2Y_f;h),Y_f)\in[w_a,\alpha).
\end{equation}

The negative half-orbit need not be at the weak-root radius when the
positive half-orbit approaches its root. To estimate it on the same
central heights, introduce a variable radius scale: let \(d>0\)
and put \(\ell_d=\ell d/a\), \(\delta_d=\kappa a/d\) and
\(a_d=a/d\), so that \(a_d/\delta_d=1/\kappa\).
In \(|\eta|=\ell_dx\), the outward drift on the negative branch is
\[
 K_d(x)=a_d+x+(\epsilon_2/a)dx^2+\hat rd^2x^3R_q(-\ell_dx),
\]
so that \(K_d(x)=K_1(dx)/d\).
Fix \(d_{\max}>0\) such that, for \(0<d\le d_{\max}\),
\begin{equation}
\label{eq:weak-root-negative-drift}
 K_d(x)\asymp a_d+x,\qquad K_d'(x)\ge\tfrac12\qquad(0\le x\le4),
\end{equation}
with constants depending on \(d_{\max}\). This holds for every small
\(d_{\max}\), and for every fixed \(d_{\max}\) when \(\epsilon_2\ge0\) and
\(\hat r\) is small.

\begin{lemma}[Two branches at the weak root]
\label{lem:weak-root-branches}
There are \(c_1,c_d>0\) and a fixed small \(Y_b>0\) such that, for
sufficiently small \(a\) and \(\kappa\), the following hold in the coordinate
\(T\).
\begin{enumerate}[label=\textup{(\roman*)}]
\item On the root interval \eqref{eq:weak-root-interval},
\(\mathfrak H_+\le-c_1/a\).
\item Let \(c_da\le d\le d_{\max}\), and let \(J\subset[1/2,x_{\max}]\) be
a compact interval, with a fixed \(x_{\max}\), such that \(K_d>0\) on
\((0,\max J)\). At the height \(\ell_d^2Y_b\), and for negative projected
states in \(J\),
\[
 d\,\mathfrak H_-\ge\min_{x\in J}\frac{K_d'(x)}x-o(1),\qquad
 |I_{-,T}|\le CE_d/d=o(1),\qquad
 E_d=\ell_d^2\bigl(z_*/(\ell_d^2Y_b)\bigr)^{-2q_2},
\]
with \(o(1)\to0\) as \(\kappa\to0\) uniformly in \(d\) and \(J\), up to a
root of \(K_d\) at \(\max J\); the constants and the smallness of
\(a,\kappa\) depend on \(x_{\max}\). In particular,
\(\mathfrak H_-\ge c/d>0\) for \(J=[0.9,4]\).
\item On \eqref{eq:weak-root-interval}, the negative graph has
\(z<(Y_b/2)|\eta|^2\) at \(|\eta|=\ell c_d\).
\end{enumerate}
\end{lemma}

\begin{proof}
(i) The boundary-flow coordinate applies with small parameter
\(\kappa\), since the error from \(B\) is \(O(\ell)=O(\kappa)\). At
\(t=\ell^2Y_f\) the coefficient ratios are \(|q_0|/t=\kappa/Y_f\),
\(q_1/\sqrt t=\kappa/\sqrt{Y_f}\) and \(|q_2|\le C_2\kappa a\). Hence
\cref{lem:boundary-flow-endpoint} gives
\(-I_{+,w}=E_aA_+wV_+(w)\), with
\(E_a=\ell^2(z_*/(\ell^2Y_f))^{-2q_2}\) and
\(V_+=1+O_{C^3}(\kappa)\). The fixed positive phase advance from
\(x_L\) to \(w_a\) allows us to apply
\cref{lem:direct-weighted-phase} in the coordinate \(\widehat T\).
Since \(K_+'(w)/w\) has a fixed negative margin on
\([w_a,\alpha)\), \cref{cor:direct-weak-root-comparison} gives (i).

(ii) By \eqref{eq:weak-root-negative-drift} and \(a_d/\delta_d=1/\kappa\ge1\),
the case of a small central coefficient in \cref{lem:direct-source-height},
with \(K_*=a_d\), gives a uniform height unit at the entry \(x_i=1/4\),
also as \(a_d\to0\).
The boundary-flow projection at \(Y_b\), on a fixed radius interval
extending \([1/4,x_{\max}]\), has inverse energy error
\(O_{C^k}(\delta_dY)\), since \(\ell_d/\delta_d=\hat rd^2\) is bounded.
By \cref{lem:direct-boundary-flow}, the projection's height bound \(Y_0\) depends
only on the lower end of that interval, so \(Y_b\) does not depend on
\(x_{\max}\).

For output states in \(J\), the advance from \(x_i\) to \(1/2\)
takes a fixed positive boundary time, since \(K_d/x\) is bounded
there. Thus \cref{lem:direct-weighted-phase} applies uniformly up to
a root of \(K_d\). At \(\ell_d^2Y_b\), the coefficient ratios are
\(\delta_da_d/Y_b\), \(\delta_d/\sqrt{Y_b}\) and
\(|q_2|\le C_2\kappa a\), and \(\ell_d\to0\).
These quantities are small because \(\delta_d\le\kappa/c_d\)
and \(a_d\le1/c_d\). \Cref{lem:boundary-flow-endpoint} with
\(\ell=\ell_d\), \(\delta=\delta_d\) and \(J_w=[1/2,x_{\max}]\)
therefore gives \(-I_{-,w}=E_dA_dw(1+O_{C^3}(\delta_d))\).

To express these estimates in the common coordinate, use
\(\widehat T_d=-\delta_d\log h=T/d\), so
\(\partial_T=d^{-1}\partial_{\widehat T_d}\). In \(\widehat T_d\),
the first two lines of \eqref{eq:direct-rate-identities}, with
\(\theta-1,\theta'=o(1)\), give \(|I_{-,\widehat T_d}|\le CE_d\) and
\(\partial_{\widehat T_d}\log(hI_{-,h})\ge\min_JK_d'(x)/x-o(1)\).
Division by \(d\) gives the two estimates. On \([0.9,4]\),
\(K_d'(x)/x\ge1/8\) by \eqref{eq:weak-root-negative-drift}.

Finally, \(|q_2|\le C_2\kappa a<1/8\), so
\eqref{eq:direct-tail-amplitude} with \(t=\ell_d^2Y_b\) gives
\(E_d\le C\ell_d^{3/2}\) and
\(E_d/d\le C(\ell/a)^{3/2}d^{1/2}\to0\), uniformly for
\(0<d\le d_{\max}\), since \(\ell/a=\hat r\kappa a\to0\).

(iii) Choose \(c_d\) so small that the negative prefix action to
\(x=c_d\) in the weak-root scale is less than one quarter of the fixed
positive phase advance \(\tau_*\) from \(x_L\) to \(w_a\). On
\eqref{eq:weak-root-interval} this advance gives
\(h/(\kappa\ell^2)\le(Y_f/\kappa)\e^{-\tau_*/\kappa}\). The negative
central-height coordinate then gives
\[
 \frac{z(-\ell c_d;h)}{(\ell c_d)^2}
 \le\frac{CY_f}{c_d^2}\,\e^{-3\tau_*/(4\kappa)}
 <\frac{Y_b}{2}
\]
for small \(\kappa\): the exponential decay in \(1/\kappa\) absorbs the
bounded height units and the polynomial prefactors.
\end{proof}

\begin{corollary}[Negative branch up to its first root]
\label{cor:weak-root-negative-first-root}
Assume \eqref{eq:weak-root-scales} and \eqref{eq:weak-root-negative-drift},
with \(Y_b\) as in \cref{lem:weak-root-branches}, and suppose that the
negative drift \(K_1\) in \(|\eta|=\ell X/a\) has a first positive root
\(\beta\le\beta^*\), with a fixed \(\beta^*\). For sufficiently small \(a\)
and \(\kappa\), every central height whose negative graph has
\(z<Y_b|\eta|^2\) at \(|\eta|=\ell_{d_{\max}}\) satisfies, in the
coordinate \(T\),
\[
 \mathfrak H_-\ge\min_{d_{\max}/2\le X\le\beta}\frac{K_1'(X)}X-o(1),
 \qquad |I_{-,T}|=o(1),
\]
where \(o(1)\to0\) as \(\kappa\to0\), uniformly as the negative projected
state at the height \(\ell_{d_{\max}}^2Y_b\) approaches the root.
\end{corollary}

\begin{proof}
By \eqref{eq:weak-root-negative-drift} at \(d=d_{\max}\) and
\(K_d(x)=K_1(dx)/d\), \(K_1>0\) on \((0,4d_{\max}]\).
The first positive root of \(K_{d_{\max}}\) is therefore
\(\beta/d_{\max}>4\). In \(|\eta|=\ell_{d_{\max}}x\),
\(z=\ell_{d_{\max}}^2Y\), the graph reaches \(x=1\) below
\(Y=Y_b\). Equation~\eqref{eq:direct-projection-energy} puts its
projected state there above \(1/2\), since \(Y_b\) and \(\kappa\)
are small. The projected state then increases; by
\cref{rem:direct-projected-continuation}, it does not cross
\(\beta/d_{\max}\), and the graph continues to \(Y_b\).

At that height the projected state lies in
\(J=[1/2,\beta/d_{\max}]\). Apply part~(ii) of
\cref{lem:weak-root-branches} with \(d=d_{\max}\), this \(J\), and
\(x_{\max}=\beta^*/d_{\max}\). The assertion follows from
\(K_d'(x)/x=d\,K_1'(dx)/(dx)\), with \(d_{\max}\) fixed.
\end{proof}

On the root interval, a negative graph is controlled by
\cref{lem:weak-root-branches} if \(z/|\eta|^2\) reaches \(Y_b\) in the
prescribed radius range. The following proposition isolates the
additional estimate needed for graphs that remain below this level.
Each application will verify that estimate on its own coefficient
scale.

\begin{proposition}[Weak root]
\label{prop:weak-root-passage}
Assume \eqref{eq:weak-root-scales} and \eqref{eq:weak-root-negative-drift},
and let \(Y_b\) be as in \cref{lem:weak-root-branches}. Suppose there are
\(C_n\) and \(\varepsilon_p\to0\) such that every central height in the
root interval \eqref{eq:weak-root-interval} whose negative graph has
\(z<Y_b|\eta|^2\) at \(|\eta|=\ell_{d_{\max}}\) satisfies
\begin{equation}
\label{eq:weak-root-outer}
 \mathfrak H_-\ge-C_n,\qquad |I_{-,T}|\le\varepsilon_p .
\end{equation}
Then there are \(a_\circ,\kappa_\circ>0\) such that, for \(0<a\le a_\circ\),
\(0<\kappa\le\kappa_\circ\) and all sufficiently small remaining parameters,
every zero of \(\mathcal L\) on \(\cH_p\) satisfies
\(\mathcal L_T>0\). Consequently \(\Delta_p\) has at most two distinct
zeros there.
\end{proposition}

\begin{proof}
\emph{The negative branch on the root interval.} By
\cref{lem:weak-root-branches}(iii), the ratio \(z/|\eta|^2\) along the
negative graph is below \(Y_b/2\) at \(|\eta|=\ell c_d\).
If the ratio reaches \(Y_b\) at a
radius \(\ell_d\) with \(c_da<d\le d_{\max}\), choose the first
such \(d\) for the central height under consideration. At
\(\ell_d^2Y_b\), the normalized radius is then one, and small
\(Y_b\) puts the projected state in \([0.9,4]\). Hold this section
fixed when varying the central height locally. The resulting
representation of \(I_-\) agrees with \(I_-(h)\), so
\cref{lem:weak-root-branches}(ii) applies as in
\cref{lem:direct-frozen-section-sign}. Otherwise, the ratio is still
below \(Y_b\) at \(|\eta|=\ell_{d_{\max}}\), and
\eqref{eq:weak-root-outer} applies. In both cases,
\(\mathfrak H_-\ge-C_n\) and \(|I_{-,T}|=o(1)\) throughout
\eqref{eq:weak-root-interval}.

\emph{The root interval.} Combine this bound for the negative branch
with \cref{lem:weak-root-branches}(i) for the positive branch.
\Cref{cor:direct-weak-root-comparison} gives \(\mathcal L_T>0\)
throughout \eqref{eq:weak-root-interval}, after \(a_\circ\) is reduced.
The constants are fixed on the preliminary small-\(a\) region.

\emph{The height cover.} Apply \cref{lem:prefix-root-cover} with
\(\delta=\kappa\), \(\epsilon_3=\ell/\kappa=\hat ra^2\), comparison drifts
\(1-\sigma x+\epsilon_2x^2\), and the radii \(w_a,R_p,\alpha\) above.
The comparison drifts have the required uniform positive margin, the
gap \(R_p-w_a\) is uniform, and the boundary-flow inverse is defined
up to \(\alpha\) at \(Y_f\). Together with \(\mathcal L_T>0\) on
the root interval, the lemma gives at most two distinct zeros of
\(\Delta_p\) on \(\cH_p\).
\end{proof}

\section{Negative offsets on the primary scale}
\label{sec:primary-regimes}

On the primary scale the cubic term remains in the limiting boundary
drift. Its contribution gives opposite signs to the curvature
coefficient \(wK''-K'\) on the two branches. When the normalized
constant coefficient is bounded away from zero, we combine this
curvature estimate on bounded radii with contraction at early heights
and large exits. In particular,
\cref{lem:curvature-criterion} gives concavity of \(\mathcal L\)
through a root-free passage and up to simple or double boundary roots, and
\cref{lem:direct-adaptive-positive} treats the positive half-orbits that
leave this scale. When the constant coefficient is small compared with
\(q_1^{3/2}\), a boundary root approaches the central line, and
\cref{prop:weak-root-passage} applies on the linear scale.

\subsection{Negative drift on the primary scale}
\label{sec:direct-primary-negative}

Assume the magnitude of the constant coefficient has a positive
normalized lower bound on the primary scale. Positive boundary roots
may still merge. At bounded radii the drift-weighted curvature estimate
continues through this fold; at early heights and large exits we use
contraction instead. We prove that these ranges overlap and cover
\(\mathcal H_p\), so that \cref{cor:ordered-cover-zero-count} applies
to a single logarithmic slope ratio on the full interval.

\begin{proposition}[Primary negative drift]
\label{prop:direct-primary-negative}
For every fixed \(a_0^{\rm pr}>0\), the return displacement has at
most three distinct zeros when
\begin{equation}
\label{eq:direct-primary-region}
 q_0=-aH^3,\quad q_1=bH^2,\quad q_2=\bar c H,\qquad
 a_0^{\rm pr}\le a\le1,\quad 0\le b\le1,\quad |\bar c|\le1,
\end{equation}
for all sufficiently small \(H>0\) and original parameters.
\end{proposition}

Write \(|\eta_\sigma|=HX,\ z=H^2Y\). The outward
boundary drifts are
\begin{equation}
\label{eq:direct-primary-drifts}
 K_\sigma(X)=-H^{-3}q(\sigma HX)
 =a-\sigma bX-\bar c X^2-\sigma X^3R_q(\sigma HX).
\end{equation}
With the central-height scale \(Q=H^3\) of \cref{tab:direct-scale-guide},
the common logarithmic height coordinate below is \(T=-H\log(h/Q)\).
The coefficient calculation in the appendix shows that the two
drifts converge, on fixed boxes and in every required state order, to
\[
 K_+^0=a-bX-\bar c X^2+X^3,\qquad
 K_-^0=a+bX-\bar c X^2-X^3.
\]
For \(X>0\), their curvature coefficients are \(b/X^2+3\) and \(-b/X^2-3\),
respectively, so the endpoint estimates below must preserve these
opposite signs uniformly as positive boundary roots merge.
The drift \(K_-^0\) has a unique, uniformly simple positive root \(\beta\)
separated from zero. Its derivative has at most one positive
zero, and
\[
 -\frac{(K_-^0)'(\beta)}{\beta}
        =2a/\beta^2+b/\beta+\beta>0 .
\]
Compactness gives the uniform bounds, which persist under the
\(O(H)\) perturbation. All positive roots of \(K_+^0\), if
present, also lie in a fixed compact interval away from zero.

On every fixed positive radius interval, the height equation is
\[
 X_{\log Y}=HF_\sigma(X)+YB(\sigma HX,H^2Y)/X,\qquad
 F_\sigma=K_\sigma/X.
\]
Thus \cref{lem:direct-boundary-flow} applies with \(\delta=H\), and
\eqref{eq:direct-projection-energy} holds for its inverse projection.
Below the first root of \(F_\sigma\), where \(F_\sigma>0\),
\cref{rem:direct-projected-continuation} applies.

\subsubsection{Early central heights}

Choose a small fixed \(d_f>0\) and put \(z_e=H^3d_f\). The present
scales are \eqref{eq:direct-quadratic-prefix-scales} with
\(\delta=\ell=H\) and \(\bar c\) replaced by \(-\bar c\). By
\cref{lem:early-central-heights}, there is a small fixed \(R>0\) such
that the set \eqref{eq:direct-primary-early}, with \(r_0=R\), is an
initial interval in \(T\) on which the return contracts.

\subsubsection{Bounded radii, including the fold}

For every remaining central height, the positive half-orbit reaches
the radius \(R\) before \(z_e\). Fix a small \(Y_f>0\) for the central
entry and the negative-root box below, and let \(M\) be large; \(M\) is
fixed after \(M_1\) in the covering step. The same \(Y_f\) serves every
such \(M\): at \(H=0\) the positive projection on \([R/2,M]\) is
\(w^2=X^2-2B_0Y\), and only its lower-radius margin matters, while a
larger \(M\) requires only a smaller final \(H\) for the derivative
bounds.
Consider those central heights with \(X_+(H^2Y_f;h)\le M\).
Their projected state \(w_+\) at this height lies in \([R/2,M]\):
its value at \(z_e\), when used, differs from \(X_+(z_e)\)
by \(O(Hd_f)\), and \(w_+\) increases.
The phase comparison below also places their negative-branch state
\(w_-\) in \([d_0,\beta)\) for a fixed \(d_0>0\).

Choose a small fixed \(x_L>0\) inside the central positive-drift margins
of both \(K_\sigma\). With
\[
 \mathcal A_\sigma(x)=\int_0^x\frac{u\,du}{K_\sigma(u)},
\]
\eqref{eq:direct-common-phase} gives
\begin{equation}
\label{eq:direct-primary-phases}
 w_\sigma=\Phi_{F_\sigma}^{\,T+C_\sigma+
                              H\Xi_\sigma(\xi_\sigma)}(x_L),
 \quad
 C_\sigma=H\log(Y_f/H)-\mathcal A_\sigma(x_L),\quad
 \xi_\sigma=\e^{\mathcal A_\sigma(x_L)/H}h/H^3 .
\end{equation}
The phase corrections \(\Xi_\sigma\) vanish at zero and have
uniformly bounded derivatives of the required orders.
Positive advance to \(R/2\) takes a fixed positive boundary
time. Take both entry actions smaller than one quarter of
this minimum time. The two phases in
\eqref{eq:direct-primary-phases} then have a common positive
lower bound. The negative scalar flow advances into a fixed
interval away from zero and remains below its first root.
Thus \(d_0\le w_-<\beta\), and the inverse projection
continues the negative graph to \(Y_f\).

At \(t=H^2Y_f\), all three coefficient ratios in
\eqref{eq:direct-tail-smallness} are \(O(H)\), and
\cref{lem:boundary-flow-endpoint} with \(\ell=\delta=H\) gives
\begin{equation}
\label{eq:direct-primary-endpoints}
 -I_{\sigma,w}=E A_\sigma^*wV_\sigma(w),\qquad
 E=H^2(z_*/(H^2Y_f))^{-2q_2}\asymp H^2,\qquad
 A_\sigma^*\asymp1,\quad V_\sigma=1+O_{C^3}(H).
\end{equation}
The negative endpoint estimate must also hold in the large-exit range,
before the positive radius cut is chosen.
The positive states range over \([R/2,M]\), so the constants of the
positive endpoint depend on \(M\). The negative states lie in
\([d_0,\beta)\), and the constants of the negative endpoint are
independent of \(M\).

The drifts \eqref{eq:direct-primary-drifts} have
\(K_\sigma'(0)=-\sigma b\) and, in the notation of
\eqref{eq:direct-curvature-coefficient-cubic},
\(R_{K_\sigma}(w)=-\sigma R_q(\sigma Hw)\), which is
\(\sigma+O_{C^2}(H)\) on the state intervals by the coefficient
calculation in the appendix. Hence
\eqref{eq:direct-curvature-coefficient-cubic} gives
\begin{equation}
\label{eq:direct-primary-opposite-curvatures}
 \frac{wK_+''-K_+'}{w^2}=\frac b{w^2}+3+O_M(H)>2,\qquad
 \frac{wK_-''-K_-'}{w^2}=-\frac b{w^2}-3+O_M(H)<-2 .
\end{equation}
The margins hold also when \(b=0\). Since the negative state stays in
\([d_0,\beta)\), \cref{lem:curvature-criterion} gives \(\mathcal L''<0\)
in \(T\), uniformly up to either root. The bounded-radius central
heights form an interval, by \cref{lem:central-height-order}.

\subsubsection{Escaping radii and their endpoint slope}

Fix \(M_0\) beyond all positive roots and critical points of
the limiting primary polynomial. For the positive radius
\(\varrho=\eta\ge M_0H\), write
\[
 \varrho^+=H+\varrho,\qquad u=H/\varrho,\qquad
 v=\varrho/(H+\varrho).
\]
Then
\[
 -q(\varrho)=\varrho^2\varrho^+\mathcal K,\qquad
 \mathcal K=v\{au^3-bu^2-\bar c u-R_q(\varrho)\}.
\]
For large \(M_0\) this \(\mathcal K\) is positive and has
uniform derivative bounds on an enlarged base box.
It equals \(1+O(u+H+\varrho)\). Thus
\cref{lem:direct-adaptive-positive} applies with \(s=\varrho_0=H\).
Its exit and landing satisfy
\(\varrho_e\le\varrho_l\le2\varrho_e\), with uniform derivatives,
before the admissible endpoint.

If the thin exit is at \(\varrho_e\ge2M_0H\), there is no positive
boundary root before \(M_0H\). Indeed at a thin hit of
\(M_0H\) the primary height is \(Y=O_{M_0}(H)\). The
projection, constructed on the fixed interval up to \(M_0\),
differs from the radius by \(O_{M_0}(H)\) there.
Since any first root is separated from \(M_0\), a projected state below
that root cannot give such a hit.

On these no-root parameters, transport the regular central
height coordinate from \(x_L\) to \(M_0\), using
\eqref{eq:direct-bottleneck-action}. The endpoint collars
have fixed positive drift margins, although the action between
them may diverge. In the height \(y=z/H^3\), this gives
\begin{equation}
\label{eq:direct-primary-bottleneck-prefix}
 \Psi_o(y_o)=\e^{\mathcal A_+(M_0)/H}\frac h{H^3}
 \Upsilon\!\left(\e^{\mathcal A_+(x_L)/H}\frac h{H^3}\right),
 \qquad \Psi_o(y)=y+O(y^2),\quad \Upsilon(0)=1 .
\end{equation}
By \eqref{eq:direct-bottleneck-inverse-composition}, the units and
height coordinates, together with their inverses, have uniform
derivative bounds in the expanded input height.

There is a fixed positive boundary-action gap from \(M_0\)
to \(2M_0\). Hence an exit beyond \(2M_0H\) makes the
height at \(M_0H\), normalized by \(H^3\), at most
\(C H^{-N}\e^{-\mathcal A_{\rm gap}/H}\). The unscaled exit radius is
bounded, so its prefactor costs only a fixed power of \(1/H\).
Differentiating the regular units in
\eqref{eq:direct-primary-bottleneck-prefix} gives
\[
 \frac{d\log\widehat z(M_0H)}{d\log h}
             =1+O(\e^{-\mathcal A_{\rm gap}/(2H)}).
\]
This is the bound on \(d\lambda_{\rm in}/d\log h\) required in
\eqref{eq:direct-adaptive-source-slope}, with \(\varrho_i=M_0H\). The tail
condition there holds at the landing: \(\varrho_l\ge M_0H\) gives
\(|q_0|/(m_{\rm land}\varrho_l^2)+q_1/\sqrt{m_{\rm land}\varrho_l^2}=O(H)\), and
\(|q_2|\le H\).

To compare with the positive endpoint slope, return to the negative
branch. It retains a fixed positive entry-time gap; its projected state lies
in \([d_0,\beta)\), regardless of the large positive exit.
The independently chosen negative endpoint bounds in
\eqref{eq:direct-primary-endpoints} therefore give
\[
 0<m_-=hI_{-,h}\le CHE\le CH^3 .
\]
Thus \cref{cor:large-exit-contraction} applies, with \(s=H\), \(C_2=1\)
and \(C_1\) the constant just obtained, to the central heights in
\(\cH_p\) with \(\varrho_e\ge2M_0H\), an interval since the exit radius
decreases with \(h\). Fix \(M_1\ge2M_0\) as required by the corollary; this gives the
large-exit interval \(\varrho_e\ge M_1H\).

\subsubsection{Covering the central heights}

The bounded-radius estimate allows any fixed \(M\), whereas the
large-exit threshold \(M_1\) uses only the independently fixed negative
endpoint bound. We may therefore choose \(M\) sufficiently large after
\(M_1\). We show that every height outside the bounded-radius range
satisfies the large-exit estimate. If the
positive half-orbit of a central height enters the adaptive range
thin and has not exited when its radius reaches \(MH\), its
eventual exit is larger.
If it has exited, its radius before cone landing is at most
\(2\varrho_e\). After landing the cone estimates give, whenever the
fixed output height is reached,
\[
 U_+(H^2Y_f)\le C\varrho_e^2+C H^2Y_f .
\]
If the half-orbit was already above the thin threshold at \(M_0H\) but
below the cone, \cref{lem:direct-adaptive-positive} bounds its landing
radius by \(2M_0H\), and the same inequality holds with
\(\varrho_e=M_0H\). If already in the cone, its continuation gives this
bound directly.
Thus, for a sufficiently large fixed \(M\), every central height
outside the bounded-radius range must have \(\varrho_e\ge M_1H\).

The early interval \eqref{eq:direct-primary-early}, the bounded-radius
interval, and the large-exit interval cover \(\cH_p\), with fixed
overlaps. On these intervals the same function \(\mathcal L\) of \(T\)
satisfies \(\mathcal L<0\) in the early and large-exit ranges, and
\(\mathcal L''<0\) on the bounded-radius interval, including the
shared cuts. \Cref{cor:ordered-cover-zero-count} gives at most three
zeros of \(\Delta_p\). This proves \cref{prop:direct-primary-negative}.

\subsection{Negative displaced contact: the weak root on the linear scale}
\label{sec:direct-negative-displaced}

The preceding estimate loses its normalized constant-coefficient
margin when that coefficient becomes small. If
\(q_2/\sqrt{q_1}\) remains bounded, the approaching positive boundary
root gives the weak-root scale. \Cref{prop:weak-root-passage} reduces
the count to an estimate for the negative branch beyond the weak-root
radius. We verify this estimate on the linear scale.

\begin{proposition}[Negative displaced contact]
\label{prop:direct-negative-displaced}
For every fixed \(A_0<\infty\), there are fixed
\(a_{\rm w},\kappa_{\rm w}>0\) such that the return displacement has
at most two distinct zeros on
\[
 q_1=\rho_1^2>0,\quad q_2=\bar c\rho_1,\quad |\bar c|\le A_0,\quad
 q_0=-a\rho_1^3,\quad 0<a\le a_{\rm w},\quad
 0<\kappa=\rho_1/a\le\kappa_{\rm w} ,
\]
for all sufficiently small \(\rho_1\) and original parameters.
\end{proposition}

\begin{proof}
To apply \cref{prop:weak-root-passage}, set
\(\ell=\rho_1a=|q_0|/q_1\). Then
\eqref{eq:weak-root-scales} holds with \(\hat r=1\) and
\(\epsilon_2=-\bar c a\), \(|\epsilon_2|\le A_0a\), and the common
logarithmic height coordinate is \(T=-\rho_1\log h=-\kappa a\log h\).
At the radius \(\ell_d=\rho_1d\) the negative drift is
\[
 K_d(x)=a_d+x-\bar c dx^2+d^2x^3R_q(-\ell_dx),
\]
so \eqref{eq:weak-root-negative-drift} holds for a small fixed
\(d_{\max}=d_0\) depending on \(A_0\).

It remains to verify \eqref{eq:weak-root-outer} for the negative graphs
with \(z<Y_b|\eta|^2\) at \(|\eta|=\rho_1d_0\). Return to the linear
scale \(\rho_1\), with \(|\eta|=\rho_1X\). Since \(\ell/a=\rho_1\),
the negative drift on this scale is \(K_1\) from
\cref{cor:weak-root-negative-first-root}:
\[
 K_-^{\rho_1}(X)=a+X-\bar c X^2+X^3R_q(-\rho_1 X).
\]
Its first positive root \(\beta\) stays in a fixed compact interval.
At \(a=\rho_1=0\), it is the positive root of
\(1-\bar c X-X^2=0\); the bound \(|\bar c|\le A_0\) gives uniform
separation and a simple-root margin. On \([d_0/2,\beta]\),
\((K_-^{\rho_1})'(X)/X\ge-C_{A_0,d_0}\).

Apply \cref{cor:weak-root-negative-first-root} with
\(d_{\max}=d_0\) and with \(\beta^*\) the right end of that compact
interval. It gives \(\mathfrak H_-\ge-C_{A_0,d_0}-1\) and
\(|I_{-,T}|=o(1)\) for these graphs, including the limiting negative
root. Thus \eqref{eq:weak-root-outer} holds.
\Cref{prop:weak-root-passage} now gives at most two distinct zeros
of \(\Delta_p\), with \(a_{\rm w}\) and \(\kappa_{\rm w}\) equal to
the constants \(a_\circ\) and \(\kappa_\circ\) of that proposition.
\end{proof}

\section{The secondary scale and positive quadratic drift}
\label{sec:quadratic-regimes}

When the normalized ratios of \(q_0\) and \(q_1\) to the scale set by
\(q_2\) are small, we use the secondary scales \(s,r\) of
\cref{sec:standard-tools}.
For \(q_2>0\), we compare the projected boundary flows near their first
roots and continue the half-orbits to the fixed output windows.
We use curvature away from the symmetric limit and the exterior
contraction in root-distance coordinates near it.
The case \(q_2<0\) is treated in
\cref{sec:direct-negative-annulus-buffer,sec:direct-negative-root-regimes}:
we first estimate the annular passage and the negative branch, then
treat the positive branch in each coefficient region.
For fixed parameters, the early, annular or root, and large-exit ranges
form overlapping subintervals of \(\mathcal H_p\). Their estimates give
crossing directions or critical-point curvature for \(\mathcal L\),
which are combined by \cref{cor:ordered-cover-zero-count}.
Thus the coefficient region selects a passage argument, while the
ordered height cover combines the estimates within that argument.
\Cref{fig:secondary-cover} shows how these regions cover the normalized
coefficients of both signs of \(q_2\).

\begin{figure}[H]
\centering
\begin{tikzpicture}[x=4.6cm,y=3.2cm,>=Stealth,font=\small,
    bar/.style={line width=2.6pt,gray!55},
    lab/.style={font=\footnotesize,inner sep=1.5pt,text=black}]
  % (a) negative quadratic coefficient
  \begin{scope}
    \draw[->] (0,0)--(1.2,0) node[right] {$a$};
    \draw[->] (0,0)--(0,1.62) node[above] {$b$};
    \draw[very thick] (0,1)--(1,1)--(1,0);
    \draw (0,1)--(-.03,1) node[left] {$1$};
    \draw (1,0)--(1,-.03) node[below] {$1$};
    \draw (.25,1)--(.25,.96) node[below,lab] {$\tfrac14$};
    \draw[bar] (.01,1.09)--(.14,1.09) node[lab,right] {weak root};
    \draw[bar] (.07,1.2)--(.22,1.2) node[lab,right] {separated roots};
    \draw[bar] (.19,1.31)--(.31,1.31) node[lab,right] {fold};
    \draw[bar] (.28,1.42)--(1,1.42);
    \node[lab,above] at (.64,1.44) {buffered};
    \draw[bar] (1.08,0)--(1.08,1);
    \node[lab,right] at (1.1,.5) {buffered};
    \node at (.5,-.25) {(a) \(q_2<0\)};
  \end{scope}
  % (b) positive quadratic coefficient
  \begin{scope}[xshift=7.4cm]
    \draw[->] (0,0)--(1.2,0) node[right] {$a$};
    \draw[->] (0,0)--(0,1.62) node[above] {$b$};
    \draw[very thick] (0,1)--(1,1)--(1,0);
    \draw (0,1)--(-.03,1) node[left] {$1$};
    \draw (1,0)--(1,-.03) node[below] {$1$};
    \draw[bar] (.01,1.09)--(.16,1.09) node[lab,right] {weak root};
    \draw[bar] (.08,1.2)--(1,1.2);
    \node[lab,above] at (.6,1.22) {curvature};
    \draw[bar] (1.08,.2)--(1.08,1);
    \draw[bar] (1.16,0)--(1.16,.4);
    \node[lab,right] at (1.1,.72) {curvature};
    \node[lab,right] at (1.18,.2) {small \(b\)};
    \node at (.5,-.25) {(b) \(q_2>0\)};
  \end{scope}
\end{tikzpicture}
\caption{Schematic of the compact coefficient covers on the secondary
charts \eqref{eq:direct-negative-secondary-scales} and
\eqref{eq:direct-positive-secondary-chart}; thresholds are not to scale.
The normalized coefficients lie on the curve \(\max\{\sqrt a,b\}=1\), and
the bars show the overlapping regions. (a) For \(q_2<0\): the buffered
region \(b^2-4a\le-\nu\) (\cref{prop:direct-buffered-negative}), the fold
\(|b^2-4a|\le2\nu\) (\cref{prop:direct-negative-fold}), separated roots
\(b^2-4a\ge\nu\), \(a\ge a_0^-\) (\cref{prop:direct-negative-separated}), and
the weak root \(a\le a_-\)
(\cref{prop:direct-negative-weak-root}). (b) For \(q_2>0\): curvature for
\(a\ge a_0^+\), \(b\ge b_{*,+}/2\), root-distance coordinates for small \(b\), and
the weak root \(a\le a_+\) (\cref{prop:direct-positive-quadratic}).
Near \(a=0\), the cover combines the weak-root regions, which require
\(s/a\le\kappa_\pm\), with the contact strip of
\cref{prop:direct-bounded-contact}.}
\label{fig:secondary-cover}
\end{figure}


\subsection{Positive quadratic-drift estimates}
\label{sec:direct-positive-quadratic}

Each branch on this chart has a first boundary root. When the central
coefficient keeps these roots away from the central line, we compare
the endpoint slopes either by curvature or, near the symmetric limit,
by distances to the roots. A further rescaling treats the root that
approaches the central line.

Suppose
\(q_2=s>0,\ q_0<0,\ q_1>0\), and write
\begin{equation}
\label{eq:direct-positive-secondary-chart}
 r=\max\{\sqrt{|q_0|/s^3},q_1/s^2\},\quad
 q_0=-as^3r^2,\quad q_1=bs^2r,\quad
 \max\{\sqrt a,b\}=1 .
\end{equation}
The small parameters \(s,r>0\) are independent. Put
\(t_0=s^2r^2,\ Q=st_0\), and \(X=|\eta|/(sr)\).
Here \(sr\) is the radius scale, \(Q=s^3r^2\) the central
height scale, and \(t_0\) the fixed-output height scale;
the common logarithmic height coordinate used below is \(T=-s\log h\).
The outward drifts are
\begin{equation}
\label{eq:direct-positive-secondary-drifts}
 K_\sigma(X)=a-\sigma bX-X^2
                 -\sigma rX^3R_q(\sigma srX),\qquad
 F_\sigma=K_\sigma/X .
\end{equation}
\begin{proposition}[Positive quadratic drift]
\label{prop:direct-positive-quadratic}
There is \(\kappa_+>0\) such that, for all sufficiently small
independent \(s,r>0\) and every \(L>1/\kappa_+\), the return
displacement has at most three distinct zeros on the part of
\eqref{eq:direct-positive-secondary-chart} outside the strip
\(|q_0|\le Lq_1^2\).
\end{proposition}

\subsubsection{Constant coefficient bounded away from zero}

First fix \(a\ge a_0^+>0\). This lower bound keeps both roots in a
compact positive radius interval, where we can prepare the endpoint
maps uniformly. The roots are simple and satisfy
\[
 \alpha_\sigma
 =\frac{\sqrt{b^2+4a}-\sigma b}{2}+O(r),\qquad
 c\le\alpha_\sigma\le C,\qquad
 K_\sigma'(\alpha_\sigma)<-c .
\]
At height \(z=t_0Y\),
\[
 X_{\log Y}=sF_\sigma(X)+YB(\sigma srX,t_0Y)/X .
\]
The boundary-flow coordinate applies with \(\delta=s\),
since the error in \(B\) is \(O(sr)=O(s)\).
At a fixed small output height \(Y_f>0\) the coefficient ratios
are \(O(s)\), and \cref{lem:boundary-flow-endpoint} with
\(\ell=sr\) and \(\delta=s\) gives
\begin{equation}
\label{eq:direct-positive-secondary-endpoints}
 -I_{\sigma,w}=EA_\sigma wV_\sigma(w),\qquad
 E=t_0(z_*/(t_0Y_f))^{-2s}\to0,\qquad
 V_\sigma=1+O_{C^3}(s),
\end{equation}
on fixed intervals of positive projected states up to the roots.
Here \(A_\sigma\asymp1\) and
\(A_-/A_+\le\exp(2C_{\rm cone}\sqrt{z_*}+Cs)\).

Use the shifted coordinate \(T_Q=T+s\log Q\), as in \eqref{eq:direct-source-clock-convention}.
The phases are
\begin{equation}
\label{eq:direct-positive-secondary-phases}
 \begin{aligned}
 w_\sigma&=\Phi_{F_\sigma}^{\,T_Q+C_\sigma+
                                  s\Xi_\sigma(\xi_\sigma)}(x_{L,\sigma}),
 &T_Q&=-s\log(h/Q),\\
 \xi_\sigma&=\e^{\mathcal A_\sigma(x_{L,\sigma})/s}h/Q,
 &C_\sigma&=s\log(Y_f/s)-\mathcal A_\sigma(x_{L,\sigma}),\\
 \mathcal A_\sigma(x)&=\int_0^x u/K_\sigma(u)\,du .
 \end{aligned}
\end{equation}
Choose both regular entry actions below one quarter of the minimum
positive terminal action. Then \(w_+\ge\alpha_+/2\) forces a
fixed advance on both legs. The negative state lies in a fixed
interval away from zero; by \cref{rem:direct-projected-continuation}
it stays below \(\alpha_-\), and the inverse projection continues the
negative graph.

Suppose also \(b\ge \underline b>0\). The drifts
\eqref{eq:direct-positive-secondary-drifts} have
\(K_\sigma'(0)=-\sigma b\) and, in the notation of
\eqref{eq:direct-curvature-coefficient-cubic},
\(R_{K_\sigma}(w)=-\sigma rR_q(\sigma srw)=O_{C^2}(r)\). Hence
\eqref{eq:direct-curvature-coefficient-cubic} gives, on these compact
intervals,
\[
 \frac{wK_+''-K_+'}{w^2}=\frac b{w^2}+O(r)>c,\qquad
 \frac{wK_-''-K_-'}{w^2}=-\frac b{w^2}+O(r)<-c .
\]
Since the negative state stays in a fixed interval away from zero,
\cref{lem:curvature-criterion} gives \(\mathcal L''<0\) on this root
interval, including all approaches to its ends.

To extend this comparison to larger central heights, join the root
interval to the regular prefix.
Apply \cref{lem:prefix-root-cover} with \(\delta=s\), \(\ell=sr\),
\(\epsilon_3=r\), comparison drifts \(a-\sigma bx-x^2\),
\(R=3\alpha_+/4\) and \(w_a=\alpha_+/2\).
The positive drift margin and the gap \(R-w_a\) are fixed, and the
boundary-flow inverse is defined up to \(\alpha_+\) at \(Y_f\).
Since \(\mathcal L_{TT}<0\) on this root range, the lemma gives at most
three zeros of \(\Delta_p\) on \(\cH_p\).

\subsubsection{Small linear coefficient: comparison in root-distance coordinates}

For small \(b\), the curvature comparison loses its fixed margin.
We use linearizing coordinates for the two scalar boundary flows to
obtain a power relation between the distances to the roots, then
combine it with the exterior contraction. Put
\[
 k_\sigma=-F_\sigma'(\alpha_\sigma)
                  =2+O(b+r)>c .
\]
There is a positive root-distance coordinate
\(\ell_\sigma(w)\), decreasing to zero, such that
\begin{equation}
\label{eq:direct-positive-root-distance}
 F_\sigma(w)\ell_\sigma'(w)=-k_\sigma\ell_\sigma(w),
 \qquad \ell_\sigma'(\alpha_\sigma)=-\alpha_\sigma .
\end{equation}
To construct this coordinate, factor out \(\alpha_\sigma-w\).
The logarithmic derivative of the remaining unit is
\(-k_\sigma/F_\sigma+1/(\alpha_\sigma-w)\), whose pole
is removable. Integration and normalization at the root
give uniform derivative bounds and a regular inverse.
Extend to the regular entry using the same scalar equation.
At \(b=r=0\) the formulas are
\[
 F_\sigma=(a-w^2)/w,\qquad k_\sigma=2,\qquad
 \ell_\sigma=(a-w^2)/2 .
\]

The constants
\[
 c_\sigma^{\rm lin}=\ell_\sigma(x_{L,\sigma})
              \exp[k_\sigma \mathcal A_\sigma(x_{L,\sigma})]
\]
are independent of the entry chosen, by
\eqref{eq:direct-positive-root-distance}.
They have positive bounded values and smooth dependence on
the compact coefficients. The phase identities give
\begin{equation}
\label{eq:direct-positive-root-distance-phase}
 \ell_\sigma(w_\sigma)
 =c_\sigma^{\rm lin}(Y_f/s)^{-sk_\sigma}
       \e^{-k_\sigma T_Q}\e^{-sk_\sigma\Xi_\sigma(\xi_\sigma)} .
\end{equation}
Put \(x=\ell_+(w_+)\), \(y=\ell_-(w_-)\), and
\(\nu_k=k_-/k_+\). Since \(\nu_k k_+=k_-\),
\[
 \frac{y}{x^{\nu_k}}
 =\frac{c_-^{\rm lin}}{(c_+^{\rm lin})^{\nu_k}}
   \exp\!\left(-sk_-[\Xi_-(\xi_-)-\Xi_+(\xi_+)]\right).
\]
The exponential factors involving \(T_Q\) and \(Y_f/s\) cancel, giving
\begin{equation}
\label{eq:direct-positive-two-root-power}
 y=u_px^{\nu_k}(1+\rho_p(x)),\qquad
 \nu_k=k_-/k_+=1+O(b+r),\qquad
 u_p=c_-^{\rm lin}/(c_+^{\rm lin})^{\nu_k}=1+O(b+r).
\end{equation}
At \(b=r=0\), \(c_+^{\rm lin}=c_-^{\rm lin}=a/2\) and
\(\nu_k=1\), so the leading unit is exactly one. Its estimate
is uniform because \(a\ge a_0^+\).

On a sufficiently small fixed collar \(0<x<x_c\), a fixed root
distance forces a positive phase gap on both legs. The argument of
\cref{lem:direct-weighted-phase}, now using
\(F_\sigma\asymp\ell_\sigma\) near the simple roots, gives
\[
 |\partial_T^j\xi_\sigma|
 \le\varepsilon_{N,j}(s)\ell_\sigma(w_\sigma)^N,
 \qquad\varepsilon_{N,j}(s)\to0
\]
for every fixed \(N,j\). The two distances are comparable through fixed
positive powers by \eqref{eq:direct-positive-root-distance-phase}, and
\(d\log x/dT=-k_+\partial_T\tau_+\) stays bounded away from zero.
Thus substitution in \eqref{eq:direct-positive-two-root-power} gives
\[
 |(x\partial_x)^j\rho_p(x)|\le C_jx^\vartheta,
 \qquad 0\le j\le2,
\]
with fixed \(\vartheta>0\), uniformly as \(h\downarrow0\).

In these root-distance coordinates the endpoints are
\[
 I_\sigma(\ell_\sigma)=b_\sigma+E\mathcal F_\sigma(\ell_\sigma),\qquad
 \mathcal F_\sigma(0)=0,\quad
 c\le\mathcal F_\sigma'\le C,\quad
 \|\mathcal F_\sigma\|_{C^3}\le C .
\]
The normalization
\(\ell_\sigma'(\alpha_\sigma)=-\alpha_\sigma\) and
\eqref{eq:direct-positive-secondary-endpoints} give
\[
 \mathcal F_-'(0)/\mathcal F_+'(0)
                   \le\exp(2C_{\rm cone}\sqrt{z_*}+C(s+b+r)).
\]
If some central height in the collar lies in \(\cH_p\), the endpoints
on the collar, including the reference endpoints \(b_\sigma\), differ
from admissible endpoints by \(O(E)\), so they lie in \(\widetilde W_\pm\)
by \cref{lem:exterior-multiplier-new}.
In these coordinates, the logarithmic slope ratio is
\[
 \mathcal L(x)
 =\log\left[\Pi_p'(I_-(y(x)))
       \frac{\mathcal F_-'(y(x))y'(x)}{\mathcal F_+'(x)}\right]
 =c_p+(\nu_k-1)\log x+e_p(x),
\]
where
\[
 c_p=\log\left[\Pi_p'(b_-)
        \frac{\mathcal F_-'(0)}{\mathcal F_+'(0)}u_p\nu_k\right]
       <-c_0,\qquad |e_p|+|xe_p'|\le Cx^\vartheta.
\]
Choose \(b_{*,+}\), then \(r,s\), for this strict margin.
The nonlinear exterior contributes \(O(Ey(x))\) with
the same Euler estimate. Since
\(dx/dh=I_{+,h}/(E\mathcal F_+'(x))>0\), the common factor
\(dx/dh\) cancels from the displayed quotient, and \(x_T<0\).
\Cref{lem:one-logarithm-zero-count} shows that every zero in a
sufficiently small fixed collar satisfies \(x\mathcal L_x<0\), and that
such zeros occur only if \(\nu_k<1\). Consequently
\(\mathcal L_T>0\) at every collar zero.

Choose \(R=(1-\epsilon_R)\alpha_+\), with fixed small \(\epsilon_R>0\), strictly inside
the prepared collar. The positive drift margin is fixed. Apply
\cref{lem:prefix-root-cover} on the same scales, choosing \(w_a<R\)
inside this collar with a fixed gap. Every collar zero has \(\mathcal L_T>0\), so
\(\Delta_p\) has at most two zeros for \(0\le b\le b_{*,+}\).
Taking \(b\ge b_{*,+}/2\) in the curvature
argument gives an overlap and covers the compact \(a\) region.

\subsubsection{The central weak root}
\label{sec:positive-quadratic-weak-root}

It remains to allow the constant coefficient to approach zero. For
small \(a>0\), normalization in
\eqref{eq:direct-positive-secondary-chart} forces \(b=1\).
Put \(\kappa=s/a\). With \(\ell=sra\), the scales
\eqref{eq:weak-root-scales} hold with \(\hat r=r\) and \(\epsilon_2=-a\).
At the radius \(\ell_d=srd\) the negative drift is
\[
 K_d(x)=a_d+x-dx^2+rd^2x^3R_q(-\ell_dx),
\]
so \eqref{eq:weak-root-negative-drift} holds for small fixed
\(d_{\max}=d_0\) and \(r\le r(d_0)\).

For \eqref{eq:weak-root-outer}, consider the negative graphs with
\(z<Y_b|\eta|^2\) at \(|\eta|=srd_0\). In the secondary radius
\(X=|\eta|/(sr)\) the negative drift is
\[
 K_-(X)=a+X-X^2+rX^3R_q(-srX),
\]
which is the drift \(K_1\) of \cref{cor:weak-root-negative-first-root},
since \(\ell/a=sr\). Its first root is \(\beta=1+O(a+r)\le2\), and
\(K_-'/X\ge-C\) on \([d_0/2,\beta]\).
\Cref{cor:weak-root-negative-first-root}, with \(d_{\max}=d_0\) and
\(\beta^*=2\), gives \(\mathfrak H_-\ge-C-1\) and \(|I_{-,T}|=o(1)\) for
these graphs, up to the limiting root. \Cref{prop:weak-root-passage} gives
at most two zeros of \(\Delta_p\) for \(0<a\le a_+\) and
\(0<\kappa\le\kappa_+\), where \(a_+\) and \(\kappa_+\) are its constants
\(a_\circ\) and \(\kappa_\circ\).

The complement of \(\kappa\le\kappa_+\) in the region \(0<a\le a_+\) has
\(|q_0|/q_1^2=a/s<1/\kappa_+\), and lies in the strip
\(|q_0|\le Lq_1^2\) for every \(L>1/\kappa_+\). The compact-\(a\)
argument, applied down to a threshold \(a_0^+<a_+\), overlaps the
weak-root region. The curvature, root-distance, and weak-root
comparisons give bounds three, two, and two, respectively, so their
maximum is three outside that strip. The thresholds for \(s,r\) are those
of the three regions and are independent of \(L\).
This proves \cref{prop:direct-positive-quadratic}.

\section{Negative quadratic drift: annulus and buffered region}
\label{sec:direct-negative-annulus-buffer}

For \(q_2<0\), a half-orbit can pass from the secondary scale to the
primary scale through an annulus. The two endpoint rates have the same
leading term there; their comparison therefore depends on the first
radial corrections and on errors that vanish with the radial variables.
We first express these rates in the common central height. We then extend
the negative-branch estimate to its primary root and, for positive
half-orbits that cross the annulus, extend the positive-branch estimate
to larger radii. The buffered region below uses these estimates to count
zeros on the whole height interval; the root regimes of
\cref{sec:direct-negative-root-regimes} use them to compare the two
branches near positive roots and their degeneracies.

\subsection{The annulus for negative quadratic drift}
\label{sec:direct-negative-annulus}

For the rest of the secondary analysis put
\begin{equation}
\label{eq:direct-negative-secondary-scales}
 q_0=-as^3r^2,\qquad q_1=bs^2r,\qquad q_2=-s,\qquad
 0<a\le1,\quad 0\le b\le1,\quad \max\{\sqrt a,b\}=1.
\end{equation}
The positive scales \(s,r\) tend to zero independently.
For a signed radius \(\eta=\sigma s\rho\), define
\[
 u=r/\rho,\quad v=\rho,\quad
 D_\sigma(\rho)=-q(\sigma s\rho)/s^3,\quad
 K_\sigma(u,v)=1+au^2-\sigma bu-\sigma vR_q(\sigma sv).
\]
The variable \(u\) decreases as the radius moves away from the secondary
scale, whereas \(v\) increases toward the primary scale. For the uniform
estimates we regard \(u,v\) as independent base variables, including at zero.
On the curve \(u=r/\rho,\ v=\rho>0\), we have
\(K_\sigma(r/\rho,\rho)=D_\sigma(\rho)/\rho^2\).
Choose a large fixed \(K_0\) and a small fixed \(v_1>0\).
On a slightly enlarged box \(0\le u\le1/K_0\),
\(0\le v\le2v_1\), the functions \(K_\sigma\) have uniform
positive bounds and uniform derivative bounds. Each branch
entering the thin annular box therefore has positive outward
radial flux. The regional arguments below identify the central heights
in \(\mathcal H_p\) whose half-orbits reach this box. The radial Euler
derivative is
\[
 \mathcal E=\rho\partial_\rho
       =-u\partial_u+v\partial_v=:\mathcal E_{uv}.
\]
Here \(\mathcal E\) acts on the exit radius, and \(\partial_T\), with
\(T=-s\log h\), on the central height.
With \(\kappa_3=-R_q(0)=1+O(s)\), Taylor's formula gives
\begin{equation}
\label{eq:direct-annulus-drift}
 K_\sigma=1-\sigma bu+\sigma\kappa_3 v
                       +O_{\mathcal E_{uv}^j}((u+v)^2).
\end{equation}
The notation \(O_{\mathcal E_{uv}^j}\) includes every fixed
derivative order required below, with uniform constants.

\subsubsection{Transport and central height}

We begin by integrating the height growth through the annulus. The
resulting first-exit equation will relate the exit radius to the central
height and determine their order.
Put \(y=z/(s^3\rho^2)\), and use the fast time
\(d\tau=d\log\rho/s\). The extended equations are
\[
 u_\tau=-su,\qquad v_\tau=sv,\qquad
 y_\tau=(K_\sigma^{-1}-2s)y+y^2\mathcal N_\sigma,
\]
\[
 \mathcal N_\sigma=-\frac{\mathcal B_\sigma}
                  {K_\sigma(K_\sigma+\mathcal B_\sigma y)},
 \quad
 \mathcal B_\sigma=B(\sigma sv,s^3v^2y).
\]
On \(0\le y\le2\delta_{\rm an}\), with \(\delta_{\rm an}>0\) fixed small,
the linear coefficient \(K_\sigma^{-1}-2s\), the logarithmic rate
\(y_\tau/y\), and the derivative \(\partial_y y_\tau\) are bounded
below by positive constants.
We take \(\delta_{\rm an}\) also so small that \(y_0=\delta_{\rm an}\) is an
admissible thin threshold in \cref{lem:direct-adaptive-positive} for every
drift factor \(\mathcal K\) with \(1/8\le\mathcal K\le2\); this smallness
depends only on these bounds and on \(B\).
Since \(y_\tau/y=1/(K_\sigma+\mathcal B_\sigma y)-2s\),
\begin{equation}
\label{eq:direct-annulus-level-crossing}
 \frac{d\log y}{d\log\rho}
   =\frac1{s(K_\sigma+\mathcal B_\sigma y)}-2 .
\end{equation}
At \(y=\delta_{\rm an}\) the right side is positive wherever
\(0<s(K_\sigma+\mathcal B_\sigma\delta_{\rm an})<1/2\), and graphs cross
this level only upward there. This holds for both signs on the annular
box, for \(\sigma=-\) also on \(\rho\le0.9\), where \(K_-\) keeps a positive margin,
and for \(\sigma=+\) wherever \(K_+\) has a positive margin and
\(sK_+=O(s+|\eta|)\) is small.
To integrate this growth exactly, we construct a height coordinate with
uniform bounds on its required derivatives,
\[
 \Psi_\sigma=y\exp\Theta_\sigma,\qquad
 \Theta_\sigma=y\widetilde\Theta_\sigma,\qquad
 \Psi_{\sigma,\tau}=(K_\sigma^{-1}-2s)\Psi_\sigma,
 \qquad c<\Psi_{\sigma,y}<C .
\]
As in \cref{lem:direct-adaptive-positive}, extend the coefficients
by bounded smooth retractions into the enlarged positive box, preserving
the height-rate margins. The backward base flow is
\((u\e^{st},v\e^{-st})\), with order-\(j\) derivatives bounded by
\(\e^{jst}\). Thus \cref{lem:direct-contracting-height-coordinate}
applies with \(c_x=s\), after reducing \(s\), and supplies all the
stated height and inverse bounds.

Consequently, for \(\widehat z=s^3\rho^2\Psi_\sigma\),
\begin{equation}
\label{eq:direct-annulus-action}
 \widehat z(\rho)=
 \exp\left(\frac1s\int_{\rho_i}^{\rho}
                   \frac{dt}{tK_\sigma(r/t,t)}\right)
 \widehat z(\rho_i),\qquad \rho_i=K_0r.
\end{equation}
Suppose the central prefix reaches \(\rho_i\) with adapted height
\begin{equation}
\label{eq:direct-annulus-prefix}
 \widehat z(\rho_i)
       =h \e^{\mathcal A_{\sigma,i}/s}\Upsilon_\sigma(\xi_\sigma),
 \qquad
 \xi_\sigma=\e^{\mathcal A_{\sigma,i}/s}
                    h/(s^3r^2),\qquad \Upsilon_\sigma(0)=1.
\end{equation}
The positive units \(\Upsilon_\sigma\), their reciprocals, and their required height
derivatives are uniformly bounded for small \(\xi_\sigma\).
After reducing this fixed interval, the maps
\(\xi\mapsto\xi\Upsilon_\sigma(\xi)\) have uniformly positive
derivatives and uniformly bounded inverse derivatives.
On regular prefixes this preparation follows from
\cref{lem:direct-source-height}. Across a bottleneck it follows from the
composition \eqref{eq:direct-bottleneck-inverse-composition}; this case
is carried out in \cref{sec:direct-no-root-fold}. To separate the possibly
unbounded prefix action \(\mathcal A_{\sigma,i}\) from the annular
increment, write
\[
 \mathcal A_\sigma(\rho)=\mathcal A_{\sigma,i}
                +\int_{\rho_i}^{\rho}
                      \frac{dt}{tK_\sigma(r/t,t)} .
\]
At the first thin exit after \(\rho_i\), \(y=\delta_{\rm an}\),
\begin{equation}
\label{eq:direct-annulus-source}
 s\log(h/s^3)=-\mathcal A_\sigma(\rho)+2s\log\rho
       +s\log\Psi_\sigma(u,v,\delta_{\rm an})
       -s\log\Upsilon_\sigma(\xi_\sigma).
\end{equation}

Take \(u\le u_0<1/K_0\), with a sufficiently large fixed
gap from the inner cut, and write
\[
 \mathcal J_\sigma(\rho)=\mathcal A_\sigma(\rho)-\mathcal A_{\sigma,i}
       =\int_{K_0r}^{\rho}\frac{dt}{tK_\sigma(r/t,t)}.
\]
This action increment satisfies $\mathcal J_\sigma\ge c\log(\rho/(K_0r))$.
The equation for the central height becomes
\[
 \xi_\sigma\Upsilon_\sigma(\xi_\sigma)=\mathfrak b_\sigma(\rho),\qquad
 \mathfrak b_\sigma(\rho)
       :=u^{-2}\e^{-\mathcal J_\sigma(\rho)/s}\Psi_\sigma(u,v,\delta_{\rm an})>0.
\]
Since $\mathcal A_{\sigma,i}$ is constant under $\mathcal E$,
for every prescribed $j\ge1$,
\[
 \mathcal E^j\mathcal J_\sigma
       =\mathcal E_{uv}^{j-1}(K_\sigma^{-1}),\qquad
 |\mathcal E^j\mathcal J_\sigma|\le C_j.
\]
Thus $\mathcal E\log\mathfrak b_\sigma
 =2-(sK_\sigma)^{-1}+\mathcal E_{uv}\log\Psi_\sigma$,
and all higher fixed Euler derivatives of its logarithm are $O_j(s^{-1})$.
Differentiating its exponential yields
\begin{equation}
\label{eq:direct-annulus-input-jets}
 |\mathcal E^j\mathfrak b_\sigma|
       \le C_js^{-j}\mathfrak b_\sigma,\qquad
 \mathfrak b_\sigma\le Cu^{c/(2s)-2},\qquad j\ge0,
\end{equation}
after decreasing the fixed $u_0$.

The uniformly regular map
$\mathcal C_\sigma(\xi)=\xi\Upsilon_\sigma(\xi)$ fixes zero.
For small $s$, $\mathfrak b_\sigma\le1$ lies in its inverse domain;
the chain rule and \eqref{eq:direct-annulus-input-jets} give
$|\mathcal E^j\mathcal C_\sigma^{-1}(\mathfrak b_\sigma)|
\le C_js^{-j}\mathfrak b_\sigma$. Since
$\log\Upsilon_\sigma(0)=0$, the same estimate holds for its composition.
Thus, for every fixed $N,j$,
after reducing $s$,
\begin{equation}
\label{eq:direct-annulus-prefix-weight}
 |\mathcal E^j\xi_\sigma|
       +|\mathcal E^j\log\Upsilon_\sigma(\xi_\sigma)|
                         \le C_{N,j}u^N .
\end{equation}
Indeed the excess power of $u\le u_0<1$ in
\eqref{eq:direct-annulus-input-jets} absorbs $s^{-j}$.
The small expanded incoming height and the outward increase of
unscaled height show that the reconstructed prefix is thin.

The event equation gives
\[
 s\partial_{\log\rho}
       \log\bigl[s^3\rho^2\Psi_\sigma(u,v,\delta_{\rm an})\bigr]
       -1/K_\sigma
 =-(y_\tau|_{y=\delta_{\rm an}})\,
                      \partial_y\log\Psi_\sigma(u,v,\delta_{\rm an})<-c.
\]
Hence the central height decreases strictly with the exit radius, with inverse
derivatives \(O(s^j)\) in logarithmic incoming height. Scalar order
makes the first-exit domains and radius subintervals into intervals of
central heights.

Define the positive central-height rate
\(G_\sigma=-s\mathcal E\log h_\sigma\), and put
\(g_0=1-2s\). From \eqref{eq:direct-annulus-source},
\begin{equation}
\label{eq:direct-annulus-source-rate}
 \begin{split}
 G_\sigma&=K_\sigma^{-1}-2s
       -s\mathcal E_{uv}\log\Psi_\sigma
       +s\mathcal E\log\Upsilon_\sigma\\
 &=g_0+\sigma bu-\sigma\kappa_3 v
        +O_{\mathcal E^j}((u+v)^2+s(u+v)),\qquad c<G_\sigma<C .
 \end{split}
\end{equation}
For \(j\ge1\), the bound
\(\mathcal E_{uv}^j\log\Psi_\sigma=O(u+v)\) follows from
uniform base derivatives and the vanishing coefficients of
\(\mathcal E_{uv}\).

\subsubsection{The endpoint slope and its vanishing error}

The first-exit equation controls differentiation in central height.
We now follow the graph from the thin exit to the fixed section,
retaining the radial weight in the slope error. This weight is needed
when the leading terms of the two rates cancel.

\begin{lemma}[Annular endpoint and central-height rates]
\label{lem:direct-annulus-rates}
Under \eqref{eq:direct-annulus-prefix}, consider a thin exit
in the smaller annular box just specified. Let \(I_\sigma\)
be its endpoint and \(M_\sigma=-\mathcal E I_\sigma\).
Then \(M_\sigma>0\), and
\begin{equation}
\label{eq:direct-annulus-endpoint-rate}
 \mathcal E\log M_\sigma=2g_0+e_\sigma,\qquad
 |\mathcal E^j e_\sigma|
       \le C_js(1+|\log s|)(u+v),\quad 0\le j\le2,\qquad
 M_\sigma\le C(s\rho)^{2g_0} .
\end{equation}
In the common logarithmic height coordinate \(T=-s\log h\), put
\(m_\sigma=hI_{\sigma,h}=sM_\sigma/G_\sigma\) and
\(\mathfrak H_\sigma=\partial_T\log m_\sigma\).
Then
\begin{equation}
\label{eq:direct-annulus-clock-rate}
 \begin{split}
 \mathfrak H_\sigma
   &=2+\sigma(-\gamma_sbu+\delta_s\kappa_3 v)+\varepsilon_\sigma,\\
 \gamma_s&=(1-4s)/g_0^2,\qquad
 \delta_s=(3-4s)/g_0^2,\\
 |\varepsilon_\sigma|+|\partial_T\varepsilon_\sigma|
   &\le C(u+v)\{u+v+s(1+|\log s|)\}.
 \end{split}
\end{equation}
The two signs may have unrelated exit radii.
\end{lemma}

\begin{proof}
After the thin exit use \(\mathcal Y=z/|\eta|^2\) as the independent variable,
and write \(\mathfrak l=\log(|\eta|/(s\rho))\).
Until a small fixed \(\mathcal Y=m_{\rm land}\), its equation is
\begin{equation}
\label{eq:direct-annulus-landing-clock}
 \frac{d\mathfrak l}{d\mathcal Y}
 =\frac{sK_\sigma(u\e^{-\mathfrak l},v\e^{\mathfrak l})+B_\sigma \mathcal Y}
        {\mathcal Y[1-2sK_\sigma(u\e^{-\mathfrak l},v\e^{\mathfrak l})-2B_\sigma \mathcal Y]},
 \quad \mathfrak l(s\delta_{\rm an})=0,\quad
 B_\sigma=B(\sigma sv\e^{\mathfrak l},s^2v^2\e^{2\mathfrak l}\mathcal Y).
\end{equation}
The bracketed factor in the denominator has a fixed positive margin.
Integration gives
\[
 0\le\mathfrak l\le C_Bm_{\rm land}
      +C_Ks\log\frac{m_{\rm land}}{s\delta_{\rm an}}.
\]
Here \(C_B,C_K\) use only the coefficient bounds on the enlarged
drift box and the denominator margin. Choose the common landing level
so that \(C_Bm_{\rm land}<\tfrac12\log(0.95/0.9)\), then reduce
\(s\), with \(\delta_{\rm an}\) fixed, so that the second term is
below the same bound. Thus \(\mathfrak l<\log(0.95/0.9)<\log2\).
A stopped bootstrap keeps the graph in the enlarged positive annulus.
For the negative-branch extension to exits \(\rho\le0.9\) below,
use the fixed enlarged band \(v<0.95\):
\(K_-=1+au^2+bu-v+O(s)\) has a positive margin there for small
\(s\), and the stricter growth bound keeps the landing inside that
band. It lands in the
cone at
\[
 t=m_{\rm land}s^2\rho^2\e^{2\mathfrak l},\qquad U(t)=t/(2m_{\rm land}).
\]
Admissible graphs land before their endpoint: that endpoint has
\(z/|\eta|^2\) bounded below, and \(m_{\rm land}\) is chosen smaller.

At \(u=v=0\), \(K=1\) and \(B_\sigma=B_0(p)\), so both signs
have the landing shift
\[
 \mathfrak l_0=\frac{s}{g_0}\log\frac{m_{\rm land}}{s\delta_{\rm an}}
       -\frac{1}{2g_0}
          \log\frac{g_0-2B_0m_{\rm land}}{g_0-2B_0s\delta_{\rm an}}.
\]
Write \(R(u,v,\mathfrak l,\mathcal Y)\) for the right side of
\eqref{eq:direct-annulus-landing-clock}, and put
\(\epsilon_s=s(1+|\log s|)\). Every positive-order mixed
\((u,v,\mathfrak l)\) derivative of \(B_\sigma\) has a factor
\(s\), while \(K_\sigma\) occurs multiplied by \(s\).
The denominator margin therefore gives
\[
 |\partial_{u,v,\mathfrak l}^{\alpha}R|
       \le C_\alpha(s/\mathcal Y+s),\qquad 1\le|\alpha|\le4.
\]
The difference from the equation at the base origin has this bound
times \(u+v\). The finite-order variational argument used for
\eqref{eq:direct-adaptive-landing} now has fixed initial height
\(s\delta_{\rm an}\), zero initial base variations, and
\[
 \int_{s\delta_{\rm an}}^{m_{\rm land}}(s/\mathcal Y+s)\,d\mathcal Y
       \le C\epsilon_s.
\]
It gives \(\|\partial_{u,v}^{\alpha}\mathfrak l\|\le
C_\alpha\epsilon_s\) for \(1\le|\alpha|\le4\), and therefore
\(|\mathfrak l-\mathfrak l_0|\le C\epsilon_s(u+v)\).
At positive orders, \(\mathcal E_{uv}^j\) is a finite sum of
operators \(c_{i_1i_2}u^{i_1}v^{i_2}\partial_u^{i_1}\partial_v^{i_2}\), with
\(1\le i_1+i_2\le j\); each coefficient vanishes at the base
origin. Thus
\begin{equation}
\label{eq:direct-annulus-landing-weight}
 |\mathcal E_{uv}^j(\mathfrak l-\mathfrak l_0)|
       \le C_js(1+|\log s|)(u+v),\qquad 0\le j\le4.
\end{equation}
In particular \(\mathcal E\log t=2+2\mathcal E_{uv}\mathfrak l>0\).

Choose \(c_{\rm land}=1/(2m_{\rm land})>4\sup B\).
The moving-cone tail in \cref{lem:direct-weighted-tail} gives
\[
 \frac{dU(z_*;c_{\rm land}t)}{dt}
       =(z_*/t)^{2s}\tilde a_\sigma(t),\qquad c<\tilde a_\sigma<C,
\]
where, for \(1\le j\le3\),
\[
 |(t\partial_t)^j\log\tilde a_\sigma(t)|
 \le C_j(\sqrt t+q_1/\sqrt t+|q_0|/t)
 \le C_js(v+bu+u^2).
\]
At this landing the coefficients satisfy the cone smallness
condition. Hence
\begin{equation}
\label{eq:direct-annulus-exact-slope}
 M_\sigma=2(1+\mathcal E_{uv}\mathfrak l)\,
             z_*^{2s}t^{\,1-2s}\tilde a_\sigma(t).
\end{equation}
Put \(\mathfrak d=\mathcal E_{uv}\mathfrak l\), \(b_t=2(1+\mathfrak d)\), and
\(A_j=(t\partial_t)^j\log\tilde a_\sigma\). Logarithmic
differentiation of \eqref{eq:direct-annulus-exact-slope}
gives
\[
 e_\sigma=2g_0\mathfrak d+\frac{\mathcal E_{uv}\mathfrak d}{1+\mathfrak d}
          +b_tA_1,
 \qquad \mathcal E A_j=b_tA_{j+1}.
\]
Here \(\mathfrak d,\mathcal E_{uv}\mathfrak d,\mathcal E_{uv}^2\mathfrak d,
 \mathcal E_{uv}^3\mathfrak d\) are \(O(\epsilon_s(u+v))\), whereas
\(A_1,A_2,A_3=O(s(u+v))\). Two further Euler derivatives
of the displayed error involve derivatives of \(\mathfrak l\)
only through order four and of \(\log\tilde a_\sigma\) only
through order three; every term still contains one of these
vanishing factors. The denominators \(1+\mathfrak d\) stay bounded
away from zero. This proves the error bounds in
\eqref{eq:direct-annulus-endpoint-rate}; the bound on \(M_\sigma\)
follows from \eqref{eq:direct-annulus-exact-slope} and
\(t\asymp(s\rho)^2\).

To express the endpoint estimate in the common central height, use
\(\partial_T=\mathcal E/G_\sigma\). Thus
\[
 \mathfrak H_\sigma
   =\frac{\mathcal E\log M_\sigma
                      -\mathcal E\log G_\sigma}{G_\sigma}.
\]
Insert \eqref{eq:direct-annulus-source-rate} and
\eqref{eq:direct-annulus-endpoint-rate}. The constant term
is exactly \(2g_0/g_0=2\); the two linear coefficients
are those in \eqref{eq:direct-annulus-clock-rate}.
Differentiation using \(u_T=-u/G_\sigma\) and
\(v_T=v/G_\sigma\) preserves the stated weighted error.
\end{proof}

For comparable exit radii \(\omega=\Omega\rho\), with
fixed positive bounds on \(\Omega\),
\begin{equation}
\label{eq:direct-annulus-slope-calibration}
 \log\frac{M_-(\omega)}{M_+(\rho)}
 \le 2g_0\log\Omega+2C_{\rm cone}\sqrt{z_*}
               +Cs(1+|\log s|)(u+v).
\end{equation}
Indeed the shifts in \eqref{eq:direct-annulus-exact-slope}
have the same base value \(\mathfrak l_0\), and the moving-section
factor is
\[
 c_{\rm land}-\mathcal G_\sigma(c_{\rm land}t,t)
       =c_{\rm land}g_0-B_0+O(s(u+v)).
\]
The remaining factor of \(\tilde a_\sigma\) is the tail variation
\eqref{eq:direct-tail-variation} divided by \(\varpi_t\), that is, the slope
\(a_{t,p}\) of \eqref{eq:direct-centered-tail} at \(\mathcal U_\circ=c_{\rm land}\).
This landing energy lies in the common cone, so on each branch
\eqref{eq:direct-tail-first-slope} bounds the logarithm of that factor by
\(C_{\rm cone}(\sqrt{z_*}+q_1/\sqrt t)\), and \(q_1/\sqrt t\le Csu\). This
proves \eqref{eq:direct-annulus-slope-calibration}.

\begin{corollary}[Two annular exits]
\label{cor:direct-annular-pair}
Under \eqref{eq:direct-annulus-prefix} on both branches, let the thin
exits from an interval of \(\mathcal H_p\) lie in the smaller annular box
\(u_\sigma\le u_0\), \(v_\sigma\le v_0\). In \(T=-s\log h\),
\begin{equation}
\label{eq:direct-annular-pair-curvature}
 \begin{split}
 (\mathfrak H_--\mathfrak H_+)'&=
 -\frac{\gamma_sb(u_-+u_+)+\delta_s\kappa_3(v_-+v_+)}{g_0}\\
 &\quad+O\left(\sum_\sigma(u_\sigma+v_\sigma)
       \{u_\sigma+v_\sigma+s(1+|\log s|)\}\right).
 \end{split}
\end{equation}
The negative annular branch alone gives
\begin{equation}
\label{eq:direct-annular-exterior-curvature}
 |\partial_T^2[(\log\Pi_p')(I_-)]|\le C(sv_-)^{2-4s}
                    \le Cs^{3/2}v_-,\qquad s<1/8,\quad v_-<1.
\end{equation}
For any fixed \(b_{\min}>0\), choosing the annular margins small relative to
\(b_{\min}\), then reducing \(s\), gives \(\mathcal L''<0\) when \(b\ge b_{\min}\).
\end{corollary}

\begin{proof}
Differentiate \eqref{eq:direct-annulus-clock-rate} using
\(\partial_T=\mathcal E_\sigma/G_\sigma\),
\(u_{\sigma,T}=-u_\sigma/G_\sigma\), \(v_{\sigma,T}=v_\sigma/G_\sigma\),
and \(G_\sigma=g_0+O(u_\sigma+v_\sigma)\); this gives
\eqref{eq:direct-annular-pair-curvature}.
By \cref{lem:direct-annulus-rates}, \(I_{-,T}=-M_-/G_-\), and
\(|I_{-,T}|+|I_{-,TT}|\le C(sv_-)^{2-4s}\).
Equation \eqref{eq:direct-exterior-curvature} therefore gives the first
bound in \eqref{eq:direct-annular-exterior-curvature}; the second follows
from \((sv_-)^{2-4s}/v_-=s^{2-4s}v_-^{1-4s}\le s^{3/2}\).
For \(b\ge b_{\min}\), the leading negative term dominates both the weighted
error and the exterior contribution, proving \(\mathcal L''<0\).
\end{proof}

\subsubsection{The negative branch beyond the annulus}
\label{sec:negative-branch-far-root}

Beyond the annulus, negative outputs may lie in compact primary
intervals or near the primary root close to \(X=1\). We continue the
annular calculation on the compact intervals. Near the root, where the
drift tends to zero, the boundary-flow coordinate gives a negative
endpoint rate and estimates its derivative with the same drift weight
as the endpoint derivatives. These weights allow the exterior
contribution to be absorbed in the regional comparisons.

For the negative branch fix the regular primary entry at
\(\rho_{\rm pr}=0.8\), take \(0.85\) as the lower bound for the projected
far-root state \(w_n\), and use \(\rho=0.9\) as the last primary
thin-exit test. These thresholds give a fixed phase advance and overlapping
ranges of central heights.
In primary coordinates \(|\eta|=sX,\ z=s^2Y\), the negative boundary
drift is
\[
 K_n(X)=ar^2+brX+X^2+X^3R_q(-sX),\qquad F_n=K_n/X .
\]
It has a simple first far root \(X_*=1+O(r+s)\),
with \(K_n'(X_*)<0\), and is positive below it.
Construct the boundary-flow coordinate on a band from
below \(0.8\) to beyond \(X_*\), with small fixed output
height \(Y_n>0\). Its far-root interval is
\begin{equation}
\label{eq:direct-negative-far-core}
 w_n=\mathcal W_n(X(s^2Y_n;h),Y_n)\ge0.85,\qquad w_n<X_* .
\end{equation}

\begin{lemma}[Negative branch]
\label{lem:negative-branch-far-root}
Fix \(M\ge K_0\). Suppose that the negative graph of a central
height in \(\cH_p\) is thin at \(\rho=Mr\), with the prefix preparation
\eqref{eq:direct-annulus-prefix} at \(\rho_i=K_0r\).
\begin{enumerate}[label=\textup{(\roman*)}]
\item If its first thin exit occurs at \(\rho\le0.9\), then \(\rho>Mr\)
and, in \(T=-s\log h\),
\begin{equation}
\label{eq:direct-negative-primary-rates}
 \begin{split}
 \mathfrak H_-&=2-3\rho+
                O(r/\rho+s(1+|\log s|)),\\
 \mathfrak H_-'&=-3\rho(1-\rho)
                +O(r+s(1+|\log s|))
 \quad\hbox{on fixed compact primary intervals},\\
 |I_{-,T}|+|I_{-,TT}|&\le C(s\rho)^{2-4s}=o(1).
 \end{split}
\end{equation}
\item If the graph is still thin at \(\rho=0.9\), it satisfies
\eqref{eq:direct-negative-far-core} with \(w_n>0.89\). For every central
height with \eqref{eq:direct-negative-far-core},
\begin{equation}
\label{eq:direct-negative-far-rate}
 \begin{gathered}
 \mathfrak H_-=2-3w_n+O(r+s+\varepsilon_{n,s})<-c_n<0,\qquad
 \mathfrak H_-'=F_n(w_n)\{Q_n(w_n)+o(1)\},\\
 Q_n\le-c_n,\qquad |I_{-,T}|+|I_{-,TT}|\le CE_nF_n(w_n),
 \end{gathered}
\end{equation}
where \(Q_n=(XK_n''-K_n')/X^2\), \(E_n=s^2(z_*/(s^2Y_n))^{2s}=O(s^2)\),
and \(\varepsilon_{n,s}\to0\), uniformly up to the far root.
\end{enumerate}
The exit radius in (i) and the state \(w_n\) in (ii) increase with \(T\).
The constants are uniform on the chart
\eqref{eq:direct-negative-secondary-scales}.
\end{lemma}

\begin{proof}
\emph{Part (i).} Since the graph is thin at \(Mr\) and, by
\eqref{eq:direct-annulus-level-crossing}, crosses \(y=\delta_{\rm an}\) only
upward on the annulus, a first thin exit before \(\rho=0.9\) has \(\rho>Mr\).
The negative annular coordinate extends to \(\rho\le0.9\):
its drift factor has a uniform positive margin there.
Its equation for the central height gives
\[
 G_-=\frac1{1-\rho}+O_{\mathcal E^j}(r/\rho+s).
\]
The endpoint argument of
\cref{lem:direct-annulus-rates} on this enlarged compact
base box gives
\[
 \mathcal E\log M_-=2-4s+
        O_{\mathcal E^j}\bigl(s(1+|\log s|)(r/\rho+\rho)\bigr).
\]
The weighted error follows from the landing equation:
the comparison at the base origin has \(K=1\), and
the derivatives remain bounded on \(0\le\rho\le0.9\).
Inserting these expressions into
\((\mathcal E\log M_--\mathcal E\log G_-)/G_-\) gives
\eqref{eq:direct-negative-primary-rates}; the last estimate there uses
\(I_{-,T}=-M_-/G_-\) and its logarithmic rate.

\emph{Part (ii).} Every branch still thin at \(0.9\) reaches the output
\(Y=Y_n\) inside the band. Its initial projected state is
\(0.9+O(s\delta_{\rm an})>0.89\) and lies below \(X_*\).
The projected flow increases toward the root, and the inverse
projection differs by \(O(Y_n)\), remaining in the band.
The inverse-flux identity gives continuation to \(Y_n\).

At \(\rho_{\rm pr}=0.8\), write
\[
 \begin{gathered}
 \mathcal J_b=\int_{K_0r}^{\rho_{\rm pr}}\frac{dt}{tK_-(r/t,t)},
 \qquad
 \mathcal A_b=\mathcal A_{-,i}+\mathcal J_b,\\
 \xi_n=\e^{\mathcal A_b/s}\frac{h}{s^3}.
 \end{gathered}
\]
The secondary and primary expanded inputs are related by
\[
 \xi_-=\beta_b\xi_n,\qquad
 \beta_b=r^{-2}\e^{-\mathcal J_b/s}.
\]
The factor \(r^{-2}\) comes from changing the scale
\(s^3r^2\) of the central height to \(s^3\). The uniform upper bound
for the annular drift gives
\(\mathcal J_b\ge c\log(\rho_{\rm pr}/(K_0r))\). After fixing smaller
upper bounds for \(r\) and then \(s\), this implies
\[
 0<\beta_b\le r^{c/(2s)-2}\le1.
\]

For \(Y_i=z/s^2\), write
\(y_i=Y_i/s=z/s^3=\rho_{\rm pr}^2y_a(\rho_{\rm pr})\), where
\(y_a=z/(s^3\rho^2)\), and set
\[
 \mathcal C_b(y_i)
   =\rho_{\rm pr}^2\Psi_-(r/\rho_{\rm pr},\rho_{\rm pr},
                         y_i/\rho_{\rm pr}^2).
\]
This height coordinate and its inverse have uniform bounds
on the required derivatives, with \(\mathcal C_b'(0)=1\).
Annular transport gives
\[
 \mathcal C_b(y_i)
     =\xi_n\Upsilon_-(\beta_b\xi_n).
\]
As in \eqref{eq:direct-bottleneck-inverse-composition}, inversion
gives \(Y_i=s\xi_nV_n(\xi_n)\), with \(V_n(0)=1\). Since
\(\beta_b\le1\), the unit, its reciprocal, and the map inverse have
uniform derivative bounds on a fixed small \(\xi_n\)-interval.

For central heights with \eqref{eq:direct-negative-far-core}, the projected
state at \(X=\rho_{\rm pr}\) satisfies
\(w_i=\mathcal W_n(\rho_{\rm pr},Y_i)\le\rho_{\rm pr}\). The boundary
flow must advance from at most \(0.8\) to at least
\(0.85\) before the output \(Y_n\). Its bounded
speed and the projection identity give
\[
 s\log(Y_n/Y_i)\ge c_b>0,\qquad
 y_i=Y_i/s\le(Y_n/s)\e^{-c_b/s}.
\]
Thus the graph is thin at \(\rho_{\rm pr}\).
By \eqref{eq:direct-annulus-level-crossing}, it has no downward crossing
of \(y_a=\delta_{\rm an}\) on the preceding annular segment, so the whole
segment is thin. Annular transport gives
\[
 \frac{\widehat z(K_0r)}{s^3r^2}
       =\beta_b\mathcal C_b(y_i).
\]
This lies in the uniform small output domain of the
inner prefix. Its inverse recovers
\(\xi_-=O(\beta_by_i)\); uniqueness of the
flow identifies the recovered central height with \(h\). Thus
\(\xi_n=\xi_-/\beta_b\le Cy_i\) lies in the required small
domain. At the fixed output \(z=s^2Y_n\), the entry
\(X=\rho_{\rm pr}=0.8\) has \(F_n(\rho_{\rm pr})\) bounded below,
and the advance to \(w_n\ge0.85\) supplies a fixed positive phase
gap. Thus \cref{lem:direct-weighted-phase} applies with \(\delta=s\)
in the same coordinate \(T=-s\log h\). We use its drift-weighted
bound \eqref{eq:direct-phase-weight} with \(N=1\) and \(j=2\)
for the curvature estimate below; in particular,
\[
 (w_n)_T=\theta_nF_n(w_n),\qquad
 |\theta_n-1|+|\theta_n'|+|\theta_n''|
                         \le\varepsilon_{n,s}\longrightarrow0 .
\]
\Cref{lem:boundary-flow-endpoint}, with \(\ell=\delta=s\), gives
\[
 -I_{-,w_n}=E_nA_nw_nU_n(w_n),\qquad
 E_n=s^2(z_*/(s^2Y_n))^{2s}=O(s^2),\qquad
 U_n=1+O_{C^3}(s).
\]
The rate identity yields the first estimate in
\eqref{eq:direct-negative-far-rate}, uniformly up to the root, since
\(K_n'(w)/w=2-3w+O(r+s)\). By
\eqref{eq:direct-curvature-coefficient-cubic}, with \(K_n'(0)=br\) and
\(R_K(X)=R_q(-sX)\), one has
\(Q_n(X)=-br/X^2+3R_q(-sX)+O(s)<-c_n\) on the band, since
\(R_q=-1+O(s)\). The weighted phase estimate and
\eqref{eq:direct-rate-curvature} give the remaining assertions of
\eqref{eq:direct-negative-far-rate}. Finally \(\rho_T=\rho/G_->0\), and
the projected far-root state increases along the boundary flow.
\end{proof}

\subsubsection{The positive branch beyond the annulus}

A positive half-orbit may leave compact \(\rho\)-intervals and reach
unscaled radii \(\varrho\gg s\), where the drift grows cubically.
The adaptive coordinate of \cref{lem:direct-adaptive-positive} follows
this transition from the inner cut and gives the endpoint rate in both
the buffered and no-root regions, using their respective preparations
at the inner cut.

The annular action controls the expanded incoming height and its
derivatives with respect to the exit radius. This control makes the
endpoint-rate comparison uniform as the positive exit moves from the
quadratic to the cubic regime.

On the positive branch write \(\varrho=\eta=s\rho>0\), \(Q=s^3r^2\), and
\begin{equation}
\label{eq:positive-adaptive-variables}
 \varrho^+=s+\varrho,\quad u=\frac{sr}\varrho,\quad
 v=\frac\varrho{s+\varrho},\quad
 K_a=1+(1-v)(au^2-bu)-v\{R_q(\varrho)+1\}.
\end{equation}
Here \(u=r/\rho\) retains its annular meaning. In \(K_a\) and the following
lemma, \(v\) is the adaptive variable of
\eqref{eq:direct-adaptive-drift}; the annular coordinate
\(\Psi_+\) still has arguments \((r/\rho,\rho)\). With
\(\varrho_i=K_0sr\), put
\[
 \mathcal J_a(\varrho)=\int_{\varrho_i}^{\varrho}\frac{x\,dx}{-q(x)},
\]
so that the action increment of the annulus is \(\mathcal J_+(\rho)=s\mathcal J_a(s\rho)\).
At every radius the annular and adaptive normalized heights satisfy
\begin{equation}
\label{eq:positive-annular-adaptive-heights}
 \frac{z}{s\varrho^2}
   =\Bigl(1+\frac\varrho s\Bigr)\frac{z}{\varrho^2(s+\varrho)} ,
\end{equation}
so, for a half-orbit thin at \(\varrho_i\), the annular thin exit at the
same threshold \(\delta_{\rm an}\) precedes the adaptive exit of part~(iii)
below.
At an exit radius \(\varrho_e\) we write \(u_e,v_e,\varrho_e^+\) for the
values of \(u,v,\varrho^+\). The operator
\(\mathcal E=\varrho_e\partial_{\varrho_e}\) varies the exit radius at fixed
parameters, and \(T=-s\log h\).

\begin{lemma}[Positive branch beyond the annulus]
\label{lem:positive-branch-adaptive-exit}
Let \(\mathcal J\subset\cH_p\) be an interval of central heights whose
positive half-orbits are thin at \(\varrho_i\), that is,
\(z(\varrho_i;h)<\delta_{\rm an}K_0^2Q\), and satisfy
\begin{equation}
\label{eq:positive-adaptive-prefix}
 z(\varrho_i;h)=Q\,\xi\mathcal V(\xi),\qquad \xi=\e^{\mathcal A_i/s}h/Q .
\end{equation}
Here \(\mathcal A_i\) depends only on \(p\), \(\xi\) lies in a fixed
interval \([0,\bar\xi]\), \(\mathcal V(0)=1\), \(\mathcal V\) and
\(1/\mathcal V\) have uniformly bounded derivatives of the required orders
there, and \(\xi\mapsto\xi\mathcal V(\xi)\) has a uniformly positive
derivative and an inverse with uniformly bounded derivatives.
\begin{enumerate}[label=\textup{(\roman*)}]
\item One has \(-q(\varrho)=\varrho^2(s+\varrho)K_a\). On the box
\(0\le u\le3/(2K_0)\), \(0\le v\le1\), with small unscaled radius and
small \(s\), \(K_a\) lies between \(1/8\) and \(2\) and has uniform
derivative bounds, and \(K_a-1\) and its fixed Euler derivatives are
\(O(u+\varrho^+)\). Thus \cref{lem:direct-adaptive-positive} applies with
\(\mathcal K=K_a\), \(\varrho_0=sr\) and thin threshold \(y_0=\delta_{\rm an}\);
let \(\Psi_a\) be its coordinate and \(\widehat z_a\) its adapted height.
\item Put \(u_i=1/K_0\), \(v_i=K_0r/(1+K_0r)\) and \(c_i=K_0^2(1+K_0r)\),
and define \(\Upsilon_a\) by
\begin{equation}
\label{eq:positive-adaptive-unit}
 \xi\Upsilon_a(\xi)
   =c_i\Psi_a\bigl(u_i,v_i,\varrho_i,\xi\mathcal V(\xi)/c_i\bigr).
\end{equation}
Then \(\Upsilon_a(0)=1\), and \(\Upsilon_a\), its reciprocal and the
inverse of \(\xi\mapsto\xi\Upsilon_a(\xi)\) have uniformly bounded
derivatives on the interval of those \(\xi\le\bar\xi\) with
\(\xi\mathcal V(\xi)\le\delta_{\rm an}K_0^2\). Moreover
\(\widehat z_a(\varrho_i)=Q\xi\Upsilon_a(\xi)\), so that
\(\lambda_{\rm in}=\log\widehat z_a(\varrho_i)\) satisfies
\(d\lambda_{\rm in}/d\log h=1+O(\xi)\).
\item The adaptive exit, the first radius \(\varrho_e>\varrho_i\) with
\(z=\delta_{\rm an}\varrho^2(s+\varrho)\), satisfies
\begin{equation}
\label{eq:direct-no-root-adaptive-input}
 \xi\Upsilon_a(\xi)=\frac{\varrho_e^2\varrho_e^+}{Q}
      \Psi_a(u_e,v_e,\varrho_e,\delta_{\rm an})\e^{-\mathcal J_a(\varrho_e)} .
\end{equation}
Its landing height \(t=m_{\rm land}\varrho_l^2\) satisfies
\(\sqrt t+q_1/\sqrt t+|q_0|/t\le C(\varrho_e+su_e)\), and
\(|q_2|=s\); in particular \eqref{eq:direct-tail-smallness} holds there
for small \(s\).
\item Suppose in addition that \(\varrho_e\ge s\rho_{\rm g}\) on
\(\mathcal J\) and \(\mathcal J_+(\rho_{\rm g})\ge c_{\rm g}\log(1/r)\), for fixed
\(\rho_{\rm g},c_{\rm g}>0\), and that \(r\le1/2\). Then, for every
fixed \(j\) and small \(s\),
\begin{equation}
\label{eq:positive-adaptive-input-weight}
 |\mathcal E^j\xi|+|\mathcal E^j\log\Upsilon_a(\xi)|
       \le C_js^{-3-j}r^{c_{\rm g}/s-2},
\end{equation}
which is \(O(s^N)\) for every fixed \(N\). The rate
\(\mathfrak g=-\varrho_e^+\mathcal E\log h\) satisfies
\begin{equation}
\label{eq:direct-no-root-adaptive-source}
 \mathfrak g=1+O(u_e+\varrho_e^+),\qquad
 \mathcal E\log\mathfrak g=O(u_e+\varrho_e^+),\qquad
 \mathcal ET=(s/\varrho_e^+)\mathfrak g>0 .
\end{equation}
For small fixed \(m_{\rm land}\) and radius box, and after reducing
\(r\) and then \(s\), the endpoint rate
\(\mathfrak H_+=\partial_T\log m_+\) satisfies
\begin{equation}
\label{eq:direct-no-root-adaptive-rate}
 \mathfrak H_+
 =\frac{\varrho_e^+}{s\mathfrak g}
        [\,2-4s+v_e+O(\epsilon_e+u_e+\varrho_e^+)\,]>c\varrho_e^+/s>0,
\end{equation}
with \(\epsilon_e\) as in \eqref{eq:direct-adaptive-landing}.
\end{enumerate}
The constants depend on \(\mathcal V\) only through its assumed bounds.
The constants in \eqref{eq:direct-no-root-adaptive-source} and
\eqref{eq:direct-no-root-adaptive-rate}, apart from terms \(O(s^N)\), are
independent of \(\mathcal V\), \(\mathcal A_i\), \(\rho_{\rm g}\) and
\(c_{\rm g}\).
\end{lemma}

\begin{proof}
\emph{(i)} Since \((s+\varrho)(1-v)=s\) and \((s+\varrho)v=\varrho\),
\[
 \varrho^2(s+\varrho)K_a
   =s\varrho^2+as^3r^2-bs^2r\varrho-\varrho^3R_q(\varrho)=-q(\varrho).
\]
Because \(\kappa_3=-R_q(0)=1+O(s)\), one has \(R_q(\varrho)+1=O(s+\varrho)\)
with its fixed Euler derivatives; here \(\mathcal Eu=-u\),
\(\mathcal Ev=v(1-v)\) and \(\mathcal E\varrho=\varrho\). On the box,
\(|(1-v)(au^2-bu)|\le u<3/4\), since \(0<a\le1\), \(0\le b\le1\) and
\(K_0>2\). This gives the bounds on \(K_a\) and on \(K_a-1\). Along the
passage from \(\varrho_i\), \(u\) decreases from \(1/K_0\) while \(v\) and
\(\varrho\) increase, so the base point stays in the box, and
\(y_0=\delta_{\rm an}\) is admissible by the choice of \(\delta_{\rm an}\)
in \cref{sec:direct-negative-annulus}.

\emph{(ii)} At the inner section,
\(\varrho_i^2(s+\varrho_i)=Qc_i\), so the adapted height is exactly
\(Q\xi\Upsilon_a(\xi)\) by \eqref{eq:positive-adaptive-unit}.
Here \(K_0^2\le c_i\le2K_0^2\), and the stated thin interval gives
\(\xi\mathcal V(\xi)/c_i\le\delta_{\rm an}\). The uniform regularity
and positive height derivative of \(\Psi_a\), together with the assumed
bounds for \(\xi\mathcal V(\xi)\), give all the unit and inverse bounds.
The normalization \(\Psi_{a,y}(\cdot,0)=\mathcal V(0)=1\) gives
\(\Upsilon_a(0)=1\), and logarithmic differentiation gives the last claim.

\emph{(iii)} Thinness at \(\varrho_i\) allows
\cref{lem:direct-adaptive-positive} to supply the exit and landing
before the admissible endpoint. Its exact action identity, evaluated
at the exit and combined with (ii), is
\eqref{eq:direct-no-root-adaptive-input}. Since
\(\varrho_e\le\varrho_l\le2\varrho_e\),
\[
 \sqrt t\le2\varrho_e,\qquad
 q_1/\sqrt t\le Cbsu_e,\qquad |q_0|/t\le Casu_e^2.
\]
With \(u_e\le1/K_0\), these are the claimed tail bounds.

\emph{(iv)} Denote the right side of
\eqref{eq:direct-no-root-adaptive-input} by \(\mathfrak b_a\), a function
of \(\varrho_e\); through this equation \(\xi\) and \(h\) are functions of
the exit radius. By (i), \(-q>0\) on the passage, so \(\mathcal J_a\) increases;
since \(s\mathcal J_a(s\rho_{\rm g})=\mathcal J_+(\rho_{\rm g})\), we get
\(\e^{-\mathcal J_a(\varrho_e)}\le r^{c_{\rm g}/s}\). Since
\(\mathcal E\mathcal J_a=1/(\varrho_e^+K_a)\), \(\varrho_e^+\ge s\), and the Euler
derivatives of \(K_a\) are uniformly bounded, each fixed radial derivative
of \(\mathfrak b_a\) costs at most another factor \(s^{-1}\):
\(|\mathcal E^j\mathfrak b_a|\le C_js^{-j}\mathfrak b_a\), as in
\eqref{eq:direct-annulus-input-jets}. In the fixed radius box
\(\mathfrak b_a\le Cs^{-3}r^{c_{\rm g}/s-2}\); since \(r\le1/2\), this is
\(O(s^N)\) for every fixed \(N\), and in particular \(\mathfrak b_a\le1\)
for small \(s\). The chain-rule step before
\eqref{eq:direct-annulus-prefix-weight}, applied to
\(\xi=(\xi\Upsilon_a)^{-1}(\mathfrak b_a)\) with the inverse of (ii),
bounds \(|\mathcal E^j\xi|+|\mathcal E^j\log\Upsilon_a(\xi)|\) by
\(C_js^{-j}\mathfrak b_a\), which gives
\eqref{eq:positive-adaptive-input-weight}.

Logarithmic differentiation of \eqref{eq:direct-no-root-adaptive-input},
with \(\log\xi=\mathcal A_i/s+\log(h/Q)\),
\(\mathcal E\log(\varrho_e^2\varrho_e^+)=2+v_e\) and
\(\varrho_e^+(2+v_e)=2s+3\varrho_e\), gives
\[
 \mathfrak g
 =K_a^{-1}-(2s+3\varrho_e)
       -\varrho_e^+\mathcal E\log\Psi_a
       +\varrho_e^+\mathcal E\log\Upsilon_a(\xi),
\]
with \(K_a\) and \(\Psi_a(\cdot,\delta_{\rm an})\) evaluated at the exit.
By (i), the uniform Euler derivatives of \(\Psi_a\), and
\eqref{eq:positive-adaptive-input-weight}, this gives
\eqref{eq:direct-no-root-adaptive-source}.

At the cone landing, \eqref{eq:direct-adaptive-landing} gives
\(\mathcal E\log t=2+O(\epsilon_e)\) and
\(\mathcal E^2\log t=O(\epsilon_e)\). By (iii), the moving-cone
part of \cref{lem:direct-weighted-tail} applies at the landing, as before
\eqref{eq:direct-adaptive-source-slope}. With \(q_2=-s\), the endpoint
energy has derivative \((z_*/t)^{2s}a(t)\), and
\eqref{eq:direct-moving-tail} gives
\[
 |t\partial_t\log a(t)|
   \le C(\sqrt t+q_1/\sqrt t+|q_0|/t)
   \le C(\varrho_e+su_e+su_e^2).
\]
Therefore \(M_a=-\mathcal EI_+\) satisfies
\[
 M_a=t(z_*/t)^{2s}a(t)\,\mathcal E\log t,\qquad
 \mathcal E\log M_a=2-4s+
              O(\epsilon_e+\varrho_e+su_e+su_e^2).
\]
Since \(m_+=hI_{+,h}=\varrho_e^+M_a/\mathfrak g\),
\(\mathcal E\log\varrho_e^+=v_e\) and \(\partial_T=\mathcal E/\mathcal ET\),
\eqref{eq:direct-no-root-adaptive-source} gives the identity in
\eqref{eq:direct-no-root-adaptive-rate}. The factor \(K_a\) and its
derivatives have the common bounds of (i), while \(\Upsilon_a\) contributes
only terms \(O(s^N)\); this gives the stated dependence of the constants.
At the exits \(u_e\le r/\rho_{\rm g}\). Hence, for small \(m_{\rm land}\)
and radius box, after reducing \(r\) and then \(s\), the bracket in
\eqref{eq:direct-no-root-adaptive-rate} exceeds one and \(\mathfrak g<2\).
\end{proof}

\subsection{Negative quadratic drift away from the secondary fold}
\label{sec:direct-buffered-negative}

We first treat the part of
\eqref{eq:direct-negative-secondary-scales} on which
\begin{equation}
\label{eq:direct-buffered-region}
 b^2-4a\le-\nu<0 .
\end{equation}
The fixed margin implies \(a\ge\nu/4\). Both secondary passages therefore
have regular prefixes. The remaining comparison must match their exits
at the same central height, even when the individual actions become
large.

\begin{proposition}[Buffered negative quadratic drift]
\label{prop:direct-buffered-negative}
For every fixed \(\nu>0\), the return displacement \(\Delta_p\) on
\eqref{eq:direct-buffered-region} has at most
three distinct zeros, for all sufficiently small independent
\(s,r>0\) and all central heights in \(\cH_p\).
Moreover, the logarithmic slope ratio \(\mathcal L\) has
at most two distinct zeros on \(\cH_p\).
\end{proposition}

\subsubsection{Actions and the regular prefix}

On a small fixed interval \(0\le t\le v_1\),
\[
 D_\sigma(t)=ar^2-\sigma brt+t^2
                    -\sigma t^3R_q(\sigma st)
                 \asymp r^2+t^2.
\]
Indeed the two quadratic forms have a uniform positive
margin by \eqref{eq:direct-buffered-region}, and the cubic
term is absorbed by reducing \(v_1\).
Set
\begin{equation}
\label{eq:direct-buffered-actions}
 \mathcal A_\sigma(\rho)=
                    \int_0^\rho t/D_\sigma(t)\,dt .
\end{equation}
The positive action is defined up to any fixed positive primary radius;
we use the negative action first on compact subintervals below its
far root near \(\rho=1\).
At \(\rho_i=K_0r\), the central scale
\(\ell=sr,\ Q=s^3r^2\) has drifts
\[
 a-\sigma bx+x^2-\sigma rx^3R_q(\sigma srx),
                 \qquad 0\le x\le K_0 .
\]
They have a fixed positive margin for sufficiently small \(r\).
The central-height construction and composition with the
annular coordinate therefore give
\eqref{eq:direct-annulus-prefix}, with
\(\mathcal A_{\sigma,i}=\mathcal A_\sigma(K_0r)\).

At a common small primary radius,
\begin{equation}
\label{eq:direct-buffered-action-difference}
 |\mathcal A_+(\rho)-\mathcal A_-(\rho)|
 \le C\int_0^\rho
          \frac{brt^2+t^4}{(r^2+t^2)^2}\,dt
 \le C(b+\rho).
\end{equation}
The substitution \(t=rx\) gives an integrable density for the first
term, and the second integrand is bounded. This bounds the difference
even when the individual actions grow like \(\log(1/r)\).

\subsubsection{The intermediate annulus}

We prove that there are fixed \(M_0>2K_0\) and \(v_0>0\) such that the
interval of central heights in \(\mathcal H_p\) with positive thin exit in
\begin{equation}
\label{eq:direct-buffered-intermediate}
 M_0r\le\rho\le v_0
\end{equation}
contains at most three zeros of \(\Delta_p\).
We first construct the negative exit for each of these central heights.
It has the form \(\omega=\Omega\rho\), where
\begin{equation}
\label{eq:direct-buffered-comparable}
 0<\Omega_0\le\Omega\le\Omega_1,\qquad
 |\log\Omega|\le C(b+\rho+s).
\end{equation}
By \eqref{eq:direct-negative-secondary-scales}, \(0\le b\le1\).
Fix the lower action-slope
margin \(c_{\rm a}>0\) on the whole regular range
\(K_0r\le t\le v_1\).  Choose fixed factors
\(0<\Omega_0<1<\Omega_1\) so far apart that
\begin{equation}
\label{eq:direct-buffered-trial-separation}
 c_{\rm a}\min\{|\log\Omega_0|,\log\Omega_1\}>8C_{\rm act},
\end{equation}
where \(C_{\rm act}\) is a fixed bound for the \(b\)-dependent part of the
full \(\Phi(1;\rho)\) comparison, including the action and bounded
height-factor terms, when \(0\le b\le1\) and \(0<\rho\le1\).
The terms proportional to \(s\log\Omega\) are handled after the fixed
trial factors have been chosen.
Then choose \(M_0>2K_0/\Omega_0\) and \(v_0>0\) so small
that \(\Omega_1v_0<v_1\).  On
\([\Omega_0\rho,\Omega_1\rho]\), with \(M_0r\le\rho\le v_0\),
the positive margin in \eqref{eq:direct-buffered-region} gives
\begin{equation}
\label{eq:direct-buffered-negative-action-slope}
 0<c_{\rm a}\le
 \frac{t^2}{D_-(t)}\le C_{\rm a},\qquad
 \left|\frac{t^2}{D_-(t)}-\frac{t^2}{D_+(t)}\right|\le C(b+t),
\end{equation}
after reducing \(v_0\).  Let \(h_\sigma(\rho)\) denote the central height
obtained from \eqref{eq:direct-annulus-source} at the first thin exit
\(\rho\) on branch \(\sigma\).  The event equation and
\eqref{eq:direct-annulus-source-rate} imply that \(h_\sigma\) is \(C^3\)
and strictly decreasing, and that
\begin{equation}
\label{eq:direct-buffered-height-comparison}
 \Phi(\Omega;\rho):=s\log\frac{h_-(\Omega\rho)}{h_+(\rho)},\qquad
 \partial_{\log\Omega}\Phi
 =-\frac{(\Omega\rho)^2}{D_-(\Omega\rho)}+O(s)\le-\frac12c_{\rm a}.
\end{equation}
At \(\Omega=1\), the action difference estimate
\eqref{eq:direct-buffered-action-difference}, together with the bounded
height factors in \eqref{eq:direct-annulus-source}, gives
\begin{equation}
\label{eq:direct-buffered-height-at-one}
 |\Phi(1;\rho)|\le C_{\rm act}+C(v_0+s),
\end{equation}
with \(C_{\rm act}\) as chosen above. First fix the trial separation in
\eqref{eq:direct-buffered-trial-separation}; then reduce \(v_0\) and
finally \(s\).  The bound in
\eqref{eq:direct-buffered-height-at-one} is then smaller than
\(\frac14c_{\rm a}\min\{|\log\Omega_0|,\log\Omega_1\}\) uniformly in
\(b\in[0,1]\).  Integrating
\eqref{eq:direct-buffered-height-comparison} then gives
\(\Phi(\Omega_0;\rho)>0>\Phi(\Omega_1;\rho)\), uniformly for
\(M_0r\le\rho\le v_0\). The intermediate value theorem therefore gives an exit
\(\omega=\Omega\rho\) with \(h_-(\omega)=h_+(\rho)\).  Strict decrease in
\eqref{eq:direct-buffered-height-comparison} gives uniqueness of the
matching radius, and each passage has a unique first exit by
\eqref{eq:direct-annulus-source-rate}.

At the matching radius, combine the equal-height equation with
\eqref{eq:direct-buffered-action-difference}. The lower bound
in \eqref{eq:direct-buffered-height-comparison} yields
\begin{equation}
\label{eq:direct-buffered-log-omega-bound}
 |\log\Omega|\le C(b+\rho+s),
\end{equation}
after absorbing the \(O(s)\) height factors, proving the second bound in
\eqref{eq:direct-buffered-comparable} for sufficiently small parameters.
Differentiation of the equation of equal central heights gives
\[
 \mathcal E\log\omega=G_+(\rho)/G_-(\omega),\qquad
 \mathcal E\log\Omega=O(u+v).
\]

For \(b\ge b_{*,-}>0\), where the threshold will be fixed by the
small-\(b\) estimate below, the preceding prefixes and comparable exit
radii meet \cref{cor:direct-annular-pair} after choosing the annular
margins small relative to \(b_{*,-}\), then reducing \(s\). It gives
\(\mathcal L''<0\) in \(T=-s\log h\), and at most three zeros of \(\Delta_p\).

For small \(b\), curvature need not have a uniform margin. We bound the
endpoint slope quotient and use the exterior contraction. The quotient is
\[
 \frac{I_{-,h}}{I_{+,h}}
       =\frac{M_-(\omega)}{M_+(\rho)}
                         \frac{G_+(\rho)}{G_-(\omega)}.
\]
The logarithm of this quotient has, by
\eqref{eq:direct-annulus-slope-calibration},
\eqref{eq:direct-buffered-comparable}, and
\eqref{eq:direct-annulus-source-rate}, the upper bound
\[
 C\{b+u_0+v_0+s(1+|\log s|)\}
       +2C_{\rm cone}\sqrt{z_*}.
\]
Choose \(b_{*,-}>0\) and the annular margins small to make this less than
\(-\log\Lambda_+\) for \(b\le b_{*,-}\), after reducing \(s\).
Then \(\Pi_p'I_{-,h}/I_{+,h}<1\), so \(\Delta_p\) has at most
one zero. Shrink the annular margins further for the curvature
argument with this same \(b_{*,-}\), then reduce \(s\).
This proves the three-zero bound on
\eqref{eq:direct-buffered-intermediate}.

\subsubsection{Exits on a compact primary interval}
\label{sec:direct-buffered-compact-primary}

Fix \(0<\rho_{\rm lo}<\rho_{\rm hi}<\infty\). Suppose the positive first exit
at \(y=z/(s|\eta|^2)=\delta_{\rm an}\) satisfies \(\rho_{\rm lo}\le\rho\le\rho_{\rm hi}\).
On this compact interval we first match the two exits, then prepare their
endpoints and the return displacement. The matching equation gives a
rational leading relation with a controlled remainder.
Choose a small fixed reference radius \(v_b>0\), and set
\[
 L_+(\rho)=\log\frac{\rho}{1+\rho},\qquad
 L_-(\rho)=\log\frac{\rho}{1-\rho},\qquad
 \mathfrak c_\sigma(p)=\mathcal A_\sigma(v_b)-L_\sigma(v_b).
\]
On compact positive intervals away from zero, and for
the negative sign away from one,
\begin{equation}
\label{eq:direct-buffered-primary-action}
 \mathcal A_\sigma(\rho)
       =\mathfrak c_\sigma+L_\sigma(\rho)+O_{C^j}(r+s).
\end{equation}
Subtract the integrands on the interval between \(v_b\)
and \(\rho\). Their denominators have fixed positive
margins, and \(R_q(\sigma st)=-1+O_{C^j}(s)\);
integration and differentiation give this estimate.
Equation \eqref{eq:direct-buffered-action-difference} implies
\[
 \Lambda=\exp(\mathfrak c_+-\mathfrak c_-),\qquad
 0<\underline\Lambda_{\rm act}\le\Lambda
     \le\overline\Lambda_{\rm act}<\infty.
\]

The annular height construction extends to any such fixed
compact primary interval with positive drift margin.
For these exits, the increase of action from \(K_0r\)
is at least \(c_1\log(\rho/(K_0r))\).
The equation for the central height therefore makes the expanded
prefix input and its first \(j\) radius derivatives
\(O(s^{-j}\e^{-c_2/s})\), uniformly in small \(r\).
Both signs have thin central prefixes: the positive action
gap can be made larger than their fixed central action
difference by reducing \(r\).
Thus the equation for the central height is
\begin{equation}
\label{eq:direct-buffered-primary-source}
 s\log(h/s^3)=-\mathcal A_\sigma(\rho_\sigma)
                 +sM_\sigma^0(\rho_\sigma),\qquad
 \|M_\sigma^0\|_{C^4}\le C .
\end{equation}
Here \(M_\sigma^0=2\log\rho_\sigma+\log\Psi_\sigma-\log\Upsilon_\sigma\).
Equating the two equations for the central height gives
\[
 L_-(\omega)-L_+(\rho)=\log\Lambda+O_{C^3}(r+s).
\]
For each value of the parameter-dependent multiplier \(\Lambda\),
the leading scalar relation has the solution
\[
 f_\Lambda(\rho)=\frac{\Lambda\rho}{1+(1+\Lambda)\rho}.
\]
For \(\rho_{\rm lo}\le\rho\le\rho_{\rm hi}\), \(f_\Lambda(\rho)\) lies
in a fixed compact subset of \((0,1)\). Comparing the central-height
equations at slightly wider negative exits gives opposite signs;
first-exit order gives the matching negative passage.
The lower bound on \(L_-'\) then gives, by implicit differentiation,
\begin{equation}
\label{eq:direct-buffered-fractional-relation}
 \omega=f_\Lambda(\rho)+O_{C^3}(r+s).
\end{equation}

The rational profile is the inverse, after a sign change and positive
rescaling, of the base-point relation of Huzak--Kristiansen
\cite[(2.13) and (6.8)]{HK2026}.
For this scalar relation, the substitutions
\[
 \rho=-\beta_{\rm HK}x_{\rm out}^{b_{\rm HK}},\qquad
 \omega=\beta_{\rm HK}x_{\rm in}^{b_{\rm HK}},\qquad
 \Lambda=\e^{-K_{\rm HK}},\qquad
 K_{\rm HK}=\frac{\lambda_1^{\rm HK}\pi}
            {\sqrt{-4\lambda_0^{\rm HK}-(\lambda_1^{\rm HK})^2}}
\]
give \(\omega=f_\Lambda(\rho)\); their domain becomes
\(\rho>0\), \(0<\omega<\Lambda/(1+\Lambda)\).
The two-scale \(C^3\) remainder in
\eqref{eq:direct-buffered-fractional-relation} follows by matching the
two central-height equations above.

For the endpoints apply \cref{lem:direct-small-height-continuation}
with \(\ell=\delta=s\), \(x=\rho_\sigma\) and
\(\zeta_0(x)=\delta_{\rm an}x^2\). The coefficient ratios are bounded,
and \(K_\sigma(x)=-s^{-3}q(\sigma sx)=D_\sigma(x)\) has a fixed
positive margin on the compact exit intervals above. The entrance
factor is positive with bounded \(C^3\) norm, including its moving
height. Choose a common \(Z_1=Z_0\) and put
\(b_\sigma^e=b_\sigma\), \(a_\sigma^e=\varpi_{s^2Z_0}A_\sigma\).
Since \(q_2=-s\), the multiplier is
\(\varpi_{s^2Z_0}=1+O(s(1+|\log s|))\). Thus \(a_\sigma^e\) has fixed
positive bounds, and
\begin{equation}
\label{eq:direct-buffered-primary-endpoints}
 I_\sigma(\rho_\sigma)=b_\sigma^e
       -\tfrac12s^2a_\sigma^e\rho_\sigma^2
       +s^2O_{C^3}(s(1+|\log s|)).
\end{equation}

The endpoint preparation now turns the matching relation into an
expansion of the return displacement. Substituting
\eqref{eq:direct-buffered-fractional-relation} and
\eqref{eq:direct-buffered-primary-endpoints} into the exterior map gives
\begin{equation}
\label{eq:direct-buffered-prepared-displacement}
 \frac{2\Delta}{s^2a_+^e}
 =D(p)+\rho^2-\lambda(p)f_\Lambda(\rho)^2
                 +O_{C^3}(r+s(1+|\log s|)),
 \qquad
 \lambda=\Pi_p'(b_-^e)a_-^e/a_+^e .
\end{equation}
Here \(D(p)\) depends only on \(p\), and \(\lambda\) has fixed positive
bounds. For every \(h\in\mathcal H_p\) whose positive first exit lies in
\([\rho_{\rm lo},\rho_{\rm hi}]\), the reference endpoints \(b_\sigma^e\)
differ from the corresponding admissible endpoints by \(O(s^2)\).
Thus, whenever this set of heights is nonempty, the reference endpoints
lie in the enlarged exterior windows. The exterior Taylor remainder adds
\(O_{C^3}(s^2)\) after normalization.
For the principal part \(\widehat\Delta_{\rm m}=D+\rho^2-\lambda f_\Lambda^2\),
\[
 \left(\frac{\widehat\Delta_{\rm m}'}{2\rho}\right)'
 =\frac{3\lambda\Lambda^2(1+\Lambda)}
             {[1+(1+\Lambda)\rho]^4}>c_*>0 .
\]
Apply \cref{lem:direct-derivative-ratio-count} with
\(\phi=\rho^2\), \(\psi=f_\Lambda^2\), and
\(\widehat\Delta=2\Delta/(s^2a_+^e)\). The compact positive radius
interval gives \(\phi'\ge2\rho_{\rm lo}\); the displayed margin
therefore persists under the remainder in
\eqref{eq:direct-buffered-prepared-displacement}. It gives at most two
zeros on this interval and
\[
 \left(\frac{\widehat\Delta_\rho}{2\rho}\right)_\rho>0.
\]
For the crossing direction, use the same endpoint identity as in
\eqref{eq:direct-quadratic-prefix-zero-orientation}, here in the exit
coordinate:
\[
 \frac{\widehat\Delta_\rho}{2\rho}
 =\frac{I_{+,\rho}}{s^2a_+^e\rho}(\e^{\mathcal L}-1),
 \qquad I_{+,\rho}<0.
\]
At a slope zero it gives \(\mathcal L_\rho<0\). Since
\eqref{eq:direct-buffered-primary-source} gives
\(T_\rho=\mathcal A_+'-s(M_+^0)'>0\), every such zero has
\(\mathcal L_T<0\).

\subsubsection{Large positive exits}

We now treat the central heights in \(\mathcal H_p\) whose positive
half-orbit is thin at \(\rho_i\), with the prefix
\eqref{eq:direct-annulus-prefix}, and whose adaptive exit, at
\(z=\delta_{\rm an}\varrho^2(s+\varrho)\) in the notation of
\cref{lem:positive-branch-adaptive-exit}, has radius
\(\varrho_e=s\rho_e\) with \(\rho_e\ge1\). The adaptive exit is the
unique upward crossing of that level, so \(\rho_e\ge1\) means
\(z(s;h)\le2\delta_{\rm an}s^3\). By \cref{lem:central-height-order},
this condition and thinness at \(\rho_i\) each define an interval;
their intersection is the interval considered here. For each such height,
the negative annular thin exit at \(z/(s|\eta|^2)=\delta_{\rm an}\) remains
in a fixed compact interval strictly below its far root, and for large
\(\rho_e\) the return contracts.

At \(\rho_i\), \eqref{eq:direct-annulus-prefix} reads
\[
 K_0^2\Psi_+(1/K_0,K_0r,y_i)=\xi_+\Upsilon_+(\xi_+),\qquad
 y_i=\frac{z(\varrho_i;h)}{K_0^2Q},\qquad \varrho_i=K_0sr .
\]
Inverting this height coordinate gives
\[
 V_{\rm an}(\xi)=\frac{K_0^2}{\xi}
   \Psi_+^{-1}\bigl(1/K_0,K_0r;\xi\Upsilon_+(\xi)/K_0^2\bigr),
 \qquad z(\varrho_i;h)/Q=\xi_+V_{\rm an}(\xi_+).
\]
Thinness keeps the inverse in its fixed domain. Its derivative at zero
is one, so \(V_{\rm an}(0)=1\), and the annular coordinate bounds
supply the unit, reciprocal and inverse bounds required by
\eqref{eq:positive-adaptive-prefix}. Thus
\cref{lem:positive-branch-adaptive-exit} applies with
\(\mathcal V=V_{\rm an}\) and \(\mathcal A_i=\mathcal A_{+,i}\).

For \(1\le\rho\le\varrho_*/s\), where \(\varrho_*\) is a small
fixed unscaled radius, subtraction of the positive
action integrands gives the value estimate
\[
 \mathcal A_+(\rho)=\mathfrak c_++L_+(\rho)
            +O\{r+s(1+\log(1+\rho))\}.
\]
Indeed \(D_+(t)\asymp t^2(1+t)\) and
\(R_q(st)+1=O(s+st)\); the difference is bounded by
\(Cr/t^2+Cs/(1+t)\) for \(t\ge v_b\).
The action increment from \(K_0r\) to \(1\) gives the gap in part
(iv) of the lemma, with \(\rho_{\rm g}=1\): since \(D_+>0\), the annular
bound \(\mathcal J_+(\rho)\ge c\log(\rho/(K_0r))\) at \(\rho=v_1\) gives
\(\mathcal J_+(1)\ge \mathcal J_+(v_1)\ge\frac c2\log(1/r)\) for small \(r\).
Taking logarithms in \eqref{eq:direct-no-root-adaptive-input}, with
\(\mathcal A_{+,i}+s\mathcal J_a(s\rho_e)=\mathcal A_+(\rho_e)\) and
\(\varrho_e^2(s+\varrho_e)/Q=\rho_e^2(1+\rho_e)/r^2\), gives the equation
for the central height at the adaptive exit:
\[
 s\log(h/s^3)=-\mathcal A_+(\rho_e)
   +s\log[\rho_e^2(1+\rho_e)\Psi_a(u_e,v_e,\varrho_e,\delta_{\rm an})]
   -s\log\Upsilon_a(\xi_+).
\]
By \eqref{eq:positive-adaptive-input-weight}, \(\xi_+\) and
\(\log\Upsilon_a(\xi_+)\) are at most \(Cs^{-3}r^{c/(2s)-2}\). After
reducing \(r\), the negative input
\(\xi_-=\e^{(\mathcal A_{-,i}-\mathcal A_{+,i})/s}\xi_+\) is exponentially
small as well, since
\eqref{eq:direct-buffered-action-difference} bounds the difference of the
central actions. The remaining height prefactor
\(\rho_e^2(1+\rho_e)\Psi_a\) is at least \(c\delta_{\rm an}\) and, since the
unscaled exit radius is bounded, at most a fixed power of \(s^{-1}\).
Thus
\[
 s\log(h/s^3)=-\mathfrak c_+-L_+(\rho_e)
                    +O(r+s(1+|\log s|)).
\]
Since \(f_\Lambda(\rho_e)\), for \(\rho_e\ge1\), lies
in a fixed compact subset of \((0,1)\), comparison with
\eqref{eq:direct-buffered-primary-source} at two wider
negative endpoints proves existence of the negative exit.
Keep these negative comparison endpoints fixed when choosing the
large-exit threshold below.

Differentiation of the equation for the central height on the
negative branch gives
\[
 \frac{d\omega}{d\log h}
   =-\frac{s}{\mathcal A_-'(\omega)
                         -s(M_-^0)'(\omega)},\qquad
 0<hI_{-,h}\le Cs^3
\]
by \eqref{eq:direct-buffered-primary-endpoints}.
On the positive branch, parts (ii) and (iv) of
\cref{lem:positive-branch-adaptive-exit} give
\(d\lambda_{\rm in}/d\log h=1+O(\xi_+)=1+O(s^N)\). By part (iii), the
landing satisfies \eqref{eq:direct-tail-smallness}. These are the
hypotheses of \eqref{eq:direct-adaptive-source-slope}. Since \(q_2=-s\),
the condition \(q_2\le C_2s\) of \cref{cor:large-exit-contraction} holds
with \(C_2=0\); take \(C_1\) the constant of the negative bound above, and
fix \(R_0\ge1\) satisfying the condition on \(M_1\) there. The corollary
gives \(\mathcal L<0\), and at most one zero of \(\Delta_p\), for
\(\rho_e\ge R_0\).
Writing \(\rho\) for the radius of the annular thin exit and \(\rho_e\)
for that of the adaptive exit, we have
\(\rho\le\rho_e\) by \eqref{eq:positive-annular-adaptive-heights}.
On the interval \(\rho\ge\rho_{\rm lo}>0\),
any zero of \(\mathcal L\) must have \(\rho_e<R_0\), because
the large-exit estimate gives \(\mathcal L<0\) otherwise.
Such a central height has \(\rho_{\rm lo}\le\rho<R_0\), so the compact-primary
comparison proves \(\mathcal L_T<0\) at that zero. Thus every
zero of \(\mathcal L\) on this entire later interval of central
heights is a strict downward crossing.

\subsubsection{Covering all central heights}

We now join the comparisons in their central-height order. On the
overlaps they concern the same slope ratio, so their crossing
directions give one zero count on the whole interval.
Use the same sufficiently small thin threshold \(\delta_{\rm an}\)
at \(z_e=s^3r^2\delta_{\rm an}\). The quadratic-prefix lemma,
with \(\delta\) there replaced by \(s\), \(\ell=sr\),
and comparison drifts \(a-\sigma bx+x^2\), gives at most two
zeros on
\[
 h\ge z_e,\qquad\hbox{or}\qquad
 h<z_e,\quad |\eta_+(z_e;h)|/(sr)\le M_0 .
\]
The compact negative inverse follows from the first assertion of
\cref{lem:prefix-root-cover},
since \(a\ge\nu/4\) and the positive margin holds up to \(M_0\).
Every other central height has its positive half-orbit below
height \(z_e\) at all radii up to \(M_0sr\). For \(K_0r\le\rho\le M_0r\),
it therefore satisfies
\[
 y(\rho)<\delta_{\rm an}(r/\rho)^2<\delta_{\rm an}.
\]
Its positive thin exit occurs beyond \(M_0r\).
Such an exit exists before the admissible endpoint: the
boundary coefficient satisfies \(q(\eta)<0\) on the positive branch throughout the
small box in the coordinates \((\eta,z)\), whereas the endpoint has
\(y\asymp s^{-1}>\delta_{\rm an}\). At a crossing of \(y=\delta_{\rm an}\),
the right side of \eqref{eq:direct-annulus-level-crossing} is positive,
because \(sK_+=O(s+|\eta|)\) past the inner cut.
Thus every crossing is upward, and the first exit is unique.

Take \(\rho_{\rm lo}=v_0/2\) in the compact-primary comparison and reduce \(r\)
so that \(M_0r<\rho_{\rm lo}\). In increasing
\(T\), the early central heights, the annular interval
\(M_0r\le\rho\le v_0\), and the later interval \(\rho\ge\rho_{\rm lo}\)
can be cut within their overlaps to give an ordered cover of
\(\mathcal H_p\). Their order
follows from the decrease of the central height with the exit radius.
The first join has an overlap: an annular exit at \(M_0r\)
has height \(s^3r^2M_0^2\delta_{\rm an}>z_e\), whereas at the height
threshold of the early interval the positive half-orbit has height
\(z_e\) at that radius. The second overlap is
\(v_0/2\le\rho\le v_0\). These assertions persist after
intersection with \(\mathcal H_p\), including empty end pieces.

On the early interval, \cref{lem:direct-quadratic-prefix} gives
\(\mathcal L_T>0\) at every zero of \(\mathcal L\).
On the middle interval, either \(\mathcal L''<0\), for
\(b\ge b_{*,-}\), or \(\mathcal L<0\), for \(b\le b_{*,-}\).
In the latter case merge the nonvanishing middle interval with the
initial interval; all zeros on their union cross upwards.
On the later interval every zero has \(\mathcal L_T<0\), by the
compact-primary and large-exit comparisons above.
\Cref{lem:direct-oriented-zero-count} therefore gives at most two
zeros of \(\mathcal L\) on the whole interval \(\mathcal H_p\), and
\cref{cor:ordered-cover-zero-count} gives at most three zeros of
\(\Delta_p\). This proves \cref{prop:direct-buffered-negative}.

\section{Fold, separated roots, and the weak root}
\label{sec:direct-negative-root-regimes}

The buffered estimate uses a fixed positive margin for the secondary
drift. Near a positive root or a minimum tending to zero, we instead
compare the drift-weighted positive passage with the negative branch
of \cref{sec:direct-negative-annulus-buffer}. At the fold this controls
bounded outputs, but on the no-root side the positive passage can
continue through the annulus after a long bottleneck. We treat that
continuation separately, then handle separated roots and the root
approaching the central line to complete the coefficient cover.

\subsection{The boundary-drift fold and bounded outputs}
\label{sec:direct-negative-fold}

Continue with \eqref{eq:direct-negative-secondary-scales}.
Take \(1/2\le b\le1\), set \(x_{\rm f}=b/2\), and choose a small
fold width \(\epsilon_{\rm f}>0\) before fixing the parameter neighborhood.
Consider
\begin{equation}
\label{eq:direct-fold-region}
 |a-b^2/4|\le\epsilon_{\rm f}^2/100,\qquad r\ll\epsilon_{\rm f}^2 .
\end{equation}
This region is specified by the quadratic comparison. The exact fold
is the parameter locus where the minimum of \(K_+\) defined below
vanishes.
In positive secondary coordinates \(x=\eta/(sr)\),
\(Z=z/(s^2r^2)\), the drift is
\[
 K_+(x)=a-bx+x^2-rx^3R_q(srx),\qquad F_+(x)=K_+(x)/x .
\]
It is strictly convex on the fixed fold box:
\(K_+''=2+O(r)>1\).
Its minimum is at \(x_{\rm f}+O(r)\). When roots are present, the first
lies in \((x_{\rm f}-\epsilon_{\rm f}/4,x_{\rm f}+\epsilon_{\rm f}/4)\) and may be double.
On \(0\le x\le x_L=x_{\rm f}-2\epsilon_{\rm f}\) there is a fixed positive
drift margin.

Fix a finite \(B_*>x_{\rm f}+\epsilon_{\rm f}\), and a small positive output
height \(Z_f\). Construct the boundary-flow coordinate
of \cref{lem:direct-boundary-flow} on a fixed positive
box from \(x_L\) to beyond \(B_*\), retaining the exact
drift \(F_+\). Put \(t_0=s^2r^2\). The bounded-output
interval of central heights is
\begin{equation}
\label{eq:direct-fold-bounded-output}
 x_{\rm f}-\epsilon_{\rm f}\le w=\mathcal W_+(x(t_0Z_f;h),Z_f)\le B_* .
\end{equation}
We consider the central heights of \(\cH_p\) in the positive-drift
component reached from \(x_L\). They form an interval, by
\cref{lem:central-height-order} and \(\mathcal W_x>0\).

\begin{proposition}[Bounded outputs at the fold]
\label{prop:direct-fold-bounded}
For each fixed \(B_*\), one can choose \(\epsilon_{\rm f}\) sufficiently
small so that \(\Delta_p\) has at most three distinct zeros on
\eqref{eq:direct-fold-bounded-output}, uniformly on both sides of
the fold. If \(K_+\) has a positive root, \(\Delta_p\) has at most
three distinct zeros on all of \(\cH_p\).
\end{proposition}

\subsubsection{Positive passage in the fold region}

Reaching the bounded-output interval forces a fixed advance through
the slow part of the boundary flow. We use this advance first to make
the entry height small, then to apply the drift-weighted phase estimate.
Since \(\mathcal W_+\le x\), for a central height in
\eqref{eq:direct-fold-bounded-output} the positive half-orbit
reaches \(x_L\) below height \(Z_f\) and then crosses the
projected interval \([x_{\rm f}-2\epsilon_{\rm f},x_{\rm f}-\epsilon_{\rm f}]\), where
\(0<F_+\le C\epsilon_{\rm f}^2\). Its incoming height at \(x=x_L\)
therefore satisfies
\begin{equation}
\label{eq:direct-fold-action-gap}
 Z_i\le Z_f\e^{-\mathcal A_{\rm gap}/s},\qquad
 \mathcal A_{\rm gap}\ge c_1/\epsilon_{\rm f},\qquad
 h/(s^3r^2)\le (Z_f/s)\e^{-\mathcal A_{\rm gap}/s}.
\end{equation}
The projected entry is at most \(x_L\), so its action is at least
the displayed increment. The central prefix is therefore thin and has the
preparation \(Z_i=s\xi \Upsilon(\xi)\), with
\(\xi=k_ph\), \(\Upsilon(0)=1\), and uniform derivative bounds.

Use \(T=-s\log h\) at the fixed output \(z=t_0Z_f\).
The regular entry margin and \eqref{eq:direct-fold-action-gap} verify
\cref{lem:direct-weighted-phase} with \(\delta=s\), \(N=1\), \(j=2\).
Thus \(w_T=\theta_+F_+(w)\), with
\(\theta_+-1,\theta_+',\theta_+''=o(1)F_+(w)\).
The constants may depend on fixed \(\epsilon_{\rm f},B_*\).
\Cref{rem:direct-projected-continuation} keeps the projected state
below its first root and the inverse graph outward, also at a double root.

The coefficient ratios at the output height \(t_0Z_f\) are
\(O(s)\), and \cref{lem:boundary-flow-endpoint} with \(\ell=sr\)
and \(\delta=s\) gives
\begin{equation}
\label{eq:direct-fold-endpoint}
 -I_{+,w}=E_f A_pw\,U_p(w),\qquad
 E_f=t_0(z_*/(t_0Z_f))^{2s}\longrightarrow0,\qquad
 U_p=1+O_{C^3}(s),\quad A_p\asymp1.
\end{equation}
Compute the rate \(\mathfrak H_+=\partial_T\log(hI_{+,h})\) and its
derivative from \eqref{eq:direct-rate-identities} with \(K=K_+\) and
\(\theta=\theta_+\).
On the fixed enlarged interval, the coefficient expansion and
\eqref{eq:direct-curvature-coefficient-cubic}, with \(K_+'(0)=-b\) and
\(R_K(x)=-rR_q(srx)\), give
\[
 \frac{K_+'(w)}w=2-\frac bw+3\kappa_3 rw+O(sr^2),\qquad
 wK_+''-K_+'=b+3\kappa_3 rw^2+O(sr^2)>b/2 .
\]
Constants in these remainders may depend on \(B_*\).
Since \(w_T=\theta_+F_+(w)>0\), the rate and its derivative
therefore satisfy
\begin{equation}
\label{eq:direct-fold-positive-rates}
 -C\epsilon_{\rm f}-o(1)\le\mathfrak H_+
       \le 2-\frac{b_{\min}}{2B_*},\qquad
 \mathfrak H_+'\ge-o(1),\qquad b_{\min}=1/2 .
\end{equation}
The constant \(C\) is independent of \(\epsilon_{\rm f}\); the errors tend to
zero after \(\epsilon_{\rm f},B_*\) are fixed, uniformly as \(K_+(w)\)
tends to zero.

\subsubsection{The negative branch}

The positive action gap also makes the negative branch thin at a
fixed secondary cut. This places it in the range of the negative-branch
estimates from the preceding section.
The negative secondary drift
\[
 a+bX+X^2+rX^3R_q(-srX)
\]
has a fixed positive margin on every fixed prefix.
Choose a large fixed \(M_n\), and then choose \(\epsilon_{\rm f}\) so
that \(\mathcal A_{\rm gap}\) in \eqref{eq:direct-fold-action-gap} exceeds
the negative action to \(X=M_n\) by a fixed positive amount.
The inverse central-height coordinate puts the negative
branch strictly below \(y=z/(s^3\rho^2)=\delta_{\rm an}\)
at \(\rho=M_nr\), with the regular prefix preparation
\eqref{eq:direct-annulus-prefix}. Hence
\cref{lem:negative-branch-far-root} applies with \(M=M_n\).

\subsubsection{Counting and the root-present cover}

Let \(\rho_{\rm f}=b_{\min}/(12B_*)\), reduced if necessary
so that \(\rho_{\rm f}<0.5\). Choose \(M_n\) large enough
to absorb the \(r/\rho\) errors, then \(\epsilon_{\rm f}\) small for
the action gap and the lower bound in
\eqref{eq:direct-fold-positive-rates}.
The positive projected output stays in the bounded interval just fixed.
We divide according to the negative output: its small-radius and far-root rates give
opposite signs for the first derivative of the slope ratio, while
the compact primary range controls its critical points.
For the slope ratio \(\mathcal L\) of \eqref{eq:direct-completed-log-ratio},
the following three intervals cover the bounded-output central heights:
\begin{enumerate}[label=(\roman*)]
\item If the negative thin exit has \(\rho\le\rho_{\rm f}\),
then \(\mathcal L'>0\), because the negative-branch rate
\(2-3\rho+o(1)\) exceeds the upper bound for the positive-branch rate in
\eqref{eq:direct-fold-positive-rates}.
\item For \(\rho_{\rm f}/2\le\rho\le0.9\),
\eqref{eq:direct-negative-primary-rates} and
\eqref{eq:direct-fold-positive-rates} give
\(\mathcal L''<0\).
\item On \eqref{eq:direct-negative-far-core},
\(\mathcal L'<0\), by
\eqref{eq:direct-negative-far-rate} and the lower bound for the
positive-branch rate. This includes every branch still thin at \(0.9\).
\end{enumerate}
The endpoint estimates and bounded derivatives of \(\Pi_p\) make the
exterior derivatives \(o(1)\) in these comparisons.
By \cref{lem:negative-branch-far-root}, the three ranges are ordered
intervals of central heights after intersection with the bounded-output
interval. The last two overlap, and the first two overlap on
\([\rho_{\rm f}/2,\rho_{\rm f}]\). Their inequalities for
\(\mathcal L(T)\) hold also at the cuts.
The first and third comparisons exclude critical points in their
respective ranges; the second controls their curvature. Thus, on the
whole bounded-output interval,
\begin{equation}
\label{eq:direct-fold-critical-maximum}
 \mathcal L'=0\quad\Longrightarrow\quad\mathcal L''<0.
\end{equation}
By \cref{cor:ordered-cover-zero-count}, \(\Delta_p\) has at most three
distinct zeros on the bounded-output interval.

To cover the early heights as well when \(K_+\) has a positive root,
apply \cref{lem:prefix-root-cover}
with \(\delta=s\), \(\ell=sr\), \(\epsilon_3=r\), comparison drifts
\(a-\sigma bx+x^2\), \(R_p=x_{\rm f}-\epsilon_{\rm f}/2\),
\(w_a=x_{\rm f}-3\epsilon_{\rm f}/4\) and the first root
\(x_*\in(x_{\rm f}-\epsilon_{\rm f}/4,x_{\rm f}+\epsilon_{\rm f}/4)\).
The positive margin up to \(R_p\) is fixed, and the inverse projection
continues in the fold box up to \(Z_f\), including when \(x_*\) is double.
Its root range lies in \eqref{eq:direct-fold-bounded-output}, where
\eqref{eq:direct-fold-critical-maximum} holds. The lemma therefore gives
the root-present bound of three.

\subsection{Passage on the no-root side of the fold}
\label{sec:direct-no-root-fold}

Suppose now that the exact minimum of \(K_+\) in
\eqref{eq:direct-fold-region} is positive. The normalized drift
\(K_+\) is then positive on the whole positive unscaled radius
interval: outside the fixed fold collar its quadratic
part has a positive margin, and its cubic term is positive.
As this minimum tends to zero, we need estimates uniform in the
possibly divergent bottleneck action. The bounded-output result
compares half-orbits whose projected states at a fixed secondary
height stay bounded. For larger projected outputs, we isolate the
bottleneck action in the exact exponential transport, prepare the
outgoing height uniformly, and compare exits on the primary and
cubic scales.

\begin{proposition}[Zero bound across the fold]
\label{prop:direct-negative-fold}
For fixed sufficiently small \(\epsilon_{\rm f}>0\), \(\Delta_p\) has at most
three distinct zeros throughout \eqref{eq:direct-fold-region}, including
the fold and both its sides, for all sufficiently small independent
\(s,r>0\).
\end{proposition}

\subsubsection{Bottleneck transport and escaped central heights}

Fix the postfold cut \(K_0>2\) from the annular construction,
and later a larger fixed \(M>K_0\). A central height is called
escaped if its positive half-orbit is still thin at \(x=M\):
\begin{equation}
\label{eq:direct-fold-escaped}
 z(Msr;h)/(s^3M^2r^2)\le\delta_{\rm an} .
\end{equation}
This is one interval of central heights.
Use a single boundary-flow coordinate on a fixed positive
box from \(x_L=x_{\rm f}-2\epsilon_{\rm f}\) to beyond \(K_0\).
The bottleneck formula \eqref{eq:direct-bottleneck-action}
gives
\[
 \Psi_{K_0}(y_{K_0})
    =\e^{\mathcal A_f/s}\Psi_L(y_L),\qquad
 \mathcal A_f=\int_{x_L}^{K_0}\frac{x\,dx}{K_+(x)} .
\]
Here \(Q=s^3r^2\) and \(y_j=z(jsr;h)/Q\) for
\(j=x_L,K_0\).
The endpoint coordinates have uniform bounds because their integrals
of \(1/F_+\) lie on regular collars. The full bottleneck action remains
in the explicit exponential factor; the composition below is
differentiated in the expanded height, where the inverse bounds stay
uniform.
The construction following \eqref{eq:direct-bottleneck-action} applies
with \(x_R=K_0\), \(\delta=s\), \(\beta=\beta_f\), and the regular
central prefix of \cref{lem:direct-source-height}, of action
\(\mathcal A_L=\int_0^{x_L}x/K_+(x)\,dx\), where
\[
 \mathcal A_{K_0}=\mathcal A_L+\mathcal A_f,\qquad
 \beta_f=\e^{-\mathcal A_f/s},\qquad
 \xi=\e^{\mathcal A_{K_0}/s}h/Q .
\]
By \eqref{eq:direct-bottleneck-inverse-composition},
\begin{equation}
\label{eq:direct-fold-inverse-composition}
 y_{K_0}=\xi V_p(\xi),\qquad V_p(0)=1,
\end{equation}
on a fixed small \(\xi\)-interval whose image contains a fixed small
output interval; the unit, its reciprocal, and the map inverse have
uniformly bounded \(\xi\)-derivatives.
Composing with the annular coordinate gives
\eqref{eq:direct-annulus-prefix} with action \(\mathcal A_{K_0}\).

We next check that every escaped central height lies in the fixed
domain of this inverse construction. For such a height, put
\(y_a(\rho)=z(s\rho;h)/(s^3\rho^2)\).
For fixed \(M>K_0\) and sufficiently small \(r\),
the segment \(K_0r\le\rho\le Mr\) lies in the
positive annular box, where \eqref{eq:direct-annulus-level-crossing}
excludes a downward crossing of \(y_a=\delta_{\rm an}\).
The escaped condition at \(Mr\) implies
\(y_a\le\delta_{\rm an}\) on the whole preceding segment.
The annular transport \eqref{eq:direct-annulus-action} yields
\[
 \begin{gathered}
 \frac{\widehat z(K_0r)}{Q}
   =\e^{-\mathcal J_{i,M}/s}\frac{\widehat z(Mr)}{Q}
   \le CM^2\delta_{\rm an}\e^{-\mathcal J_{i,M}/s},\\
 \mathcal J_{i,M}=\int_{K_0r}^{Mr}
       \frac{dt}{tK_+(r/t,t)}
       \ge c\log(M/K_0)>0.
 \end{gathered}
\]
At the inner cut, \(y_{K_0}=K_0^2y_a(K_0r)\) is
uniformly comparable to \(\widehat z(K_0r)/Q\).
Thus \(y_{K_0}\) lies in the small output domain of the bottleneck
and central-prefix inverse coordinates. Their composition gives a
small \(\xi\), and flow uniqueness identifies the recovered central
height with \(h\). Hence
\eqref{eq:direct-annulus-prefix} applies throughout the interval of
escaped central heights.

For every escaped central height, the positive half-orbit crosses
the fold collar below a bounded secondary height. The action gap gives
\begin{equation}
\label{eq:direct-escaped-source-gap}
 h/(s^3r^2)\le C_{M,\delta_{\rm an}}\e^{-\mathcal A_{\rm gap}/s},\qquad
 \mathcal A_{\rm gap}\ge c/\epsilon_{\rm f}.
\end{equation}
Fix the negative secondary cuts needed below, then choose
\(\epsilon_{\rm f}\) small so that \(\mathcal A_{\rm gap}\) exceeds their
regular actions by a fixed gap.
At each such cut the expanded negative input is exponentially small
by \eqref{eq:direct-escaped-source-gap}. The regular inverse
central-height coordinate therefore constructs a thin negative prefix;
flow uniqueness identifies it with the branch from the original \(h\).
Both signs now satisfy the hypotheses of
\cref{lem:direct-annulus-rates} whenever their exits
have small primary radius, and the negative branch satisfies those of
\cref{lem:negative-branch-far-root}.

\subsubsection{Both exits on the small primary annulus}

Take small fixed annular margins \(u_0,v_0>0\),
with \(u_0\) eventually much smaller than \(v_0\).
Choose \(M\) and the action gap so that both inner ratios
\(u_\sigma=r/\rho_\sigma\) are at most \(u_0\).
If both thin exits satisfy \(v_\sigma=\rho_\sigma\le v_0\), the prefixes
constructed above meet \cref{cor:direct-annular-pair}. Since \(b\ge1/2\),
small annular margins followed by small \(s\) give \(\mathcal L''<0\)
in \(T=-s\log h\).

\subsubsection{The negative primary and far-root outputs}

When the positive exit is at most \(v_0\), its annular
rate satisfies
\[
 \mathfrak H_+\ge
       2-Cu_0-Cv_0^2-Cs(1+|\log s|).
\]
For a negative exit in \(v_0/2\le\rho_-\le0.9\),
\eqref{eq:direct-negative-primary-rates} gives
\(\mathfrak H_-=2-3\rho_-+o(1)\).
First choose \(v_0\) small, then \(u_0\) much smaller
than \(v_0\), and finally reduce \(r,s\).
The rate difference \(\mathfrak H_--\mathfrak H_+\) is strictly negative.
The same sign
holds if the negative graph is in its far-root interval,
by \eqref{eq:direct-negative-far-rate}.
The exterior contribution to \(\mathcal L'\) tends to
zero with the negative endpoint derivative.
Hence \(\mathcal L'<0\) on both ranges.

It remains to locate the negative branch when the positive exit has
left the small annulus. For an escaped central height, a positive
annular thin exit with \(\rho_+\ge v_0\) places the negative graph in
the far-root interval, as the following comparison of central heights
shows.
Choose first the compact negative comparison cut \(0.9\) and the
annular cut \(v_0\), and let \(C_{\rm th}\) bound all fixed endpoint and
prefix terms in the two central-height equations.  Then reduce the fold
width \(\epsilon_{\rm f}\), after the preceding choices, so that
\(\mathcal A_{\rm gap}>2C_{\rm th}\).  Compare the central height for
which the positive exit is at \(v_0\) with the central height for which
the negative height at \(0.9\) equals the thin threshold.  Their actions satisfy
\begin{equation}
\label{eq:direct-fold-far-threshold}
 \begin{split}
 \mathcal A_+(v_0)
    &=\mathcal A_{K_0}+\log\frac{v_0}{K_0r}
                                   +O(K_0^{-1}+v_0),\\
 \mathcal A_-(0.9)&=\log(1/r)+O(1).
 \end{split}
\end{equation}
The first estimate follows by integrating the annular factor
\(K_+(r/\rho,\rho)^{-1}-1=O(r/\rho+\rho)\) after \(K_0r\).
For the second, the negative secondary prefix has no
bottleneck, and the remaining primary action has fixed
endpoint margins.
Since \(\mathcal A_{K_0}\ge\mathcal A_{\rm gap}\), the displayed estimates
and the order of choices give, uniformly for small \(r\),
\begin{equation}
\label{eq:direct-fold-far-action-gap}
 \mathcal A_+(v_0)-\mathcal A_-(0.9)
 \ge \mathcal A_{\rm gap}-C_{\rm th}\ge C_{\rm th}>0.
\end{equation}
Writing \(h_+^{v_0}\) and \(h_-^{0.9}\) for the two central heights
(with \(h_-^{0.9}\) defined by the thin threshold at the fixed radius
\(0.9\)),
\eqref{eq:direct-annulus-source} and the uniform boundedness of its
height factors give
\begin{equation}
\label{eq:direct-fold-far-height-order}
 s\log\frac{h_+^{v_0}}{h_-^{0.9}}
 =-\bigl[\mathcal A_+(v_0)-\mathcal A_-(0.9)\bigr]+O(s)<0
\end{equation}
after reducing \(s\).  Thus \(h_+^{v_0}<h_-^{0.9}\).  The exit-height
maps are strictly decreasing in their radii by the event equation, so
for any \(\rho_+\ge v_0\) its central height satisfies
\(h\le h_+^{v_0}<h_-^{0.9}\). The negative graph at that same \(h\)
is consequently still thin at \(0.9\): by
\cref{lem:central-height-order}, its height at this fixed radius is
smaller than the thin threshold attained by \(h_-^{0.9}\).
Part (ii) of \cref{lem:negative-branch-far-root} now gives
\(w_n>0.89\), hence \eqref{eq:direct-negative-far-core}.

\subsubsection{Positive continuation through the cubic regime}

For escaped central heights whose positive annular exit is at least
\(v_0\), which form an interval, we apply
\cref{lem:positive-branch-adaptive-exit}. Its prefix hypothesis
\eqref{eq:positive-adaptive-prefix} is
\eqref{eq:direct-fold-inverse-composition}, with \(\mathcal V=V_p\) and
\(\mathcal A_i=\mathcal A_{K_0}\). Escaped central heights are thin
at \(K_0r\), and their inverse estimate makes \(\xi\) uniformly small,
so it lies in the domain of \(V_p\) after reducing \(s\). By
\eqref{eq:positive-annular-adaptive-heights}, the adaptive exit follows
the annular thin exit; hence \(\varrho_e\ge sv_0\). By the first estimate in
\eqref{eq:direct-fold-far-threshold},
\(\mathcal J_+(v_0)=\mathcal A_+(v_0)-\mathcal A_{K_0}\ge\log(1/r)-C\), which is at
least \(\frac12\log(1/r)\) after fixing a sufficiently small upper bound
for \(r\). Thus part (iv) of the lemma applies with \(\rho_{\rm g}=v_0\)
and \(c_{\rm g}=1/2\). It gives \eqref{eq:direct-no-root-adaptive-rate},
\(\mathfrak H_+>c\varrho_e^+/s\), at these exits. The bounds on \(V_p\)
in \eqref{eq:direct-bottleneck-inverse-composition} make this estimate
uniform in the bottleneck action.
The negative graph is in the far-root interval by
\eqref{eq:direct-fold-far-threshold}. Its strictly
negative rate and the small exterior term imply
\(\mathcal L'<0\) here as well.

We now combine the exit comparisons on the interval of escaped heights.
On the interval where both exits are at most \(v_0\), we have
\(\mathcal L''<0\).
Within the escaped central heights, the condition
\(y_-(v_0/2;h)\le\delta_{\rm an}\) defines a second interval,
because the height at a fixed radius increases with \(h\).
The negative primary, negative far-root, and positive adaptive
estimates cover this interval, including their fixed overlaps.
Thus \(\mathcal L'<0\) throughout the second interval.
These two intervals cover
\eqref{eq:direct-fold-escaped}: a positive exit at least
\(v_0\) forces the negative far-root interval; otherwise
either both exits are small or the negative graph is
already in the second interval.
The estimates for \(\mathcal L(T)\) include both cuts
\(v_0/2\) and \(v_0\).
Consequently, on the entire escaped interval,
\begin{equation}
\label{eq:direct-escaped-critical-maximum}
 \mathcal L'=0\quad\Longrightarrow\quad\mathcal L''<0.
\end{equation}

\subsubsection{Covering the height interval}

Set \(B_*=4M\) in \cref{prop:direct-fold-bounded}.
Its bounded-output estimates give
\eqref{eq:direct-fold-critical-maximum}.
The quadratic-prefix lemma applies on the early interval
\[
 h\ge z_e=s^3r^2\delta_e,\qquad\hbox{or}\qquad
 h<z_e,\quad x_+(z_e;h)\le x_{\rm f}-\epsilon_{\rm f}/2 .
\]
Its regular-prefix hypotheses hold on the no-root side. In \(T\),
the early heights form an initial interval, and each zero of \(\mathcal L\) on it
satisfies \(\mathcal L'>0\).

For every other central height, the projected state of the positive
half-orbit at its hit of \(x=x_{\rm f}-\epsilon_{\rm f}/2\) is greater than
\(x_{\rm f}-3\epsilon_{\rm f}/4\).
Follow the half-orbit to the first event \(Z=Z_f\) or \(x=3M\).
If the height is reached first, the projected output
lies between \(x_{\rm f}-\epsilon_{\rm f}\) and \(B_*\), and is counted by
the bounded-output proposition.
If \(x=3M\) is reached first, the central height satisfies
\eqref{eq:direct-fold-escaped}. Indeed, between \(x=M\) and that event
the half-orbit lies in the coordinate box of
\cref{sec:direct-negative-fold}, which extends beyond \(B_*=4M\), and
\(B_0Z_f\le M^2\) since \(Z_f\) is small. For small \(s\),
\cref{lem:direct-radial-advance} with \(\delta=s\), \(X_1=M\), \(X_2=3M\)
and \(C_F=\sup|F_+|\) on that box bounds the height at \(x=M\) by
\(Z_f\e^{-M/(C_Fs)}<sM^2\delta_{\rm an}\).
Every admissible half-orbit under consideration reaches one of these
events before its fixed endpoint.

The complement of the early set is the single interval
\[
 J=\{T:h\in\cH_p,\ h<z_e,\quad
                      x_+(z_e;h)>x_{\rm f}-\epsilon_{\rm f}/2\}.
\]
Every point of \(J\) satisfies a bounded-output or escaped estimate.
Their common conclusion is the critical-point inequality
\(\mathcal L'=0\Rightarrow\mathcal L''<0\), by
\eqref{eq:direct-fold-critical-maximum} and
\eqref{eq:direct-escaped-critical-maximum}, for the same smooth function
on \(J\). The early cut
belongs to the initial interval, where
\eqref{eq:direct-quadratic-prefix-zero-orientation} includes the
terminal radius. Thus the initial interval and \(J\) cover the
whole height interval \(\cH_p\).
Apply \cref{cor:ordered-cover-zero-count}, with \(J\) as the
middle piece and an empty final piece. It gives the bound of three on
the no-root side. The same bound on the
root-present side follows from \cref{prop:direct-fold-bounded}.

\subsection{Separated roots on the negative quadratic-drift scale}
\label{sec:direct-negative-separated}

Even for separated roots, the limiting endpoint rates can agree.
We compare derivatives of these rates in the common logarithmic height
coordinate. Retaining their drift weights makes the curvature
comparison strict at every interior height.

Fix \(a_0^-,\underline b,\nu>0\), and consider
\begin{equation}
\label{eq:direct-negative-separated-region}
 a\ge a_0^-,\qquad b\ge \underline b,\qquad b^2-4a\ge\nu
\end{equation}
in \eqref{eq:direct-negative-secondary-scales}.
The first positive secondary root is simple:
\[
 \alpha=\frac{b-\sqrt{b^2-4a}}2+O(r),\qquad
 c<\alpha<C,\qquad K_+'(\alpha)<-c.
\]
Set \(x_L=\alpha/4,\ w_a=\alpha/2,\ R_p=3\alpha/4\).
The exact drift and its quadratic comparison have fixed
positive margins up to \(R_p\).

\begin{proposition}[Separated secondary roots]
\label{prop:direct-negative-separated}
In \eqref{eq:direct-negative-separated-region}, \(\Delta_p\)
has at most three zeros for all sufficiently
small independent \(s,r>0\), including parameters at
which the two limiting characteristic rates agree.
\end{proposition}

\begin{proof}
Construct the positive boundary-flow coordinate at a
fixed small output \(z=t_0Z_f,\ t_0=s^2r^2\).
The root-approach interval of central heights is
\begin{equation}
\label{eq:direct-negative-separated-root-cell}
 w=\mathcal W_+(x(t_0Z_f;h),Z_f)\in[w_a,\alpha).
\end{equation}
Advancing from a projected entry at most \(x_L\) to \(w_a\)
requires a fixed positive \(F_+\)-flow time \(\tau_*\).
Thus each central height in this interval satisfies
\[
 h/(s^3r^2)\le (Z_f/s)\e^{-\tau_*/s}.
\]
The regular prefix gives \eqref{eq:direct-common-phase} at the
fixed output \(z=t_0Z_f\). The positive drift margin at \(x_L\)
and the advance \(\tau_*\) verify
\cref{lem:direct-weighted-phase} with \(\delta=s\), \(N=1\) and
\(j=2\). Thus, in \(T=-s\log h\), the phase errors are
\(o(1)F_+(w)\) through two \(T\)-derivatives.
\Cref{lem:boundary-flow-endpoint}, with \(\ell=sr\) and
\(\delta=s\), gives \(-I_{+,w}=E_fA_pw(1+O_{C^3}(s))\), with
\(E_f=t_0(z_*/(t_0Z_f))^{2s}\to0\).
Apply \eqref{eq:direct-rate-curvature}:
\[
 \mathfrak H_+'=
 F_+(w)\{Q_+(w)+o(1)\},\qquad
 Q_+(w)=\frac{wK_+''-K_+'}{w^2}.
\]
By \eqref{eq:direct-curvature-coefficient-cubic}, with \(K_+'(0)=-b\) and
\(R_K(x)=-rR_q(srx)\), one has \(Q_+(w)=b/w^2+O(r)>c_Q>0\), since
\(b\ge \underline b\) and \(w<\alpha<C\).
Consequently \(\mathfrak H_+'\ge cF_+(w)>0\) on the root interval.

The negative output at the same central height may lie on either
scale. We follow it from bounded secondary outputs through the annulus
and the compact primary interval to its far root, retaining the
endpoint estimates needed to absorb the exterior curvature in each
range.
Choose a small fixed \(x_c>0\) so that the negative boundary action
to \(2x_c\) is less than \(\tau_*/4\). This is possible
because the central coefficient is at least \(a_0^-\).
The negative prefix coordinate gives exponentially small height
at this cut.
The negative secondary drift is
\[
 K_b(X)=a+bX+X^2+rX^3R_q(-srX),\qquad F_b=K_b/X .
\]
Choose a large fixed \(M\), and a small fixed output
height \(Z_b>0\). The boundary-flow coordinate on
a box containing \([x_c/4,5M]\) shows that every
negative output with \(X\le4M\) at \(t_0Z_b\)
has projected state \(w_b\in[x_c,4M]\).
At the regular entry \(x_c/2\), the drift \(F_b\) has a fixed
positive lower bound. The advance from that entry to
\(w_b\ge x_c\), at the fixed output \(z=t_0Z_b\), has a fixed
positive phase gap. Hence \cref{lem:direct-weighted-phase}
applies with \(\delta=s\), \(N=1\) and \(j=2\), in the same
\(T\), giving phase errors \(o(1)F_b(w_b)\) through two
\(T\)-derivatives.
By \cref{lem:boundary-flow-endpoint}, its endpoint has amplitude
\(E_b=t_0(z_*/(t_0Z_b))^{2s}\to0\). By
\eqref{eq:direct-curvature-coefficient-cubic}, with \(K_b'(0)=b\) and
\(R_K(X)=rR_q(-srX)\), one has \(Q_b(X)=-b/X^2+O(r)<-c_M<0\) for
\(x_c\le X\le4M\). The endpoint rate identities and the weighted phase
estimate then give
\[
 \mathfrak H_-'\le-cF_b(w_b),\qquad
 |I_{-,T}|+|I_{-,TT}|\le CE_bF_b(w_b).
\]

If the radius \(X\) of the negative graph at \(t_0Z_b\) exceeds
\(4M\), the graph is thin at \(X=M\). Indeed, the outward graph then
reaches \(X=4M\) below \(Z_b\), it lies in the coordinate box between
\(X=M\) and \(X=4M\), and \(8B_0Z_b\le15M^2\) since \(Z_b\) is small.
For small \(s\), \cref{lem:direct-radial-advance} with \(\delta=s\),
\(X_1=M\), \(X_2=4M\) and \(C_F=\sup|F_b|\) on the box bounds its
secondary height at \(X=M\) by
\(Z_b\e^{-3M/(2C_Fs)}<sM^2\delta_{\rm an}\).

For a thin exit in \(Mr\le\rho\le v_0\), the negative regular
prefix allows us to apply \cref{lem:direct-annulus-rates}, with
\(u=r/\rho\), \(v=\rho\). Choose \(M\) large and \(v_0,s\) small.
Since \(b\ge \underline b\), \eqref{eq:direct-annulus-clock-rate} gives
\[
 \mathfrak H_-'\le-c(u+v),\qquad
 |I_{-,T}|+|I_{-,TT}|\le C(s\rho)^{2-4s}.
\]
For \(v_0/2\le\rho\le0.9\),
\cref{lem:negative-branch-far-root}(i) instead gives
\(\mathfrak H_-'<-c_{v_0}<0\), with the same
endpoint size estimate.

Every graph still thin at \(0.9\) satisfies
\eqref{eq:direct-negative-far-core}, and
\cref{lem:negative-branch-far-root}(ii) gives, uniformly up to the
far root,
\[
 \mathfrak H_-'\le-cF_n(w_n),\qquad
 |I_{-,T}|+|I_{-,TT}|\le CE_nF_n(w_n).
\]
These ranges cover every negative-branch output. We now compare the
exterior term with the curvature margin in each range. Equation
\eqref{eq:direct-exterior-curvature} gives exterior contributions
\(o(1)F_b\) and \(o(1)F_n\) on the bounded secondary and far-root
intervals and \(o(1)\) on the compact primary interval. On the small
primary annulus, \eqref{eq:direct-annular-exterior-curvature} gives
\(o(\rho)\).
Each exterior contribution is absorbed by its negative curvature
margin. Thus
\[
 \mathcal L''=\mathfrak H_-'-\mathfrak H_+'
                      +\partial_T^2[(\log\Pi_p')(I_-)]<0
\]
throughout \eqref{eq:direct-negative-separated-root-cell}.
The positive drift weights give strict negativity at every interior
central height, although both may tend to zero at the endpoint.
Thus the curvature remains strictly negative at each interior height
even when the limiting rates agree. In the leading model the
characteristic rates agree at
\(\alpha=b/3,\ a=2b^2/9\), which is included here.

Finally apply \cref{lem:prefix-root-cover} with \(\delta=s\),
\(\ell=sr\), \(\epsilon_3=r\), comparison drifts \(a-\sigma bx+x^2\)
and the radii \(w_a,R_p,\alpha\) fixed above. The positive margin
and the gap \(R_p-w_a\) are fixed, and the inverse boundary-flow
coordinate is defined up to \(\alpha\) at \(Z_f\).
The root range is \eqref{eq:direct-negative-separated-root-cell},
where \(\mathcal L_{TT}<0\). Hence \(\Delta_p\) has at most three
zeros on \(\cH_p\).
\end{proof}

\subsubsection{The compact secondary coefficient cover}

For any fixed \(a_0^->0\), the normalized coefficients
with \(a\ge a_0^-\) now have a bound of three.
Choose the fixed fold width \(\epsilon_{\rm f}\) from
\cref{prop:direct-negative-fold}, and let \(\nu=\epsilon_{\rm f}^2/50\).
The three overlapping regions are
\[
 b^2-4a\le-\nu,\qquad |b^2-4a|\le2\nu,\qquad b^2-4a\ge\nu .
\]
They are covered respectively by
\cref{prop:direct-buffered-negative,prop:direct-negative-fold,prop:direct-negative-separated}.
For the last two regions normalization forces \(b=1\):
otherwise \(a=1\) would give \(b^2-4a\le-3\).
On the middle region
\mbox{\(|a-b^2/4|\le\epsilon_{\rm f}^2/100\)}, as required for the fold.
Each of the three propositions bounds the zeros on the whole height
interval. The minimum of their scale thresholds therefore gives a
uniform bound of three on this coefficient set.

\subsection{The secondary root approaching the central line}
\label{sec:direct-negative-weak-root}

The remaining part of the coefficient cover has \(a\downarrow0\).
Normalization in \eqref{eq:direct-negative-secondary-scales} then
forces \(b=1\). To treat the positive root approaching the central
line, apply \cref{prop:weak-root-passage} with
\[
 \kappa=s/a,\qquad \ell=sra,\qquad \hat r=r,\qquad \epsilon_2=a.
\]
These choices give \eqref{eq:weak-root-scales} and the logarithmic
height coordinate \(T=-s\log h\). The remaining hypothesis concerns
the negative branch outside the local weak-root region. We verify it
by separating bounded secondary outputs from graphs that reach the
annulus while still thin. The first group has a positive endpoint
rate; for the second, the annular and far-root estimates give the
uniform lower bound required by the weak-root comparison.

\begin{proposition}[Negative quadratic drift at the weak root]
\label{prop:direct-negative-weak-root}
There are fixed \(a_-,\kappa_->0\) such that \(\Delta_p\) has at
most two zeros on
\[
 0<a\le a_-,\qquad 0<s/a\le\kappa_-
\]
for all sufficiently small independent \(s,r>0\). Consequently, for all
sufficiently small independent \(s,r>0\) and every \(L>1/\kappa_-\),
\(\Delta_p\) has at most three zeros on the part of
\eqref{eq:direct-negative-secondary-scales} outside the strip
\(|q_0|\le Lq_1^2\).
\end{proposition}

\begin{proof}
At the radius \(\ell_d=srd\), the outward drift on the negative branch is
\[
 K_d(x)=a_d+x+dx^2+rd^2x^3R_q(-\ell_dx).
\]
For a fixed large \(M\) and \(r\le r(M)\),
\eqref{eq:weak-root-negative-drift} holds with \(d_{\max}=M\), with
constants depending on \(M\). We verify
\eqref{eq:weak-root-outer} for the negative graphs with
\(z<Y_b|\eta|^2\) at \(|\eta|=srM\).

Let \(X=|\eta|/(sr)\), and follow such a graph to the height
\(\ell_M^2Y_b\). At \(d=M\), the local radius in the weak-root estimate
is \(x=X/M\). If \(X\le4M\) at this height, the corresponding projected
state lies in \([0.9,4]\), and
\cref{lem:weak-root-branches}(ii) with \(d=M\) gives
\(\mathfrak H_-\ge c/M\) and \(|I_{-,T}|=o(1)\).

If \(X>4M\) at that height, the graph is thin at the annular cut \(X=M\):
\[
 z(-srM;h)/(s^3r^2M^2)<\delta_{\rm an},
\]
where \(\delta_{\rm an}\) is the thin threshold of
\cref{sec:direct-negative-annulus}. To obtain this estimate, use
\(|\eta|=\ell_Mx\), \(z=\ell_M^2Y\) with the same \(d=M\).
The outward graph reaches \(x=4\) below \(Y_b\) and lies in the box between
\(x=1\) and \(x=4\). Also, \(8B_0Y_b\le15\) since \(Y_b\) is
small. For small \(s\), apply \cref{lem:direct-radial-advance} in
\(x\), with \(\delta=s/M\), \(X_1=1\), \(X_2=4\) and
\(C_F=\sup|K_M/x|\) on the box. The height at \(x=1\) is at most
\(Y_b\e^{-3M/(2C_Fs)}<s\delta_{\rm an}\), as required.

For this thin input, take \(M\) large enough for the annular margins
on the negative branch. In the secondary coordinate \(X=|\eta|/(sr)\),
the negative drift is comparable to \(a+X\) and its derivative is
bounded below on this fixed box. Here
\(\delta=s\) and the weak-root condition gives \(a/s\ge1/\kappa_-\);
\cref{lem:direct-source-height} supplies the height unit at a fixed
secondary cut below \(M\), uniformly as \(a\to0\).
\Cref{lem:negative-branch-far-root}, whose constants are uniform down
to \(a=0\), gives, for exits of the negative branch up to \(\rho=0.9\),
\[
 \mathfrak H_-=2-3\rho+
                O(r/\rho+s(1+|\log s|))\ge-C_n,\qquad
 |I_{-,T}|\le C(s\rho)^{2-4s}=o(1).
\]
Graphs still thin at \(\rho=0.9\) satisfy \(\mathfrak H_-\ge-C_n\) and
\(|I_{-,T}|\le Cs^2\). Thus \eqref{eq:weak-root-outer} holds in both
cases. \Cref{prop:weak-root-passage} gives at most two zeros of
\(\Delta_p\), with \(a_-\) and \(\kappa_-\) equal to its constants
\(a_\circ\) and \(\kappa_\circ\).

To obtain the assertion outside the contact strip, combine the weak-root
bound with the compact coefficient cover of
\cref{sec:direct-negative-separated}, which gives the bound three on
\(a\ge a_-/2\), overlapping the weak-root region. On the remaining
small-\(a\) part, where \(s/a>\kappa_-\), normalization gives
\(b=1\) and \(|q_0|/q_1^2=a/s<1/\kappa_-\). This part lies in
\(|q_0|\le Lq_1^2\) for every \(L>1/\kappa_-\).
Hence the weak-root and compact-coefficient regions cover the parameter
set outside the contact strip. Their bounds on the whole interval
\(\mathcal H_p\) give at most three zeros there.
The thresholds for \(s,r\) are those of these two regions and are
independent of \(L\).
\end{proof}

\section{The finite parameter cover and the geometric conclusion}
\label{sec:direct-parameter-cover}

The regional estimates bound the displacement on the whole interval
of admissible central heights. It remains to show that their parameter
regions cover a full neighborhood of the organizing value, with common
sections and compatible constant choices. The resulting zero bound is
the maximum of the regional bounds. Geometric capture then gives the
upper bound for periodic orbits. The table below summarizes the regions;
the proof fixes their thresholds and checks their overlaps and boundary
cases.

\begin{table}[H]
\centering
\small
\caption{Parameter regions, the comparison that bounds the zeros of
\(\Delta_p\) on the whole height interval \(\cH_p\), and the resulting
bound. After the first row, \(q_1>0\). The constants
\(L,\kappa_{\rm p},m_{\rm cov},A_0,a_{\rm w},\kappa_{\rm w},r_{\rm d}\) are
chosen in the proof below, in \cref{sec:direct-parameter-cover}, which
also checks the overlaps and the boundary axes.}
\label{tab:direct-cover-guide}
\label{tab:direct-parameter-cover}
\begin{tabular}{@{}>{\raggedright\arraybackslash}p{\dimexpr.31\textwidth-10pt}
                >{\raggedright\arraybackslash}p{\dimexpr.24\textwidth-10pt}
                >{\raggedright\arraybackslash}p{\dimexpr.27\textwidth-10pt}
                >{\raggedright\arraybackslash}p{\dimexpr.1\textwidth-10pt}
                >{\raggedleft\arraybackslash}p{\dimexpr.08\textwidth}@{}}
\hline
Parameter region & Obstruction & Comparison & Section & Zero bound\\
\hline
\(q_1\le0\) & Vanishing radial flux at the central point
  & Signed paired contraction & \ref{sec:paired-variation} & 1\\
\(|q_0|\le Lq_1^2\), arbitrary \(q_2\) & Contact and low flux
  & Scattering ratio and endpoint bounds
  & \ref{sec:direct-bounded-contact} & 2\\
\(q_0>0\), \(q_1/\sqrt{q_0}\le\kappa_{\rm p}\) & Interior saddle
  & Dulac slopes and exterior contraction
  & \ref{sec:direct-positive-offset} & 2\\
\(q_0<0\), \(|q_0|/H^3\ge m_{\rm cov}\) & Merging boundary roots
  & Drift-weighted curvature; early central heights
  & \ref{sec:direct-primary-negative} & 3\\
\(q_0<0\), \(|q_2|/\sqrt{q_1}\le A_0\), \(|q_0|/q_1^{3/2}\le a_{\rm w}\),
\(q_1^2/|q_0|\le\kappa_{\rm w}\) & Weak boundary root
  & Weak-root comparison & \ref{sec:direct-negative-displaced} & 2\\
\(q_0<0\), \(q_2>0\), \(r<r_{\rm d}\), outside the contact strip
  & Near-symmetric roots & Curvature or power comparison
  & \ref{sec:direct-positive-quadratic} & 3\\
\(q_0<0\), \(q_2<0\), \(r<r_{\rm d}\), outside the contact strip
  & Unequal exits; long bottleneck
  & Radially weighted rates and exact boundary action
  & \ref{sec:direct-negative-annulus-buffer}--\ref{sec:direct-negative-root-regimes}
  & 3\\
\hline
\end{tabular}
\end{table}


\begin{proposition}[Uniform zero bound]
\label{prop:direct-uniform-displacement}
For each fixed \(C_0>1\), there is a full parameter
neighborhood of \((1,C_0,0,0,0)\) such that, with
\[
 N_{\rm dir}=3,
\]
\(\Delta_p\) has at most \(N_{\rm dir}\)
distinct zeros on the height interval \(\cH_p\).
\end{proposition}

\begin{proof}
\emph{Common sections and the contraction margin.}
First fix a preliminary box in the coordinates \((\eta,z)\) for this
\(C_0\), the common cone of \cref{sec:direct-passage-estimates}, and the
fixed adaptive and annular landing levels. The organizing endpoints satisfy
\(\eta_\pm^2/(2z_*)\to1\) as \(z_*\downarrow0\), by
\cref{lem:exterior-multiplier-new}. Hence their windows can be
chosen with a fixed lower bound \(z_*/\eta_\pm^2\ge c_{\rm end}>0\)
for all sufficiently small sections; choose each landing level
below \(c_{\rm end}/2\). The other landing restrictions use the enlarged
compact drift domains before the later coefficient and height thresholds.

Choose \(\Lambda_*\) with \(\Lambda_0<\Lambda_*<1\), and reserve a
fixed fraction of \(-\log\Lambda_*\) for tail errors.
By \cref{lem:exterior-multiplier-new} with \(\Lambda_+=\Lambda_*\),
we can then fix \(z_*\) small and choose the enlarged exterior
windows \(\widetilde W_\pm\) of that lemma. Thus
\eqref{eq:exterior-uniform-contraction-new} holds with
\(\Lambda_+=\Lambda_*\) in every regional estimate.

The same tail bound must fit the reserved contraction margin in every
region. The half-orbits enter the common cone after a finite scaled
continuation, as in \cref{lem:early-central-heights},
in \cref{sec:direct-bounded-contact,sec:direct-positive-offset}, and
by \eqref{eq:direct-finite-common-cone-entry}. From a common height
\(t\) in that cone, \eqref{eq:direct-tail-first-slope} bounds the
tail contribution to the logarithm of the endpoint derivative ratio by
\(2C_{\rm cone}(\sqrt{z_*}+|q_1|/\sqrt t)\), as in
\cref{lem:boundary-flow-endpoint}. The coefficient of \(\sqrt{z_*}\)
in every such upper bound is therefore at most \(2C_{\rm cone}\),
independently of \(L\), of coefficient margins such as \(\nu\), and
of the regional cuts.

The continuation constants and higher exterior derivatives may depend
on the cuts and on \(z_*\), but enter only the final parameter
thresholds. Choose \(z_*\) also so that
\(2C_{\rm cone}\sqrt{z_*}\) is smaller than the reserved fraction.
This choice is made after the common cone and before all later cuts
and thresholds, and meets every regional smallness requirement on
\(z_*\). If \(q_1\le0\), \cref{prop:signed-contraction-new} gives
at most one zero. Henceforth assume \(q_1>0\).

\emph{Secondary negative offsets.}
For \(q_0<0,\ q_2\ne0\), let \(s,r\) be the scales of
\cref{sec:standard-tools}. We first fix thresholds for these scales;
the contact width will be chosen after the remaining regions have
been compared.
\Cref{prop:direct-positive-quadratic} for \(q_2>0\) and the second
assertion of \cref{prop:direct-negative-weak-root} for \(q_2<0\) give
constants \(\kappa_+\) and \(\kappa_-\), respectively, and thresholds for
\(s,r\) independent of \(L\). Below these thresholds, for every
\(L>\max\{\kappa_+^{-1},\kappa_-^{-1}\}\), \(\Delta_p\) has at most three
zeros outside the strip \(|q_0|\le Lq_1^2\), for either sign of \(q_2\).
Consequently a single secondary threshold
\(0<r_{\rm d}<1/4\) for \(r\) can be fixed for both signs
before the final contact constant \(L\) is selected.

\emph{The primary and linear scales.}
For the remaining negative offsets, including \(q_2=0\), the primary
and linear scales must overlap the secondary region. We choose the
thresholds so that small normalized constant coefficient, outside the
secondary region, forces the normalized linear coefficient to stay
bounded away from zero. Put
\[
 b_{\rm cov}:=r_{\rm d}/2,\qquad A_0:=2/\sqrt{b_{\rm cov}}.
\]
For this fixed \(A_0\), the primary displaced-contact
estimate \cref{prop:direct-negative-displaced} gives
constants \(a_{\rm w},\kappa_{\rm w}>0\).
Choose \(m_{\rm cov}>0\) so that
\begin{equation}
\label{eq:direct-cover-constants}
 2m_{\rm cov}<1,\qquad \sqrt{2m_{\rm cov}}<b_{\rm cov},
             \qquad 2m_{\rm cov}/b_{\rm cov}^{3/2}<a_{\rm w}.
\end{equation}
With \(H\) and \(\rho_1\) of \cref{sec:standard-tools}, put
\(\bar a=|q_0|/H^3\) and \(\bar b=q_1/H^2\).
If \(\bar a\ge m_{\rm cov}\),
\cref{prop:direct-primary-negative} gives at most three zeros.
On the remaining region \(\bar a<m_{\rm cov}\) (and hence also
\(\bar a\le2m_{\rm cov}\)), use the secondary estimates when
\(q_2\ne0\) and \(r<r_{\rm d}\). For the complementary part,
it suffices to consider \(q_2=0\) or
\(r\ge r_{\rm d}/2=b_{\rm cov}\). On this part we claim
\begin{equation}
\label{eq:direct-cover-primary-overlap}
 \bar b\ge b_{\rm cov},\qquad
 |q_2|/\rho_1\le1/\sqrt{b_{\rm cov}}<A_0,\qquad
 |q_0|/\rho_1^3\le2m_{\rm cov}/b_{\rm cov}^{3/2}<a_{\rm w}.
\end{equation}
To prove the overlap, consider which term realizes the maximum
defining \(H\). If \(H=\rho_1\), the first assertion is immediate.
The possibility \(H=|q_0|^{1/3}\) would give
\(\bar a=1\), excluded by \eqref{eq:direct-cover-constants}.
In the remaining case \(H=|q_2|\),
\[
 r=\max\{\sqrt{\bar a},\bar b\}.
\]
Because \(\sqrt{\bar a}<b_{\rm cov}\) and \(r\ge b_{\rm cov}\),
one has \(\bar b\ge b_{\rm cov}\).
The other assertions in
\eqref{eq:direct-cover-primary-overlap} follow from
\(\rho_1=H\sqrt{\bar b}\).

If also \(q_1^2/|q_0|\le\kappa_{\rm w}\), these inequalities put
the parameter in \cref{prop:direct-negative-displaced}, giving at
most two zeros. Within this part of coefficient space, the
complementary case satisfies \(|q_0|/q_1^2<1/\kappa_{\rm w}\)
and lies in bounded contact. The bands
\(m_{\rm cov}\le\bar a\le2m_{\rm cov}\) and
\(r_{\rm d}/2\le r<r_{\rm d}\) provide fixed overlaps of these
coefficient regions.

\emph{Positive offsets and the contact strip.}
If \(q_0>0\), \cref{prop:direct-positive-offset} gives at most two
zeros whenever \(q_1/\sqrt{q_0}\le\kappa_{\rm p}\), for a fixed
\(\kappa_{\rm p}>0\). Its complement satisfies
\(q_0/q_1^2<\kappa_{\rm p}^{-2}\). We can now choose the contact
width to include all the remaining parameters:
\begin{equation}
\label{eq:direct-final-contact-constant}
 L>2\max\{\kappa_{\rm p}^{-2},\kappa_{\rm w}^{-1},
                         \kappa_+^{-1},\kappa_-^{-1},1\}.
\end{equation}
\Cref{prop:direct-bounded-contact} gives at most two zeros on this
strip for sufficiently small parameters.

\Cref{tab:direct-constant-order} records the global order of choices.
Within each regional proof, its fixed coefficient and height cuts precede
the final scale reduction, as specified there.
\begin{table}[ht]
\centering
\small
\caption{Global order of fixed choices. Each row may depend on preceding
rows; the original parameter neighborhood is reduced last.}
\label{tab:direct-constant-order}
\begin{tabular}{@{}>{\raggedright\arraybackslash}p{\dimexpr.44\textwidth-6pt}
                >{\raggedright\arraybackslash}p{\dimexpr.56\textwidth-6pt}@{}}
\hline
Choice & Purpose and permitted dependence\\
\hline
Fixed \(C_0\); preliminary box, common cone and landing level
\(m_{\rm land}\); \(\Lambda_*=\Lambda_+\), then \(z_*\), exterior windows
and capture data
  & Reserve the exterior margin using \(\Lambda_0\), and fix
    \(2C_{\rm cone}\sqrt{z_*}\) below a fixed fraction of it, before any
    coefficient or height threshold. The landing level is small for
    \cref{lem:direct-adaptive-positive} and
    \eqref{eq:direct-no-root-adaptive-rate}, which use only bounds of the
    drift factors.\\
For \(q_2>0\): \(a_+,\kappa_+\), then the compact thresholds. For
\(q_2<0\): the fold width \(\epsilon_{\rm f}\) and
\(\nu=\epsilon_{\rm f}^2/50\), \(a_-,\kappa_-\), then the compact
thresholds. Then \(r_{\rm d}\)
  & Both signs of \(q_2\); the complements of the weak-root regions lie
    in bounded contact strips.\\
\(b_{\rm cov}=r_{\rm d}/2\), \(A_0\), \(a_{\rm w}\), \(\kappa_{\rm w}\),
\(m_{\rm cov}\)
  & Fixed primary--secondary overlaps in
    \eqref{eq:direct-cover-constants}--\eqref{eq:direct-cover-primary-overlap}.\\
Positive-offset threshold \(\kappa_{\rm p}\), then contact width \(L\)
  & \eqref{eq:direct-final-contact-constant} includes the complements of
    all weak-root and saddle regions.\\
Contact height thresholds, compact passing sets and matching height;
final original parameter neighborhood
  & Cut-dependent finite-height errors and exterior derivative bounds
    are absorbed by parameter reduction with \(z_*\) fixed.\\
\hline
\end{tabular}

\end{table}

\emph{Boundary cases and a full neighborhood.}
The regions now have fixed overlaps. We check that they also include
their boundary cases. If \(q_1=0\), signed contraction applies, including
\(q_0=q_2=0\) and the axes \(H=0\) or \(r=0\). If \(q_1>0\)
and \(q_0=0\), the contact row applies.

If \(q_0<0\) and \(q_2=0\), then \(\bar a<m_{\rm cov}\) puts
\(H=\rho_1\), so \(\bar b=1\) and the primary--contact overlap
above applies. Its weak-root condition puts the point in the displaced
row; the complementary condition puts it in the contact row.

For \(q_0<0\), \(\bar a<m_{\rm cov}\) and \(q_2\ne0\), the
cases \(r<r_{\rm d}/2\), \(r_{\rm d}/2\le r<r_{\rm d}\), and
\(r\ge r_{\rm d}\) are covered by the secondary row, the doubled
overlap, and the primary--contact overlap. The primary row includes
\(\bar a=m_{\rm cov}\); the overlap and secondary rows include
\(r=r_{\rm d}/2\), and the overlap includes \(r=r_{\rm d}\).

Equality \(|q_0|=Lq_1^2\) belongs to contact. For \(q_0>0\), the
saddle row and its contact complement are exhaustive. Equality in the
saddle condition is assigned to the saddle row, and both rows allow
\(q_2=0\). Thus every parameter belongs to at least one row of
\cref{tab:direct-parameter-cover}. The maximum of the listed zero
bounds is \(N_{\rm dir}=3\).

Finally, choose the smallest of the finitely many remaining parameter
thresholds, keeping the sections fixed. The coefficient comparison in
Appendix~\ref{app:coordinates} gives a full neighborhood in the original
five parameters on which all these conditions hold. All estimates are
uniform for \((A,C)\) in the fixed neighborhood, independently of
\(H,s,r\) and \(\rho_1\). This proves the proposition.
\end{proof}

\begin{proof}[Proof of \cref{thm:direct-finite-cyclicity}]
Fix \(C_0>1\) and use the fixed sections, exterior windows,
and full parameter neighborhood constructed in the proof of
\cref{prop:direct-uniform-displacement}. Choose the
whole-graphic capture data of
\cref{prop:isolation-capture-new} for these windows, reduced as in
the paragraph before \cref{lem:exterior-multiplier-new} and in
\eqref{eq:small-central-height-new}.
Let \(V\) be the intersection of the parameter neighborhood
with \(V_{\mathrm{cap}}\), and choose
\(\varepsilon\le\rho_{\mathrm{cap}}\); by
\cref{prop:isolation-capture-new}(ii), every periodic orbit within
Hausdorff distance \(\varepsilon\) of \(\GammaI\) is then captured.

Let \(p\in V\), and let \(\gamma\) be a periodic orbit with
\(\distH(\overline\gamma,\GammaI)<\varepsilon\). By
\cref{lem:unique-central-crossing-new}, it has exactly one transverse
central crossing. The interval construction in
\cref{prop:central-displacement-interval-new} assigns it a zero of
\(\Delta_p\), injectively. The uniform zero bound therefore gives
at most \(N_{\rm dir}=3\) such orbits.

It remains to check isolation. The interior-window construction in
\cref{prop:central-displacement-interval-new} places the height of
every captured orbit in \(\operatorname{int}\mathcal H_p\).
The finite zero bound makes its displacement zero isolated in an
open central-height neighborhood. The correspondence in that
proposition preserves isolation, so every periodic orbit counted
above is a limit cycle.
Consequently \(N(V,\varepsilon)\le3\) in
\eqref{eq:cyclicity-definition}, and \(\Cycl(\GammaI)\le3\).
\end{proof}

\section{Sharpness and three hyperbolic limit cycles}
\label{sec:sharpness}

To attain the upper bound, we need three displacement zeros whose
orbits approach the whole graphic. Two arise on a fixed compact
interval of normalized exit radii: along a buffered coefficient path,
the difference of the boundary actions gives a profile with a strict
negative minimum, and varying \(A\) sets its additive constant to a
small positive value. This produces two zeros. The third comes from
the organizing contraction at a larger central height, beyond the
compact exit interval. We must then show that this third height also
tends to zero and that all three orbits converge to the whole graphic.
These are the final steps of the proof below.

\subsection{A buffered path and its action difference}

Choose the coefficient path
\begin{equation}
\label{eq:sharpness-coefficient-path}
 \nu=\min\{1,(C_0-1)/64\},\qquad a=(1+\nu)/4,\qquad
 q_0=-as^5,\quad q_1=s^3,\quad q_2=-s,\qquad s>0.
\end{equation}
It lies in the buffered region \eqref{eq:direct-buffered-region}
with \(b=1\), \(r=s\), and \(b^2-4a=-\nu<0\).
The coefficient change \eqref{eq:direct-coefficient-bridge}
determines \(\mu(s)=(\mu_1(s),\mu_2(s),\mu_3(s))\to0\).
We keep \(C=C_0\) and write
\[
 p_s(A)=(A,C_0,\mu(s)).
\]
This path fixes the boundary drift through \(\mu\), while \(A\)
remains free to adjust the connection splitting. In particular, the
boundary actions are independent of this adjustment.

Besides the limiting exterior multiplier \(\Lambda_0\) in
\eqref{eq:exterior-multiplier-new}, put
\[
 \Lambda_{\rm b}=\exp(\pi/\sqrt\nu).
\]
The choice of \(\nu\) gives the multiplier margin
\begin{equation}
\label{eq:sharpness-multiplier-margin}
 \Lambda_0\Lambda_{\rm b}^2\ge \e^{3\pi/2}>100.
\end{equation}
Indeed, the logarithm of the left-hand side is
\(12\pi/\sqrt{C_0-1}\) when \(C_0-1\le64\), and
\(2\pi-4\pi/\sqrt{C_0-1}\) otherwise.

The following lemma identifies \(\Lambda_{\rm b}\) as the limit of the
boundary-action multiplier. In the comparison with
\cite[(2.13) and (6.8)]{HK2026} in
\cref{sec:direct-buffered-compact-primary}, their
\((\lambda_0^{\rm HK},\lambda_1^{\rm HK})=(-a,-1)\) give
\(\Lambda_{\rm b}=\e^{-K_{\rm HK}}\).

\begin{samepage}
\begin{lemma}[Buffered action limit]
\label{lem:buffered-action-limit}
Along \eqref{eq:sharpness-coefficient-path}, the constants
\(\mathfrak c_\sigma\) in \eqref{eq:direct-buffered-primary-action} satisfy
\[
 \mathfrak c_+(s)-\mathfrak c_-(s)\longrightarrow\frac{\pi}{\sqrt\nu},
 \qquad
 \Lambda_s:=\e^{\mathfrak c_+(s)-\mathfrak c_-(s)}\longrightarrow\Lambda_{\rm b}.
\]
They are independent of \(A\).
\end{lemma}
\end{samepage}

\begin{proof}
Here \(D_\sigma\), \(\mathcal A_\sigma\), \(L_\sigma\), \(\mathfrak c_\sigma\) and the
reference radius \(v_b\) are those of \cref{sec:direct-buffered-negative}
with \(r=s\) and \(b=1\); thus
\(D_\sigma(t)=as^2-\sigma st+t^2-\sigma t^3R_q(\sigma st)\).
The fixed positive margin \(\nu\) gives
\(D_\sigma(t)\ge c_\nu(s^2+t^2)\) for \(0\le t\le v_b\).
The same bound holds after replacing \(R_q\) by \(-1\).
Since \(R_q(\sigma st)=-1+O(s)\), this replacement changes
each action by at most
\[
 Cs\int_0^{v_b}\frac{t^4}{(s^2+t^2)^2}\dd t=O(s).
\]

We separate the remaining integral at a fixed
\(0<\varepsilon<v_b\) to distinguish the action accumulated near
the central section from the outer reference terms. On
\([0,\varepsilon]\), deleting the
remaining cubic term changes each action by at most \(C\varepsilon\).
With \(t=sx\), the difference of the resulting inner actions tends to
\begin{align}
\label{eq:sharpness-action-integral}
 c_\infty
 &=\int_0^\infty x
   \left(\frac1{x^2-x+a}-\frac1{x^2+x+a}\right)\dd x \notag\\
 &=\int_0^\infty
       \frac{2x^2}{x^4+(2a-1)x^2+a^2}\dd x
   =\frac{\pi}{\sqrt{4a-1}}.
\end{align}
For the last equality, substitute \(x=\sqrt a\,u\), use
\(u\mapsto1/u\) to equate the integrals with numerators \(u^2\)
and \(1\), and apply \(y=u-1/u\) to their sum.
Both ends of the integral converge.

On \([\varepsilon,v_b]\), the action difference converges to
\[
 (L_+-L_-)(v_b)-(L_+-L_-)(\varepsilon).
\]
Subtract the reference terms in \(\mathfrak c_+-\mathfrak c_-\),
first let \(s\to0\), and then let \(\varepsilon\to0\).
The remaining outer term vanishes because
\((L_+-L_-)(\varepsilon)=\log((1-\varepsilon)/(1+\varepsilon))\to0\).
This proves the result.
\end{proof}

\subsection{Half-orbits with a prescribed exit}

The action comparison must be realized by half-orbits with a common
central height. We first construct each half-orbit backward from its
exit; monotonicity of the resulting height then gives the matching.
Fix \(J_+=[\rho_{\rm lo},\rho_{\rm hi}]\Subset(0,\infty)\) and
\(J_-\Subset(0,1)\) whose interior contains
\(f_{\Lambda_{\rm b}}(J_+)\). For a terminal radius
\(x_e\in J_\sigma\), prescribe
\(\eta=\sigma sx_e\), \(z=s^3\delta_{\rm an}x_e^2\) and integrate
toward the central section. In the local variables \(t=|\eta|/s\)
and \(y=z/s^3\), the exact radial equation is
\[
 y_t=\frac{ty}{s\{D_\sigma(t)+B(\sigma st,s^3y,p_s(A))y\}},
 \qquad y(x_e)=\delta_{\rm an}x_e^2.
\]
The quadratic margin and the upper end of \(J_-\) give
\(c(s^2+t^2)\le D_\sigma(t)\le C(s^2+t^2)\) on these fixed
intervals. The backward solution remains positive and bounded,
and reaches \(t=0\) with \(h_\sigma(x_e)=s^3y(0)>0\).
Its radius and height remain in the normal-rate box for small \(s\),
and \(q-zB<0\) identifies it with the corresponding half-orbit
of the original field.

The correction to its boundary action is exactly
\[
 \mathscr E_\sigma(x_e)
 =\int_0^{x_e}\frac{tBy}{sD_\sigma(D_\sigma+By)}\,dt
 =\int_{y(0)}^{y(x_e)}\frac{B}{D_\sigma}\,dy\ge0,
\]
with the coefficients evaluated along the half-orbit.
Choose a fixed \(\kappa>0\) with \(2\kappa<\inf(J_+\cup J_-)\).
The logarithmic growth rate on \([\kappa,2\kappa]\) is at least
\(c/s\), so \(y(\kappa)\le C\e^{-c/s}\).
The action correction up to \(\kappa\) is at most
\(Cs^{-2}y(\kappa)\); thereafter it is at most \(Cy(x_e)\),
since \(D_\sigma\ge c\kappa^2\). Thus
\(\mathscr E_\sigma\) is uniformly bounded and
\begin{equation}
\label{eq:sharpness-common-source}
 s\log(h_\sigma(x_e)/s^3)=-\mathcal A_\sigma(x_e)
                   +sM_\sigma^0(x_e),\qquad
 M_\sigma^0=\log(\delta_{\rm an}x_e^2)+\mathscr E_\sigma,
 \qquad \|M_\sigma^0\|_{C^0}\le C.
\end{equation}
All constants may depend on \(J_\pm,C_0\) and the fixed \(A\)-interval;
\(\delta_{\rm an}\) is held fixed before reducing \(s\).

At \(t=K_0s\) the normalized height
\(y/t^2\) is exponentially small. At any later crossing of
\(y/t^2=\delta_{\rm an}\),
\[
 \frac{d\log(y/t^2)}{d\log t}
 =\frac{t^2}{s(D_\sigma+B\delta_{\rm an}t^2)}-2>0
\]
for small \(s\), so the prescribed point is the first upward crossing
after \(t=K_0s\).
The central height decreases strictly with its exit radius:
the radial solution has positive derivative in its initial height,
whereas its terminal radial derivative exceeds the derivative of
\(\delta_{\rm an}x_e^2\) by the same inequality.

On \(J_\pm\), subtraction of the action integrands from their limits
gives \(\mathcal A_\sigma=\mathfrak c_\sigma+L_\sigma+O(s)\).
The strict endpoint margins of \(J_-\),
\cref{lem:buffered-action-limit}, and
\eqref{eq:sharpness-common-source} therefore bracket
\(h_+(\rho)\) between the two negative endpoint heights.
Continuity and strict decrease give a unique negative exit \(\omega\)
with the same central height. Equating their actions and using the
positive lower bound on \(L_-'\) gives
\begin{equation}
\label{eq:sharpness-negative-exit}
 \omega=f_{\Lambda_s}(\rho)+O_{C^0}(s),\qquad
 f_\Lambda(\rho)=\frac{\Lambda\rho}{1+(1+\Lambda)\rho}.
\end{equation}
Uniqueness at the regular central point makes these two half-orbits
parts of one orbit. The continuation below places their endpoints
in the fixed windows.

\subsection{A transverse organizing splitting}

Write \(p_0(A)=(A,C_0,0,0,0)\). The tuning argument requires a
strict sign for the derivative of the organizing splitting. We first
construct the two boundary graphs and their dependence on \(A\).
The variational equation along the parabola gives the sign of their
splitting in an affine coordinate; the section change has positive
derivative and preserves this sign in the endpoint coordinates.

\begin{lemma}[Transversality of the organizing splitting]
\label{lem:organizing-splitting-transversality}
For \(A\) close to \(1\), the graph equation
\eqref{eq:physical-graph-equation-new} at \(p_0(A)\) has unique
solutions
\[
 U_\sigma^0(z;A)=z+O(z^{3/2}),\qquad \sigma=\pm,
\]
which depend \(C^1\) on \(A\). Set
\[
 b_\sigma^0(A)=z_*-U_\sigma^0(z_*;A),\qquad
 \mathcal S_0(A)=\Pi_{p_0(A)}(b_-^0(A))-b_+^0(A).
\]
Then \(\mathcal S_0(1)=0\) and \(\mathcal S_0'(1)<0\).
\end{lemma}

\begin{proof}
\emph{Boundary graphs and their variations.}
At \(\mu=0\), both \(q\) and \(\eta=v+v^2+C_0vz\) are independent of \(A\).
Differentiating the exact projective field gives
\begin{equation}
\label{eq:sharpness-B-variation}
 B_A=-v(1+2v+C_0z),\qquad B(0,z,p_0(A))=1+C_0z.
\end{equation}
On \(U\asymp z\), uniformly for \(A\) near \(1\), the graph
right side \(\mathcal G_{\sigma,p_0(A)}\) therefore satisfies
\[
 \mathcal G_{\sigma,p_0(A)}=1+O(\sqrt z),\qquad
 |\mathcal G_{\sigma,U}|\le Cz^{-1/2},\qquad
 |\mathcal G_{\sigma,A}|\le C\sqrt z,
\]
where the middle bound is \eqref{eq:direct-tail-coefficients} at
\(\mu=0\).
The Volterra operator
\[
 \mathcal T_AU(z)=\int_0^z\mathcal G_{\sigma,p_0(A)}(U(t),t)\dd t
\]
maps \(|U(z)/z-1|\le1/2\) into itself. In the norm
\(\|U\|_*=\sup_{0<z\le z_*}|U(z)|/z\), its Lipschitz
constant is at most \(C\sqrt{z_*}<1\).
The operator is \(C^1\) in \(U,A\) on a slightly wider positive
cone, so differentiation of the convergent iteration gives
the asserted \(C^1\) fixed point. Its parameter derivative satisfies
\begin{equation}
\label{eq:sharpness-origin-variation}
 \begin{gathered}
 (\partial_AU_\sigma^0)'=\mathcal G_U\partial_AU_\sigma^0+B_A,\qquad
 \partial_AU_\sigma^0(0)=0,\\
 \partial_AU_\sigma^0(z)=O(z^{3/2}).
 \end{gathered}
\end{equation}

\emph{Comparison along the organizing parabola.}
For the entire \(p_0(A)\)-family, the function \(H_0\) of
\eqref{eq:parabola-invariance-new} satisfies the exact identity
\begin{equation}
\label{eq:sharpness-affine-variation}
 \dot H_0=2xH_0-(A-1)x^2.
\end{equation}
Continue the left boundary graph forward from the negative section,
and the right boundary graph backward from the positive section,
to the same point \(x\) on the organizing parabola. Denote their
\(A\)-derivatives at \(A=1\), measured in the \(H_0\)-coordinate
at fixed \(x\), by \(\mathfrak k_{\rm L}\) and
\(\mathfrak k_{\rm R}\), respectively. Their difference is the
signed variation of the splitting at that point. Changing from
fixed-\(x\) to fixed-\(z\) sections adds a tangential variation
along the organizing parabola, whose contribution to the variation
of \(H_0\) at \(A=1\) is zero.

Use \(x\) as time, as in the proof of
\cref{lem:exterior-multiplier-new}. The last term of
\eqref{eq:sharpness-affine-variation} gives the forcing in the
variational equation. With \(P_{\rm aff}(x)=x^2+2x+C_0>0\), both
variations satisfy
\[
 \mathfrak k'=\frac{4x}{P_{\rm aff}(x)}\mathfrak k-\frac{2x^2}{P_{\rm aff}(x)}.
\]
The boundary conditions follow from
\eqref{eq:sharpness-origin-variation}: at fixed \(z\),
\[
 \partial_Av
   =\frac{\sigma\partial_AU_\sigma^0}
       {\sqrt{2U_\sigma^0}(1+2v+C_0z)}=O(z),\qquad
 \partial_AH_0=-\frac{v\,\partial_Av}{z^2}
       =O(z^{-1/2})=O(|x|).
\]
With
\[
 W(x)=\exp\left(-\int_0^x\frac{4t}{P_{\rm aff}(t)}\dd t\right)>0,
\]
direct integration gives \(W(x)\asymp|x|^{-4}\) at both ends.
Hence \(W\mathfrak k_{\rm L}\to0\) as \(x\to-\infty\) and
\(W\mathfrak k_{\rm R}\to0\) as \(x\to+\infty\). Variation of constants
with these two endpoint conditions yields
\begin{equation}
\label{eq:sharpness-positive-integral}
 \mathfrak k_{\rm L}(x)-\mathfrak k_{\rm R}(x)
 =-W(x)^{-1}\int_{-\infty}^{\infty}
                  \frac{2W(t)t^2}{P_{\rm aff}(t)}\dd t<0.
\end{equation}
The integral is absolutely convergent and strictly positive.

\emph{The splitting in endpoint coordinates.}
To relate this variation to the displacement, let \(\chi_\pm\)
be the \(H_0\)-coordinates on the two fixed-\(z\) sections.
The section changes are regular, so \(\chi_\pm'(I_\pm^0)>0\), and
the chain rule at the positive section gives
\begin{equation}
\label{eq:sharpness-splitting-bridge}
 \mathcal S_0'(1)
 =\frac{1}{\chi_+'(I_+^0)}
   \bigl(\mathfrak k_{\rm L}(x_+)-\mathfrak k_{\rm R}(x_+)\bigr),
 \qquad \chi_+'(I_+^0)>0 .
\end{equation}
Here \(x_+\) is the organizing point on the positive section, where
both variations are evaluated. The finite exterior propagation is
included in \(\mathfrak k_{\rm L}(x_+)\), and
\(1/\chi_+'(I_+^0)\) is the positive \(H_0\)-to-\(I\) endpoint
conversion factor computed in \cref{lem:exterior-multiplier-new}.
Thus \eqref{eq:sharpness-positive-integral} and
\eqref{eq:sharpness-splitting-bridge} give \(\mathcal S_0'(1)<0\).
\end{proof}

\subsection{Tuning the splitting and obtaining three cycles}

\begin{proof}[Proof of \cref{prop:three-cycle-lower-bound}]
\emph{Endpoint preparation.}
We first continue the half-orbits from the compact exit intervals to
the fixed sections, where their endpoint maps can be composed with
the exterior transition. Use
\cref{lem:direct-small-height-continuation} with
\(\ell=\delta=s\) and \(\zeta_0(\rho_\sigma)=\delta_{\rm an}\rho_\sigma^2\),
as in \eqref{eq:direct-buffered-primary-endpoints}. On the path
\eqref{eq:sharpness-coefficient-path}, the coefficient ratios are
bounded, and the fixed compact exit intervals have positive drift
margins, uniformly for \(A\) near \(1\).
Choose the large height \(Z_1=Z_0\) of that lemma uniformly in \(A\),
and put \(t_s=s^2Z_0\). The continued and reference energies lie in
the common cone. Let \(U_\sigma^{\rm ref}\) be the tail with initial
energy \(U_\sigma^{\rm ref}(t_s)=B_0(p)t_s\), and define
\begin{equation}
\label{eq:sharpness-continuous-tail}
 \begin{aligned}
 b_\sigma^e(p)&=z_*-U_\sigma^{\rm ref}(z_*),\\
 a_\sigma^e(p)&=
 \left.\partial_uU_\sigma(z_*;t_s,u,p)\right|_{u=B_0(p)t_s}>0.
 \end{aligned}
\end{equation}
The reference endpoint and the tail-energy derivative are smooth in \(A\)
for fixed \(s>0\). They are the lemma's \(b_\sigma\) and
\(\varpi_{t_s}A_\sigma\), respectively.
Since \(\varpi_{t_s}=1+O(s(1+|\log s|))\), the coefficients
\(a_\sigma^e\) have uniform positive bounds, and
\eqref{eq:direct-small-height-endpoint} gives
\begin{equation}
\label{eq:sharpness-prepared-endpoints}
 I_\sigma(\rho_\sigma;p)
 =b_\sigma^e-\tfrac12s^2a_\sigma^e\rho_\sigma^2
                +s^2O_{C^3}(s(1+|\log s|)).
\end{equation}

\emph{Domains of the exterior transition.}
To place the reference and continued endpoints in the exterior windows
for small \(s\), uniformly on a fixed small \(A\)-interval, compare
their tails to \(U_\sigma^0(z;A)\). On the common cone, the difference
of the graph right sides at \(p_s(A)\) and \(p_0(A)\) is bounded by
\[
 C(s+s^5/z+s^3/\sqrt z),
\]
whereas, by \eqref{eq:direct-tail-coefficients}, a Lipschitz
coefficient for the difference is
\begin{equation}
\label{eq:sharpness-tail-comparison}
 2s/z+C(z^{-1/2}+s^3z^{-3/2}).
\end{equation}
Its integral over \([t_s,z_*]\) is uniformly bounded.
The initial energy difference is \(O(s^2)\), and the integrated
forcing is \(O(s)\). Gronwall's inequality gives, uniformly in
\(A\) and the fixed exit ranges,
\begin{equation}
\label{eq:sharpness-endpoint-limit}
 b_\sigma^e(A,s)\longrightarrow b_\sigma^0(A),\qquad
 I_\sigma(\rho_\sigma;A,s)\longrightarrow b_\sigma^0(A).
\end{equation}
First restrict \(A\) so that \(b_\sigma^0(A)\in\operatorname{int}W_\sigma\).
Uniform convergence then puts, for small \(s\), all the endpoints in
\eqref{eq:sharpness-endpoint-limit}, including those at the wider
comparison exits, in \(\operatorname{int}W_\sigma\). The continued
half-orbits thus give admissible central heights in
\(\mathcal H_p\). All these endpoints also lie in the enlarged
windows, where the exterior map and its derivatives are defined for
the displacement expansion and splitting comparison below.

\emph{The normalized displacement.}
We separate the constant term, which will be adjusted by \(A\), from
the coefficient of the compact exit profile. Define the continuous
splitting and slope coefficient
\begin{equation}
\label{eq:sharpness-actual-splitting}
 \mathcal S_s(A)=\Pi_{p_s(A)}(b_-^e)-b_+^e,\qquad
 \lambda_s(A)=\Pi_{p_s(A)}'(b_-^e)\frac{a_-^e}{a_+^e}.
\end{equation}
By \eqref{eq:sharpness-negative-exit} and
\eqref{eq:sharpness-prepared-endpoints}, the prepared displacement
\eqref{eq:direct-buffered-prepared-displacement} holds with \(r=s\),
\(\Lambda=\Lambda_s\), \(\lambda(p)=\lambda_s(A)\) and
\(D(p)=2\mathcal S_s(A)/(s^2a_+^e)\):
\begin{equation}
\label{eq:sharpness-prepared-displacement}
 \frac{2\Delta_{p_s(A)}(h(\rho;p_s(A)))}{s^2a_+^e}
 =\frac{2\mathcal S_s(A)}{s^2a_+^e}
      +g_{\lambda_s(A),\Lambda_s}(\rho)
      +O_{C^0}(s(1+|\log s|)),
\end{equation}
where
\[
 g_{\lambda,\Lambda}(\rho)
   =\rho^2-\lambda f_\Lambda(\rho)^2
   =\rho^2-\frac{\lambda\Lambda^2\rho^2}
                         {[1+(1+\Lambda)\rho]^2}.
\]

The three test radii must work throughout the tuning interval in
\(A\). We therefore bound the slope coefficient uniformly first.
The common factor
\(\varpi_{t_s}\) cancels from the quotient, so
\eqref{eq:direct-tail-first-slope} gives
\(\log(a_-^e/a_+^e)=O(\sqrt{z_*})+o_s(1)\). Also,
\(\Pi_0'(I_-^0)=\Lambda_0\{1+O(\sqrt{z_*})\}\) by
\cref{lem:exterior-multiplier-new}. As in
\cref{sec:direct-passage-estimates}, the coefficient of \(\sqrt{z_*}\)
depends only on the common cone; constants depending on the fixed
exit ranges enter only the \(o_s(1)\) term. Thus we may first choose
the fixed section, then restrict \(A\) and \(s\), so that
\begin{equation}
\label{eq:sharpness-lambda-range}
 \lambda_{\rm lo}:=\Lambda_0/2
 \le\lambda_s(A)\le
 \lambda_{\rm hi}:=(1+\Lambda_0)/2<1.
\end{equation}

\emph{Three test radii.}
The multiplier margin makes the model negative at an interior
minimum; a small positive additive constant will give the two outer
positive signs. To choose radii with fixed sign margins, put
\[
 M_{\rm b}=\sqrt{\lambda_{\rm lo}}\Lambda_{\rm b}>1,\qquad
 K_{\rm b}=1+\Lambda_{\rm b}.
\]
The strict negative minimum of
\(g_{\lambda_{\rm lo},\Lambda_{\rm b}}\) occurs at
\begin{equation}
\label{eq:sharpness-negative-minimum}
 \rho_{\rm b}=\frac{M_{\rm b}^{2/3}-1}{K_{\rm b}},\qquad
 g_{\lambda_{\rm lo},\Lambda_{\rm b}}(\rho_{\rm b})
       =-\frac{(M_{\rm b}^{2/3}-1)^3}{K_{\rm b}^2}<0.
\end{equation}
This follows directly from
\[
 g_{\lambda,\Lambda}'(\rho)
   =2\rho\left(1-\frac{\lambda\Lambda^2}
                            {[1+(1+\Lambda)\rho]^3}\right).
\]
Choose \(0<D_*<-g_{\lambda_{\rm lo},\Lambda_{\rm b}}(\rho_{\rm b})/4\).
Since \(g\) decreases in \(\lambda\),
\(D_*+g_{\lambda,\Lambda_{\rm b}}(\rho_{\rm b})<-3D_*\)
throughout \eqref{eq:sharpness-lambda-range}.
Take \(0<\rho_{\rm a}<\rho_{\rm b}\) so small that
\[
 \sup_{\lambda_{\rm lo}\le\lambda\le\lambda_{\rm hi}}
       |g_{\lambda,\Lambda_{\rm b}}(\rho_{\rm a})|<D_*/4,
\]
and then take \(\rho_{\rm c}>\rho_{\rm b}\) with
\(\rho_{\rm c}^2>\lambda_{\rm hi}+1\).
Since \(0<f_\Lambda<1\), the last choice gives
\(g_{\lambda,\Lambda}(\rho_{\rm c})>1\).
By \cref{lem:buffered-action-limit}, these choices imply,
uniformly in \(A\), for small \(s\),
\begin{equation}
\label{eq:sharpness-three-model-signs}
 \begin{aligned}
 D_*+g_{\lambda_s(A),\Lambda_s}(\rho_{\rm a})&>D_*/2,\\
 D_*+g_{\lambda_s(A),\Lambda_s}(\rho_{\rm b})&<-2D_*,\\
 D_*+g_{\lambda_s(A),\Lambda_s}(\rho_{\rm c})&>1+D_*.
 \end{aligned}
\end{equation}
All three radii and their negative comparison intervals are
fixed before the final reduction of \(s\). Apply the compact construction
above with \(J_+\) containing \([\rho_{\rm a},\rho_{\rm c}]\) and
\(f_{\Lambda_{\rm b}}(J_+)\subset\operatorname{int}J_-\).
\Cref{fig:sharpness-model}
shows the model profile and the three radii for \(C_0=65\).

\begin{figure}[t]
\centering
% Generated by scripts/figure_sharpness_model.py; do not edit.
% C_0 = 65, lambda = Lambda_0 = 0.207880, Lambda = Lambda_b = 23.140693,
% rho_a = 0.00251189, rho_b = 0.116736, rho_c = 1.3, D_* = 0.0076805.
\newcommand{\figModelFull}[1]{\draw[#1] (0.000,0.013) -- (0.014,0.011) -- (0.027,0.007) -- (0.041,0.001) -- (0.054,-0.006) -- (0.068,-0.014) -- (0.081,-0.021) -- (0.095,-0.029) -- (0.108,-0.036) -- (0.121,-0.043) -- (0.135,-0.050) -- (0.149,-0.057) -- (0.162,-0.064) -- (0.176,-0.070) -- (0.189,-0.075) -- (0.203,-0.081) -- (0.216,-0.086) -- (0.230,-0.091) -- (0.243,-0.095) -- (0.257,-0.100) -- (0.270,-0.104) -- (0.284,-0.108) -- (0.297,-0.111) -- (0.310,-0.115) -- (0.324,-0.118) -- (0.338,-0.121) -- (0.351,-0.124) -- (0.365,-0.126) -- (0.378,-0.129) -- (0.392,-0.131) -- (0.405,-0.133) -- (0.419,-0.135) -- (0.432,-0.137) -- (0.446,-0.138) -- (0.459,-0.140) -- (0.472,-0.141) -- (0.486,-0.142) -- (0.500,-0.144) -- (0.513,-0.145) -- (0.527,-0.145) -- (0.540,-0.146) -- (0.553,-0.147) -- (0.567,-0.147) -- (0.581,-0.148) -- (0.594,-0.148) -- (0.608,-0.148) -- (0.621,-0.148) -- (0.635,-0.148) -- (0.648,-0.148) -- (0.662,-0.148) -- (0.675,-0.148) -- (0.689,-0.148) -- (0.702,-0.147) -- (0.716,-0.147) -- (0.729,-0.146) -- (0.743,-0.145) -- (0.756,-0.145) -- (0.770,-0.144) -- (0.783,-0.143) -- (0.797,-0.142) -- (0.810,-0.141) -- (0.824,-0.140) -- (0.837,-0.139) -- (0.851,-0.138) -- (0.864,-0.136) -- (0.877,-0.135) -- (0.891,-0.133) -- (0.905,-0.132) -- (0.918,-0.130) -- (0.932,-0.129) -- (0.945,-0.127) -- (0.959,-0.125) -- (0.972,-0.124) -- (0.986,-0.122) -- (0.999,-0.120) -- (1.012,-0.118) -- (1.026,-0.116) -- (1.040,-0.114) -- (1.053,-0.112) -- (1.067,-0.109) -- (1.080,-0.107) -- (1.094,-0.105) -- (1.107,-0.102) -- (1.121,-0.100) -- (1.134,-0.098) -- (1.148,-0.095) -- (1.161,-0.092) -- (1.175,-0.090) -- (1.188,-0.087) -- (1.202,-0.084) -- (1.215,-0.082) -- (1.229,-0.079) -- (1.242,-0.076) -- (1.256,-0.073) -- (1.269,-0.070) -- (1.282,-0.067) -- (1.296,-0.064) -- (1.310,-0.061) -- (1.323,-0.057) -- (1.337,-0.054) -- (1.350,-0.051) -- (1.364,-0.048) -- (1.377,-0.044) -- (1.391,-0.041) -- (1.404,-0.037) -- (1.417,-0.034) -- (1.431,-0.030) -- (1.445,-0.027) -- (1.458,-0.023) -- (1.472,-0.019) -- (1.485,-0.015) -- (1.499,-0.012) -- (1.512,-0.008) -- (1.526,-0.004) -- (1.539,0.000) -- (1.552,0.004) -- (1.566,0.008) -- (1.580,0.012) -- (1.593,0.016) -- (1.607,0.020) -- (1.620,0.025) -- (1.634,0.029) -- (1.647,0.033) -- (1.661,0.038) -- (1.674,0.042) -- (1.688,0.046) -- (1.701,0.051) -- (1.715,0.055) -- (1.728,0.060) -- (1.742,0.065) -- (1.755,0.069) -- (1.769,0.074) -- (1.782,0.079) -- (1.796,0.083) -- (1.809,0.088) -- (1.823,0.093) -- (1.836,0.098) -- (1.850,0.103) -- (1.863,0.108) -- (1.877,0.113) -- (1.890,0.118) -- (1.904,0.123) -- (1.917,0.128) -- (1.931,0.134) -- (1.944,0.139) -- (1.958,0.144) -- (1.971,0.150) -- (1.985,0.155) -- (1.998,0.160) -- (2.011,0.166) -- (2.025,0.171) -- (2.039,0.177) -- (2.052,0.182) -- (2.066,0.188) -- (2.079,0.194) -- (2.092,0.199) -- (2.106,0.205) -- (2.120,0.211) -- (2.133,0.217) -- (2.147,0.223) -- (2.160,0.229) -- (2.174,0.235) -- (2.187,0.241) -- (2.200,0.247) -- (2.214,0.253) -- (2.228,0.259) -- (2.241,0.265) -- (2.255,0.271) -- (2.268,0.278) -- (2.281,0.284) -- (2.295,0.290) -- (2.309,0.297) -- (2.322,0.303) -- (2.336,0.309) -- (2.349,0.316) -- (2.363,0.323) -- (2.376,0.329) -- (2.389,0.336) -- (2.403,0.342) -- (2.417,0.349) -- (2.430,0.356) -- (2.444,0.363) -- (2.457,0.369) -- (2.470,0.376) -- (2.484,0.383) -- (2.498,0.390) -- (2.511,0.397) -- (2.525,0.404) -- (2.538,0.411) -- (2.551,0.418) -- (2.565,0.426) -- (2.579,0.433) -- (2.592,0.440) -- (2.606,0.447) -- (2.619,0.455) -- (2.632,0.462) -- (2.646,0.469) -- (2.660,0.477) -- (2.673,0.484) -- (2.687,0.492) -- (2.700,0.499) -- (2.714,0.507) -- (2.727,0.515) -- (2.740,0.522) -- (2.754,0.530) -- (2.768,0.538) -- (2.781,0.546) -- (2.795,0.553) -- (2.808,0.561) -- (2.822,0.569) -- (2.835,0.577) -- (2.849,0.585) -- (2.862,0.593) -- (2.876,0.601) -- (2.889,0.609) -- (2.902,0.618) -- (2.916,0.626) -- (2.930,0.634) -- (2.943,0.642) -- (2.957,0.651) -- (2.970,0.659) -- (2.984,0.667) -- (2.997,0.676) -- (3.010,0.684) -- (3.024,0.693) -- (3.038,0.701) -- (3.051,0.710) -- (3.065,0.719) -- (3.078,0.727) -- (3.092,0.736) -- (3.105,0.745) -- (3.119,0.754) -- (3.132,0.762) -- (3.146,0.771) -- (3.159,0.780) -- (3.172,0.789) -- (3.186,0.798) -- (3.200,0.807) -- (3.213,0.816) -- (3.227,0.825) -- (3.240,0.835) -- (3.254,0.844) -- (3.267,0.853) -- (3.280,0.862) -- (3.294,0.872) -- (3.308,0.881) -- (3.321,0.890) -- (3.335,0.900) -- (3.348,0.909) -- (3.362,0.919) -- (3.375,0.928) -- (3.389,0.938) -- (3.402,0.947) -- (3.416,0.957) -- (3.429,0.967) -- (3.442,0.977) -- (3.456,0.986) -- (3.470,0.996) -- (3.483,1.006) -- (3.497,1.016) -- (3.510,1.026) -- (3.524,1.036) -- (3.537,1.046) -- (3.550,1.056) -- (3.564,1.066) -- (3.578,1.076) -- (3.591,1.086) -- (3.605,1.097) -- (3.618,1.107) -- (3.632,1.117) -- (3.645,1.128) -- (3.659,1.138) -- (3.672,1.148) -- (3.686,1.159) -- (3.699,1.169) -- (3.712,1.180) -- (3.726,1.190) -- (3.740,1.201) -- (3.753,1.212) -- (3.767,1.222) -- (3.780,1.233) -- (3.794,1.244) -- (3.807,1.255) -- (3.821,1.266) -- (3.834,1.277) -- (3.848,1.287) -- (3.861,1.298) -- (3.875,1.309) -- (3.888,1.320) -- (3.902,1.332) -- (3.915,1.343) -- (3.929,1.354) -- (3.942,1.365) -- (3.956,1.376) -- (3.969,1.388) -- (3.982,1.399) -- (3.996,1.410) -- (4.010,1.422) -- (4.023,1.433) -- (4.037,1.445) -- (4.050,1.456) -- (4.064,1.468) -- (4.077,1.479) -- (4.091,1.491) -- (4.104,1.503) -- (4.117,1.514) -- (4.131,1.526) -- (4.145,1.538) -- (4.158,1.550) -- (4.171,1.562) -- (4.185,1.573) -- (4.199,1.585) -- (4.212,1.597) -- (4.226,1.609) -- (4.239,1.622) -- (4.253,1.634) -- (4.266,1.646) -- (4.280,1.658) -- (4.293,1.670) -- (4.307,1.682) -- (4.320,1.695) -- (4.333,1.707) -- (4.347,1.719) -- (4.361,1.732) -- (4.374,1.744) -- (4.388,1.757) -- (4.401,1.769) -- (4.415,1.782) -- (4.428,1.795) -- (4.442,1.807) -- (4.455,1.820) -- (4.469,1.833) -- (4.482,1.845) -- (4.495,1.858) -- (4.509,1.871) -- (4.523,1.884) -- (4.536,1.897) -- (4.550,1.910) -- (4.563,1.923) -- (4.577,1.936) -- (4.590,1.949) -- (4.604,1.962) -- (4.617,1.975) -- (4.631,1.989) -- (4.644,2.002) -- (4.658,2.015) -- (4.671,2.028) -- (4.685,2.042) -- (4.698,2.055) -- (4.712,2.069) -- (4.725,2.082) -- (4.739,2.096) -- (4.752,2.109) -- (4.766,2.123) -- (4.779,2.136) -- (4.793,2.150) -- (4.806,2.164) -- (4.820,2.177) -- (4.833,2.191) -- (4.847,2.205) -- (4.860,2.219) -- (4.873,2.233) -- (4.887,2.247) -- (4.901,2.261) -- (4.914,2.275) -- (4.928,2.289) -- (4.941,2.303) -- (4.955,2.317) -- (4.968,2.331) -- (4.982,2.346) -- (4.995,2.360) -- (5.009,2.374) -- (5.022,2.388) -- (5.035,2.403) -- (5.049,2.417) -- (5.063,2.432) -- (5.076,2.446) -- (5.090,2.461) -- (5.103,2.475) -- (5.117,2.490) -- (5.130,2.505) -- (5.144,2.519) -- (5.157,2.534) -- (5.171,2.549) -- (5.184,2.563) -- (5.197,2.578) -- (5.211,2.593) -- (5.225,2.608) -- (5.238,2.623) -- (5.252,2.638) -- (5.265,2.653) -- (5.279,2.668) -- (5.292,2.683) -- (5.306,2.698) -- (5.319,2.714) -- (5.332,2.729) -- (5.346,2.744) -- (5.360,2.759) -- (5.373,2.775) -- (5.387,2.790) -- (5.400,2.806);}
\newcommand{\figModelZoom}[1]{\draw[#1] (7.200,0.230) -- (7.214,0.230) -- (7.228,0.230) -- (7.242,0.230) -- (7.256,0.230) -- (7.270,0.230) -- (7.284,0.230) -- (7.298,0.230) -- (7.312,0.230) -- (7.326,0.230) -- (7.340,0.230) -- (7.354,0.230) -- (7.368,0.230) -- (7.382,0.230) -- (7.396,0.230) -- (7.410,0.230) -- (7.424,0.230) -- (7.438,0.230) -- (7.452,0.230) -- (7.466,0.230) -- (7.480,0.230) -- (7.494,0.230) -- (7.508,0.230) -- (7.522,0.230) -- (7.536,0.230) -- (7.550,0.230) -- (7.564,0.230) -- (7.578,0.230) -- (7.592,0.230) -- (7.606,0.229) -- (7.620,0.229) -- (7.634,0.229) -- (7.648,0.229) -- (7.662,0.229) -- (7.676,0.229) -- (7.690,0.229) -- (7.704,0.229) -- (7.718,0.229) -- (7.732,0.229) -- (7.746,0.229) -- (7.760,0.229) -- (7.774,0.229) -- (7.788,0.229) -- (7.802,0.229) -- (7.816,0.229) -- (7.830,0.229) -- (7.844,0.229) -- (7.858,0.229) -- (7.872,0.229) -- (7.886,0.228) -- (7.900,0.228) -- (7.914,0.228) -- (7.928,0.228) -- (7.942,0.228) -- (7.956,0.228) -- (7.970,0.228) -- (7.984,0.228) -- (7.998,0.228) -- (8.012,0.228) -- (8.026,0.228) -- (8.040,0.228) -- (8.054,0.227) -- (8.068,0.227) -- (8.082,0.227) -- (8.096,0.227) -- (8.110,0.227) -- (8.124,0.227) -- (8.138,0.227) -- (8.152,0.227) -- (8.166,0.226) -- (8.180,0.226) -- (8.194,0.226) -- (8.208,0.226) -- (8.222,0.226) -- (8.236,0.226) -- (8.250,0.225) -- (8.264,0.225) -- (8.278,0.225) -- (8.292,0.225) -- (8.306,0.225) -- (8.320,0.225) -- (8.334,0.224) -- (8.348,0.224) -- (8.362,0.224) -- (8.376,0.224) -- (8.390,0.223) -- (8.404,0.223) -- (8.418,0.223) -- (8.432,0.223) -- (8.446,0.222) -- (8.460,0.222) -- (8.474,0.222) -- (8.488,0.221) -- (8.502,0.221) -- (8.516,0.221) -- (8.530,0.220) -- (8.544,0.220) -- (8.558,0.220) -- (8.572,0.219) -- (8.586,0.219) -- (8.600,0.218) -- (8.614,0.218) -- (8.628,0.218) -- (8.642,0.217) -- (8.656,0.217) -- (8.670,0.216) -- (8.684,0.216) -- (8.698,0.215) -- (8.712,0.215) -- (8.726,0.214) -- (8.740,0.213) -- (8.754,0.213) -- (8.768,0.212) -- (8.782,0.212) -- (8.796,0.211) -- (8.810,0.210) -- (8.824,0.209) -- (8.838,0.209) -- (8.852,0.208) -- (8.866,0.207) -- (8.880,0.206) -- (8.894,0.206) -- (8.908,0.205) -- (8.922,0.204) -- (8.936,0.203) -- (8.950,0.202) -- (8.964,0.201) -- (8.978,0.200) -- (8.992,0.199) -- (9.006,0.198) -- (9.020,0.197) -- (9.034,0.195) -- (9.048,0.194) -- (9.062,0.193) -- (9.076,0.192) -- (9.090,0.190) -- (9.104,0.189) -- (9.118,0.188) -- (9.132,0.186) -- (9.146,0.185) -- (9.160,0.183) -- (9.174,0.181) -- (9.188,0.180) -- (9.202,0.178) -- (9.216,0.176) -- (9.230,0.174) -- (9.244,0.173) -- (9.258,0.171) -- (9.272,0.169) -- (9.286,0.167) -- (9.300,0.164) -- (9.314,0.162) -- (9.328,0.160) -- (9.342,0.158) -- (9.356,0.155) -- (9.370,0.153) -- (9.384,0.150) -- (9.398,0.148) -- (9.412,0.145) -- (9.426,0.142) -- (9.440,0.139) -- (9.454,0.136) -- (9.468,0.133) -- (9.482,0.130) -- (9.496,0.127) -- (9.510,0.123) -- (9.524,0.120) -- (9.538,0.116) -- (9.552,0.113) -- (9.566,0.109) -- (9.580,0.105) -- (9.594,0.101) -- (9.608,0.097) -- (9.622,0.093) -- (9.636,0.088) -- (9.650,0.084) -- (9.664,0.079) -- (9.678,0.075) -- (9.692,0.070) -- (9.706,0.065) -- (9.720,0.060) -- (9.734,0.054) -- (9.748,0.049) -- (9.762,0.044) -- (9.776,0.038) -- (9.790,0.032) -- (9.804,0.026) -- (9.818,0.020) -- (9.832,0.014) -- (9.846,0.007) -- (9.860,0.000) -- (9.874,-0.006) -- (9.888,-0.013) -- (9.902,-0.021) -- (9.916,-0.028) -- (9.930,-0.036) -- (9.944,-0.043) -- (9.958,-0.051) -- (9.972,-0.059) -- (9.986,-0.068) -- (10.000,-0.076) -- (10.014,-0.085) -- (10.028,-0.094) -- (10.042,-0.103) -- (10.056,-0.112) -- (10.070,-0.122) -- (10.084,-0.132) -- (10.098,-0.142) -- (10.112,-0.152) -- (10.126,-0.162) -- (10.140,-0.173) -- (10.154,-0.184) -- (10.168,-0.195) -- (10.182,-0.207) -- (10.196,-0.218) -- (10.210,-0.230) -- (10.224,-0.242) -- (10.238,-0.255) -- (10.252,-0.267) -- (10.266,-0.280) -- (10.280,-0.293) -- (10.294,-0.307) -- (10.308,-0.321) -- (10.322,-0.335) -- (10.336,-0.349) -- (10.350,-0.363) -- (10.364,-0.378) -- (10.378,-0.393) -- (10.392,-0.408) -- (10.406,-0.424) -- (10.420,-0.440) -- (10.434,-0.456) -- (10.448,-0.472) -- (10.462,-0.489) -- (10.476,-0.506) -- (10.490,-0.523) -- (10.504,-0.540) -- (10.518,-0.558) -- (10.532,-0.576) -- (10.546,-0.595) -- (10.560,-0.613) -- (10.574,-0.632) -- (10.588,-0.651) -- (10.602,-0.670) -- (10.616,-0.690) -- (10.630,-0.710) -- (10.644,-0.730) -- (10.658,-0.751) -- (10.672,-0.771) -- (10.686,-0.792) -- (10.700,-0.814) -- (10.714,-0.835) -- (10.728,-0.857) -- (10.742,-0.879) -- (10.756,-0.901) -- (10.770,-0.923) -- (10.784,-0.946) -- (10.798,-0.969) -- (10.812,-0.992) -- (10.826,-1.015) -- (10.840,-1.039) -- (10.854,-1.062) -- (10.868,-1.086) -- (10.882,-1.110) -- (10.896,-1.135) -- (10.910,-1.159) -- (10.924,-1.184) -- (10.938,-1.209) -- (10.952,-1.234) -- (10.966,-1.259) -- (10.980,-1.284) -- (10.994,-1.309) -- (11.008,-1.335) -- (11.022,-1.360) -- (11.036,-1.386) -- (11.050,-1.412) -- (11.064,-1.438) -- (11.078,-1.464) -- (11.092,-1.490) -- (11.106,-1.516) -- (11.120,-1.542) -- (11.134,-1.568) -- (11.148,-1.594) -- (11.162,-1.621) -- (11.176,-1.647) -- (11.190,-1.673) -- (11.204,-1.699) -- (11.218,-1.725) -- (11.232,-1.751) -- (11.246,-1.778) -- (11.260,-1.803) -- (11.274,-1.829) -- (11.288,-1.855) -- (11.302,-1.881) -- (11.316,-1.906) -- (11.330,-1.932) -- (11.344,-1.957) -- (11.358,-1.982) -- (11.372,-2.007) -- (11.386,-2.031) -- (11.400,-2.056) -- (11.414,-2.080) -- (11.428,-2.104) -- (11.442,-2.127) -- (11.456,-2.151) -- (11.470,-2.174) -- (11.484,-2.196) -- (11.498,-2.219) -- (11.512,-2.241) -- (11.526,-2.262) -- (11.540,-2.284) -- (11.554,-2.304) -- (11.568,-2.325) -- (11.582,-2.345) -- (11.596,-2.364) -- (11.610,-2.383) -- (11.624,-2.401) -- (11.638,-2.419) -- (11.652,-2.436) -- (11.666,-2.453) -- (11.680,-2.469) -- (11.694,-2.484) -- (11.708,-2.499) -- (11.722,-2.513) -- (11.736,-2.526) -- (11.750,-2.538) -- (11.764,-2.550) -- (11.778,-2.561) -- (11.792,-2.571) -- (11.806,-2.580) -- (11.820,-2.588) -- (11.834,-2.596) -- (11.848,-2.602) -- (11.862,-2.607) -- (11.876,-2.611) -- (11.890,-2.615) -- (11.904,-2.617) -- (11.918,-2.618) -- (11.932,-2.617) -- (11.946,-2.616) -- (11.960,-2.613) -- (11.974,-2.609) -- (11.988,-2.604) -- (12.002,-2.597) -- (12.016,-2.588) -- (12.030,-2.578) -- (12.044,-2.567) -- (12.058,-2.554) -- (12.072,-2.539) -- (12.086,-2.523) -- (12.100,-2.505) -- (12.114,-2.485) -- (12.128,-2.463) -- (12.142,-2.439) -- (12.156,-2.413) -- (12.170,-2.385) -- (12.184,-2.355) -- (12.198,-2.322) -- (12.212,-2.287) -- (12.226,-2.250) -- (12.240,-2.210) -- (12.254,-2.168) -- (12.268,-2.123) -- (12.282,-2.075) -- (12.296,-2.025) -- (12.310,-1.971) -- (12.324,-1.915) -- (12.338,-1.855) -- (12.352,-1.792) -- (12.366,-1.726) -- (12.380,-1.656) -- (12.394,-1.582) -- (12.408,-1.505) -- (12.422,-1.424) -- (12.436,-1.339) -- (12.450,-1.249) -- (12.464,-1.156) -- (12.478,-1.058) -- (12.492,-0.955) -- (12.506,-0.848) -- (12.520,-0.735) -- (12.534,-0.618) -- (12.548,-0.495) -- (12.562,-0.367) -- (12.576,-0.233) -- (12.590,-0.093) -- (12.604,0.053) -- (12.618,0.205) -- (12.632,0.364) -- (12.646,0.530) -- (12.660,0.702) -- (12.674,0.882) -- (12.688,1.070) -- (12.702,1.265) -- (12.716,1.469) -- (12.730,1.681) -- (12.744,1.901) -- (12.758,2.131) -- (12.772,2.370) -- (12.786,2.618) -- (12.800,2.877);}
\newcommand{\figModelCoordinates}{\coordinate (fB) at (0.467,-0.141); \coordinate (fC) at (5.200,2.581); \coordinate (zA) at (8.775,0.212); \coordinate (zB) at (11.693,-2.483);}
\newcommand{\figModelZoomTickThree}{8.075}
\newcommand{\figModelZoomTickTwo}{9.825}
\newcommand{\figModelZoomTickOne}{11.575}

\begin{tikzpicture}[>=Stealth,font=\small]
  \figModelCoordinates
  % (a) linear scale in rho
  \draw[->,gray] (0,0)--(5.8,0) node[right] {$\rho$};
  \draw[->,gray] (0,-0.45)--(0,3.35);
  \figModelFull{thick}
  \draw[gray] (fB |- 0,0.06)--(fB |- 0,-0.06);
  \node[below] at (fB |- 0,-0.2) {$\rho_{\rm b}$};
  \draw[gray] (fC |- 0,0.06)--(fC |- 0,-0.06) node[below] {$\rho_{\rm c}$};
  \draw[gray,dashed] (fC)--(fC |- 0,0);
  \fill (fB) circle (1.5pt) (fC) circle (1.5pt);
  \node[left] at (fC) {$+$};
  \node at (2.9,-0.8) {(a)};
  % (b) logarithmic scale in rho, near rho_a and rho_b
  \begin{scope}
    \clip (7.2,-2.85) rectangle (12.9,1.05);
    \figModelZoom{thick}
  \end{scope}
  \draw[->,gray] (7.2,0)--(13.0,0) node[right] {$\rho$};
  \draw[->,gray] (7.2,-2.85)--(7.2,1.2);
  \foreach \x/\l in {\figModelZoomTickThree/{10^{-3}},\figModelZoomTickTwo/{10^{-2}},
                     \figModelZoomTickOne/{10^{-1}}}
    {\draw[gray] (\x,0.06)--(\x,-0.06) node[below,font=\scriptsize] {$\l$};}
  \draw[gray,dashed] (zA)--(zA |- 0,0);
  \node[above,font=\scriptsize] at (zA |- 0,0.35) {$\rho_{\rm a}$};
  \draw[gray] (zB |- 0,0.06)--(zB |- 0,-0.06);
  \node[above,font=\scriptsize] at (zB |- 0,0.1) {$\rho_{\rm b}$};
  \fill (zA) circle (1.5pt) (zB) circle (1.5pt);
  \node[above right] at (zA) {$+$};
  \node[below] at (zB) {$-$};
  \node at (10.05,-3.3) {(b)};
\end{tikzpicture}
\caption{The model profile \(D_*+g_{\lambda,\Lambda}(\rho)\) of
\eqref{eq:sharpness-prepared-displacement} for \(C_0=65\), computed from
its formula with \(\lambda=\Lambda_0=\e^{-\pi/2}\) and
\(\Lambda=\Lambda_{\rm b}=\e^{\pi}\). The radii
\(\rho_{\rm a}<\rho_{\rm b}<\rho_{\rm c}\) and the constant \(D_*\) are
chosen as in the proof. (a) Linear scale in \(\rho\). (b) Logarithmic scale
in \(\rho\) near \(\rho_{\rm a}\) and \(\rho_{\rm b}\), with the vertical scale
enlarged. The two model zeros on the compact exit interval lie between
the three marked radii. The fourth sign at a larger central height in
\eqref{eq:sharpness-four-source-signs} gives the third displacement zero.}
\label{fig:sharpness-model}
\end{figure}


\emph{Realizing the additive constant.}
The choices above fix the desired normalized constant \(D_*\).
To realize it, we apply the intermediate value theorem to the
unscaled splitting. By
\eqref{eq:sharpness-endpoint-limit},
\[
 \mathcal S_s(A)\longrightarrow \mathcal S_0(A)
\]
uniformly on a fixed small \(A\)-interval. By
\cref{lem:organizing-splitting-transversality}, choose its endpoints
\(A_{\rm lo}<1<A_{\rm hi}\) so that \(\mathcal S_0\) has opposite
signs there and its only zero in the interval is \(1\).
The continuous function
\[
 \Phi_s(A)=\mathcal S_s(A)-\tfrac12s^2a_+^e(A,s)D_*
\]
converges uniformly to \(\mathcal S_0\). The intermediate value theorem
therefore gives \(A_s\in(A_{\rm lo},A_{\rm hi})\) with
\begin{equation}
\label{eq:sharpness-tuned-splitting}
 \mathcal S_s(A_s)=\tfrac12s^2a_+^e(A_s,s)D_*.
\end{equation}
Every accumulation point of any such choices is a zero of
\(\mathcal S_0\), so \(A_s\to1\). Put \(p_s=p_s(A_s)\).

\emph{Three isolated zeros.}
Let \(h_s(\rho)\) be the common central height of
\eqref{eq:sharpness-common-source} at \(p_s\).
By \eqref{eq:sharpness-prepared-displacement},
\eqref{eq:sharpness-three-model-signs}, and
\eqref{eq:sharpness-tuned-splitting}, with the remainder made
smaller than the fixed margins, \(\Delta_{p_s}\) is positive at
\(h_s(\rho_{\rm a})\) and \(h_s(\rho_{\rm c})\), and negative at
\(h_s(\rho_{\rm b})\). On the fixed radius interval, \(h_s\) is
strictly decreasing by the radial construction above.
Moreover \(\mathcal A_+(\rho)\ge c\log(1/s)-C\), so these three
central heights tend to zero.

The three test radii give two zeros. For a third zero, we need a
negative displacement at a larger central height. We obtain it from
the organizing contraction at a fixed height.
At the organizing parameter, \cref{prop:signed-contraction-new}
gives \(\Delta_0'(h)<0\), while
\(\Delta_0(h)\to0\) as \(h\downarrow0\). Choose a fixed small
\(\bar h>0\) whose endpoints lie in \(\operatorname{int}W_\pm\). Then
\(\Delta_0<0\) on \((0,\bar h]\). For each fixed
\(0<\widetilde h<\bar h\), regular dependence on the parameters at
positive central heights (\cref{lem:regular-transitions}, applied to the
field extended by \(\dot p=0\)) gives
\(\Delta_{p_s}\to\Delta_0\) uniformly on \([\widetilde h,\bar h]\).
Hence \(\Delta_{p_s}<0\) on \([\widetilde h,\bar h]\) for small \(s\);
in particular \(\Delta_{p_s}(\bar h)<0\).

All four central heights belong to \(\cH_{p_s}\), by
\cref{prop:central-displacement-interval-new}, and satisfy
\begin{equation}
\label{eq:sharpness-four-source-signs}
 h_s(\rho_{\rm c})<h_s(\rho_{\rm b})<h_s(\rho_{\rm a})<\bar h,
 \qquad
 \operatorname{sgn}\Delta_{p_s}=(+,-,+,-).
\end{equation}
Since the admissible heights form an interval, there is a zero in
each of the three intervening open intervals.

By \cref{prop:central-displacement-interval-new},
\(\Delta_{p_s}\) is real analytic on
\(\operatorname{int}\cH_{p_s}\), and the four signs in
\eqref{eq:sharpness-four-source-signs} exclude an identically zero
displacement. The three selected zeros are therefore isolated and,
by the correspondence in that proposition, give three distinct
isolated periodic orbits.

Let their central heights be \(h_{j,s}\), \(j=1,2,3\).
The first two are trapped between test heights tending to zero.
For the third, negativity of \(\Delta_{p_s}\) on every fixed
\([\widetilde h,\bar h]\) for small \(s\) excludes a positive
limiting height. Thus all three heights tend to zero, including the
zero between \(h_s(\rho_{\rm a})\) and the fixed height \(\bar h\).

\emph{Convergence of the orbit segments.}
To pass from vanishing central heights to convergence of the orbits,
we compare their graphs near infinity and along the two organizing
branches. Regular flow dependence then treats the finite connecting
arc.
The half-orbits of the smallest test central height
\(h_{{\rm low},s}=h_s(\rho_{\rm c})\) have radius \(O(s)\) in the
thin passage, energy \(O(s^2)\) at \(t_s\), and thereafter lie in
the common cone. Hence, for each fixed small \(\zeta>0\), their
graphs satisfy
\[
 0\le U_\sigma(z;h_{{\rm low},s})\le C\zeta
       \quad(h_{{\rm low},s}\le z\le\zeta),
\]
and converge uniformly to \(U_\sigma^0(z;1)\) on
\([\zeta,z_*]\), by \eqref{eq:sharpness-tail-comparison}.
Every selected zero lies above \(h_{{\rm low},s}\). The exact graph
order \eqref{eq:central-graph-order-new} therefore bounds its energy
above by the graph for \(h_{{\rm low},s}\). Its portion with
\(z\le\zeta\) lies in a neighborhood shrinking to the point at
infinity as \(\zeta\to0\).

To obtain convergence along both organizing branches away from
infinity, fix \(0<\varepsilon<\zeta\). Once
\(h_{j,s}<\varepsilon\), the same order gives, for
\(\zeta\le z\le z_*\),
\[
 U_\sigma(z;\varepsilon,p_s)
 \le U_\sigma(z;h_{j,s},p_s)
 \le U_\sigma(z;h_{{\rm low},s},p_s).
\]
By regular dependence, the left graph converges to the
organizing-parameter graph with central height \(\varepsilon\).
As \(\varepsilon\to0\), that graph converges uniformly on
\([\zeta,z_*]\) to \(U_\sigma^0(z;1)\): at
\(z=K\varepsilon\), for fixed large \(K\), both graphs lie in the
common cone with energy difference \(O(\varepsilon)\), and the
integrable variation coefficient \(O(z^{-1/2})\) propagates this
difference to every fixed height. The two bounds give uniform
convergence of both branches of each periodic orbit away from the
point at infinity.

In particular, the endpoints converge to the organizing endpoints.
The exterior passage lies in finitely many regular flow boxes, so
endpoint convergence and regular dependence of the flow give uniform
convergence of the entire finite connecting arc in the compactified
phase space.

\emph{Hausdorff convergence to the whole graphic.}
We now express these segment estimates in the fixed metric on
\(\mathbb S^2\). Let \(\mathscr P_p\) denote the compactification
map from the parameter-dependent \((\eta,z)\)-chart to \(\mathbb S^2\),
and let \(q_\infty=\mathscr P_{p_0}(0,0)\) be the nilpotent point of the
graphic. Uniform smoothness of these charts gives, on a smaller chart, a
constant \(C_{\mathbb S}\) such that
\begin{equation}
\label{eq:sharpness-sphere-tubular-estimate}
 d_{\mathbb S}\bigl(\mathscr P_p(\eta,z),q_\infty\bigr)
 \le C_{\mathbb S}\bigl(|\eta|+|z|+|p-p_0|\bigr).
\end{equation}
Because \(\eta_\sigma=\sigma\sqrt{2U_\sigma}\), the portions of the
three orbits with \(0\le z\le\zeta\) are at distance at most
\(C_{\mathbb S}'\sqrt{\zeta}+o_s(1)\) from \(q_\infty\), uniformly in
\(s\) after the fixed \(\zeta\) is chosen. The same estimate applies to
the corresponding part of \(\GammaI\) (with \(p=p_0\)).

For the reverse approximation near infinity, the central section point
\(\mathscr P_{p_s}(0,h_{j,s})\) belongs to each periodic orbit and tends to
\(q_\infty\). Thus every point of \(\GammaI\) with
\(0\le z\le\zeta\) is within
\(C_{\mathbb S}'\sqrt{\zeta}+o_s(1)\) of that orbit.
On the remaining portions, the graph convergence on \([\zeta,z_*]\)
and the regular-flow-box convergence of the finite connecting arcs
give both inclusions in the Hausdorff estimate. Therefore, for every
fixed \(\zeta>0\),
\[
 \limsup_{s\downarrow0}
 \distH(\overline{\gamma_{j,s}},\GammaI)
 \le C_{\mathbb S}'\sqrt{\zeta},\qquad j=1,2,3.
\]
Letting \(\zeta\downarrow0\) proves
\[
 \distH(\overline{\gamma_{j,s}},\GammaI)\longrightarrow0,
 \qquad j=1,2,3.
\]
Taking any sequence \(s_n\downarrow0\) completes the proof.
\end{proof}

\begin{proof}[Proof of \cref{cor:exact-local-cyclicity-three}]
By \cref{thm:direct-finite-cyclicity}, \(\Cycl(\GammaI)\le3\).
Let \(V\) be a neighborhood of \((1,C_0,0,0,0)\) and let
\(\varepsilon>0\). For large \(n\), the parameter \(p_n\) of
\cref{prop:three-cycle-lower-bound} lies in \(V\), and its three
limit cycles \(\gamma_{j,n}\) satisfy
\(\distH(\overline{\gamma_{j,n}},\GammaI)<\varepsilon\). Hence
\(N(V,\varepsilon)\ge3\) in \eqref{eq:cyclicity-definition}, and
\(\Cycl(\GammaI)\ge3\).
\end{proof}

\subsection{Hyperbolicity and separation of central heights}

\begin{corollary}[Three hyperbolic limit cycles]
\label{cor:hyperbolic-triple}
For sufficiently small positive \(s\), the three periodic
orbits constructed above are hyperbolic. In increasing order
of their central heights, they are attracting, repelling, and
attracting. For each such \(s\), all three persist, with these
stabilities, on an open neighborhood of \(p_s\) in the original
five-parameter space. Such open sets occur arbitrarily close
to the organizing parameter, with all three orbits arbitrarily
close to the whole graphic.
\end{corollary}

\begin{proof}
We obtain simplicity from the zero bound for the logarithmic slope
ratio: the two gaps between the three displacement zeros already
require two slope zeros. By \cref{prop:direct-buffered-negative},
which applies to \(p_s\) by \eqref{eq:sharpness-coefficient-path},
\(\mathcal L\) has at most two distinct zeros on all of \(\cH_{p_s}\).
By \eqref{eq:direct-completed-log-ratio}, where \(I_{+,h}>0\),
\(\Delta_{p_s,h}\) has the sign of \(\e^{\mathcal L}-1\).
Write the three displacement zeros as \(h_1<h_2<h_3\).
Rolle's theorem gives a zero of \(\mathcal L\) in each of
\((h_1,h_2)\) and \((h_2,h_3)\). If
\(\Delta_{p_s,h}(h_i)=0\) for any \(i\), then \(h_i\) would be a
third distinct zero of \(\mathcal L\). Thus all three zeros are
simple. The same argument excludes a fourth displacement zero.

The signs in \eqref{eq:sharpness-four-source-signs} now imply
\[
 \Delta_{p_s,h}(h_1)<0,\qquad
 \Delta_{p_s,h}(h_2)>0,\qquad
 \Delta_{p_s,h}(h_3)<0.
\]
At each fixed point, differentiation of
\(\cP_{p_s}=I_+^{-1}\circ\Pi_{p_s}\circ I_-\) gives
\begin{equation}
\label{eq:sharpness-return-multipliers}
 \cP_{p_s}'(h_i)=\e^{\mathcal L(h_i)}.
\end{equation}
The two negative signs give multipliers in \((0,1)\), and the
positive sign gives a multiplier greater than one. This proves
hyperbolicity and the stated stabilities.

For fixed \(s>0\), regular dependence
(\cref{lem:regular-transitions}, applied to the field extended by
\(\dot p=0\)) makes these finite transverse returns smooth in all
five parameters, with central height as the section coordinate.
Apply the implicit function theorem near the three simple fixed
points and intersect the resulting parameter neighborhoods.
This gives an open set \(V_s\ni p_s\) on which the three distinct
hyperbolic orbits persist. Their stability is unchanged after
reducing \(V_s\).

Given fixed small parameter and whole-graphic neighborhoods,
first take \(s\) small by the lower-bound construction, then
reduce \(V_s\) to retain the geometric margins.
Within the neighborhood and the Hausdorff condition
of \cref{thm:direct-finite-cyclicity}, these are all the
periodic orbits counted there.
\end{proof}

\begin{corollary}[Separation of the third central height]
\label{cor:third-height-separation}
For the parameters \(p_s\) constructed in the proof of
\cref{prop:three-cycle-lower-bound}, write the three central
heights as \(h_{1,s}<h_{2,s}<h_{3,s}\). Then
\begin{equation}
\label{eq:third-height-separation}
 \lim_{s\downarrow0}s\log\frac{h_{3,s}}{h_s(\rho_{\rm a})}=+\infty.
\end{equation}
The same limit holds with \(h_s(\rho_{\rm a})\) replaced by
\(h_{1,s}\) or \(h_{2,s}\).
\end{corollary}

\begin{proof}
We first place the third zero outside each fixed compact exit interval.
Fix \(0<\rho_{\rm lo}<\rho_{\rm a}\). The central-height construction
\eqref{eq:sharpness-common-source} extends to
\([\rho_{\rm lo},\rho_{\rm c}]\); it is strictly decreasing there and
\(h_s(\rho_{\rm lo})\to0\). Hence, for sufficiently small \(s\),
\[
 h_s(\rho_{\rm c})<h_s(\rho_{\rm lo})<\bar h.
\]
Both outer central heights belong to \(\cH_{p_s}\), as shown in the
proof of \cref{prop:three-cycle-lower-bound}. Since this is an interval,
it contains the image of \([\rho_{\rm lo},\rho_{\rm c}]\) under \(h_s\).

By \eqref{eq:sharpness-coefficient-path}, the compact-primary estimate
in \cref{sec:direct-buffered-compact-primary} applies to the chosen
\(p_s\) on \([\rho_{\rm lo},\rho_{\rm c}]\), after reducing \(s\)
with \(\rho_{\rm lo}\) fixed. It gives at most two zeros on the
corresponding height interval. That interval contains \(h_{1,s}\)
and \(h_{2,s}\), so
\[
 h_{3,s}>h_s(\rho_{\rm lo}).
\]
The count includes endpoints; equivalently, one can first apply it
on a slightly larger fixed compact interval.

We now quantify this separation. Subtracting the equations for the
two central heights and using
\eqref{eq:direct-buffered-primary-action} gives
\[
 s\log\frac{h_s(\rho_{\rm lo})}{h_s(\rho_{\rm a})}
  =L_+(\rho_{\rm a})-L_+(\rho_{\rm lo})+O_{\rho_{\rm lo}}(s).
\]
The common action constant cancels, and the terms \(sM_+^0\)
contribute \(O_{\rho_{\rm lo}}(s)\). Given \(K>0\), first fix
\(\rho_{\rm lo}\) so that
\(L_+(\rho_{\rm a})-L_+(\rho_{\rm lo})>K+1\), and then take \(s\)
small. The preceding inequalities give
\(s\log(h_{3,s}/h_s(\rho_{\rm a}))>K\), proving
\eqref{eq:third-height-separation}. Finally,
\(h_{1,s},h_{2,s}<h_s(\rho_{\rm a})\), which proves the last assertion.
\end{proof}

\begin{remark}[A numerical example]
\label{rem:numerical-three-cycles}
Take \(C_0=65\), so that \(\nu=1\) and \(a=1/2\) in
\eqref{eq:sharpness-coefficient-path}, and take \(s=0.05\) and \(A=1\),
with \(\mu\) given by \eqref{eq:direct-coefficient-bridge}. We computed the
return map to the central section numerically, by an adaptive Runge--Kutta
integration of the field in the chart \eqref{eq:projective-coordinates-new},
with \(\log z\) as height variable, and of \eqref{eq:quadratic-family}
between two crossings of a fixed section, with relative tolerance
\(10^{-11}\).

The accompanying
\href{https://github.com/h-lu/I21-reproof/raw/ccbc54f203ef41d891f4e2c890b595955c58b244/supplementary/three-cycles.zip}{numerical supplement}
contains a standalone Python script, its fixed parameters and tolerances,
execution instructions, and the reference table. It uses a Dormand--Prince
5(4) integrator, logarithmic-height bisection, and the same affine field
between section crossings. For central heights with
\(-120\le\log h\le-3.5\), the computation finds three fixed points:
\begin{center}
% Generated by scripts/numerics_three_cycles.py; do not edit.
\begin{tabular}{@{}ccc@{}}
\hline
\(j\) & \(\log h_j\) & multiplier\\
\hline
\(1\) & \(-93.45\) & \(0.095\)\\
\(2\) & \(-34.94\) & \(9.25\)\\
\(3\) & \(-11.39\) & \(0.289\)\\
\hline
\end{tabular}

\end{center}
Here the multiplier is the derivative of the return map at the fixed point.
In increasing central height the orbits are attracting, repelling and
attracting, as in \cref{cor:hyperbolic-triple}, and the largest central
height is much larger than the other two; compare
\cref{cor:third-height-separation}. The computation illustrates the
construction at these moderate parameter values.
\end{remark}


\appendix
\section{Coefficient scales}
\label{app:coordinates}

The parameter cover is expressed in boundary coefficients, so we need
both a local coordinate change from \(\mu_1,\mu_2,\mu_3\) to
\(q_0,q_1,q_2\) and a comparison of their weighted scales. We compute
the coordinate change first and then prove the scale bounds directly.

On \(z=0\), let \(v_{\rm bd}\) be the small root of
\(\mu_3+v_{\rm bd}+v_{\rm bd}^2=0\), and put
\(t=1+2v_{\rm bd}\). Writing \(v=v_{\rm bd}+u\) gives
\[
 \eta=tu+u^2,\qquad
 u=\eta/t-\eta^2/t^3+O(\eta^3).
\]
The boundary identity
\eqref{eq:boundary-drift-conjugacy-new} reads
\[
 q(\eta)=(t+2u)\{P(v_{\rm bd})+P'(v_{\rm bd})u-3v_{\rm bd}u^2-u^3\}.
\]
It follows that
\begin{equation}
\label{eq:direct-coefficient-bridge}
 q_0=tP(v_{\rm bd}),\qquad q_1=P'(v_{\rm bd})+2P(v_{\rm bd})/t,\qquad
 q_2=-3v_{\rm bd}/t+q_1/t^2-4q_0/t^4.
\end{equation}
Differentiating \eqref{eq:direct-coefficient-bridge} at the organizing
point gives the Jacobian and its inverse:
\[
 D_{\mu}(q_0,q_1,q_2)\big|_0
 =\begin{pmatrix}1&0&0\\[1mm]2&1&0\\[1mm]-2&1&3\end{pmatrix},
 \qquad
 D_q(\mu_1,\mu_2,\mu_3)\big|_0
 =\begin{pmatrix}1&0&0\\[1mm]-2&1&0\\[1mm]4/3&-1/3&1/3\end{pmatrix}.
\]
Thus the coefficient change is locally invertible.

Set
\[
 H=\max\{|q_0|^{1/3},\sqrt{|q_1|},|q_2|\},\qquad
 H_\mu=\max\{|\mu_1|^{1/3},\sqrt{|\mu_2|},|\mu_3|\}.
\]
These scales are comparable on a fixed neighborhood.
Indeed \eqref{eq:direct-coefficient-bridge} gives \(|v_{\rm bd}|\le CH\)
when \(1/2<t<3/2\), and the inverse identities
\[
 \mu_2=q_1-2q_0/t^2+3v_{\rm bd}^2,\qquad
 \mu_1=q_0/t-\mu_2v_{\rm bd}+v_{\rm bd}^3,\qquad
 \mu_3=-v_{\rm bd}-v_{\rm bd}^2
\]
give \(H_\mu\le CH\). Conversely \(v_{\rm bd}=O(H_\mu)\),
\(P(v_{\rm bd})=O(H_\mu^3)\), \(P'(v_{\rm bd})=O(H_\mu^2)\), and
\eqref{eq:direct-coefficient-bridge} imply \(H\le CH_\mu\).

\begin{samepage}
For \(H>0\), put \(\bar q_j=q_j/H^{3-j}\). Uniformly on this
compact normalized coefficient set,
\[
 \frac {v_{\rm bd}}H=-\frac{\bar q_2}{3}+O(H),\qquad
 \frac{\mu_3}{H}=\frac{\bar q_2}{3}+O(H),\qquad
 \frac{\mu_2}{H^2}=\bar q_1+\frac{\bar q_2^2}{3}+O(H),
\]
\[
 \frac{\mu_1}{H^3}
   =\bar q_0+\frac{\bar q_1\bar q_2}{3}+\frac{2\bar q_2^3}{27}+O(H),\qquad
 H^{-3}q(HX)=\bar q_0+\bar q_1X+\bar q_2X^2-X^3+O_{C^k}(H).
\]
\end{samepage}
The last estimate, for every fixed \(k\) on every fixed \(X\)-interval,
follows from analyticity and the fact that the boundary drift depends only on
\((\mu_1,\mu_2,\mu_3)\); \(A,C\) remain independent parameters
in the positive-height coefficient \(B\).

\section{Geometric capture and saddle preparation}
\label{app:geometric-saddle-details}

\subsection{The isolating block and whole-graphic capture}
\label{app:graphic-capture-details}

We prove \cref{prop:isolation-capture-new} in two steps. An outward
transverse boundary around the finite equilibrium gives the isolating
block. Near the graphic, flow boxes force crossings of local section
segments, and intersection numbers in a tubular annulus show that
each section is crossed exactly once.

At the finite equilibrium, the linear part of the organizing field \(\mathcal X_0\) is
\[
 \mathsf L_0=\begin{pmatrix}1&-1\\ C_0&0\end{pmatrix}.
\]
The symmetric matrix
\[
 \mathsf Q_0=\begin{pmatrix}
 (C_0+1)/2&-1/2\\[1mm]
 -1/2&(C_0+2)/(2C_0)
 \end{pmatrix}
\]
is positive definite and satisfies
$\mathsf L_0^T\mathsf Q_0+\mathsf Q_0\mathsf L_0=\mathsf I$.  Hence, for
$\mathsf Q(u)=u^T\mathsf Q_0u$,
\[
 \dot{\mathsf Q}(u)=|u|^2+O(|u|^3),
\]
where the dot denotes the derivative along $\mathcal X_0$.
Choose a sufficiently small disk $B_{\rm eq}=\{\mathsf Q\le\rho_{\rm eq}^2\}$,
disjoint from $\GammaI$, for which the vector field points strictly out of
$B_{\rm eq}$.  Then $B_{\mathrm{iso}}=\mathbb S^2\setminus\operatorname{int}B_{\rm eq}$
has a boundary along which the field points strictly inward. This
property persists in a full parameter neighborhood and proves (i).

The periodic orbits of \eqref{eq:quadratic-family} lie in the open upper
hemisphere of $\mathbb S^2$. The graphic
$\GammaI$ is a smooth simple closed curve in $\mathbb S^2$. Away from the nilpotent
point it is the parabola \eqref{eq:organizing-parabola}. In the chart
\eqref{eq:projective-coordinates-new}, whose coordinate $z$ takes both signs
on the sphere, the part of the parabola with $y>0$ is the graph
\[
 z=\frac{\sqrt{1+C_0v^2}-1}{C_0}=\frac{v^2}{2}+O(v^4),\qquad v\ne0,
\]
and its closure passes smoothly through the nilpotent point $(v,z)=(0,0)$.
Hence $\GammaI$ has closed tubular neighborhoods in
$\operatorname{int}B_{\mathrm{iso}}$; each of them is a compact annulus,
since $\mathbb S^2$ is orientable.

Take \(\Sigma_0\) on the affine orbit and the two projective faces
\(\Sigma_{\rm in},\Sigma_{\rm out}\) on its regular branches.
Choose a closed tubular annulus
\(B_{\rm cap}\Subset\operatorname{int}B_{\rm iso}\) so thin that
\[
 S_j=\Sigma_j\cap B_{\rm cap},\qquad
 j\in\{0,\mathrm{in},\mathrm{out}\},
\]
are disjoint proper spanning arcs in the cyclic order of \(\GammaI\).
Choose a smaller tube
\(\overline{\mathcal N'}\Subset B_{\rm cap}\) for which
\(T_j=S_j\cap\overline{\mathcal N'}\) are compact transverse cores.
Around the marked points \(\xi_j=\GammaI\cap S_j\), choose nested,
pairwise disjoint flow boxes \(F'_j\Subset F_j\Subset B_{\rm cap}\),
with \(T_j\subset F_j\), such that every trajectory through \(F'_j\)
meets \(T_j\) before leaving \(F_j\), in positive or negative time.
All crossings of each \(T_j\) have the organizing orientation.
These properties persist for nearby parameters by transversality
and smooth dependence of the flow boxes.

For sufficiently small \(\delta>0\),
\[
 N_\delta(\GammaI)\subset\mathcal N',\qquad
 N_\delta(\xi_j)\subset F'_j,
\]
where \(N_\delta\) denotes the metric neighborhood. If
\(\distH(\overline\gamma,\GammaI)<\delta\), one Hausdorff inclusion
puts \(\gamma\) in \(\mathcal N'\), and the other supplies a point in
each \(F'_j\). Hence \(\gamma\) crosses every \(T_j\) at least once,
and all its intersections with \(S_j\) lie in that core and have the
same sign.

Let \(\pi:B_{\rm cap}\to\GammaI\) be the tubular projection.
The algebraic intersection with a spanning arc is
\(I(\gamma,S_j)=\deg(\pi|_\gamma)\), up to orientation, and is nonzero.
Thus \(\gamma\), an embedded circle when taken with minimal period,
is essential in the annulus. Such a circle represents a primitive
generator of \(H_1(B_{\rm cap};\mathbb Z)\); its intersection number
with each spanning arc has absolute value one. Since all crossings
have the same sign, there is exactly one per arc. Disjoint spanning
arcs are met in their cyclic order, and the flow fixes that order to
be the direction of \(\GammaI\). This proves (ii).

\subsection{Uniform preparation of the interior saddle}
\label{app:positive-saddle-preparation}

For \eqref{eq:direct-positive-scaled-field}, we first straighten the
stable and unstable graphs with uniform finite-order bounds, then
normalize the equations on the two axes. The parameter-dependent
stable manifold theorem \cite[Theorem~9.6]{Teschl2012} supplies the
graphs through the following amplitude comparison.

Fix a finite order \(k\), with extra derivatives for the coordinate
changes below. In the scaled parameters
\(\widehat p=(A,C,\kappa,q_2,\ell)\), where
\(q_0=\ell^2\) and \(q_1=\kappa\ell\), the field extends smoothly
to a neighborhood of a compact parameter box including \(\ell=0\).
Factoring \(B(\ell X,\ell^2S,p)-B(0,0,p)\) by \(\ell\) writes it as
\[
 \mathcal X_{\widehat p}=\mathcal X^0_{B_0(p)}
   +\kappa\mathcal X_1+q_2\mathcal X_2+\ell\mathcal X_{3,\widehat p},
 \qquad \mathcal X_1=(X,0),\quad \mathcal X_2=(X^2,0),
\]
where \(\mathcal X^0_{B_0}\) is \eqref{eq:direct-positive-comparison} and
\(\mathcal X_{3,\widehat p}\) is smooth with uniformly bounded state derivatives on
the fixed state box. Replace the displayed coefficients
\(\kappa,q_2,\ell\) by independent small amplitudes
\(\alpha=(\alpha_1,\alpha_2,\alpha_3)\), keeping \(\widehat p\) fixed.
At \(\alpha=0\), the comparison
saddle has eigenvalues \(1,-1\); reducing the amplitude box gives
a common positive spectral margin. The cited theorem on a finite
cover of the compact \((\widehat p,\alpha)\) set gives stable and unstable
graphs on a common smaller state box, with bounded state derivatives
and one additional amplitude derivative. The mean value theorem from
\(\alpha=0\) to \(\alpha=(\kappa,q_2,\ell)\) gives graph differences
\(O_{C^k}(\kappa+|q_2|+\ell)\). Smooth eigen-coordinates centered
at each saddle have the same bounds; straightening the graphs gives
\[
 x'=xg_{\rm u}(x,y),\qquad y'=-yg_{\rm s}(x,y),\qquad g_{\rm u},g_{\rm s}>c>0.
\]
An orbital time change makes \(x'=x,\ y'=-yg(x,y)\),
where \(g(0,0)=\nu_p\).

Normalize the \(x=0\) equation by the regular coordinate
\[
 \psi(y)=y\exp\int_0^y
           \left(\frac{\nu_p}{g(0,t)}-1\right)\frac{dt}{t}.
\]
The new equation is \(y'=-y[\nu_p+xc_1(x,y)]\).
Multiplication of \(y\) by
\(\exp\int_0^xc_1(t,0)\,dt\) then gives exactly
the prepared equation \eqref{eq:direct-positive-saddle-preparation}.

\section{Notation and terminology index}
\label{app:notation-terminology}
\begingroup
% Use nonshrinking vertical fill in the split table output.
\let\vss\vfil
\small
\begin{longtable}{@{}>{\raggedright\arraybackslash}p{\dimexpr.3\textwidth-6pt}
                >{\raggedright\arraybackslash}p{\dimexpr.7\textwidth-6pt}@{}}
\caption{Notation used in more than one section.}
\label{tab:notation-guide}\\
\hline
Symbol & Meaning\\
\hline
\endfirsthead
\caption[]{Notation used in more than one section (continued).}\\
\hline
Symbol & Meaning\\
\hline
\endhead
\(C_0\); \(A,C,\mu_1,\mu_2,\mu_3\); \(p\), \(\mathcal X_p\)
  & Fixed organizing value; parameters of \eqref{eq:quadratic-family};
    the parameter point \(p=(A,C,\mu_1,\mu_2,\mu_3)\) and its vector field.\\
\(\GammaI\), \(\mathbb S^2\), \(\distH\)
  & The graphic \(I_2^1\); the Poincar\'e sphere; Hausdorff distance on
    \(\mathbb S^2\).\\
\(\varepsilon\), \(V\), \(N(V,\varepsilon)\)
  & Hausdorff radius, parameter neighborhood and the supremum over
    \(p\in V\) of the number of limit cycles satisfying the Hausdorff
    condition in \eqref{eq:cyclicity-definition}.\\
\((v,z)\), \(P\), \(\eta\)
  & Projective chart \eqref{eq:projective-coordinates-new};
    \(P(v)=\mu_1+\mu_2v-v^3\); normal-rate coordinate
    \eqref{eq:normal-rate-coordinate-new}.\\
\(q\), \(q_0,q_1,q_2\), \(R_q\); \(B\), \(B_0\), \(b_0,b_1,b_2\)
  & Boundary drift \(q=q_0+q_1\eta+q_2\eta^2+\eta^3R_q\); positive
    coefficient in \eqref{eq:physical-normal-form-new}, with
    \(B_0=B(0,0,p)\), and the margins
    \(B\ge b_0\), \(B_\eta\ge b_1\),
    \(3R_q+\eta\partial_\eta R_q\le-b_2\) in
    \eqref{eq:normal-sign-margins-new}.\\
\(z_*\), \(W_\pm\), \(\widetilde W_\pm\), \(\Pi_p\), \(I\)
  & Height of the fixed section; capture windows and the enlarged windows
    of \cref{lem:exterior-multiplier-new}; exterior map; section coordinate
    \(I=z_*-\eta^2/2\).\\
\(\Lambda_0\), \(\Lambda_\pm\), \(\Lambda_*\)
  & Limit of the organizing exterior multiplier as \(z_*\downarrow0\);
    bounds \(\Lambda_-\le\Pi_p'\le\Lambda_+\)
    of \cref{lem:exterior-multiplier-new}; the value of \(\Lambda_+\) fixed
    in \cref{sec:direct-parameter-cover}.\\
\(h\), \(U_\sigma(z;h)\), \(\eta_\sigma\), \(\mathcal G_{\sigma,p}\)
  & Central height; energy \(\eta_\sigma^2/2\) of the half-orbit of sign
    \(\sigma\); right side of \eqref{eq:physical-graph-equation-new}.\\
\(I_\sigma\), \(\cH_p\), \(\Delta_p\), \(\cD_p\)
  & Endpoints; interval of admissible central heights; displacement;
    domain of the central return.\\
\(H\), \(s\), \(r\), \(\rho_1\)
  & Coefficient scales of \cref{tab:direct-scale-guide};
    \(\rho_1=\sqrt{q_1}\) is distinct from the annular exit radius
    \(\rho\).\\
\(a\), \(b\), \(\bar c\)
  & Normalized constant, linear and quadratic coefficients of the drift
    in each chart.\\
\(\mathcal L\), \(m_\sigma\), \(M_\sigma\)
  & Logarithmic slope ratio \eqref{eq:direct-completed-log-ratio};
    endpoint slopes \(m_\sigma=hI_{\sigma,h}\) and, in the annulus,
    \(M_\sigma=-\mathcal E I_\sigma\), where \(\mathcal E=\rho\partial_\rho\)
    and \(m_\sigma=sM_\sigma/G_\sigma\) with
    \(G_\sigma=-s\mathcal E\log h_\sigma>0\).\\
\(T\), \(T_Q\), \(c_T\)
  & Logarithmic height coordinate
    \eqref{eq:direct-source-clock-convention}, with all parameter scales
    fixed; \(T=-c_T\log(h/h_{\rm ref})\), \(T_Q=T+s\log Q\) in the positive
    quadratic region.\\
\(c_\pm\), \(C_{\rm cone}\), \(\varpi_t\)
  & Common cone and its first-variation constant; tail multiplier
    \(\varpi_t(z)=(z/t)^{-2q_2}\) of \cref{lem:direct-weighted-tail}.\\
\(c_\sigma^{\rm lin}\), \(\mathfrak c_\sigma(p)\)
  & Root-distance normalization constants in
    \eqref{eq:direct-positive-root-distance-phase}; boundary-action
    constants in \eqref{eq:direct-buffered-primary-action}.\\
\(\delta\), \(\epsilon_3\), \(F\), \(K\)
  & Small parameter and drift in \eqref{eq:direct-height-system}; outward
    boundary drift \(K\), with \(F=K/X\);
    \(\epsilon_3=\ell/\delta\) on the quadratic-prefix scales.\\
\(\mathcal W\), \(w\), \(\chi\)
  & Boundary-flow coordinate, projected state and inverse projection of
    \cref{lem:direct-boundary-flow}.\\
\(\mathcal U\), \(\Theta\), \(\Upsilon\)
  & Scaled energy; logarithmic height correction; positive prefix unit,
    with their scales and arguments specified in each construction.\\
\(\ell\), \(E\), \(A_p\), \(U_p\)
  & Radius scale, amplitude, amplitude constant and dimensionless
    positive unit factor in an endpoint derivative,
    \cref{lem:boundary-flow-endpoint}; in
    \cref{sec:direct-positive-offset}, \(\ell=\sqrt{q_0}\).\\
\(\mathcal A\), \(\mathfrak H\), \(Q_\sigma\)
  & Boundary action integral; logarithmic endpoint rate in a common
    logarithmic height coordinate; curvature coefficient
    \((wK_\sigma''-K_\sigma')/w^2\).\\
\(\widehat h\), \(\widehat h_+\)
  & Rescaled central height \(h/q_1^2\) in
    \cref{sec:direct-bounded-contact}; \(h/q_0\) in
    \cref{sec:direct-positive-offset}.\\
\(\kappa\), \(a_\circ\), \(\kappa_\circ\)
  & Weak-root scale of \eqref{eq:weak-root-scales}; the constants of
    \cref{prop:weak-root-passage}.\\
\(\varrho\), \(\rho\)
  & Positive unscaled radius; normalized annular radius, with
    \(\varrho=s\rho\) as a coordinate relation in the negative quadratic
    chart.\\
\(\varrho_0\), \(\varrho^+\); \(\rho_e\), \(\varrho_e\)
  & Fixed coefficient in \(u=\varrho_0/\varrho\); \(\varrho^+=s+\varrho\);
    normalized and unscaled adaptive exit radii, \(\varrho_e=s\rho_e\).
    The annular exit has unscaled radius \(s\rho\); at the same thin
    threshold, \(\rho\le\rho_e\) in the buffered comparison of
    \cref{sec:direct-buffered-negative}.\\
\(\delta_{\rm an}\), \(\nu\), \(\epsilon_{\rm f}\)
  & Thin threshold of \cref{sec:direct-negative-annulus}; margin of the
    buffered region \eqref{eq:direct-buffered-region}; fold width in
    \eqref{eq:direct-fold-region}.\\
\(\mathcal Y\), \(m_{\rm land}\)
  & Height-to-radius ratio \(\mathcal Y=z/\varrho^2\); fixed geometric
    landing level.\\
\(a_0^{\rm pr},a_0^+,a_0^-\), \(b_{*,+},b_{*,-}\), \(\underline b\)
  & Regional coefficient thresholds in the primary, positive-secondary
    and negative-secondary arguments; see \cref{tab:direct-constant-order}.\\
\(\kappa_{\rm p}\); \(a_{\rm w},\kappa_{\rm w}\); \(a_\pm,\kappa_\pm\);
\(r_{\rm d}\), \(b_{\rm cov}\), \(A_0\), \(m_{\rm cov}\); \(L\)
  & Fixed thresholds, margins and overlaps of the parameter cover;
    their choice order is given in \cref{tab:direct-constant-order}.\\
\(\Lambda\), \(\Lambda_s\), \(\lambda_s\)
  & Boundary-action multiplier; its value on the path of
    \cref{sec:sharpness}; slope coefficient in the normalized
    displacement \eqref{eq:sharpness-prepared-displacement}.\\
\(c\), \(C\)
  & Positive constants whose values may change from line to line, with
    or without indices. The family parameter \(C\) appears only in
    \eqref{eq:quadratic-family}, in \(p\), in the formulas of
    \cref{sec:graphic-central-return}, and as \((A,C)\) or
    \(C=C_0\).\\
\hline
\end{longtable}
\endgroup

\paragraph{Passage terminology.}
\label{tab:terminology}
A central height is \emph{admissible} if it belongs to \(\cH_p\);
the associated half-orbits reach the windows \(W_\pm\).
A \emph{captured} periodic orbit is one satisfying the Hausdorff
condition of \cref{prop:isolation-capture-new}. Its central height
is admissible by \cref{prop:central-displacement-interval-new}.

A \emph{prefix} runs from the central point to a prescribed radius,
and a \emph{tail} from an incoming height to \(z_*\).
The \emph{projected state} is the boundary-flow coordinate \(w\).
A half-orbit is \emph{thin} while its stated normalized height is
below the chosen threshold; its \emph{thin exit} is the first outward
crossing of that threshold. The subsequent \emph{landing} is the first
hit of \(z=m_{\rm land}|\eta|^2\).

A \emph{weak root} approaches the central line relative to the current
coefficient scale. The \emph{fold} is a double root of the positive
secondary drift; a small positive minimum on its no-root side is the
\emph{bottleneck}. The \emph{far root} is the negative branch's first
primary root \(X_*=1+O(r+s)\) in
\cref{sec:negative-branch-far-root}.

\section*{Acknowledgments and AI tool disclosure}
OpenAI’s ChatGPT and Codex were used during the development and preparation
of this work. These tools assisted with exploring proof strategies,
developing and checking mathematical arguments, preparing verification
scripts, and revising the exposition and LaTeX source. They were also used
to help identify relevant literature and potential gaps in intermediate
drafts. The author is responsible for the mathematical claims, references,
and final presentation of the paper.
\bibliographystyle{amsplain}
\bibliography{references}
\end{document}
